\documentclass[10pt,openany]{book}
% One editorial system for the five existing Studium collections.
% Load after \documentclass, before book-specific configuration.
\usepackage[utf8]{inputenc}
\usepackage[T1]{fontenc}
\usepackage{lmodern}
\DeclareUnicodeCharacter{2212}{\ensuremath{-}}
\DeclareUnicodeCharacter{2192}{\ensuremath{\to}}
\DeclareUnicodeCharacter{03B1}{\ensuremath{\alpha}}
\usepackage[spanish,es-nodecimaldot]{babel}
\spanishplainpercent
\AtBeginDocument{\spanishdeactivate{"~<>}}
\usepackage[a4paper,margin=2.5cm,headheight=14pt]{geometry}
\usepackage{amsmath,amssymb,mathtools,booktabs,array,longtable,graphicx,xcolor,tikz,enumitem,microtype,titlesec,fancyhdr,listings,etoolbox}
\usetikzlibrary{arrows.meta,calc,positioning,patterns}
\usepackage[breakable,skins]{tcolorbox}
\usepackage[labelfont=bf]{caption}
\usepackage[hyphens]{url}
\usepackage[hidelinks,unicode,hypertexnames=false]{hyperref}
\hypersetup{pdfauthor={Daniel Miguel Tejedor}}
\addto\captionsspanish{\renewcommand{\contentsname}{Índice general}\renewcommand{\tablename}{Cuadro}}
\AtBeginDocument{\normalfont\DeclareFontShape{T1}{lmr}{bx}{sc}{<->ssub*lmr/m/sc}{}}
\definecolor{studiumInk}{HTML}{142C40}
\definecolor{studiumRule}{HTML}{142C40}
\definecolor{studiumBlue}{RGB}{24,74,130}
\definecolor{studiumCyan}{RGB}{0,184,205}
\definecolor{studiumLight}{RGB}{232,240,250}
\definecolor{studiumTipBack}{RGB}{232,244,236}
\definecolor{studiumTipFrame}{RGB}{27,107,58}
\definecolor{studiumDefBack}{RGB}{232,240,250}
\definecolor{studiumDefFrame}{RGB}{24,74,130}
\definecolor{studiumWorkedBack}{RGB}{255,244,230}
\definecolor{studiumWorkedFrame}{RGB}{138,78,12}
\definecolor{studiumCheckBack}{RGB}{243,236,248}
\definecolor{studiumCheckFrame}{RGB}{96,48,130}
\definecolor{studiumCodeBack}{RGB}{245,245,242}
\definecolor{studiumCodeRule}{RGB}{48,48,52}
\titleformat{\chapter}[display]{\normalfont\filright}{\large\mdseries\scshape\chaptertitlename\ \thechapter}{1ex}{\normalfont\huge\bfseries}
\titlespacing*{\chapter}{0pt}{2.5ex plus 1ex minus .2ex}{2.3ex}
\titleformat{name=\chapter,numberless}[display]{\normalfont\filright}{}{0pt}{\normalfont\huge\bfseries}
\titleformat{\section}{\normalfont\Large\bfseries}{\thesection}{.8em}{}
\titleformat{\subsection}{\normalfont\large\bfseries}{\thesubsection}{.7em}{}
\titleformat{\subsubsection}{\normalfont\normalsize\bfseries}{\thesubsubsection}{.7em}{}
\makeatletter
\renewcommand*\l@section{\@dottedtocline{1}{1.5em}{3.5em}}
\@openrightfalse
\let\cleardoublepage\clearpage
\makeatother
\setcounter{secnumdepth}{2}\setcounter{tocdepth}{1}
\setlength{\parindent}{1.5em}\setlength{\parskip}{3pt plus .5pt}
\setlength{\emergencystretch}{2em}\raggedbottom
\setlist{itemsep=2pt,topsep=4pt,leftmargin=*}
\setlist[description]{style=nextline,leftmargin=0pt,labelwidth=0pt,labelindent=0pt}
\allowdisplaybreaks[2]
\newcounter{studiumfigure}
\makeatletter\@addtoreset{studiumfigure}{chapter}\makeatother
\renewcommand{\thestudiumfigure}{\thechapter.\arabic{studiumfigure}}
\newcommand{\studiumcaption}[2]{\refstepcounter{studiumfigure}\label{#1}\addcontentsline{lof}{figure}{\protect\numberline{\thestudiumfigure}#2}{\small\textbf{Figura~\thestudiumfigure.} #2\par}}
\newcommand{\Raire}{287.05287}
\newcommand{\gzero}{9.80665}
\tcbset{studiumcommon/.style={breakable,boxrule=.65pt,arc=.7pt,left=1.7mm,right=1.7mm,top=1.3mm,bottom=1.3mm,before skip=10pt,after skip=10pt,fonttitle=\bfseries,coltitle=white,coltext=black},
 studiumtip/.style={studiumcommon,colback=studiumTipBack,colframe=studiumTipFrame,colbacktitle=studiumTipFrame},
 studiumdef/.style={studiumcommon,colback=studiumDefBack,colframe=studiumDefFrame,colbacktitle=studiumDefFrame},
 studiumworked/.style={studiumcommon,colback=studiumWorkedBack,colframe=studiumWorkedFrame,colbacktitle=studiumWorkedFrame},
 studiumcheck/.style={studiumcommon,colback=studiumCheckBack,colframe=studiumCheckFrame,colbacktitle=studiumCheckFrame}}
\newtcolorbox{studiumworked}[1]{studiumworked,title={#1}}
\newcommand{\studiumchapterguide}[1]{\par\noindent\textcolor{studiumInk}{\textbf{Orientación del capítulo.}}\enspace\textit{#1}\par\medskip}
\newcommand{\objetivos}[1]{\studiumchapterguide{#1}}
\newcommand{\caja}[2]{\begin{tcolorbox}[studiumdef,title={#1}]#2\end{tcolorbox}}
\newcommand{\resumen}[1]{\begin{tcolorbox}[studiumcheck,title={Síntesis del capítulo}]#1\end{tcolorbox}}
\newcommand{\studiumchaptercheck}{\begin{tcolorbox}[studiumcheck,title={Comprobación final}]
\begin{enumerate}
\item Explica el modelo del capítulo y las hipótesis que permiten utilizarlo.
\item Reproduce un desarrollo o un ejemplo sin consultar su solución y comprueba unidades, signos y condiciones de validez.
\item Resuelve las preguntas o ejercicios disponibles y distingue el resultado calculado de su interpretación física.
\end{enumerate}\end{tcolorbox}}
\newcommand{\figura}[2]{\begin{figure}[tbp]\centering\includegraphics[width=.83\linewidth]{#1.pdf}\caption{#2}\label{fig:#1}\end{figure}\noindent La figura~\ref{fig:#1} ofrece una representación del modelo discutido.}
\newsavebox{\studiumMathBox}
\newcommand{\fitmath}[1]{\sbox{\studiumMathBox}{$\displaystyle #1$}\ifdim\wd\studiumMathBox>\linewidth\resizebox{\linewidth}{!}{\usebox{\studiumMathBox}}\else\usebox{\studiumMathBox}\fi}
\pagestyle{fancy}\fancyhf{}
\fancyhead[LE]{\footnotesize\thepage\hspace{1.5em}\mdseries\scshape\nouppercase{\leftmark}}
\fancyhead[RO]{\footnotesize\mdseries\scshape\nouppercase{\leftmark}\hspace{1.5em}\thepage}
\renewcommand{\headrulewidth}{.4pt}\renewcommand{\footrulewidth}{0pt}
\fancypagestyle{plain}{\fancyhf{}\fancyfoot[C]{\small\thepage}\renewcommand{\headrulewidth}{0pt}}
\lstdefinestyle{studium}{basicstyle=\ttfamily\small,columns=fullflexible,keepspaces=true,showstringspaces=false,breaklines=true,breakatwhitespace=false,frame=leftline,framerule=1.15pt,rulecolor=\color{studiumCodeRule},backgroundcolor=\color{studiumCodeBack},xleftmargin=1.15em,framexleftmargin=.75em,aboveskip=.9em,belowskip=.7em,tabsize=4,keywordstyle=\ttfamily\bfseries,commentstyle=\ttfamily\itshape,stringstyle=\ttfamily}
\lstset{style=studium}
% Shared cover: every letter remains editable vector text; photographs keep original pixels.
\providecommand{\studiumcoverphoto}{}
\providecommand{\studiumcoveredition}{Primera edición}
\providecommand{\studiumcovercategory}{MANUAL UNIVERSITARIO}
\providecommand{\studiumphotocredit}{Crédito y procedencia en el proyecto del libro.}
\providecommand{\studiumphototop}{0}
\newsavebox{\studiumPhotoBox}
\newcommand{\studiumcover}[3]{%
\hypersetup{pageanchor=false}\begin{titlepage}\thispagestyle{empty}
\begin{tikzpicture}[remember picture,overlay,x=1pt,y=-1pt]
\begin{scope}[shift={(current page.north west)}]
\fill[studiumInk] (0,0) rectangle (595.276,841.89);
\begin{scope}
\clip (0,89.5) rectangle (595.276,585.5);
\sbox{\studiumPhotoBox}{\includegraphics[width=\paperwidth]{\studiumcoverphoto}}%
\ifdim\ht\studiumPhotoBox<496pt
 \node[anchor=north,inner sep=0pt] at (297.638,89.5){\includegraphics[height=496pt]{\studiumcoverphoto}};
\else
 \node[anchor=north,inner sep=0pt] at (297.638,\studiumphototop){\usebox{\studiumPhotoBox}};
\fi
\end{scope}
\fill[white] (0,0) rectangle (595.276,89.5);
\node[anchor=west,text=studiumInk,font=\sffamily\bfseries\fontsize{18}{20}\selectfont,inner sep=0pt] at (51.223,38.5){STUDIUM};
\fill[studiumCyan] (51.223,61.5) rectangle (80.8,63.7);
\node[anchor=west,text=studiumInk,font=\sffamily\fontsize{9}{11}\selectfont,inner sep=0pt] at (89.49,65.5){BIBLIOTECA DE INGENIERÍA};
\node[anchor=west,text=studiumCyan,font=\sffamily\bfseries\fontsize{10}{12}\selectfont,inner sep=0pt] at (51.223,608){#3};
\node[anchor=north west,text=white,text width=493pt,font=\rmfamily\bfseries\fontsize{32}{35}\selectfont,inner sep=0pt] at (51.223,632){#1};
\node[anchor=north west,text=white!82!studiumInk,text width=493pt,font=\sffamily\fontsize{11}{14}\selectfont,inner sep=0pt] at (51.223,722){#2};
\node[anchor=west,text=white,font=\sffamily\bfseries\fontsize{10}{12}\selectfont,inner sep=0pt] at (51.223,763){GRADO EN INGENIERÍA AEROESPACIAL};
\node[anchor=west,text=white!82!studiumInk,font=\sffamily\fontsize{10}{12}\selectfont,inner sep=0pt] at (51.223,778){Universidad de León};
\draw[white!75!studiumInk,line width=.4pt] (51.223,792) -- (544.05,792);
\node[anchor=west,text=white!82!studiumInk,font=\sffamily\fontsize{10}{12}\selectfont,inner sep=0pt] at (51.223,815){2026};
\node[anchor=east,text=white,font=\sffamily\bfseries\fontsize{10}{12}\selectfont,inner sep=0pt] at (544.05,815){\studiumcoveredition};
\end{scope}\end{tikzpicture}\null\end{titlepage}\hypersetup{pageanchor=true}}
\newcommand{\portada}[2]{\studiumcover{#1}{#2}{\studiumcovercategory}}

\providecommand{\studiumvolumetype}{manual}
\newcommand{\studiumreadingguide}{%
\chapter*{Cómo utilizar este volumen}\addcontentsline{toc}{chapter}{Cómo utilizar este volumen}
\ifdefstring{\studiumvolumetype}{manual}{\noindent\textbf{Ruta de estudio.} Empieza por la orientación del capítulo y sus conceptos. Escribe las hipótesis antes de utilizar las ecuaciones y sigue los desarrollos con sus unidades. Después intenta los problemas y las preguntas de repaso sin consultar las soluciones. Cierra cada capítulo con la comprobación final.}{%
\ifdefstring{\studiumvolumetype}{problemas}{\noindent\textbf{Ruta de trabajo.} Lee el enunciado y copia sus datos con unidades. Define las incógnitas, el modelo y las hipótesis. Intenta resolver el problema antes de consultar el desarrollo. Compara después el resultado, la interpretación física y las comprobaciones.}{\noindent\textbf{Ruta de entrenamiento.} Lee las instrucciones particulares de cada simulacro y reserva el tiempo propuesto. Trabaja primero con los enunciados. Después utiliza las soluciones indicadas en el volumen para revisar el planteamiento, los cálculos y las comprobaciones.}}

\noindent\textbf{Recuadros.} El azul identifica definiciones y resultados; el verde, consejos; el naranja, problemas resueltos; y el violeta, síntesis y autoevaluación. Los apartados técnicos conservan los nombres propios de cada asignatura.

\noindent\textbf{Problemas y simulacros.} Separa enunciado, datos, hipótesis, modelo, desarrollo, resultado y comprobación. En los simulacros trabaja primero con los enunciados y consulta después las soluciones. Sus tiempos y puntuaciones son propuestas de entrenamiento cuando así se indican, sin sustituir las instrucciones oficiales del curso.

\noindent\textbf{Consulta.} El índice general permite localizar los temas. La notación, los formularios, las soluciones y la bibliografía mantienen sus apartados al final del manual cuando forman parte de él.
\section*{Créditos y alcance editorial}\addcontentsline{toc}{section}{Créditos y alcance editorial}
\noindent Daniel Miguel Tejedor. Colección Studium. Edición de estudio independiente para el Grado en Ingeniería Aeroespacial, Universidad de León.

\noindent La compilación y la unificación editorial no añaden una aprobación institucional ni una nueva verificación científica. Cada libro conserva las hipótesis, límites y comprobaciones de sus modelos; los informes registran su alcance y sus pendientes.

\noindent\textbf{Imagen de portada.} \studiumphotocredit
}

% Book-specific metadata; all visual rules live in shared/latex/.
\renewcommand{\studiumcoverphoto}{../figures/cover.jpeg}
\renewcommand{\studiumcoveredition}{Primera edición}
\renewcommand{\studiumphotocredit}{Fotografía conservada del PDF de clase aportado por el usuario. La identificación y el crédito de origen siguen pendientes.}
\renewcommand{\studiumphototop}{0}
\graphicspath{{../figures/}}

\hypersetup{pdftitle={Turbomáquinas}}
\renewcommand{\studiumvolumetype}{problemas}
\begin{document}
\studiumcover{TURBOMÁQUINAS}{Problemas resueltos y comprobaciones}{PROBLEMAS RESUELTOS}
\frontmatter
\studiumreadingguide
\tableofcontents
\mainmatter
\chapter{Fundamentos, clasificación y estados del gas}
\begin{studiumworked}{Problema resuelto C01-01-R}


\noindent\textbf{Enunciado}
\(C01-01-R \cdot\) Estado estático desde remanso. Una estación de ensayo tiene los datos totales y Mach indicados. Obtenga temperatura, presión y densidad estáticas; la incógnita final es la densidad. Datos: \(T_0=318\, \mathrm{K}\); \(p_0=154000\, \mathrm{Pa}\); \(\mathrm{Ma}=0.62\); \(\gamma=1.4\); \(R_g=287\, \mathrm{J}\,\mathrm{kg}^{-1}\,\mathrm{K}^{-1}\).

\noindent\textbf{Datos}
\begin{itemize}
\item $T_0 = 318\,\mathrm{K}$ — temperatura total
\item $p_0 = 154000\,\mathrm{Pa}$ — presión total
\item $\mathrm{Ma} = 0.62$ — Mach absoluto
\item $\gamma = 1.4$ — razón de calores
\item $R_g = 287\,\mathrm{J}\,\mathrm{kg}^{-1}\,\mathrm{K}^{-1}$ — constante del gas
\end{itemize}

\noindent\textbf{Incógnita}
$\rho$, en $\mathrm{kg\,m^{-3}}$; también las magnitudes intermedias desarrolladas.

\noindent\textbf{Modelo físico}

\noindent\textit{Estado estático desde remanso}

\[
\fitmath{\rho=\frac{p_0 \left(\frac{\mathrm{Ma}^{2} \left(\gamma - 1\right)}{2} + 1\right)^{- \frac{1}{\gamma - 1}}}{R_g T_0}}
\]

\noindent\textbf{Hipótesis}
\begin{itemize}
\item Gas perfecto y propiedades constantes.
\item Remanso definido localmente; no se presume isentropía entre estaciones.
\end{itemize}

\noindent\textbf{Desarrollo paso a paso}

\begin{enumerate}

\item El remanso local fija la temperatura estática.

\[\fitmath{T=\frac{T_0}{1+\frac{\left(\gamma-1\right)\cdot {\mathrm{Ma}}^{2}}{2}}}\]

\[\fitmath{T=\frac{(318)}{1+\frac{\left((1.4)-1\right)\cdot {(0.62)}^{2}}{2}}\approx 295.2975\,\mathrm{K}}\]

\item Se aplica el frenado isentrópico local para recuperar la presión.

\[\fitmath{p=\frac{p_0}{{\left(1+\frac{\left(\gamma-1\right)\cdot {\mathrm{Ma}}^{2}}{2}\right)}^{\frac{\gamma}{\gamma-1}}}}\]

\[\fitmath{p=\frac{(154000)}{{\left(1+\frac{\left((1.4)-1\right)\cdot {(0.62)}^{2}}{2}\right)}^{\frac{(1.4)}{(1.4)-1}}}\approx 118832.5\,\mathrm{Pa}}\]

\item La ecuación de estado usa las dos magnitudes estáticas.

\[\fitmath{\rho=\frac{p_0}{R_g\cdot T_0\cdot {\left(1+\frac{\left(\gamma-1\right)\cdot {\mathrm{Ma}}^{2}}{2}\right)}^{\frac{1}{\gamma-1}}}}\]

\[\fitmath{\rho=\frac{(154000)}{(287)\cdot (318)\cdot {\left(1+\frac{\left((1.4)-1\right)\cdot {(0.62)}^{2}}{2}\right)}^{\frac{1}{(1.4)-1}}}\approx 1.402147\,\mathrm{kg/m^{3}}}\]

\end{enumerate}


\noindent\textbf{Sustitución numérica}

\[
\fitmath{\frac{\left(154000\,\mathrm{Pa}\right) \left(\frac{\left(0.62\right)^{2} \left(\left(1.4\right) - 1\right)}{2} + 1\right)^{- \frac{1}{\left(1.4\right) - 1}}}{\left(287\,\mathrm{J}\,\mathrm{kg}^{-1}\,\mathrm{K}^{-1}\right) \left(318\,\mathrm{K}\right)}}
\]

\noindent\textbf{Resultado}

\[
\fitmath{\rho=1.402146609\times10^{0}\,\mathrm{kg}\,\mathrm{m}^{-3}}
\]

\noindent\textbf{Interpretación}
La densidad total no sustituye a la estática en el área de paso. Verifique que T y p son menores que sus valores totales.

\noindent\textbf{Comprobación de los pasos}
\begin{itemize}
\item Paso 2: $295.298$ — reproducido
\item Paso 4: $118832$ — reproducido
\item Paso 6: $1.40215$ — reproducido
\end{itemize}

\noindent\textbf{Comprobación de la respuesta} la aritmética del último paso reproduce la respuesta enunciada a la precisión con que está escrita. Es una reproducción aritmética; no valida el modelo físico ni sus hipótesis.

\noindent\textbf{Limitaciones}
\begin{itemize}
\item Modelo educativo; la aplicabilidad a una máquina real requiere revisión experta.
\end{itemize}

\noindent\textbf{Errores frecuentes}
\begin{itemize}
\item Introducir Celsius en un cociente.
\item Usar presión total y temperatura estática en la misma ecuación de estado.
\end{itemize}
% studium:end-record eyJpZCI6ICJQUkItMDAwMSIsICJzZWN0aW9uIjogImMwMSJ9
\end{studiumworked}

\begin{studiumworked}{Problema resuelto C01-01-P}


\noindent\textbf{Enunciado}
\(C01-01-P \cdot\) Estado estático desde remanso. Ejercicio de transferencia. Una estación de ensayo tiene los datos totales y Mach indicados. Obtenga temperatura, presión y densidad estáticas; la incógnita final es la densidad. Antes de consultar la solución, prediga el efecto de cambiar el primer dato respecto del caso resuelto de esta familia. Justifique la tendencia y explique qué restricción podría convertirse en limitante. Datos: \(T_0=338.67\, \mathrm{K}\); \(p_0=154000\, \mathrm{Pa}\); \(\mathrm{Ma}=0.62\); \(\gamma=1.4\); \(R_g=287\, \mathrm{J}\,\mathrm{kg}^{-1}\,\mathrm{K}^{-1}\).

\noindent\textbf{Datos}
\begin{itemize}
\item $T_0 = 338.67\,\mathrm{K}$ — temperatura total
\item $p_0 = 154000\,\mathrm{Pa}$ — presión total
\item $\mathrm{Ma} = 0.62$ — Mach absoluto
\item $\gamma = 1.4$ — razón de calores
\item $R_g = 287\,\mathrm{J}\,\mathrm{kg}^{-1}\,\mathrm{K}^{-1}$ — constante del gas
\end{itemize}

\noindent\textbf{Incógnita}
$\rho$, en $\mathrm{kg\,m^{-3}}$; también las magnitudes intermedias desarrolladas.

\noindent\textbf{Modelo físico}

\noindent\textit{Estado estático desde remanso}

\[
\fitmath{\rho=\frac{p_0 \left(\frac{\mathrm{Ma}^{2} \left(\gamma - 1\right)}{2} + 1\right)^{- \frac{1}{\gamma - 1}}}{R_g T_0}}
\]

\noindent\textbf{Hipótesis}
\begin{itemize}
\item Gas perfecto y propiedades constantes.
\item Remanso definido localmente; no se presume isentropía entre estaciones.
\end{itemize}

\noindent\textbf{Desarrollo paso a paso}

\begin{enumerate}

\item El remanso local fija la temperatura estática.

\[\fitmath{T=\frac{T_0}{1+\frac{\left(\gamma-1\right)\cdot {\mathrm{Ma}}^{2}}{2}}}\]

\[\fitmath{T=\frac{(338.67)}{1+\frac{\left((1.4)-1\right)\cdot {(0.62)}^{2}}{2}}\approx 314.4919\,\mathrm{K}}\]

\item Se aplica el frenado isentrópico local para recuperar la presión.

\[\fitmath{p=\frac{p_0}{{\left(1+\frac{\left(\gamma-1\right)\cdot {\mathrm{Ma}}^{2}}{2}\right)}^{\frac{\gamma}{\gamma-1}}}}\]

\[\fitmath{p=\frac{(154000)}{{\left(1+\frac{\left((1.4)-1\right)\cdot {(0.62)}^{2}}{2}\right)}^{\frac{(1.4)}{(1.4)-1}}}\approx 118832.5\,\mathrm{Pa}}\]

\item La ecuación de estado usa las dos magnitudes estáticas.

\[\fitmath{\rho=\frac{p_0}{R_g\cdot T_0\cdot {\left(1+\frac{\left(\gamma-1\right)\cdot {\mathrm{Ma}}^{2}}{2}\right)}^{\frac{1}{\gamma-1}}}}\]

\[\fitmath{\rho=\frac{(154000)}{(287)\cdot (338.67)\cdot {\left(1+\frac{\left((1.4)-1\right)\cdot {(0.62)}^{2}}{2}\right)}^{\frac{1}{(1.4)-1}}}\approx 1.31657\,\mathrm{kg/m^{3}}}\]

\end{enumerate}


\noindent\textbf{Sustitución numérica}

\[
\fitmath{\frac{\left(154000\,\mathrm{Pa}\right) \left(\frac{\left(0.62\right)^{2} \left(\left(1.4\right) - 1\right)}{2} + 1\right)^{- \frac{1}{\left(1.4\right) - 1}}}{\left(287\,\mathrm{J}\,\mathrm{kg}^{-1}\,\mathrm{K}^{-1}\right) \left(338.67\,\mathrm{K}\right)}}
\]

\noindent\textbf{Resultado}

\[
\fitmath{\rho=1.316569586\times10^{0}\,\mathrm{kg}\,\mathrm{m}^{-3}}
\]

\noindent\textbf{Interpretación}
La densidad total no sustituye a la estática en el área de paso. Verifique que T y p son menores que sus valores totales.

\noindent\textbf{Comprobación de los pasos}
\begin{itemize}
\item Paso 2: $314.492$ — reproducido
\item Paso 4: $118832$ — reproducido
\item Paso 6: $1.31657$ — reproducido
\end{itemize}

\noindent\textbf{Comprobación de la respuesta} la aritmética del último paso reproduce la respuesta enunciada a la precisión con que está escrita. Es una reproducción aritmética; no valida el modelo físico ni sus hipótesis.

\noindent\textbf{Limitaciones}
\begin{itemize}
\item Modelo educativo; la aplicabilidad a una máquina real requiere revisión experta.
\end{itemize}

\noindent\textbf{Errores frecuentes}
\begin{itemize}
\item Introducir Celsius en un cociente.
\item Usar presión total y temperatura estática en la misma ecuación de estado.
\end{itemize}
% studium:end-record eyJpZCI6ICJQUkItMDAwMiIsICJzZWN0aW9uIjogImMwMSJ9
\end{studiumworked}

\begin{studiumworked}{Problema resuelto C01-02-R}


\noindent\textbf{Enunciado}
\(C01-02-R \cdot\) Difusor adiabático. Un difusor reduce el módulo de velocidad. Halle el aumento de temperatura estática y la temperatura de salida. Datos: \(T_e=291\, \mathrm{K}\); \(V_e=218\, \mathrm{m}\,\mathrm{s}^{-1}\); \(V_s=94\, \mathrm{m}\,\mathrm{s}^{-1}\); \(c_p=1004.5\, \mathrm{J}\,\mathrm{kg}^{-1}\,\mathrm{K}^{-1}\).

\noindent\textbf{Datos}
\begin{itemize}
\item $T_e = 291\,\mathrm{K}$ — temperatura estática de entrada
\item $V_e = 218\,\mathrm{m}\,\mathrm{s}^{-1}$ — velocidad de entrada
\item $V_s = 94\,\mathrm{m}\,\mathrm{s}^{-1}$ — velocidad de salida
\item $c_p = 1004.5\,\mathrm{J}\,\mathrm{kg}^{-1}\,\mathrm{K}^{-1}$ — calor específico
\end{itemize}

\noindent\textbf{Incógnita}
$T_s$, en $\mathrm{K}$; también las magnitudes intermedias desarrolladas.

\noindent\textbf{Modelo físico}

\noindent\textit{Difusor adiabático}

\[
\fitmath{T_s=T_e + \frac{V_e^{2} - V_s^{2}}{2 c_p}}
\]

\noindent\textbf{Hipótesis}
\begin{itemize}
\item Estacionario, adiabático, sin trabajo ni cambio de cota.
\end{itemize}

\noindent\textbf{Desarrollo paso a paso}

\begin{enumerate}

\item Sin trabajo ni calor, la entalpía total permanece constante.

\[\fitmath{\mathrm{dT}=\frac{{\mathrm{Ve}}^{2}-{\mathrm{Vs}}^{2}}{2\cdot \mathrm{cp}}}\]

\[\fitmath{\mathrm{dT}=\frac{{(218)}^{2}-{(94)}^{2}}{2\cdot (1004.5)}\approx 19.25734\,\mathrm{K}}\]

\item Se suma el aumento estático al estado de entrada.

\[\fitmath{\mathrm{Ts}=\mathrm{Te}+\frac{{\mathrm{Ve}}^{2}-{\mathrm{Vs}}^{2}}{2\cdot \mathrm{cp}}}\]

\[\fitmath{\mathrm{Ts}=(291)+\frac{{(218)}^{2}-{(94)}^{2}}{2\cdot (1004.5)}\approx 310.2573\,\mathrm{K}}\]

\end{enumerate}


\noindent\textbf{Sustitución numérica}

\[
\fitmath{\left(291\,\mathrm{K}\right) + \frac{\left(\left(218\,\mathrm{m}\,\mathrm{s}^{-1}\right)\right)^{2} - \left(\left(94\,\mathrm{m}\,\mathrm{s}^{-1}\right)\right)^{2}}{2 \left(1004.5\,\mathrm{J}\,\mathrm{kg}^{-1}\,\mathrm{K}^{-1}\right)}}
\]

\noindent\textbf{Resultado}

\[
\fitmath{T_s=3.102573420\times10^{2}\,\mathrm{K}}
\]

\noindent\textbf{Interpretación}
La temperatura estática aumenta aunque no se suministra trabajo al gas. La presión final no puede calcularse sin una condición de pérdida.

\noindent\textbf{Comprobación de los pasos}
\begin{itemize}
\item Paso 2: $19.2573$ — reproducido
\item Paso 4: $310.257$ — reproducido
\end{itemize}

\noindent\textbf{Comprobación de la respuesta} la aritmética del último paso reproduce la respuesta enunciada a la precisión con que está escrita. Es una reproducción aritmética; no valida el modelo físico ni sus hipótesis.

\noindent\textbf{Limitaciones}
\begin{itemize}
\item Modelo educativo; la aplicabilidad a una máquina real requiere revisión experta.
\end{itemize}

\noindent\textbf{Errores frecuentes}
\begin{itemize}
\item Interpretar el aumento estático como trabajo de eje.
\end{itemize}
% studium:end-record eyJpZCI6ICJQUkItMDAwMyIsICJzZWN0aW9uIjogImMwMSJ9
\end{studiumworked}

\begin{studiumworked}{Problema resuelto C01-02-P}


\noindent\textbf{Enunciado}
\(C01-02-P \cdot\) Difusor adiabático. Ejercicio de transferencia. Un difusor reduce el módulo de velocidad. Halle el aumento de temperatura estática y la temperatura de salida. Antes de consultar la solución, prediga el efecto de cambiar el primer dato respecto del caso resuelto de esta familia. Justifique la tendencia y explique qué restricción podría convertirse en limitante. Datos: \(T_e=309.915\, \mathrm{K}\); \(V_e=218\, \mathrm{m}\,\mathrm{s}^{-1}\); \(V_s=94\, \mathrm{m}\,\mathrm{s}^{-1}\); \(c_p=1004.5\, \mathrm{J}\,\mathrm{kg}^{-1}\,\mathrm{K}^{-1}\).

\noindent\textbf{Datos}
\begin{itemize}
\item $T_e = 309.915\,\mathrm{K}$ — temperatura estática de entrada
\item $V_e = 218\,\mathrm{m}\,\mathrm{s}^{-1}$ — velocidad de entrada
\item $V_s = 94\,\mathrm{m}\,\mathrm{s}^{-1}$ — velocidad de salida
\item $c_p = 1004.5\,\mathrm{J}\,\mathrm{kg}^{-1}\,\mathrm{K}^{-1}$ — calor específico
\end{itemize}

\noindent\textbf{Incógnita}
$T_s$, en $\mathrm{K}$; también las magnitudes intermedias desarrolladas.

\noindent\textbf{Modelo físico}

\noindent\textit{Difusor adiabático}

\[
\fitmath{T_s=T_e + \frac{V_e^{2} - V_s^{2}}{2 c_p}}
\]

\noindent\textbf{Hipótesis}
\begin{itemize}
\item Estacionario, adiabático, sin trabajo ni cambio de cota.
\end{itemize}

\noindent\textbf{Desarrollo paso a paso}

\begin{enumerate}

\item Sin trabajo ni calor, la entalpía total permanece constante.

\[\fitmath{\mathrm{dT}=\frac{{\mathrm{Ve}}^{2}-{\mathrm{Vs}}^{2}}{2\cdot \mathrm{cp}}}\]

\[\fitmath{\mathrm{dT}=\frac{{(218)}^{2}-{(94)}^{2}}{2\cdot (1004.5)}\approx 19.25734\,\mathrm{K}}\]

\item Se suma el aumento estático al estado de entrada.

\[\fitmath{\mathrm{Ts}=\mathrm{Te}+\frac{{\mathrm{Ve}}^{2}-{\mathrm{Vs}}^{2}}{2\cdot \mathrm{cp}}}\]

\[\fitmath{\mathrm{Ts}=(309.915)+\frac{{(218)}^{2}-{(94)}^{2}}{2\cdot (1004.5)}\approx 329.1723\,\mathrm{K}}\]

\end{enumerate}


\noindent\textbf{Sustitución numérica}

\[
\fitmath{\left(309.915\,\mathrm{K}\right) + \frac{\left(\left(218\,\mathrm{m}\,\mathrm{s}^{-1}\right)\right)^{2} - \left(\left(94\,\mathrm{m}\,\mathrm{s}^{-1}\right)\right)^{2}}{2 \left(1004.5\,\mathrm{J}\,\mathrm{kg}^{-1}\,\mathrm{K}^{-1}\right)}}
\]

\noindent\textbf{Resultado}

\[
\fitmath{T_s=3.291723420\times10^{2}\,\mathrm{K}}
\]

\noindent\textbf{Interpretación}
La temperatura estática aumenta aunque no se suministra trabajo al gas. La presión final no puede calcularse sin una condición de pérdida.

\noindent\textbf{Comprobación de los pasos}
\begin{itemize}
\item Paso 2: $19.2573$ — reproducido
\item Paso 4: $329.172$ — reproducido
\end{itemize}

\noindent\textbf{Comprobación de la respuesta} la aritmética del último paso reproduce la respuesta enunciada a la precisión con que está escrita. Es una reproducción aritmética; no valida el modelo físico ni sus hipótesis.

\noindent\textbf{Limitaciones}
\begin{itemize}
\item Modelo educativo; la aplicabilidad a una máquina real requiere revisión experta.
\end{itemize}

\noindent\textbf{Errores frecuentes}
\begin{itemize}
\item Interpretar el aumento estático como trabajo de eje.
\end{itemize}
% studium:end-record eyJpZCI6ICJQUkItMDAwNCIsICJzZWN0aW9uIjogImMwMSJ9
\end{studiumworked}

\begin{studiumworked}{Problema resuelto C01-03-R}


\noindent\textbf{Enunciado}
\(C01-03-R \cdot\) Entropía en fila fija. Una fila fija conserva temperatura total y reduce la presión total con la razón indicada. Halle la generación específica de entropía. Datos: \(\frac{p_{0s}}{p_{0e}}=0.963\); \(R_g=287\, \mathrm{J}\,\mathrm{kg}^{-1}\,\mathrm{K}^{-1}\).

\noindent\textbf{Datos}
\begin{itemize}
\item $\frac{p_{0s}}{p_{0e}} = 0.963$ — presión total salida/entrada
\item $R_g = 287\,\mathrm{J}\,\mathrm{kg}^{-1}\,\mathrm{K}^{-1}$ — constante del gas
\end{itemize}

\noindent\textbf{Incógnita}
$\Delta s$, en $\mathrm{J\,kg^{-1}\,K^{-1}}$; también las magnitudes intermedias desarrolladas.

\noindent\textbf{Modelo físico}

\noindent\textit{Entropía en fila fija}

\[
\fitmath{\Delta s=- R_g \log{\left(\frac{p_{0s}}{p_{0e}} \right)}}
\]

\noindent\textbf{Hipótesis}
\begin{itemize}
\item Gas perfecto, adiabático, sin trabajo ni mezcla de especies.
\end{itemize}

\noindent\textbf{Desarrollo paso a paso}

\begin{enumerate}

\item La igualdad de temperaturas totales elimina el término térmico de la entropía.

\[\fitmath{L=-\ln\left(\mathrm{ratio}\right)}\]

\[\fitmath{L=-\ln\left((0.963)\right)\approx 0.03770187\,}\]

\item El descenso de presión produce generación positiva.

\[\fitmath{\mathrm{ds}=-R_g\cdot \ln\left(\mathrm{ratio}\right)}\]

\[\fitmath{\mathrm{ds}=-(287)\cdot \ln\left((0.963)\right)\approx 10.82044\,\mathrm{J/(kg\,K)}}\]

\end{enumerate}


\noindent\textbf{Sustitución numérica}

\[
\fitmath{- \left(287\,\mathrm{J}\,\mathrm{kg}^{-1}\,\mathrm{K}^{-1}\right) \log{\left(\left(0.963\right) \right)}}
\]

\noindent\textbf{Resultado}

\[
\fitmath{\Delta s=1.082043588\times10^{1}\,\mathrm{J}\,\mathrm{kg}^{-1}\,\mathrm{K}^{-1}}
\]

\noindent\textbf{Interpretación}
El signo debe ser positivo para la fila adiabática pasiva. No se ha perdido energía total hacia el ambiente.

\noindent\textbf{Comprobación de los pasos}
\begin{itemize}
\item Paso 2: $0.0377019$ — reproducido
\item Paso 4: $10.8204$ — reproducido
\end{itemize}

\noindent\textbf{Comprobación de la respuesta} la aritmética del último paso reproduce la respuesta enunciada a la precisión con que está escrita. Es una reproducción aritmética; no valida el modelo físico ni sus hipótesis.

\noindent\textbf{Limitaciones}
\begin{itemize}
\item Modelo educativo; la aplicabilidad a una máquina real requiere revisión experta.
\end{itemize}

\noindent\textbf{Errores frecuentes}
\begin{itemize}
\item Invertir la razón y obtener entropía negativa.
\end{itemize}
% studium:end-record eyJpZCI6ICJQUkItMDAwNSIsICJzZWN0aW9uIjogImMwMSJ9
\end{studiumworked}

\begin{studiumworked}{Problema resuelto C01-03-P}


\noindent\textbf{Enunciado}
\(C01-03-P \cdot\) Entropía en fila fija. Ejercicio de transferencia. Una fila fija conserva temperatura total y reduce la presión total con la razón indicada. Halle la generación específica de entropía. Antes de consultar la solución, prediga el efecto de cambiar el primer dato respecto del caso resuelto de esta familia. Justifique la tendencia y explique qué restricción podría convertirse en limitante. Datos: \(\frac{p_{0s}}{p_{0e}}=0.947\); \(R_g=287\, \mathrm{J}\,\mathrm{kg}^{-1}\,\mathrm{K}^{-1}\).

\noindent\textbf{Datos}
\begin{itemize}
\item $\frac{p_{0s}}{p_{0e}} = 0.947$ — presión total salida/entrada
\item $R_g = 287\,\mathrm{J}\,\mathrm{kg}^{-1}\,\mathrm{K}^{-1}$ — constante del gas
\end{itemize}

\noindent\textbf{Incógnita}
$\Delta s$, en $\mathrm{J\,kg^{-1}\,K^{-1}}$; también las magnitudes intermedias desarrolladas.

\noindent\textbf{Modelo físico}

\noindent\textit{Entropía en fila fija}

\[
\fitmath{\Delta s=- R_g \log{\left(\frac{p_{0s}}{p_{0e}} \right)}}
\]

\noindent\textbf{Hipótesis}
\begin{itemize}
\item Gas perfecto, adiabático, sin trabajo ni mezcla de especies.
\end{itemize}

\noindent\textbf{Desarrollo paso a paso}

\begin{enumerate}

\item La igualdad de temperaturas totales elimina el término térmico de la entropía.

\[\fitmath{L=-\ln\left(\mathrm{ratio}\right)}\]

\[\fitmath{L=-\ln\left((0.947)\right)\approx 0.05445619\,}\]

\item El descenso de presión produce generación positiva.

\[\fitmath{\mathrm{ds}=-R_g\cdot \ln\left(\mathrm{ratio}\right)}\]

\[\fitmath{\mathrm{ds}=-(287)\cdot \ln\left((0.947)\right)\approx 15.62893\,\mathrm{J/(kg\,K)}}\]

\end{enumerate}


\noindent\textbf{Sustitución numérica}

\[
\fitmath{- \left(287\,\mathrm{J}\,\mathrm{kg}^{-1}\,\mathrm{K}^{-1}\right) \log{\left(\left(0.947\right) \right)}}
\]

\noindent\textbf{Resultado}

\[
\fitmath{\Delta s=1.562892532\times10^{1}\,\mathrm{J}\,\mathrm{kg}^{-1}\,\mathrm{K}^{-1}}
\]

\noindent\textbf{Interpretación}
El signo debe ser positivo para la fila adiabática pasiva. No se ha perdido energía total hacia el ambiente.

\noindent\textbf{Comprobación de los pasos}
\begin{itemize}
\item Paso 2: $0.0544562$ — reproducido
\item Paso 4: $15.6289$ — reproducido
\end{itemize}

\noindent\textbf{Comprobación de la respuesta} la aritmética del último paso reproduce la respuesta enunciada a la precisión con que está escrita. Es una reproducción aritmética; no valida el modelo físico ni sus hipótesis.

\noindent\textbf{Limitaciones}
\begin{itemize}
\item Modelo educativo; la aplicabilidad a una máquina real requiere revisión experta.
\end{itemize}

\noindent\textbf{Errores frecuentes}
\begin{itemize}
\item Invertir la razón y obtener entropía negativa.
\end{itemize}
% studium:end-record eyJpZCI6ICJQUkItMDAwNiIsICJzZWN0aW9uIjogImMwMSJ9
\end{studiumworked}

\begin{studiumworked}{Problema resuelto C01-04-R}


\noindent\textbf{Enunciado}
\(C01-04-R \cdot\) Anillo con remolino. Dimensione el área frontal anular a partir de un módulo absoluto y un ángulo con el eje. Halle también la altura a radio medio fijado; esa altura es la incógnita final. Datos: \(\dot m=19.6\, \mathrm{kg}\,\mathrm{s}^{-1}\); \(\rho=1.38\, \mathrm{kg}\,\mathrm{m}^{-3}\); \(C_e=186\, \mathrm{m}\,\mathrm{s}^{-1}\); \(\alpha_e=0.5585054\, \mathrm{rad}\); \(r_m=0.29\, \mathrm{m}\).

\noindent\textbf{Datos}
\begin{itemize}
\item $\dot m = 19.6\,\mathrm{kg}\,\mathrm{s}^{-1}$ — caudal
\item $\rho = 1.38\,\mathrm{kg}\,\mathrm{m}^{-3}$ — densidad estática
\item $C_e = 186\,\mathrm{m}\,\mathrm{s}^{-1}$ — módulo absoluto
\item $\alpha_e = 0.55850536063819\,\mathrm{rad}$ — ángulo absoluto con el eje
\item $r_m = 0.29\,\mathrm{m}$ — radio medio aritmético
\end{itemize}

\noindent\textbf{Incógnita}
$b$, en $\mathrm{m}$; también las magnitudes intermedias desarrolladas.

\noindent\textbf{Modelo físico}

\noindent\textit{Anillo con remolino}

\[
\fitmath{b=\frac{\dot m}{2 \pi C_e \rho r_m \cos{\left(\alpha_e \right)}}}
\]

\noindent\textbf{Hipótesis}
\begin{itemize}
\item Flujo axial medio uniforme, bloqueo despreciable.
\end{itemize}

\noindent\textbf{Desarrollo paso a paso}

\begin{enumerate}

\item La proyección axial atraviesa el área frontal.

\[\fitmath{\mathrm{Cx}=\mathrm{Vel}\cdot \cos\left(\mathrm{ang}\right)}\]

\[\fitmath{\mathrm{Cx}=(186)\cdot \cos\left((0.5585053606)\right)\approx 157.7369\,\mathrm{m/s}}\]

\item La continuidad fija el área anular.

\[\fitmath{\mathrm{Aa}=\frac{\mathrm{md}}{\rho\cdot \mathrm{Vel}\cdot \cos\left(\mathrm{ang}\right)}}\]

\[\fitmath{\mathrm{Aa}=\frac{(19.6)}{(1.38)\cdot (186)\cdot \cos\left((0.5585053606)\right)}\approx 0.09004167\,\mathrm{m^{2}}}\]

\item La identidad geométrica A=2 pi rm b fija la altura.

\[\fitmath{b=\frac{\mathrm{md}}{\rho\cdot \mathrm{Vel}\cdot \cos\left(\mathrm{ang}\right)\cdot 2\cdot \pi\cdot \mathrm{rm}}}\]

\[\fitmath{b=\frac{(19.6)}{(1.38)\cdot (186)\cdot \cos\left((0.5585053606)\right)\cdot 2\cdot \pi\cdot (0.29)}\approx 0.04941578\,\mathrm{m}}\]

\end{enumerate}


\noindent\textbf{Sustitución numérica}

\[
\fitmath{\frac{\left(19.6\,\mathrm{kg}\,\mathrm{s}^{-1}\right)}{2 \pi \left(186\,\mathrm{m}\,\mathrm{s}^{-1}\right) \left(1.38\,\mathrm{kg}\,\mathrm{m}^{-3}\right) \left(0.29\,\mathrm{m}\right) \cos{\left(\left(0.55850536063819\,\mathrm{rad}\right) \right)}}}
\]

\noindent\textbf{Resultado}

\[
\fitmath{b=4.941578411\times10^{-2}\,\mathrm{m}}
\]

\noindent\textbf{Interpretación}
Usar el módulo de velocidad reduciría indebidamente el área. Reconstruya el caudal con el área obtenida.

\noindent\textbf{Comprobación de los pasos}
\begin{itemize}
\item Paso 2: $157.737$ — reproducido
\item Paso 4: $0.0900417$ — reproducido
\item Paso 6: $0.0494158$ — reproducido
\end{itemize}

\noindent\textbf{Comprobación de la respuesta} la aritmética del último paso reproduce la respuesta enunciada a la precisión con que está escrita. Es una reproducción aritmética; no valida el modelo físico ni sus hipótesis.

\noindent\textbf{Limitaciones}
\begin{itemize}
\item Modelo educativo; la aplicabilidad a una máquina real requiere revisión experta.
\end{itemize}

\noindent\textbf{Errores frecuentes}
\begin{itemize}
\item Usar el seno del ángulo axial.
\item Confundir altura y diámetro.
\end{itemize}
% studium:end-record eyJpZCI6ICJQUkItMDAwNyIsICJzZWN0aW9uIjogImMwMSJ9
\end{studiumworked}

\begin{studiumworked}{Problema resuelto C01-04-P}


\noindent\textbf{Enunciado}
\(C01-04-P \cdot\) Anillo con remolino. Ejercicio de transferencia. Dimensione el área frontal anular a partir de un módulo absoluto y un ángulo con el eje. Halle también la altura a radio medio fijado; esa altura es la incógnita final. Antes de consultar la solución, prediga el efecto de cambiar el primer dato respecto del caso resuelto de esta familia. Justifique la tendencia y explique qué restricción podría convertirse en limitante. Datos: \(\dot m=20.874\, \mathrm{kg}\,\mathrm{s}^{-1}\); \(\rho=1.38\, \mathrm{kg}\,\mathrm{m}^{-3}\); \(C_e=186\, \mathrm{m}\,\mathrm{s}^{-1}\); \(\alpha_e=0.5585054\, \mathrm{rad}\); \(r_m=0.29\, \mathrm{m}\).

\noindent\textbf{Datos}
\begin{itemize}
\item $\dot m = 20.874\,\mathrm{kg}\,\mathrm{s}^{-1}$ — caudal
\item $\rho = 1.38\,\mathrm{kg}\,\mathrm{m}^{-3}$ — densidad estática
\item $C_e = 186\,\mathrm{m}\,\mathrm{s}^{-1}$ — módulo absoluto
\item $\alpha_e = 0.55850536063819\,\mathrm{rad}$ — ángulo absoluto con el eje
\item $r_m = 0.29\,\mathrm{m}$ — radio medio aritmético
\end{itemize}

\noindent\textbf{Incógnita}
$b$, en $\mathrm{m}$; también las magnitudes intermedias desarrolladas.

\noindent\textbf{Modelo físico}

\noindent\textit{Anillo con remolino}

\[
\fitmath{b=\frac{\dot m}{2 \pi C_e \rho r_m \cos{\left(\alpha_e \right)}}}
\]

\noindent\textbf{Hipótesis}
\begin{itemize}
\item Flujo axial medio uniforme, bloqueo despreciable.
\end{itemize}

\noindent\textbf{Desarrollo paso a paso}

\begin{enumerate}

\item La proyección axial atraviesa el área frontal.

\[\fitmath{\mathrm{Cx}=\mathrm{Vel}\cdot \cos\left(\mathrm{ang}\right)}\]

\[\fitmath{\mathrm{Cx}=(186)\cdot \cos\left((0.5585053606)\right)\approx 157.7369\,\mathrm{m/s}}\]

\item La continuidad fija el área anular.

\[\fitmath{\mathrm{Aa}=\frac{\mathrm{md}}{\rho\cdot \mathrm{Vel}\cdot \cos\left(\mathrm{ang}\right)}}\]

\[\fitmath{\mathrm{Aa}=\frac{(20.874)}{(1.38)\cdot (186)\cdot \cos\left((0.5585053606)\right)}\approx 0.09589438\,\mathrm{m^{2}}}\]

\item La identidad geométrica A=2 pi rm b fija la altura.

\[\fitmath{b=\frac{\mathrm{md}}{\rho\cdot \mathrm{Vel}\cdot \cos\left(\mathrm{ang}\right)\cdot 2\cdot \pi\cdot \mathrm{rm}}}\]

\[\fitmath{b=\frac{(20.874)}{(1.38)\cdot (186)\cdot \cos\left((0.5585053606)\right)\cdot 2\cdot \pi\cdot (0.29)}\approx 0.05262781\,\mathrm{m}}\]

\end{enumerate}


\noindent\textbf{Sustitución numérica}

\[
\fitmath{\frac{\left(20.874\,\mathrm{kg}\,\mathrm{s}^{-1}\right)}{2 \pi \left(186\,\mathrm{m}\,\mathrm{s}^{-1}\right) \left(1.38\,\mathrm{kg}\,\mathrm{m}^{-3}\right) \left(0.29\,\mathrm{m}\right) \cos{\left(\left(0.55850536063819\,\mathrm{rad}\right) \right)}}}
\]

\noindent\textbf{Resultado}

\[
\fitmath{b=5.262781008\times10^{-2}\,\mathrm{m}}
\]

\noindent\textbf{Interpretación}
Usar el módulo de velocidad reduciría indebidamente el área. Reconstruya el caudal con el área obtenida.

\noindent\textbf{Comprobación de los pasos}
\begin{itemize}
\item Paso 2: $157.737$ — reproducido
\item Paso 4: $0.0958944$ — reproducido
\item Paso 6: $0.0526278$ — reproducido
\end{itemize}

\noindent\textbf{Comprobación de la respuesta} la aritmética del último paso reproduce la respuesta enunciada a la precisión con que está escrita. Es una reproducción aritmética; no valida el modelo físico ni sus hipótesis.

\noindent\textbf{Limitaciones}
\begin{itemize}
\item Modelo educativo; la aplicabilidad a una máquina real requiere revisión experta.
\end{itemize}

\noindent\textbf{Errores frecuentes}
\begin{itemize}
\item Usar el seno del ángulo axial.
\item Confundir altura y diámetro.
\end{itemize}
% studium:end-record eyJpZCI6ICJQUkItMDAwOCIsICJzZWN0aW9uIjogImMwMSJ9
\end{studiumworked}

\chapter{Momento angular, Euler y triángulos de velocidades}
\begin{studiumworked}{Problema resuelto C02-01-R}


\noindent\textbf{Enunciado}
\(C02-01-R \cdot\) Par de un compresor. Calcule el par comunicado al fluido y la potencia ideal de un rotor a radio constante. La incógnita final es la potencia. Datos: \(\dot m=22\, \mathrm{kg}\,\mathrm{s}^{-1}\); \(r=0.31\, \mathrm{m}\); \(\Omega=970\, \mathrm{s}^{-1}\); \(C_{\theta e}=24\, \mathrm{m}\,\mathrm{s}^{-1}\); \(C_{\theta s}=109\, \mathrm{m}\,\mathrm{s}^{-1}\).

\noindent\textbf{Datos}
\begin{itemize}
\item $\dot m = 22\,\mathrm{kg}\,\mathrm{s}^{-1}$ — caudal
\item $r = 0.31\,\mathrm{m}$ — radio
\item $\Omega = 970\,\mathrm{s}^{-1}$ — velocidad angular en rad/s
\item $C_{\theta e} = 24\,\mathrm{m}\,\mathrm{s}^{-1}$ — remolino de entrada
\item $C_{\theta s} = 109\,\mathrm{m}\,\mathrm{s}^{-1}$ — remolino de salida
\end{itemize}

\noindent\textbf{Incógnita}
$\dot{W}$, en $\mathrm{W}$; también las magnitudes intermedias desarrolladas.

\noindent\textbf{Modelo físico}

\noindent\textit{Par de un compresor}

\[
\fitmath{\dot{W}=\Omega \dot m r \left(- C_{\theta e} + C_{\theta s}\right)}
\]

\noindent\textbf{Hipótesis}
\begin{itemize}
\item Estacionario y sin pares adicionales.
\item Radio constante y flujo uniforme en los cortes.
\end{itemize}

\noindent\textbf{Desarrollo paso a paso}

\begin{enumerate}

\item Se calcula el flujo neto de momento angular.

\[\fitmath{\tau=\mathrm{md}\cdot \mathrm{rad}\cdot \left(\mathrm{cs}-\mathrm{ce}\right)}\]

\[\fitmath{\tau=(22)\cdot (0.31)\cdot \left((109)-(24)\right)\approx 579.7\,\mathrm{N\,m}}\]

\item Se multiplica el par por la velocidad angular.

\[\fitmath{\mathrm{Pow}=\mathrm{md}\cdot \mathrm{rad}\cdot \left(\mathrm{cs}-\mathrm{ce}\right)\cdot \Omega}\]

\[\fitmath{\mathrm{Pow}=(22)\cdot (0.31)\cdot \left((109)-(24)\right)\cdot (970)\approx 562309\,\mathrm{W}}\]

\end{enumerate}


\noindent\textbf{Sustitución numérica}

\[
\fitmath{\left(970\,\mathrm{s}^{-1}\right) \left(22\,\mathrm{kg}\,\mathrm{s}^{-1}\right) \left(0.31\,\mathrm{m}\right) \left(- \left(24\,\mathrm{m}\,\mathrm{s}^{-1}\right) + \left(109\,\mathrm{m}\,\mathrm{s}^{-1}\right)\right)}
\]

\noindent\textbf{Resultado}

\[
\fitmath{\dot{W}=5.623090000\times10^{5}\,\mathrm{W}}
\]

\noindent\textbf{Interpretación}
La potencia sobre el fluido es positiva; la reacción sobre el rotor es negativa.

\noindent\textbf{Comprobación de los pasos}
\begin{itemize}
\item Paso 2: $579.7$ — reproducido
\item Paso 4: $562309$ — reproducido
\end{itemize}

\noindent\textbf{Comprobación de la respuesta} la aritmética del último paso reproduce la respuesta enunciada a la precisión con que está escrita. Es una reproducción aritmética; no valida el modelo físico ni sus hipótesis.

\noindent\textbf{Limitaciones}
\begin{itemize}
\item Modelo educativo; la aplicabilidad a una máquina real requiere revisión experta.
\end{itemize}

\noindent\textbf{Errores frecuentes}
\begin{itemize}
\item Usar velocidades relativas en el volumen fijo.
\end{itemize}
% studium:end-record eyJpZCI6ICJQUkItMDAwOSIsICJzZWN0aW9uIjogImMwMiJ9
\end{studiumworked}

\begin{studiumworked}{Problema resuelto C02-01-P}


\noindent\textbf{Enunciado}
\(C02-01-P \cdot\) Par de un compresor. Ejercicio de transferencia. Calcule el par comunicado al fluido y la potencia ideal de un rotor a radio constante. La incógnita final es la potencia. Antes de consultar la solución, prediga el efecto de cambiar el primer dato respecto del caso resuelto de esta familia. Justifique la tendencia y explique qué restricción podría convertirse en limitante. Datos: \(\dot m=23.43\, \mathrm{kg}\,\mathrm{s}^{-1}\); \(r=0.31\, \mathrm{m}\); \(\Omega=970\, \mathrm{s}^{-1}\); \(C_{\theta e}=24\, \mathrm{m}\,\mathrm{s}^{-1}\); \(C_{\theta s}=109\, \mathrm{m}\,\mathrm{s}^{-1}\).

\noindent\textbf{Datos}
\begin{itemize}
\item $\dot m = 23.43\,\mathrm{kg}\,\mathrm{s}^{-1}$ — caudal
\item $r = 0.31\,\mathrm{m}$ — radio
\item $\Omega = 970\,\mathrm{s}^{-1}$ — velocidad angular en rad/s
\item $C_{\theta e} = 24\,\mathrm{m}\,\mathrm{s}^{-1}$ — remolino de entrada
\item $C_{\theta s} = 109\,\mathrm{m}\,\mathrm{s}^{-1}$ — remolino de salida
\end{itemize}

\noindent\textbf{Incógnita}
$\dot{W}$, en $\mathrm{W}$; también las magnitudes intermedias desarrolladas.

\noindent\textbf{Modelo físico}

\noindent\textit{Par de un compresor}

\[
\fitmath{\dot{W}=\Omega \dot m r \left(- C_{\theta e} + C_{\theta s}\right)}
\]

\noindent\textbf{Hipótesis}
\begin{itemize}
\item Estacionario y sin pares adicionales.
\item Radio constante y flujo uniforme en los cortes.
\end{itemize}

\noindent\textbf{Desarrollo paso a paso}

\begin{enumerate}

\item Se calcula el flujo neto de momento angular.

\[\fitmath{\tau=\mathrm{md}\cdot \mathrm{rad}\cdot \left(\mathrm{cs}-\mathrm{ce}\right)}\]

\[\fitmath{\tau=(23.43)\cdot (0.31)\cdot \left((109)-(24)\right)\approx 617.3805\,\mathrm{N\,m}}\]

\item Se multiplica el par por la velocidad angular.

\[\fitmath{\mathrm{Pow}=\mathrm{md}\cdot \mathrm{rad}\cdot \left(\mathrm{cs}-\mathrm{ce}\right)\cdot \Omega}\]

\[\fitmath{\mathrm{Pow}=(23.43)\cdot (0.31)\cdot \left((109)-(24)\right)\cdot (970)\approx 598859.1\,\mathrm{W}}\]

\end{enumerate}


\noindent\textbf{Sustitución numérica}

\[
\fitmath{\left(970\,\mathrm{s}^{-1}\right) \left(23.43\,\mathrm{kg}\,\mathrm{s}^{-1}\right) \left(0.31\,\mathrm{m}\right) \left(- \left(24\,\mathrm{m}\,\mathrm{s}^{-1}\right) + \left(109\,\mathrm{m}\,\mathrm{s}^{-1}\right)\right)}
\]

\noindent\textbf{Resultado}

\[
\fitmath{\dot{W}=5.988590850\times10^{5}\,\mathrm{W}}
\]

\noindent\textbf{Interpretación}
La potencia sobre el fluido es positiva; la reacción sobre el rotor es negativa.

\noindent\textbf{Comprobación de los pasos}
\begin{itemize}
\item Paso 2: $617.38$ — reproducido
\item Paso 4: $598859$ — reproducido
\end{itemize}

\noindent\textbf{Comprobación de la respuesta} la aritmética del último paso reproduce la respuesta enunciada a la precisión con que está escrita. Es una reproducción aritmética; no valida el modelo físico ni sus hipótesis.

\noindent\textbf{Limitaciones}
\begin{itemize}
\item Modelo educativo; la aplicabilidad a una máquina real requiere revisión experta.
\end{itemize}

\noindent\textbf{Errores frecuentes}
\begin{itemize}
\item Usar velocidades relativas en el volumen fijo.
\end{itemize}
% studium:end-record eyJpZCI6ICJQUkItMDAxMCIsICJzZWN0aW9uIjogImMwMiJ9
\end{studiumworked}

\begin{studiumworked}{Problema resuelto C02-02-R}


\noindent\textbf{Enunciado}
\(C02-02-R \cdot\) Triángulo y referencia relativa. Una estación tiene las componentes absolutas y velocidad periférica indicadas. Obtenga la componente relativa, el módulo relativo y la temperatura total relativa. Datos: \(C_x=158\, \mathrm{m}\,\mathrm{s}^{-1}\); \(C_\theta=76\, \mathrm{m}\,\mathrm{s}^{-1}\); \(U=284\, \mathrm{m}\,\mathrm{s}^{-1}\); \(T_e=306\, \mathrm{K}\); \(c_p=1004.5\, \mathrm{J}\,\mathrm{kg}^{-1}\,\mathrm{K}^{-1}\).

\noindent\textbf{Datos}
\begin{itemize}
\item $C_x = 158\,\mathrm{m}\,\mathrm{s}^{-1}$ — componente axial
\item $C_\theta = 76\,\mathrm{m}\,\mathrm{s}^{-1}$ — componente tangencial absoluta
\item $U = 284\,\mathrm{m}\,\mathrm{s}^{-1}$ — velocidad periférica
\item $T_e = 306\,\mathrm{K}$ — temperatura estática
\item $c_p = 1004.5\,\mathrm{J}\,\mathrm{kg}^{-1}\,\mathrm{K}^{-1}$ — calor específico
\end{itemize}

\noindent\textbf{Incógnita}
$T_{0r}$, en $\mathrm{K}$; también las magnitudes intermedias desarrolladas.

\noindent\textbf{Modelo físico}

\noindent\textit{Triángulo y referencia relativa}

\[
\fitmath{T_{0r}=T_e + \frac{C_x^{2} + \left(C_\theta - U\right)^{2}}{2 c_p}}
\]

\noindent\textbf{Hipótesis}
\begin{itemize}
\item Mismo punto y referencia de rotor con U tangencial.
\end{itemize}

\noindent\textbf{Desarrollo paso a paso}

\begin{enumerate}

\item La transformación vectorial resta U a la componente tangencial.

\[\fitmath{\mathrm{wt}=\mathrm{ct}-\mathrm{uu}}\]

\[\fitmath{\mathrm{wt}=(76)-(284)\approx -208\,\mathrm{m/s}}\]

\item El módulo usa componentes ortogonales.

\[\fitmath{\mathrm{Wr}=\sqrt{{\mathrm{cx}}^{2}+{\left(\mathrm{ct}-\mathrm{uu}\right)}^{2}}}\]

\[\fitmath{\mathrm{Wr}=\sqrt{{(158)}^{2}+{\left((76)-(284)\right)}^{2}}\approx 261.2049\,\mathrm{m/s}}\]

\item El remanso relativo se construye con W.

\[\fitmath{\mathrm{Tzr}=\mathrm{Te}+\frac{{\mathrm{cx}}^{2}+{\left(\mathrm{ct}-\mathrm{uu}\right)}^{2}}{2\cdot \mathrm{cp}}}\]

\[\fitmath{\mathrm{Tzr}=(306)+\frac{{(158)}^{2}+{\left((76)-(284)\right)}^{2}}{2\cdot (1004.5)}\approx 339.9612\,\mathrm{K}}\]

\end{enumerate}


\noindent\textbf{Sustitución numérica}

\[
\fitmath{\left(306\,\mathrm{K}\right) + \frac{\left(\left(158\,\mathrm{m}\,\mathrm{s}^{-1}\right)\right)^{2} + \left(\left(76\,\mathrm{m}\,\mathrm{s}^{-1}\right) - \left(284\,\mathrm{m}\,\mathrm{s}^{-1}\right)\right)^{2}}{2 \left(1004.5\,\mathrm{J}\,\mathrm{kg}^{-1}\,\mathrm{K}^{-1}\right)}}
\]

\noindent\textbf{Resultado}

\[
\fitmath{T_{0r}=3.399611747\times10^{2}\,\mathrm{K}}
\]

\noindent\textbf{Interpretación}
El beta relativo es negativo. La temperatura total relativa no tiene por qué coincidir con la absoluta.

\noindent\textbf{Comprobación de los pasos}
\begin{itemize}
\item Paso 2: $-208$ — reproducido
\item Paso 4: $261.205$ — reproducido
\item Paso 6: $339.961$ — reproducido
\end{itemize}

\noindent\textbf{Comprobación de la respuesta} la aritmética del último paso reproduce la respuesta enunciada a la precisión con que está escrita. Es una reproducción aritmética; no valida el modelo físico ni sus hipótesis.

\noindent\textbf{Limitaciones}
\begin{itemize}
\item Modelo educativo; la aplicabilidad a una máquina real requiere revisión experta.
\end{itemize}

\noindent\textbf{Errores frecuentes}
\begin{itemize}
\item Restar módulos de velocidades.
\end{itemize}
% studium:end-record eyJpZCI6ICJQUkItMDAxMSIsICJzZWN0aW9uIjogImMwMiJ9
\end{studiumworked}

\begin{studiumworked}{Problema resuelto C02-02-P}


\noindent\textbf{Enunciado}
\(C02-02-P \cdot\) Triángulo y referencia relativa. Ejercicio de transferencia. Una estación tiene las componentes absolutas y velocidad periférica indicadas. Obtenga la componente relativa, el módulo relativo y la temperatura total relativa. Antes de consultar la solución, prediga el efecto de cambiar el primer dato respecto del caso resuelto de esta familia. Justifique la tendencia y explique qué restricción podría convertirse en limitante. Datos: \(C_x=168.27\, \mathrm{m}\,\mathrm{s}^{-1}\); \(C_\theta=76\, \mathrm{m}\,\mathrm{s}^{-1}\); \(U=284\, \mathrm{m}\,\mathrm{s}^{-1}\); \(T_e=306\, \mathrm{K}\); \(c_p=1004.5\, \mathrm{J}\,\mathrm{kg}^{-1}\,\mathrm{K}^{-1}\).

\noindent\textbf{Datos}
\begin{itemize}
\item $C_x = 168.27\,\mathrm{m}\,\mathrm{s}^{-1}$ — componente axial
\item $C_\theta = 76\,\mathrm{m}\,\mathrm{s}^{-1}$ — componente tangencial absoluta
\item $U = 284\,\mathrm{m}\,\mathrm{s}^{-1}$ — velocidad periférica
\item $T_e = 306\,\mathrm{K}$ — temperatura estática
\item $c_p = 1004.5\,\mathrm{J}\,\mathrm{kg}^{-1}\,\mathrm{K}^{-1}$ — calor específico
\end{itemize}

\noindent\textbf{Incógnita}
$T_{0r}$, en $\mathrm{K}$; también las magnitudes intermedias desarrolladas.

\noindent\textbf{Modelo físico}

\noindent\textit{Triángulo y referencia relativa}

\[
\fitmath{T_{0r}=T_e + \frac{C_x^{2} + \left(C_\theta - U\right)^{2}}{2 c_p}}
\]

\noindent\textbf{Hipótesis}
\begin{itemize}
\item Mismo punto y referencia de rotor con U tangencial.
\end{itemize}

\noindent\textbf{Desarrollo paso a paso}

\begin{enumerate}

\item La transformación vectorial resta U a la componente tangencial.

\[\fitmath{\mathrm{wt}=\mathrm{ct}-\mathrm{uu}}\]

\[\fitmath{\mathrm{wt}=(76)-(284)\approx -208\,\mathrm{m/s}}\]

\item El módulo usa componentes ortogonales.

\[\fitmath{\mathrm{Wr}=\sqrt{{\mathrm{cx}}^{2}+{\left(\mathrm{ct}-\mathrm{uu}\right)}^{2}}}\]

\[\fitmath{\mathrm{Wr}=\sqrt{{(168.27)}^{2}+{\left((76)-(284)\right)}^{2}}\approx 267.5421\,\mathrm{m/s}}\]

\item El remanso relativo se construye con W.

\[\fitmath{\mathrm{Tzr}=\mathrm{Te}+\frac{{\mathrm{cx}}^{2}+{\left(\mathrm{ct}-\mathrm{uu}\right)}^{2}}{2\cdot \mathrm{cp}}}\]

\[\fitmath{\mathrm{Tzr}=(306)+\frac{{(168.27)}^{2}+{\left((76)-(284)\right)}^{2}}{2\cdot (1004.5)}\approx 341.6291\,\mathrm{K}}\]

\end{enumerate}


\noindent\textbf{Sustitución numérica}

\[
\fitmath{\left(306\,\mathrm{K}\right) + \frac{\left(\left(168.27\,\mathrm{m}\,\mathrm{s}^{-1}\right)\right)^{2} + \left(\left(76\,\mathrm{m}\,\mathrm{s}^{-1}\right) - \left(284\,\mathrm{m}\,\mathrm{s}^{-1}\right)\right)^{2}}{2 \left(1004.5\,\mathrm{J}\,\mathrm{kg}^{-1}\,\mathrm{K}^{-1}\right)}}
\]

\noindent\textbf{Resultado}

\[
\fitmath{T_{0r}=3.416290657\times10^{2}\,\mathrm{K}}
\]

\noindent\textbf{Interpretación}
El beta relativo es negativo. La temperatura total relativa no tiene por qué coincidir con la absoluta.

\noindent\textbf{Comprobación de los pasos}
\begin{itemize}
\item Paso 2: $-208$ — reproducido
\item Paso 4: $267.542$ — reproducido
\item Paso 6: $341.629$ — reproducido
\end{itemize}

\noindent\textbf{Comprobación de la respuesta} la aritmética del último paso reproduce la respuesta enunciada a la precisión con que está escrita. Es una reproducción aritmética; no valida el modelo físico ni sus hipótesis.

\noindent\textbf{Limitaciones}
\begin{itemize}
\item Modelo educativo; la aplicabilidad a una máquina real requiere revisión experta.
\end{itemize}

\noindent\textbf{Errores frecuentes}
\begin{itemize}
\item Restar módulos de velocidades.
\end{itemize}
% studium:end-record eyJpZCI6ICJQUkItMDAxMiIsICJzZWN0aW9uIjogImMwMiJ9
\end{studiumworked}

\begin{studiumworked}{Problema resuelto C02-03-R}


\noindent\textbf{Enunciado}
\(C02-03-R \cdot\) Euler con radios distintos. Un rotor radial trabaja con prerrotación no nula. Calcule el trabajo comunicado al fluido y el incremento de temperatura total. Datos: \(\Omega=1650\, \mathrm{s}^{-1}\); \(r_e=0.075\, \mathrm{m}\); \(r_s=0.19\, \mathrm{m}\); \(C_{\theta e}=31\, \mathrm{m}\,\mathrm{s}^{-1}\); \(C_{\theta s}=247\, \mathrm{m}\,\mathrm{s}^{-1}\); \(c_p=1004.5\, \mathrm{J}\,\mathrm{kg}^{-1}\,\mathrm{K}^{-1}\).

\noindent\textbf{Datos}
\begin{itemize}
\item $\Omega = 1650\,\mathrm{s}^{-1}$ — velocidad angular
\item $r_e = 0.075\,\mathrm{m}$ — radio entrada
\item $r_s = 0.19\,\mathrm{m}$ — radio salida
\item $C_{\theta e} = 31\,\mathrm{m}\,\mathrm{s}^{-1}$ — remolino entrada
\item $C_{\theta s} = 247\,\mathrm{m}\,\mathrm{s}^{-1}$ — remolino salida
\item $c_p = 1004.5\,\mathrm{J}\,\mathrm{kg}^{-1}\,\mathrm{K}^{-1}$ — calor específico
\end{itemize}

\noindent\textbf{Incógnita}
$\Delta T_0$, en $\mathrm{K}$; también las magnitudes intermedias desarrolladas.

\noindent\textbf{Modelo físico}

\noindent\textit{Euler con radios distintos}

\[
\fitmath{\Delta T_0=\frac{\Omega \left(- C_{\theta e} r_e + C_{\theta s} r_s\right)}{c_p}}
\]

\noindent\textbf{Hipótesis}
\begin{itemize}
\item Adiabático, propiedades constantes.
\item Valores medios compatibles con el flujo de momento angular.
\end{itemize}

\noindent\textbf{Desarrollo paso a paso}

\begin{enumerate}

\item La velocidad periférica de salida depende de su radio.

\[\fitmath{\mathrm{Us}=\Omega\cdot \mathrm{rs}}\]

\[\fitmath{\mathrm{Us}=(1650)\cdot (0.19)\approx 313.5\,\mathrm{m/s}}\]

\item Se conserva el término de prerrotación en Euler.

\[\fitmath{\mathrm{wc}=\Omega\cdot \left(\mathrm{rs}\cdot \mathrm{cs}-\mathrm{re}\cdot \mathrm{ce}\right)}\]

\[\fitmath{\mathrm{wc}=(1650)\cdot \left((0.19)\cdot (247)-(0.075)\cdot (31)\right)\approx 73598.25\,\mathrm{J/kg}}\]

\item Se transforma trabajo en aumento de remanso.

\[\fitmath{\mathrm{dTz}=\frac{\Omega\cdot \left(\mathrm{rs}\cdot \mathrm{cs}-\mathrm{re}\cdot \mathrm{ce}\right)}{\mathrm{cp}}}\]

\[\fitmath{\mathrm{dTz}=\frac{(1650)\cdot \left((0.19)\cdot (247)-(0.075)\cdot (31)\right)}{(1004.5)}\approx 73.26854\,\mathrm{K}}\]

\end{enumerate}


\noindent\textbf{Sustitución numérica}

\[
\fitmath{\frac{\left(1650\,\mathrm{s}^{-1}\right) \left(- \left(31\,\mathrm{m}\,\mathrm{s}^{-1}\right) \left(0.075\,\mathrm{m}\right) + \left(247\,\mathrm{m}\,\mathrm{s}^{-1}\right) \left(0.19\,\mathrm{m}\right)\right)}{\left(1004.5\,\mathrm{J}\,\mathrm{kg}^{-1}\,\mathrm{K}^{-1}\right)}}
\]

\noindent\textbf{Resultado}

\[
\fitmath{\Delta T_0=7.326854156\times10^{1}\,\mathrm{K}}
\]

\noindent\textbf{Interpretación}
Omitir el remolino de entrada sobreestima el trabajo en este caso. Compare ambas evaluaciones.

\noindent\textbf{Comprobación de los pasos}
\begin{itemize}
\item Paso 2: $313.5$ — reproducido
\item Paso 4: $73598.2$ — reproducido
\item Paso 6: $73.2685$ — reproducido
\end{itemize}

\noindent\textbf{Comprobación de la respuesta} la aritmética del último paso reproduce la respuesta enunciada a la precisión con que está escrita. Es una reproducción aritmética; no valida el modelo físico ni sus hipótesis.

\noindent\textbf{Limitaciones}
\begin{itemize}
\item Modelo educativo; la aplicabilidad a una máquina real requiere revisión experta.
\end{itemize}

\noindent\textbf{Errores frecuentes}
\begin{itemize}
\item Usar el radio de salida para ambos extremos.
\end{itemize}
% studium:end-record eyJpZCI6ICJQUkItMDAxMyIsICJzZWN0aW9uIjogImMwMiJ9
\end{studiumworked}

\begin{studiumworked}{Problema resuelto C02-03-P}


\noindent\textbf{Enunciado}
\(C02-03-P \cdot\) Euler con radios distintos. Ejercicio de transferencia. Un rotor radial trabaja con prerrotación no nula. Calcule el trabajo comunicado al fluido y el incremento de temperatura total. Antes de consultar la solución, prediga el efecto de cambiar el primer dato respecto del caso resuelto de esta familia. Justifique la tendencia y explique qué restricción podría convertirse en limitante. Datos: \(\Omega=1757.25\, \mathrm{s}^{-1}\); \(r_e=0.075\, \mathrm{m}\); \(r_s=0.19\, \mathrm{m}\); \(C_{\theta e}=31\, \mathrm{m}\,\mathrm{s}^{-1}\); \(C_{\theta s}=247\, \mathrm{m}\,\mathrm{s}^{-1}\); \(c_p=1004.5\, \mathrm{J}\,\mathrm{kg}^{-1}\,\mathrm{K}^{-1}\).

\noindent\textbf{Datos}
\begin{itemize}
\item $\Omega = 1757.25\,\mathrm{s}^{-1}$ — velocidad angular
\item $r_e = 0.075\,\mathrm{m}$ — radio entrada
\item $r_s = 0.19\,\mathrm{m}$ — radio salida
\item $C_{\theta e} = 31\,\mathrm{m}\,\mathrm{s}^{-1}$ — remolino entrada
\item $C_{\theta s} = 247\,\mathrm{m}\,\mathrm{s}^{-1}$ — remolino salida
\item $c_p = 1004.5\,\mathrm{J}\,\mathrm{kg}^{-1}\,\mathrm{K}^{-1}$ — calor específico
\end{itemize}

\noindent\textbf{Incógnita}
$\Delta T_0$, en $\mathrm{K}$; también las magnitudes intermedias desarrolladas.

\noindent\textbf{Modelo físico}

\noindent\textit{Euler con radios distintos}

\[
\fitmath{\Delta T_0=\frac{\Omega \left(- C_{\theta e} r_e + C_{\theta s} r_s\right)}{c_p}}
\]

\noindent\textbf{Hipótesis}
\begin{itemize}
\item Adiabático, propiedades constantes.
\item Valores medios compatibles con el flujo de momento angular.
\end{itemize}

\noindent\textbf{Desarrollo paso a paso}

\begin{enumerate}

\item La velocidad periférica de salida depende de su radio.

\[\fitmath{\mathrm{Us}=\Omega\cdot \mathrm{rs}}\]

\[\fitmath{\mathrm{Us}=(1757.25)\cdot (0.19)\approx 333.8775\,\mathrm{m/s}}\]

\item Se conserva el término de prerrotación en Euler.

\[\fitmath{\mathrm{wc}=\Omega\cdot \left(\mathrm{rs}\cdot \mathrm{cs}-\mathrm{re}\cdot \mathrm{ce}\right)}\]

\[\fitmath{\mathrm{wc}=(1757.25)\cdot \left((0.19)\cdot (247)-(0.075)\cdot (31)\right)\approx 78382.14\,\mathrm{J/kg}}\]

\item Se transforma trabajo en aumento de remanso.

\[\fitmath{\mathrm{dTz}=\frac{\Omega\cdot \left(\mathrm{rs}\cdot \mathrm{cs}-\mathrm{re}\cdot \mathrm{ce}\right)}{\mathrm{cp}}}\]

\[\fitmath{\mathrm{dTz}=\frac{(1757.25)\cdot \left((0.19)\cdot (247)-(0.075)\cdot (31)\right)}{(1004.5)}\approx 78.031\,\mathrm{K}}\]

\end{enumerate}


\noindent\textbf{Sustitución numérica}

\[
\fitmath{\frac{\left(1757.25\,\mathrm{s}^{-1}\right) \left(- \left(31\,\mathrm{m}\,\mathrm{s}^{-1}\right) \left(0.075\,\mathrm{m}\right) + \left(247\,\mathrm{m}\,\mathrm{s}^{-1}\right) \left(0.19\,\mathrm{m}\right)\right)}{\left(1004.5\,\mathrm{J}\,\mathrm{kg}^{-1}\,\mathrm{K}^{-1}\right)}}
\]

\noindent\textbf{Resultado}

\[
\fitmath{\Delta T_0=7.803099676\times10^{1}\,\mathrm{K}}
\]

\noindent\textbf{Interpretación}
Omitir el remolino de entrada sobreestima el trabajo en este caso. Compare ambas evaluaciones.

\noindent\textbf{Comprobación de los pasos}
\begin{itemize}
\item Paso 2: $333.877$ — reproducido
\item Paso 4: $78382.1$ — reproducido
\item Paso 6: $78.031$ — reproducido
\end{itemize}

\noindent\textbf{Comprobación de la respuesta} la aritmética del último paso reproduce la respuesta enunciada a la precisión con que está escrita. Es una reproducción aritmética; no valida el modelo físico ni sus hipótesis.

\noindent\textbf{Limitaciones}
\begin{itemize}
\item Modelo educativo; la aplicabilidad a una máquina real requiere revisión experta.
\end{itemize}

\noindent\textbf{Errores frecuentes}
\begin{itemize}
\item Usar el radio de salida para ambos extremos.
\end{itemize}
% studium:end-record eyJpZCI6ICJQUkItMDAxNCIsICJzZWN0aW9uIjogImMwMiJ9
\end{studiumworked}

\begin{studiumworked}{Problema resuelto C02-04-R}


\noindent\textbf{Enunciado}
\(C02-04-R \cdot\) Comprobación de rotalpía. Para un rotor se conocen estados cinemáticos de entrada y salida. Calcule el cambio de entalpía estática exigido por rotalpía. Datos: \(U_e=192\, \mathrm{m}\,\mathrm{s}^{-1}\); \(U_s=331\, \mathrm{m}\,\mathrm{s}^{-1}\); \(W_e=246\, \mathrm{m}\,\mathrm{s}^{-1}\); \(W_s=214\, \mathrm{m}\,\mathrm{s}^{-1}\).

\noindent\textbf{Datos}
\begin{itemize}
\item $U_e = 192\,\mathrm{m}\,\mathrm{s}^{-1}$ — velocidad periférica entrada
\item $U_s = 331\,\mathrm{m}\,\mathrm{s}^{-1}$ — velocidad periférica salida
\item $W_e = 246\,\mathrm{m}\,\mathrm{s}^{-1}$ — velocidad relativa entrada
\item $W_s = 214\,\mathrm{m}\,\mathrm{s}^{-1}$ — velocidad relativa salida
\end{itemize}

\noindent\textbf{Incógnita}
$\Delta h$, en $\mathrm{J\,kg^{-1}}$; también las magnitudes intermedias desarrolladas.

\noindent\textbf{Modelo físico}

\noindent\textit{Comprobación de rotalpía}

\[
\fitmath{\Delta h=- \frac{U_e^{2}}{2} + \frac{U_s^{2}}{2} + \frac{W_e^{2}}{2} - \frac{W_s^{2}}{2}}
\]

\noindent\textbf{Hipótesis}
\begin{itemize}
\item Rotor adiabático, estacionario en referencia móvil, sin cambio de cota.
\end{itemize}

\noindent\textbf{Desarrollo paso a paso}

\begin{enumerate}

\item Se calcula la contribución de cambio de radio.

\[\fitmath{\mathrm{du}=\frac{{\mathrm{us}}^{2}-{\mathrm{ue}}^{2}}{2}}\]

\[\fitmath{\mathrm{du}=\frac{{(331)}^{2}-{(192)}^{2}}{2}\approx 36348.5\,\mathrm{J/kg}}\]

\item Se evalúa por separado la variación de energía cinética relativa.

\[\fitmath{\mathrm{ref}=\frac{{\mathrm{ws}}^{2}-{\mathrm{we}}^{2}}{2}}\]

\[\fitmath{\mathrm{ref}=\frac{{(214)}^{2}-{(246)}^{2}}{2}\approx -7360\,\mathrm{J/kg}}\]

\item De I constante se despeja la variación estática.

\[\fitmath{\mathrm{dh}=\frac{{\mathrm{us}}^{2}-{\mathrm{ue}}^{2}-{\mathrm{ws}}^{2}+{\mathrm{we}}^{2}}{2}}\]

\[\fitmath{\mathrm{dh}=\frac{{(331)}^{2}-{(192)}^{2}-{(214)}^{2}+{(246)}^{2}}{2}\approx 43708.5\,\mathrm{J/kg}}\]

\end{enumerate}


\noindent\textbf{Sustitución numérica}

\[
\fitmath{- \frac{\left(\left(192\,\mathrm{m}\,\mathrm{s}^{-1}\right)\right)^{2}}{2} + \frac{\left(\left(331\,\mathrm{m}\,\mathrm{s}^{-1}\right)\right)^{2}}{2} + \frac{\left(\left(246\,\mathrm{m}\,\mathrm{s}^{-1}\right)\right)^{2}}{2} - \frac{\left(\left(214\,\mathrm{m}\,\mathrm{s}^{-1}\right)\right)^{2}}{2}}
\]

\noindent\textbf{Resultado}

\[
\fitmath{\Delta h=4.370850000\times10^{4}\,\mathrm{J}\,\mathrm{kg}^{-1}}
\]

\noindent\textbf{Interpretación}
El término centrífugo no puede omitirse cuando U cambia. Reintroduzca dh en la identidad de rotalpía y compruebe residuo cero.

\noindent\textbf{Comprobación de los pasos}
\begin{itemize}
\item Paso 2: $36348.5$ — reproducido
\item Paso 4: $-7360$ — reproducido
\item Paso 6: $43708.5$ — reproducido
\end{itemize}

\noindent\textbf{Comprobación de la respuesta} la aritmética del último paso reproduce la respuesta enunciada a la precisión con que está escrita. Es una reproducción aritmética; no valida el modelo físico ni sus hipótesis.

\noindent\textbf{Limitaciones}
\begin{itemize}
\item Modelo educativo; la aplicabilidad a una máquina real requiere revisión experta.
\end{itemize}

\noindent\textbf{Errores frecuentes}
\begin{itemize}
\item Conservar h+W al cuadrado/2 a radio variable.
\end{itemize}
% studium:end-record eyJpZCI6ICJQUkItMDAxNSIsICJzZWN0aW9uIjogImMwMiJ9
\end{studiumworked}

\begin{studiumworked}{Problema resuelto C02-04-P}


\noindent\textbf{Enunciado}
\(C02-04-P \cdot\) Comprobación de rotalpía. Ejercicio de transferencia. Para un rotor se conocen estados cinemáticos de entrada y salida. Calcule el cambio de entalpía estática exigido por rotalpía. Antes de consultar la solución, prediga el efecto de cambiar el primer dato respecto del caso resuelto de esta familia. Justifique la tendencia y explique qué restricción podría convertirse en limitante. Datos: \(U_e=204.48\, \mathrm{m}\,\mathrm{s}^{-1}\); \(U_s=331\, \mathrm{m}\,\mathrm{s}^{-1}\); \(W_e=246\, \mathrm{m}\,\mathrm{s}^{-1}\); \(W_s=214\, \mathrm{m}\,\mathrm{s}^{-1}\).

\noindent\textbf{Datos}
\begin{itemize}
\item $U_e = 204.48\,\mathrm{m}\,\mathrm{s}^{-1}$ — velocidad periférica entrada
\item $U_s = 331\,\mathrm{m}\,\mathrm{s}^{-1}$ — velocidad periférica salida
\item $W_e = 246\,\mathrm{m}\,\mathrm{s}^{-1}$ — velocidad relativa entrada
\item $W_s = 214\,\mathrm{m}\,\mathrm{s}^{-1}$ — velocidad relativa salida
\end{itemize}

\noindent\textbf{Incógnita}
$\Delta h$, en $\mathrm{J\,kg^{-1}}$; también las magnitudes intermedias desarrolladas.

\noindent\textbf{Modelo físico}

\noindent\textit{Comprobación de rotalpía}

\[
\fitmath{\Delta h=- \frac{U_e^{2}}{2} + \frac{U_s^{2}}{2} + \frac{W_e^{2}}{2} - \frac{W_s^{2}}{2}}
\]

\noindent\textbf{Hipótesis}
\begin{itemize}
\item Rotor adiabático, estacionario en referencia móvil, sin cambio de cota.
\end{itemize}

\noindent\textbf{Desarrollo paso a paso}

\begin{enumerate}

\item Se calcula la contribución de cambio de radio.

\[\fitmath{\mathrm{du}=\frac{{\mathrm{us}}^{2}-{\mathrm{ue}}^{2}}{2}}\]

\[\fitmath{\mathrm{du}=\frac{{(331)}^{2}-{(204.48)}^{2}}{2}\approx 33874.46\,\mathrm{J/kg}}\]

\item Se evalúa por separado la variación de energía cinética relativa.

\[\fitmath{\mathrm{ref}=\frac{{\mathrm{ws}}^{2}-{\mathrm{we}}^{2}}{2}}\]

\[\fitmath{\mathrm{ref}=\frac{{(214)}^{2}-{(246)}^{2}}{2}\approx -7360\,\mathrm{J/kg}}\]

\item De I constante se despeja la variación estática.

\[\fitmath{\mathrm{dh}=\frac{{\mathrm{us}}^{2}-{\mathrm{ue}}^{2}-{\mathrm{ws}}^{2}+{\mathrm{we}}^{2}}{2}}\]

\[\fitmath{\mathrm{dh}=\frac{{(331)}^{2}-{(204.48)}^{2}-{(214)}^{2}+{(246)}^{2}}{2}\approx 41234.46\,\mathrm{J/kg}}\]

\end{enumerate}


\noindent\textbf{Sustitución numérica}

\[
\fitmath{- \frac{\left(\left(204.48\,\mathrm{m}\,\mathrm{s}^{-1}\right)\right)^{2}}{2} + \frac{\left(\left(331\,\mathrm{m}\,\mathrm{s}^{-1}\right)\right)^{2}}{2} + \frac{\left(\left(246\,\mathrm{m}\,\mathrm{s}^{-1}\right)\right)^{2}}{2} - \frac{\left(\left(214\,\mathrm{m}\,\mathrm{s}^{-1}\right)\right)^{2}}{2}}
\]

\noindent\textbf{Resultado}

\[
\fitmath{\Delta h=4.123446480\times10^{4}\,\mathrm{J}\,\mathrm{kg}^{-1}}
\]

\noindent\textbf{Interpretación}
El término centrífugo no puede omitirse cuando U cambia. Reintroduzca dh en la identidad de rotalpía y compruebe residuo cero.

\noindent\textbf{Comprobación de los pasos}
\begin{itemize}
\item Paso 2: $33874.5$ — reproducido
\item Paso 4: $-7360$ — reproducido
\item Paso 6: $41234.5$ — reproducido
\end{itemize}

\noindent\textbf{Comprobación de la respuesta} la aritmética del último paso reproduce la respuesta enunciada a la precisión con que está escrita. Es una reproducción aritmética; no valida el modelo físico ni sus hipótesis.

\noindent\textbf{Limitaciones}
\begin{itemize}
\item Modelo educativo; la aplicabilidad a una máquina real requiere revisión experta.
\end{itemize}

\noindent\textbf{Errores frecuentes}
\begin{itemize}
\item Conservar h+W al cuadrado/2 a radio variable.
\end{itemize}
% studium:end-record eyJpZCI6ICJQUkItMDAxNiIsICJzZWN0aW9uIjogImMwMiJ9
\end{studiumworked}

\chapter{Rendimientos, pérdidas y semejanza}
\begin{studiumworked}{Problema resuelto C03-01-R}


\noindent\textbf{Enunciado}
\(C03-01-R \cdot\) Compresión total-total. Calcule temperatura total de salida y trabajo de una compresión con relación de presión y rendimiento dados. Datos: \(T_{0e}=294\, \mathrm{K}\); \(\pi_c=4.6\); \(\eta=0.84\); \(\gamma=1.4\); \(c_p=1004.5\, \mathrm{J}\,\mathrm{kg}^{-1}\,\mathrm{K}^{-1}\).

\noindent\textbf{Datos}
\begin{itemize}
\item $T_{0e} = 294\,\mathrm{K}$ — temperatura total entrada
\item $\pi_c = 4.6$ — relación presión total
\item $\eta = 0.84$ — rendimiento total-total
\item $\gamma = 1.4$ — razón de calores
\item $c_p = 1004.5\,\mathrm{J}\,\mathrm{kg}^{-1}\,\mathrm{K}^{-1}$ — calor específico
\end{itemize}

\noindent\textbf{Incógnita}
$w_c$, en $\mathrm{J\,kg^{-1}}$; también las magnitudes intermedias desarrolladas.

\noindent\textbf{Modelo físico}

\noindent\textit{Compresión total-total}

\[
\fitmath{w_c=\frac{T_{0e} c_p \left(\pi_c^{\frac{\gamma - 1}{\gamma}} - 1\right)}{\eta}}
\]

\noindent\textbf{Hipótesis}
\begin{itemize}
\item Gas caloríficamente perfecto y proceso adiabático.
\end{itemize}

\noindent\textbf{Desarrollo paso a paso}

\begin{enumerate}

\item Temperatura total de salida

\[\fitmath{\mathrm{Tout}=\mathrm{Tin}\cdot \left(1+\frac{{\mathrm{pr}}^{\frac{\gamma-1}{\gamma}}-1}{\eta}\right)}\]

\[\fitmath{\mathrm{Tout}=(294)\cdot \left(1+\frac{{(4.6)}^{\frac{(1.4)-1}{(1.4)}}-1}{(0.84)}\right)\approx 485.2868\,\mathrm{K}}\]

\item La referencia isentrópica utiliza la misma presión de salida.

\[\fitmath{\mathrm{dTs}=\mathrm{Tin}\cdot \left({\mathrm{pr}}^{\frac{\gamma-1}{\gamma}}-1\right)}\]

\[\fitmath{\mathrm{dTs}=(294)\cdot \left({(4.6)}^{\frac{(1.4)-1}{(1.4)}}-1\right)\approx 160.6809\,\mathrm{K}}\]

\item El trabajo real es mayor que el ideal.

\[\fitmath{\mathrm{wc}=\frac{\mathrm{cp}\cdot \mathrm{Tin}\cdot \left({\mathrm{pr}}^{\frac{\gamma-1}{\gamma}}-1\right)}{\eta}}\]

\[\fitmath{\mathrm{wc}=\frac{(1004.5)\cdot (294)\cdot \left({(4.6)}^{\frac{(1.4)-1}{(1.4)}}-1\right)}{(0.84)}\approx 192147.6\,\mathrm{J/kg}}\]

\end{enumerate}


\noindent\textbf{Sustitución numérica}

\[
\fitmath{\frac{\left(294\,\mathrm{K}\right) \left(1004.5\,\mathrm{J}\,\mathrm{kg}^{-1}\,\mathrm{K}^{-1}\right) \left(\left(4.6\right)^{\frac{\left(1.4\right) - 1}{\left(1.4\right)}} - 1\right)}{\left(0.84\right)}}
\]

\noindent\textbf{Resultado}

\[
\fitmath{w_c=1.921475762\times10^{5}\,\mathrm{J}\,\mathrm{kg}^{-1}}
\]

\noindent\textbf{Temperatura total de salida}
\[
\fitmath{T_{0s}=T_{0e}+w_c/c_p=485.2867856\,\mathrm{K}}
\]

\noindent\textbf{Comprobación energética.} $c_p(T_{0s}-T_{0e})=w_c$; no confundir el incremento isentrópico con la temperatura real de salida.

\noindent\textbf{Interpretación}
La eficiencia divide el trabajo ideal en compresión. El resultado debe superar el coste isentrópico.

\noindent\textbf{Comprobación de los pasos}
\begin{itemize}
\item Paso 2: $160.681$ — reproducido
\item Paso 4: $192148$ — reproducido
\end{itemize}

\noindent\textbf{Comprobación de la respuesta} la aritmética del último paso reproduce la respuesta enunciada a la precisión con que está escrita. Es una reproducción aritmética; no valida el modelo físico ni sus hipótesis.

\noindent\textbf{Limitaciones}
\begin{itemize}
\item Modelo educativo; la aplicabilidad a una máquina real requiere revisión experta.
\end{itemize}

\noindent\textbf{Errores frecuentes}
\begin{itemize}
\item Multiplicar por eta para calcular el trabajo real.
\end{itemize}
% studium:end-record eyJpZCI6ICJQUkItMDAxNyIsICJzZWN0aW9uIjogImMwMyJ9
\end{studiumworked}

\begin{studiumworked}{Problema resuelto C03-01-P}


\noindent\textbf{Enunciado}
\(C03-01-P \cdot\) Compresión total-total. Ejercicio de transferencia. Calcule temperatura total de salida y trabajo de una compresión con relación de presión y rendimiento dados. Antes de consultar la solución, prediga el efecto de cambiar el primer dato respecto del caso resuelto de esta familia. Justifique la tendencia y explique qué restricción podría convertirse en limitante. Datos: \(T_{0e}=313.11\, \mathrm{K}\); \(\pi_c=4.6\); \(\eta=0.84\); \(\gamma=1.4\); \(c_p=1004.5\, \mathrm{J}\,\mathrm{kg}^{-1}\,\mathrm{K}^{-1}\).

\noindent\textbf{Datos}
\begin{itemize}
\item $T_{0e} = 313.11\,\mathrm{K}$ — temperatura total entrada
\item $\pi_c = 4.6$ — relación presión total
\item $\eta = 0.84$ — rendimiento total-total
\item $\gamma = 1.4$ — razón de calores
\item $c_p = 1004.5\,\mathrm{J}\,\mathrm{kg}^{-1}\,\mathrm{K}^{-1}$ — calor específico
\end{itemize}

\noindent\textbf{Incógnita}
$w_c$, en $\mathrm{J\,kg^{-1}}$; también las magnitudes intermedias desarrolladas.

\noindent\textbf{Modelo físico}

\noindent\textit{Compresión total-total}

\[
\fitmath{w_c=\frac{T_{0e} c_p \left(\pi_c^{\frac{\gamma - 1}{\gamma}} - 1\right)}{\eta}}
\]

\noindent\textbf{Hipótesis}
\begin{itemize}
\item Gas caloríficamente perfecto y proceso adiabático.
\end{itemize}

\noindent\textbf{Desarrollo paso a paso}

\begin{enumerate}

\item Temperatura total de salida

\[\fitmath{\mathrm{Tout}=\mathrm{Tin}\cdot \left(1+\frac{{\mathrm{pr}}^{\frac{\gamma-1}{\gamma}}-1}{\mathrm{eta}}\right)}\]

\[\fitmath{\mathrm{Tout}=(313.11)\cdot \left(1+\frac{{(4.6)}^{\frac{(1.4)-1}{(1.4)}}-1}{(0.84)}\right)\approx 516.8304\,\mathrm{K}}\]

\item La referencia isentrópica utiliza la misma presión de salida.

\[\fitmath{\mathrm{dTs}=\mathrm{Tin}\cdot \left({\mathrm{pr}}^{\frac{\gamma-1}{\gamma}}-1\right)}\]

\[\fitmath{\mathrm{dTs}=(313.11)\cdot \left({(4.6)}^{\frac{(1.4)-1}{(1.4)}}-1\right)\approx 171.1252\,\mathrm{K}}\]

\item El trabajo real es mayor que el ideal.

\[\fitmath{\mathrm{wc}=\frac{\mathrm{cp}\cdot \mathrm{Tin}\cdot \left({\mathrm{pr}}^{\frac{\gamma-1}{\gamma}}-1\right)}{\mathrm{eta}}}\]

\[\fitmath{\mathrm{wc}=\frac{(1004.5)\cdot (313.11)\cdot \left({(4.6)}^{\frac{(1.4)-1}{(1.4)}}-1\right)}{(0.84)}\approx 204637.2\,\mathrm{J/kg}}\]

\end{enumerate}


\noindent\textbf{Sustitución numérica}

\[
\fitmath{\frac{\left(313.11\,\mathrm{K}\right) \left(1004.5\,\mathrm{J}\,\mathrm{kg}^{-1}\,\mathrm{K}^{-1}\right) \left(\left(4.6\right)^{\frac{\left(1.4\right) - 1}{\left(1.4\right)}} - 1\right)}{\left(0.84\right)}}
\]

\noindent\textbf{Resultado}

\[
\fitmath{w_c=2.046371686\times10^{5}\,\mathrm{J}\,\mathrm{kg}^{-1}}
\]

\noindent\textbf{Temperatura total de salida}
\[
\fitmath{T_{0s}=T_{0e}+w_c/c_p=516.8304267\,\mathrm{K}}
\]

\noindent\textbf{Comprobación energética.} $c_p(T_{0s}-T_{0e})=w_c$; no confundir el incremento isentrópico con la temperatura real de salida.

\noindent\textbf{Interpretación}
La eficiencia divide el trabajo ideal en compresión. El resultado debe superar el coste isentrópico.

\noindent\textbf{Comprobación de los pasos}
\begin{itemize}
\item Paso 2: $171.125$ — reproducido
\item Paso 4: $204637$ — reproducido
\end{itemize}

\noindent\textbf{Comprobación de la respuesta} la aritmética del último paso reproduce la respuesta enunciada a la precisión con que está escrita. Es una reproducción aritmética; no valida el modelo físico ni sus hipótesis.

\noindent\textbf{Limitaciones}
\begin{itemize}
\item Modelo educativo; la aplicabilidad a una máquina real requiere revisión experta.
\end{itemize}

\noindent\textbf{Errores frecuentes}
\begin{itemize}
\item Multiplicar por eta para calcular el trabajo real.
\end{itemize}
% studium:end-record eyJpZCI6ICJQUkItMDAxOCIsICJzZWN0aW9uIjogImMwMyJ9
\end{studiumworked}

\begin{studiumworked}{Problema resuelto C03-02-R}


\noindent\textbf{Enunciado}
\(C03-02-R \cdot\) Expansión total-total. Determine el trabajo positivo extraído por una turbina con razón de presión salida/entrada dada. Datos: \(T_{0e}=1435\, \mathrm{K}\); \(\pi_t=0.43\); \(\eta=0.89\); \(\gamma=1.33\); \(c_p=1156.697\, \mathrm{J}\,\mathrm{kg}^{-1}\,\mathrm{K}^{-1}\).

\noindent\textbf{Datos}
\begin{itemize}
\item $T_{0e} = 1435\,\mathrm{K}$ — temperatura total entrada
\item $\pi_t = 0.43$ — presión total salida/entrada
\item $\eta = 0.89$ — rendimiento total-total
\item $\gamma = 1.33$ — razón de calores
\item $c_p = 1156.696969697\,\mathrm{J}\,\mathrm{kg}^{-1}\,\mathrm{K}^{-1}$ — calor específico compatible con R=287
\end{itemize}

\noindent\textbf{Incógnita}
$w_t$, en $\mathrm{J\,kg^{-1}}$; también las magnitudes intermedias desarrolladas.

\noindent\textbf{Modelo físico}

\noindent\textit{Expansión total-total}

\[
\fitmath{w_t=T_{0e} c_p \eta \left(1 - \pi_t^{\frac{\gamma - 1}{\gamma}}\right)}
\]

\noindent\textbf{Hipótesis}
\begin{itemize}
\item Sin sangrados ni refrigeración, gas perfecto constante.
\end{itemize}

\noindent\textbf{Desarrollo paso a paso}

\begin{enumerate}

\item La caída ideal se refiere a igual presión final.

\[\fitmath{\mathrm{dTs}=\mathrm{Tin}\cdot \left(1-{\mathrm{pr}}^{\frac{\gamma-1}{\gamma}}\right)}\]

\[\fitmath{\mathrm{dTs}=(1435)\cdot \left(1-{(0.43)}^{\frac{(1.33)-1}{(1.33)}}\right)\approx 271.1206\,\mathrm{K}}\]

\item La caída real es una fracción de la ideal.

\[\fitmath{\mathrm{wt}=\eta\cdot \mathrm{cp}\cdot \mathrm{Tin}\cdot \left(1-{\mathrm{pr}}^{\frac{\gamma-1}{\gamma}}\right)}\]

\[\fitmath{\mathrm{wt}=(0.89)\cdot (1156.69697)\cdot (1435)\cdot \left(1-{(0.43)}^{\frac{(1.33)-1}{(1.33)}}\right)\approx 279107.9\,\mathrm{J/kg}}\]

\end{enumerate}


\noindent\textbf{Sustitución numérica}

\[
\fitmath{\left(1435\,\mathrm{K}\right) \left(1156.696969697\,\mathrm{J}\,\mathrm{kg}^{-1}\,\mathrm{K}^{-1}\right) \left(0.89\right) \left(1 - \left(0.43\right)^{\frac{\left(1.33\right) - 1}{\left(1.33\right)}}\right)}
\]

\noindent\textbf{Resultado}

\[
\fitmath{w_t=2.791078880\times10^{5}\,\mathrm{J}\,\mathrm{kg}^{-1}}
\]

\noindent\textbf{Interpretación}
El trabajo útil idealizado es positivo y menor que la caída isentrópica disponible.

\noindent\textbf{Comprobación de los pasos}
\begin{itemize}
\item Paso 2: $271.121$ — reproducido
\item Paso 4: $279108$ — reproducido
\end{itemize}

\noindent\textbf{Comprobación de la respuesta} la aritmética del último paso reproduce la respuesta enunciada a la precisión con que está escrita. Es una reproducción aritmética; no valida el modelo físico ni sus hipótesis.

\noindent\textbf{Limitaciones}
\begin{itemize}
\item Modelo educativo; la aplicabilidad a una máquina real requiere revisión experta.
\end{itemize}

\noindent\textbf{Errores frecuentes}
\begin{itemize}
\item Usar presión entrada/salida sin transformar la fórmula.
\end{itemize}
% studium:end-record eyJpZCI6ICJQUkItMDAxOSIsICJzZWN0aW9uIjogImMwMyJ9
\end{studiumworked}

\begin{studiumworked}{Problema resuelto C03-02-P}


\noindent\textbf{Enunciado}
\(C03-02-P \cdot\) Expansión total-total. Ejercicio de transferencia. Determine el trabajo positivo extraído por una turbina con razón de presión salida/entrada dada. Antes de consultar la solución, prediga el efecto de cambiar el primer dato respecto del caso resuelto de esta familia. Justifique la tendencia y explique qué restricción podría convertirse en limitante. Datos: \(T_{0e}=1528.275\, \mathrm{K}\); \(\pi_t=0.43\); \(\eta=0.89\); \(\gamma=1.33\); \(c_p=1156.697\, \mathrm{J}\,\mathrm{kg}^{-1}\,\mathrm{K}^{-1}\).

\noindent\textbf{Datos}
\begin{itemize}
\item $T_{0e} = 1528.275\,\mathrm{K}$ — temperatura total entrada
\item $\pi_t = 0.43$ — presión total salida/entrada
\item $\eta = 0.89$ — rendimiento total-total
\item $\gamma = 1.33$ — razón de calores
\item $c_p = 1156.696969697\,\mathrm{J}\,\mathrm{kg}^{-1}\,\mathrm{K}^{-1}$ — calor específico compatible con R=287
\end{itemize}

\noindent\textbf{Incógnita}
$w_t$, en $\mathrm{J\,kg^{-1}}$; también las magnitudes intermedias desarrolladas.

\noindent\textbf{Modelo físico}

\noindent\textit{Expansión total-total}

\[
\fitmath{w_t=T_{0e} c_p \eta \left(1 - \pi_t^{\frac{\gamma - 1}{\gamma}}\right)}
\]

\noindent\textbf{Hipótesis}
\begin{itemize}
\item Sin sangrados ni refrigeración, gas perfecto constante.
\end{itemize}

\noindent\textbf{Desarrollo paso a paso}

\begin{enumerate}

\item La caída ideal se refiere a igual presión final.

\[\fitmath{\mathrm{dTs}=\mathrm{Tin}\cdot \left(1-{\mathrm{pr}}^{\frac{\gamma-1}{\gamma}}\right)}\]

\[\fitmath{\mathrm{dTs}=(1528.275)\cdot \left(1-{(0.43)}^{\frac{(1.33)-1}{(1.33)}}\right)\approx 288.7434\,\mathrm{K}}\]

\item La caída real es una fracción de la ideal.

\[\fitmath{\mathrm{wt}=\mathrm{eta}\cdot \mathrm{cp}\cdot \mathrm{Tin}\cdot \left(1-{\mathrm{pr}}^{\frac{\gamma-1}{\gamma}}\right)}\]

\[\fitmath{\mathrm{wt}=(0.89)\cdot (1156.69697)\cdot (1528.275)\cdot \left(1-{(0.43)}^{\frac{(1.33)-1}{(1.33)}}\right)\approx 297249.9\,\mathrm{J/kg}}\]

\end{enumerate}


\noindent\textbf{Sustitución numérica}

\[
\fitmath{\left(1528.275\,\mathrm{K}\right) \left(1156.696969697\,\mathrm{J}\,\mathrm{kg}^{-1}\,\mathrm{K}^{-1}\right) \left(0.89\right) \left(1 - \left(0.43\right)^{\frac{\left(1.33\right) - 1}{\left(1.33\right)}}\right)}
\]

\noindent\textbf{Resultado}

\[
\fitmath{w_t=2.972499007\times10^{5}\,\mathrm{J}\,\mathrm{kg}^{-1}}
\]

\noindent\textbf{Interpretación}
El trabajo útil idealizado es positivo y menor que la caída isentrópica disponible.

\noindent\textbf{Comprobación de los pasos}
\begin{itemize}
\item Paso 2: $288.743$ — reproducido
\item Paso 4: $297250$ — reproducido
\end{itemize}

\noindent\textbf{Comprobación de la respuesta} la aritmética del último paso reproduce la respuesta enunciada a la precisión con que está escrita. Es una reproducción aritmética; no valida el modelo físico ni sus hipótesis.

\noindent\textbf{Limitaciones}
\begin{itemize}
\item Modelo educativo; la aplicabilidad a una máquina real requiere revisión experta.
\end{itemize}

\noindent\textbf{Errores frecuentes}
\begin{itemize}
\item Usar presión entrada/salida sin transformar la fórmula.
\end{itemize}
% studium:end-record eyJpZCI6ICJQUkItMDAyMCIsICJzZWN0aW9uIjogImMwMyJ9
\end{studiumworked}

\begin{studiumworked}{Problema resuelto C03-03-R}


\noindent\textbf{Enunciado}
\(C03-03-R \cdot\) Conversión politrópica. Un compresor tiene eficiencia politrópica uniforme. Calcule la razón de temperaturas y su eficiencia isentrópica global. Datos: \(\pi_c=12.4\); \(\eta_p=0.905\); \(\gamma=1.4\).

\noindent\textbf{Datos}
\begin{itemize}
\item $\pi_c = 12.4$ — relación presión total
\item $\eta_p = 0.905$ — eficiencia politrópica
\item $\gamma = 1.4$ — razón de calores
\end{itemize}

\noindent\textbf{Incógnita}
$\eta$ (adimensional); también las magnitudes intermedias desarrolladas.

\noindent\textbf{Modelo físico}

\noindent\textit{Conversión politrópica}

\[
\fitmath{\eta=\frac{\pi_c^{\frac{\gamma - 1}{\gamma}} - 1}{\pi_c^{\frac{\gamma - 1}{\eta_p \gamma}} - 1}}
\]

\noindent\textbf{Hipótesis}
\begin{itemize}
\item Eficiencia politrópica y propiedades constantes.
\end{itemize}

\noindent\textbf{Desarrollo paso a paso}

\begin{enumerate}

\item Se integra la relación diferencial politrópica.

\[\fitmath{\mathrm{Tr}={\mathrm{pr}}^{\frac{\gamma-1}{\gamma\cdot \mathrm{ep}}}}\]

\[\fitmath{\mathrm{Tr}={(12.4)}^{\frac{(1.4)-1}{(1.4)\cdot (0.905)}}\approx 2.214115\,}\]

\item Se compara el trabajo isentrópico global con el real.

\[\fitmath{\mathrm{eta}=\frac{{\mathrm{pr}}^{\frac{\gamma-1}{\gamma}}-1}{{\mathrm{pr}}^{\frac{\gamma-1}{\gamma\cdot \mathrm{ep}}}-1}}\]

\[\fitmath{\mathrm{eta}=\frac{{(12.4)}^{\frac{(1.4)-1}{(1.4)}}-1}{{(12.4)}^{\frac{(1.4)-1}{(1.4)\cdot (0.905)}}-1}\approx 0.8673654\,}\]

\end{enumerate}


\noindent\textbf{Sustitución numérica}

\[
\fitmath{\frac{\left(12.4\right)^{\frac{\left(1.4\right) - 1}{\left(1.4\right)}} - 1}{\left(12.4\right)^{\frac{\left(1.4\right) - 1}{\left(0.905\right) \left(1.4\right)}} - 1}}
\]

\noindent\textbf{Resultado}

\[
\fitmath{\eta=8.673653999\times10^{-1}}
\]

\noindent\textbf{Interpretación}
La eficiencia global es menor que la politrópica en este modelo de compresión.

\noindent\textbf{Comprobación de los pasos}
\begin{itemize}
\item Paso 2: $2.21412$ — reproducido
\item Paso 4: $0.867365$ — reproducido
\end{itemize}

\noindent\textbf{Comprobación de la respuesta} la aritmética del último paso reproduce la respuesta enunciada a la precisión con que está escrita. Es una reproducción aritmética; no valida el modelo físico ni sus hipótesis.

\noindent\textbf{Limitaciones}
\begin{itemize}
\item Modelo educativo; la aplicabilidad a una máquina real requiere revisión experta.
\end{itemize}

\noindent\textbf{Errores frecuentes}
\begin{itemize}
\item Elevar la eficiencia de escalón al número de escalones.
\end{itemize}
% studium:end-record eyJpZCI6ICJQUkItMDAyMSIsICJzZWN0aW9uIjogImMwMyJ9
\end{studiumworked}

\begin{studiumworked}{Problema resuelto C03-03-P}


\noindent\textbf{Enunciado}
\(C03-03-P \cdot\) Conversión politrópica. Ejercicio de transferencia. Un compresor tiene eficiencia politrópica uniforme. Calcule la razón de temperaturas y su eficiencia isentrópica global. Antes de consultar la solución, prediga el efecto de cambiar el primer dato respecto del caso resuelto de esta familia. Justifique la tendencia y explique qué restricción podría convertirse en limitante. Datos: \(\pi_c=13.206\); \(\eta_p=0.905\); \(\gamma=1.4\).

\noindent\textbf{Datos}
\begin{itemize}
\item $\pi_c = 13.206$ — relación presión total
\item $\eta_p = 0.905$ — eficiencia politrópica
\item $\gamma = 1.4$ — razón de calores
\end{itemize}

\noindent\textbf{Incógnita}
$\eta$ (adimensional); también las magnitudes intermedias desarrolladas.

\noindent\textbf{Modelo físico}

\noindent\textit{Conversión politrópica}

\[
\fitmath{\eta=\frac{\pi_c^{\frac{\gamma - 1}{\gamma}} - 1}{\pi_c^{\frac{\gamma - 1}{\eta_p \gamma}} - 1}}
\]

\noindent\textbf{Hipótesis}
\begin{itemize}
\item Eficiencia politrópica y propiedades constantes.
\end{itemize}

\noindent\textbf{Desarrollo paso a paso}

\begin{enumerate}

\item Se integra la relación diferencial politrópica.

\[\fitmath{\mathrm{Tr}={\mathrm{pr}}^{\frac{\gamma-1}{\gamma\cdot \mathrm{ep}}}}\]

\[\fitmath{\mathrm{Tr}={(13.206)}^{\frac{(1.4)-1}{(1.4)\cdot (0.905)}}\approx 2.258576\,}\]

\item Se compara el trabajo isentrópico global con el real.

\[\fitmath{\eta=\frac{{\mathrm{pr}}^{\frac{\gamma-1}{\gamma}}-1}{{\mathrm{pr}}^{\frac{\gamma-1}{\gamma\cdot \mathrm{ep}}}-1}}\]

\[\fitmath{\eta=\frac{{(13.206)}^{\frac{(1.4)-1}{(1.4)}}-1}{{(13.206)}^{\frac{(1.4)-1}{(1.4)\cdot (0.905)}}-1}\approx 0.8663416\,}\]

\end{enumerate}


\noindent\textbf{Sustitución numérica}

\[
\fitmath{\frac{\left(13.206\right)^{\frac{\left(1.4\right) - 1}{\left(1.4\right)}} - 1}{\left(13.206\right)^{\frac{\left(1.4\right) - 1}{\left(0.905\right) \left(1.4\right)}} - 1}}
\]

\noindent\textbf{Resultado}

\[
\fitmath{\eta=8.663416108\times10^{-1}}
\]

\noindent\textbf{Interpretación}
La eficiencia global es menor que la politrópica en este modelo de compresión.

\noindent\textbf{Comprobación de los pasos}
\begin{itemize}
\item Paso 2: $2.25858$ — reproducido
\item Paso 4: $0.866342$ — reproducido
\end{itemize}

\noindent\textbf{Comprobación de la respuesta} la aritmética del último paso reproduce la respuesta enunciada a la precisión con que está escrita. Es una reproducción aritmética; no valida el modelo físico ni sus hipótesis.

\noindent\textbf{Limitaciones}
\begin{itemize}
\item Modelo educativo; la aplicabilidad a una máquina real requiere revisión experta.
\end{itemize}

\noindent\textbf{Errores frecuentes}
\begin{itemize}
\item Elevar la eficiencia de escalón al número de escalones.
\end{itemize}
% studium:end-record eyJpZCI6ICJQUkItMDAyMiIsICJzZWN0aW9uIjogImMwMyJ9
\end{studiumworked}

\begin{studiumworked}{Problema resuelto C03-04-R}


\noindent\textbf{Enunciado}
\(C03-04-R \cdot\) Corrección de caudal. Convierta el caudal medido a las referencias de presión y temperatura especificadas. Datos: \(\dot m=31.7\, \mathrm{kg}\,\mathrm{s}^{-1}\); \(T_{0e}=271\, \mathrm{K}\); \(p_{0e}=82500\, \mathrm{Pa}\); \(T_{\mathrm{ref}}=288.15\, \mathrm{K}\); \(p_{\mathrm{ref}}=101325\, \mathrm{Pa}\).

\noindent\textbf{Datos}
\begin{itemize}
\item $\dot m = 31.7\,\mathrm{kg}\,\mathrm{s}^{-1}$ — caudal real
\item $T_{0e} = 271\,\mathrm{K}$ — temperatura total entrada
\item $p_{0e} = 82500\,\mathrm{Pa}$ — presión total entrada
\item $T_{\mathrm{ref}} = 288.15\,\mathrm{K}$ — temperatura referencia
\item $p_{\mathrm{ref}} = 101325\,\mathrm{Pa}$ — presión referencia
\end{itemize}

\noindent\textbf{Incógnita}
$\dot{m}_c$, en $\mathrm{kg\,s^{-1}}$; también las magnitudes intermedias desarrolladas.

\noindent\textbf{Modelo físico}

\noindent\textit{Corrección de caudal}

\[
\fitmath{\dot{m}_c=\frac{\dot m p_{\mathrm{ref}} \sqrt{\frac{T_{0e}}{T_{\mathrm{ref}}}}}{p_{0e}}}
\]

\noindent\textbf{Hipótesis}
\begin{itemize}
\item Mismo gas y referencias declaradas.
\end{itemize}

\noindent\textbf{Desarrollo paso a paso}

\begin{enumerate}

\item La corrección térmica es la raíz de la razón de temperaturas.

\[\fitmath{\theta=\sqrt{\frac{\mathrm{Tin}}{\mathrm{Tr}}}}\]

\[\fitmath{\theta=\sqrt{\frac{(271)}{(288.15)}}\approx 0.9697847\,}\]

\item La presión corrige la capacidad de paso del gas.

\[\fitmath{\mathrm{mc}=\frac{\mathrm{md}\cdot \sqrt{\frac{\mathrm{Tin}}{\mathrm{Tr}}}}{\frac{\mathrm{pin}}{\mathrm{pr}}}}\]

\[\fitmath{\mathrm{mc}=\frac{(31.7)\cdot \sqrt{\frac{(271)}{(288.15)}}}{\frac{(82500)}{(101325)}}\approx 37.75698\,\mathrm{kg/s}}\]

\end{enumerate}


\noindent\textbf{Sustitución numérica}

\[
\fitmath{\frac{\left(31.7\,\mathrm{kg}\,\mathrm{s}^{-1}\right) \left(101325\,\mathrm{Pa}\right) \sqrt{\frac{\left(271\,\mathrm{K}\right)}{\left(288.15\,\mathrm{K}\right)}}}{\left(82500\,\mathrm{Pa}\right)}}
\]

\noindent\textbf{Resultado}

\[
\fitmath{\dot{m}_c=3.775698084\times10^{1}\,\mathrm{kg}\,\mathrm{s}^{-1}}
\]

\noindent\textbf{Interpretación}
El caudal corregido es una variable de comparación. No es una nueva medición de masa real.

\noindent\textbf{Comprobación de los pasos}
\begin{itemize}
\item Paso 2: $0.969785$ — reproducido
\item Paso 4: $37.757$ — reproducido
\end{itemize}

\noindent\textbf{Comprobación de la respuesta} la aritmética del último paso reproduce la respuesta enunciada a la precisión con que está escrita. Es una reproducción aritmética; no valida el modelo físico ni sus hipótesis.

\noindent\textbf{Limitaciones}
\begin{itemize}
\item Modelo educativo; la aplicabilidad a una máquina real requiere revisión experta.
\end{itemize}

\noindent\textbf{Errores frecuentes}
\begin{itemize}
\item Invertir la corrección de presión.
\end{itemize}
% studium:end-record eyJpZCI6ICJQUkItMDAyMyIsICJzZWN0aW9uIjogImMwMyJ9
\end{studiumworked}

\begin{studiumworked}{Problema resuelto C03-04-P}


\noindent\textbf{Enunciado}
\(C03-04-P \cdot\) Corrección de caudal. Ejercicio de transferencia. Convierta el caudal medido a las referencias de presión y temperatura especificadas. Antes de consultar la solución, prediga el efecto de cambiar el primer dato respecto del caso resuelto de esta familia. Justifique la tendencia y explique qué restricción podría convertirse en limitante. Datos: \(\dot m=33.7605\, \mathrm{kg}\,\mathrm{s}^{-1}\); \(T_{0e}=271\, \mathrm{K}\); \(p_{0e}=82500\, \mathrm{Pa}\); \(T_{\mathrm{ref}}=288.15\, \mathrm{K}\); \(p_{\mathrm{ref}}=101325\, \mathrm{Pa}\).

\noindent\textbf{Datos}
\begin{itemize}
\item $\dot m = 33.7605\,\mathrm{kg}\,\mathrm{s}^{-1}$ — caudal real
\item $T_{0e} = 271\,\mathrm{K}$ — temperatura total entrada
\item $p_{0e} = 82500\,\mathrm{Pa}$ — presión total entrada
\item $T_{\mathrm{ref}} = 288.15\,\mathrm{K}$ — temperatura referencia
\item $p_{\mathrm{ref}} = 101325\,\mathrm{Pa}$ — presión referencia
\end{itemize}

\noindent\textbf{Incógnita}
$\dot{m}_c$, en $\mathrm{kg\,s^{-1}}$; también las magnitudes intermedias desarrolladas.

\noindent\textbf{Modelo físico}

\noindent\textit{Corrección de caudal}

\[
\fitmath{\dot{m}_c=\frac{\dot m p_{\mathrm{ref}} \sqrt{\frac{T_{0e}}{T_{\mathrm{ref}}}}}{p_{0e}}}
\]

\noindent\textbf{Hipótesis}
\begin{itemize}
\item Mismo gas y referencias declaradas.
\end{itemize}

\noindent\textbf{Desarrollo paso a paso}

\begin{enumerate}

\item La corrección térmica es la raíz de la razón de temperaturas.

\[\fitmath{\theta=\sqrt{\frac{\mathrm{Tin}}{\mathrm{Tr}}}}\]

\[\fitmath{\theta=\sqrt{\frac{(271)}{(288.15)}}\approx 0.9697847\,}\]

\item La presión corrige la capacidad de paso del gas.

\[\fitmath{\mathrm{mc}=\frac{\mathrm{md}\cdot \sqrt{\frac{\mathrm{Tin}}{\mathrm{Tr}}}}{\frac{\mathrm{pin}}{\mathrm{pr}}}}\]

\[\fitmath{\mathrm{mc}=\frac{(33.7605)\cdot \sqrt{\frac{(271)}{(288.15)}}}{\frac{(82500)}{(101325)}}\approx 40.21118\,\mathrm{kg/s}}\]

\end{enumerate}


\noindent\textbf{Sustitución numérica}

\[
\fitmath{\frac{\left(33.7605\,\mathrm{kg}\,\mathrm{s}^{-1}\right) \left(101325\,\mathrm{Pa}\right) \sqrt{\frac{\left(271\,\mathrm{K}\right)}{\left(288.15\,\mathrm{K}\right)}}}{\left(82500\,\mathrm{Pa}\right)}}
\]

\noindent\textbf{Resultado}

\[
\fitmath{\dot{m}_c=4.021118460\times10^{1}\,\mathrm{kg}\,\mathrm{s}^{-1}}
\]

\noindent\textbf{Interpretación}
El caudal corregido es una variable de comparación. No es una nueva medición de masa real.

\noindent\textbf{Comprobación de los pasos}
\begin{itemize}
\item Paso 2: $0.969785$ — reproducido
\item Paso 4: $40.2112$ — reproducido
\end{itemize}

\noindent\textbf{Comprobación de la respuesta} la aritmética del último paso reproduce la respuesta enunciada a la precisión con que está escrita. Es una reproducción aritmética; no valida el modelo físico ni sus hipótesis.

\noindent\textbf{Limitaciones}
\begin{itemize}
\item Modelo educativo; la aplicabilidad a una máquina real requiere revisión experta.
\end{itemize}

\noindent\textbf{Errores frecuentes}
\begin{itemize}
\item Invertir la corrección de presión.
\end{itemize}
% studium:end-record eyJpZCI6ICJQUkItMDAyNCIsICJzZWN0aW9uIjogImMwMyJ9
\end{studiumworked}

\chapter{Turbinas axiales de acción}
\begin{studiumworked}{Problema resuelto C04-01-R}


\noindent\textbf{Enunciado}
\(C04-01-R \cdot\) Acción ideal. Una turbina de acción ideal recibe las componentes y velocidad periférica dadas. Halle remolino de salida, trabajo y descenso de temperatura total. Datos: \(C_\theta=462\, \mathrm{m}\,\mathrm{s}^{-1}\); \(U=211\, \mathrm{m}\,\mathrm{s}^{-1}\); \(c_p=1156.697\, \mathrm{J}\,\mathrm{kg}^{-1}\,\mathrm{K}^{-1}\).

\noindent\textbf{Datos}
\begin{itemize}
\item $C_\theta = 462\,\mathrm{m}\,\mathrm{s}^{-1}$ — remolino entrada rotor
\item $U = 211\,\mathrm{m}\,\mathrm{s}^{-1}$ — velocidad periférica
\item $c_p = 1156.696969697\,\mathrm{J}\,\mathrm{kg}^{-1}\,\mathrm{K}^{-1}$ — calor específico
\end{itemize}

\noindent\textbf{Incógnita}
$\Delta T_0$, en $\mathrm{K}$; también las magnitudes intermedias desarrolladas.

\noindent\textbf{Modelo físico}

\noindent\textit{Acción ideal}

\[
\fitmath{\Delta T_0=\frac{2 U \left(C_\theta - U\right)}{c_p}}
\]

\noindent\textbf{Hipótesis}
\begin{itemize}
\item Acción ideal, radio y componente axial constantes.
\end{itemize}

\noindent\textbf{Desarrollo paso a paso}

\begin{enumerate}

\item La simetría relativa da la componente tangencial de salida.

\[\fitmath{\mathrm{cts}=2\cdot \mathrm{uu}-\mathrm{ct}}\]

\[\fitmath{\mathrm{cts}=2\cdot (211)-(462)\approx -40\,\mathrm{m/s}}\]

\item Euler proporciona trabajo positivo de turbina.

\[\fitmath{\mathrm{wt}=2\cdot \mathrm{uu}\cdot \left(\mathrm{ct}-\mathrm{uu}\right)}\]

\[\fitmath{\mathrm{wt}=2\cdot (211)\cdot \left((462)-(211)\right)\approx 105922\,\mathrm{J/kg}}\]

\item Se divide por el calor específico.

\[\fitmath{\mathrm{dTz}=\frac{2\cdot \mathrm{uu}\cdot \left(\mathrm{ct}-\mathrm{uu}\right)}{\mathrm{cp}}}\]

\[\fitmath{\mathrm{dTz}=\frac{2\cdot (211)\cdot \left((462)-(211)\right)}{(1156.69697)}\approx 91.57282\,\mathrm{K}}\]

\end{enumerate}


\noindent\textbf{Sustitución numérica}

\[
\fitmath{\frac{2 \left(211\,\mathrm{m}\,\mathrm{s}^{-1}\right) \left(\left(462\,\mathrm{m}\,\mathrm{s}^{-1}\right) - \left(211\,\mathrm{m}\,\mathrm{s}^{-1}\right)\right)}{\left(1156.696969697\,\mathrm{J}\,\mathrm{kg}^{-1}\,\mathrm{K}^{-1}\right)}}
\]

\noindent\textbf{Resultado}

\[
\fitmath{\Delta T_0=9.157281706\times10^{1}\,\mathrm{K}}
\]

\noindent\textbf{Interpretación}
El remolino de salida es negativo con estos datos. Mantener su signo es necesario para cerrar Euler.

\noindent\textbf{Comprobación de los pasos}
\begin{itemize}
\item Paso 2: $-40$ — reproducido
\item Paso 4: $105922$ — reproducido
\item Paso 6: $91.5728$ — reproducido
\end{itemize}

\noindent\textbf{Comprobación de la respuesta} la aritmética del último paso reproduce la respuesta enunciada a la precisión con que está escrita. Es una reproducción aritmética; no valida el modelo físico ni sus hipótesis.

\noindent\textbf{Limitaciones}
\begin{itemize}
\item Modelo educativo; la aplicabilidad a una máquina real requiere revisión experta.
\end{itemize}

\noindent\textbf{Errores frecuentes}
\begin{itemize}
\item Tomar valor absoluto del remolino antes de aplicar Euler.
\end{itemize}
% studium:end-record eyJpZCI6ICJQUkItMDAyNSIsICJzZWN0aW9uIjogImMwNCJ9
\end{studiumworked}

\begin{studiumworked}{Problema resuelto C04-01-P}


\noindent\textbf{Enunciado}
\(C04-01-P \cdot\) Acción ideal. Ejercicio de transferencia. Una turbina de acción ideal recibe las componentes y velocidad periférica dadas. Halle remolino de salida, trabajo y descenso de temperatura total. Antes de consultar la solución, prediga el efecto de cambiar el primer dato respecto del caso resuelto de esta familia. Justifique la tendencia y explique qué restricción podría convertirse en limitante. Datos: \(C_\theta=492.03\, \mathrm{m}\,\mathrm{s}^{-1}\); \(U=211\, \mathrm{m}\,\mathrm{s}^{-1}\); \(c_p=1156.697\, \mathrm{J}\,\mathrm{kg}^{-1}\,\mathrm{K}^{-1}\).

\noindent\textbf{Datos}
\begin{itemize}
\item $C_\theta = 492.03\,\mathrm{m}\,\mathrm{s}^{-1}$ — remolino entrada rotor
\item $U = 211\,\mathrm{m}\,\mathrm{s}^{-1}$ — velocidad periférica
\item $c_p = 1156.696969697\,\mathrm{J}\,\mathrm{kg}^{-1}\,\mathrm{K}^{-1}$ — calor específico
\end{itemize}

\noindent\textbf{Incógnita}
$\Delta T_0$, en $\mathrm{K}$; también las magnitudes intermedias desarrolladas.

\noindent\textbf{Modelo físico}

\noindent\textit{Acción ideal}

\[
\fitmath{\Delta T_0=\frac{2 U \left(C_\theta - U\right)}{c_p}}
\]

\noindent\textbf{Hipótesis}
\begin{itemize}
\item Acción ideal, radio y componente axial constantes.
\end{itemize}

\noindent\textbf{Desarrollo paso a paso}

\begin{enumerate}

\item La simetría relativa da la componente tangencial de salida.

\[\fitmath{\mathrm{cts}=2\cdot \mathrm{uu}-\mathrm{ct}}\]

\[\fitmath{\mathrm{cts}=2\cdot (211)-(492.03)\approx -70.03\,\mathrm{m/s}}\]

\item Euler proporciona trabajo positivo de turbina.

\[\fitmath{\mathrm{wt}=2\cdot \mathrm{uu}\cdot \left(\mathrm{ct}-\mathrm{uu}\right)}\]

\[\fitmath{\mathrm{wt}=2\cdot (211)\cdot \left((492.03)-(211)\right)\approx 118594.7\,\mathrm{J/kg}}\]

\item Se divide por el calor específico.

\[\fitmath{\mathrm{dTz}=\frac{2\cdot \mathrm{uu}\cdot \left(\mathrm{ct}-\mathrm{uu}\right)}{\mathrm{cp}}}\]

\[\fitmath{\mathrm{dTz}=\frac{2\cdot (211)\cdot \left((492.03)-(211)\right)}{(1156.69697)}\approx 102.5287\,\mathrm{K}}\]

\end{enumerate}


\noindent\textbf{Sustitución numérica}

\[
\fitmath{\frac{2 \left(211\,\mathrm{m}\,\mathrm{s}^{-1}\right) \left(\left(492.03\,\mathrm{m}\,\mathrm{s}^{-1}\right) - \left(211\,\mathrm{m}\,\mathrm{s}^{-1}\right)\right)}{\left(1156.696969697\,\mathrm{J}\,\mathrm{kg}^{-1}\,\mathrm{K}^{-1}\right)}}
\]

\noindent\textbf{Resultado}

\[
\fitmath{\Delta T_0=1.025287202\times10^{2}\,\mathrm{K}}
\]

\noindent\textbf{Interpretación}
El remolino de salida es negativo con estos datos. Mantener su signo es necesario para cerrar Euler.

\noindent\textbf{Comprobación de los pasos}
\begin{itemize}
\item Paso 2: $-70.03$ — reproducido
\item Paso 4: $118595$ — reproducido
\item Paso 6: $102.529$ — reproducido
\end{itemize}

\noindent\textbf{Comprobación de la respuesta} la aritmética del último paso reproduce la respuesta enunciada a la precisión con que está escrita. Es una reproducción aritmética; no valida el modelo físico ni sus hipótesis.

\noindent\textbf{Limitaciones}
\begin{itemize}
\item Modelo educativo; la aplicabilidad a una máquina real requiere revisión experta.
\end{itemize}

\noindent\textbf{Errores frecuentes}
\begin{itemize}
\item Tomar valor absoluto del remolino antes de aplicar Euler.
\end{itemize}
% studium:end-record eyJpZCI6ICJQUkItMDAyNiIsICJzZWN0aW9uIjogImMwNCJ9
\end{studiumworked}

\begin{studiumworked}{Problema resuelto C04-02-R}


\noindent\textbf{Enunciado}
\(C04-02-R \cdot\) Óptimo de utilización. Para una velocidad de entrada al rotor y ángulo fijados, obtenga la velocidad periférica de máximo trabajo ideal y ese trabajo. Datos: \(C_e=590\, \mathrm{m}\,\mathrm{s}^{-1}\); \(\alpha_e=1.117011\, \mathrm{rad}\).

\noindent\textbf{Datos}
\begin{itemize}
\item $C_e = 590\,\mathrm{m}\,\mathrm{s}^{-1}$ — módulo entrada rotor
\item $\alpha_e = 1.1170107212764\,\mathrm{rad}$ — ángulo absoluto axial
\end{itemize}

\noindent\textbf{Incógnita}
$w_{\max}$, en $\mathrm{J\,kg^{-1}}$; también las magnitudes intermedias desarrolladas.

\noindent\textbf{Modelo físico}

\noindent\textit{Óptimo de utilización}

\[
\fitmath{w_{\max}=\frac{C_e^{2} \sin^{2}{\left(\alpha_e \right)}}{2}}
\]

\noindent\textbf{Hipótesis}
\begin{itemize}
\item Velocidad y orientación incidentes fijadas; rotor de acción ideal.
\end{itemize}

\noindent\textbf{Desarrollo paso a paso}

\begin{enumerate}

\item Se proyecta la velocidad tangencial.

\[\fitmath{\mathrm{ct}=\mathrm{Vel}\cdot \sin\left(\mathrm{ang}\right)}\]

\[\fitmath{\mathrm{ct}=(590)\cdot \sin\left((1.117010721)\right)\approx 530.2885\,\mathrm{m/s}}\]

\item El óptimo de la parábola está en U=Ctheta/2.

\[\fitmath{\mathrm{uopt}=\frac{\mathrm{Vel}\cdot \sin\left(\mathrm{ang}\right)}{2}}\]

\[\fitmath{\mathrm{uopt}=\frac{(590)\cdot \sin\left((1.117010721)\right)}{2}\approx 265.1442\,\mathrm{m/s}}\]

\item Se evalúa el máximo ideal.

\[\fitmath{\mathrm{wmax}=\frac{{\mathrm{Vel}}^{2}\cdot {\sin\left(\mathrm{ang}\right)}^{2}}{2}}\]

\[\fitmath{\mathrm{wmax}=\frac{{(590)}^{2}\cdot {\sin\left((1.117010721)\right)}^{2}}{2}\approx 140602.9\,\mathrm{J/kg}}\]

\end{enumerate}


\noindent\textbf{Sustitución numérica}

\[
\fitmath{\frac{\left(\left(590\,\mathrm{m}\,\mathrm{s}^{-1}\right)\right)^{2} \sin^{2}{\left(\left(1.1170107212764\,\mathrm{rad}\right) \right)}}{2}}
\]

\noindent\textbf{Resultado}

\[
\fitmath{w_{\max}=1.406029399\times10^{5}\,\mathrm{J}\,\mathrm{kg}^{-1}}
\]

\noindent\textbf{Interpretación}
Queda velocidad axial a la salida incluso en el óptimo. La utilización cinética no es una eficiencia total-total.

\noindent\textbf{Comprobación de los pasos}
\begin{itemize}
\item Paso 2: $530.288$ — reproducido
\item Paso 4: $265.144$ — reproducido
\item Paso 6: $140603$ — reproducido
\end{itemize}

\noindent\textbf{Comprobación de la respuesta} la aritmética del último paso reproduce la respuesta enunciada a la precisión con que está escrita. Es una reproducción aritmética; no valida el modelo físico ni sus hipótesis.

\noindent\textbf{Limitaciones}
\begin{itemize}
\item Modelo educativo; la aplicabilidad a una máquina real requiere revisión experta.
\end{itemize}

\noindent\textbf{Errores frecuentes}
\begin{itemize}
\item Suponer que el óptimo elimina toda velocidad de salida.
\end{itemize}
% studium:end-record eyJpZCI6ICJQUkItMDAyNyIsICJzZWN0aW9uIjogImMwNCJ9
\end{studiumworked}

\begin{studiumworked}{Problema resuelto C04-02-P}


\noindent\textbf{Enunciado}
\(C04-02-P \cdot\) Óptimo de utilización. Ejercicio de transferencia. Para una velocidad de entrada al rotor y ángulo fijados, obtenga la velocidad periférica de máximo trabajo ideal y ese trabajo. Antes de consultar la solución, prediga el efecto de cambiar el primer dato respecto del caso resuelto de esta familia. Justifique la tendencia y explique qué restricción podría convertirse en limitante. Datos: \(C_e=628.35\, \mathrm{m}\,\mathrm{s}^{-1}\); \(\alpha_e=1.117011\, \mathrm{rad}\).

\noindent\textbf{Datos}
\begin{itemize}
\item $C_e = 628.35\,\mathrm{m}\,\mathrm{s}^{-1}$ — módulo entrada rotor
\item $\alpha_e = 1.1170107212764\,\mathrm{rad}$ — ángulo absoluto axial
\end{itemize}

\noindent\textbf{Incógnita}
$w_{\max}$, en $\mathrm{J\,kg^{-1}}$; también las magnitudes intermedias desarrolladas.

\noindent\textbf{Modelo físico}

\noindent\textit{Óptimo de utilización}

\[
\fitmath{w_{\max}=\frac{C_e^{2} \sin^{2}{\left(\alpha_e \right)}}{2}}
\]

\noindent\textbf{Hipótesis}
\begin{itemize}
\item Velocidad y orientación incidentes fijadas; rotor de acción ideal.
\end{itemize}

\noindent\textbf{Desarrollo paso a paso}

\begin{enumerate}

\item Se proyecta la velocidad tangencial.

\[\fitmath{\mathrm{ct}=\mathrm{Vel}\cdot \sin\left(\mathrm{ang}\right)}\]

\[\fitmath{\mathrm{ct}=(628.35)\cdot \sin\left((1.117010721)\right)\approx 564.7572\,\mathrm{m/s}}\]

\item El óptimo de la parábola está en U=Ctheta/2.

\[\fitmath{\mathrm{uopt}=\frac{\mathrm{Vel}\cdot \sin\left(\mathrm{ang}\right)}{2}}\]

\[\fitmath{\mathrm{uopt}=\frac{(628.35)\cdot \sin\left((1.117010721)\right)}{2}\approx 282.3786\,\mathrm{m/s}}\]

\item Se evalúa el máximo ideal.

\[\fitmath{\mathrm{wmax}=\frac{{\mathrm{Vel}}^{2}\cdot {\sin\left(\mathrm{ang}\right)}^{2}}{2}}\]

\[\fitmath{\mathrm{wmax}=\frac{{(628.35)}^{2}\cdot {\sin\left((1.117010721)\right)}^{2}}{2}\approx 159475.4\,\mathrm{J/kg}}\]

\end{enumerate}


\noindent\textbf{Sustitución numérica}

\[
\fitmath{\frac{\left(\left(628.35\,\mathrm{m}\,\mathrm{s}^{-1}\right)\right)^{2} \sin^{2}{\left(\left(1.1170107212764\,\mathrm{rad}\right) \right)}}{2}}
\]

\noindent\textbf{Resultado}

\[
\fitmath{w_{\max}=1.594753695\times10^{5}\,\mathrm{J}\,\mathrm{kg}^{-1}}
\]

\noindent\textbf{Interpretación}
Queda velocidad axial a la salida incluso en el óptimo. La utilización cinética no es una eficiencia total-total.

\noindent\textbf{Comprobación de los pasos}
\begin{itemize}
\item Paso 2: $564.757$ — reproducido
\item Paso 4: $282.379$ — reproducido
\item Paso 6: $159475$ — reproducido
\end{itemize}

\noindent\textbf{Comprobación de la respuesta} la aritmética del último paso reproduce la respuesta enunciada a la precisión con que está escrita. Es una reproducción aritmética; no valida el modelo físico ni sus hipótesis.

\noindent\textbf{Limitaciones}
\begin{itemize}
\item Modelo educativo; la aplicabilidad a una máquina real requiere revisión experta.
\end{itemize}

\noindent\textbf{Errores frecuentes}
\begin{itemize}
\item Suponer que el óptimo elimina toda velocidad de salida.
\end{itemize}
% studium:end-record eyJpZCI6ICJQUkItMDAyOCIsICJzZWN0aW9uIjogImMwNCJ9
\end{studiumworked}

\begin{studiumworked}{Problema resuelto C04-03-R}


\noindent\textbf{Enunciado}
\(C04-03-R \cdot\) Rotor con reducción de velocidad relativa. Un rotor recibe W tangencial de entrada positiva y descarga con dirección relativa simétrica y módulo reducido por k. Halle la potencia extraída. Datos: \(C_x=176\, \mathrm{m}\,\mathrm{s}^{-1}\); \(C_\theta=493\, \mathrm{m}\,\mathrm{s}^{-1}\); \(U=238\, \mathrm{m}\,\mathrm{s}^{-1}\); \(k_W=0.92\); \(\dot m=17.8\, \mathrm{kg}\,\mathrm{s}^{-1}\).

\noindent\textbf{Datos}
\begin{itemize}
\item $C_x = 176\,\mathrm{m}\,\mathrm{s}^{-1}$ — componente axial entrada
\item $C_\theta = 493\,\mathrm{m}\,\mathrm{s}^{-1}$ — remolino entrada
\item $U = 238\,\mathrm{m}\,\mathrm{s}^{-1}$ — velocidad periférica
\item $k_W = 0.92$ — cociente Wsal/Went
\item $\dot m = 17.8\,\mathrm{kg}\,\mathrm{s}^{-1}$ — caudal
\end{itemize}

\noindent\textbf{Incógnita}
$\dot{W}$, en $\mathrm{W}$; también las magnitudes intermedias desarrolladas.

\noindent\textbf{Modelo físico}

\noindent\textit{Rotor con reducción de velocidad relativa}

\[
\fitmath{\dot{W}=\dot m U \left(C_\theta + k_W \left(C_\theta - U\right) - U\right)}
\]

\noindent\textbf{Hipótesis}
\begin{itemize}
\item Radio constante; k y dirección relativa son datos del modelo.
\end{itemize}

\noindent\textbf{Desarrollo paso a paso}

\begin{enumerate}

\item La salida tangencial absoluta combina U y la componente relativa reducida.

\[\fitmath{\mathrm{cts}=\mathrm{uu}-\mathrm{kk}\cdot \left(\mathrm{ct}-\mathrm{uu}\right)}\]

\[\fitmath{\mathrm{cts}=(238)-(0.92)\cdot \left((493)-(238)\right)\approx 3.4\,\mathrm{m/s}}\]

\item Se calcula el trabajo de Euler.

\[\fitmath{\mathrm{wt}=\mathrm{uu}\cdot \left(\mathrm{ct}-\left(\mathrm{uu}-\mathrm{kk}\cdot \left(\mathrm{ct}-\mathrm{uu}\right)\right)\right)}\]

\[\fitmath{\mathrm{wt}=(238)\cdot \left((493)-\left((238)-(0.92)\cdot \left((493)-(238)\right)\right)\right)\approx 116524.8\,\mathrm{J/kg}}\]

\item El caudal convierte trabajo en potencia.

\[\fitmath{\mathrm{Pow}=\mathrm{md}\cdot \mathrm{uu}\cdot \left(\mathrm{ct}-\left(\mathrm{uu}-\mathrm{kk}\cdot \left(\mathrm{ct}-\mathrm{uu}\right)\right)\right)}\]

\[\fitmath{\mathrm{Pow}=(17.8)\cdot (238)\cdot \left((493)-\left((238)-(0.92)\cdot \left((493)-(238)\right)\right)\right)\approx 2074141\,\mathrm{W}}\]

\end{enumerate}


\noindent\textbf{Sustitución numérica}

\[
\fitmath{\left(17.8\,\mathrm{kg}\,\mathrm{s}^{-1}\right) \left(238\,\mathrm{m}\,\mathrm{s}^{-1}\right) \left(\left(493\,\mathrm{m}\,\mathrm{s}^{-1}\right) + \left(0.92\right) \left(\left(493\,\mathrm{m}\,\mathrm{s}^{-1}\right) - \left(238\,\mathrm{m}\,\mathrm{s}^{-1}\right)\right) - \left(238\,\mathrm{m}\,\mathrm{s}^{-1}\right)\right)}
\]

\noindent\textbf{Resultado}

\[
\fitmath{\dot{W}=2.074141440\times10^{6}\,\mathrm{W}}
\]

\noindent\textbf{Interpretación}
La componente axial de salida es k veces la de entrada al mantener el ángulo relativo simétrico. No se impone además acción estática perfecta.

\noindent\textbf{Comprobación de los pasos}
\begin{itemize}
\item Paso 2: $3.4$ — reproducido
\item Paso 4: $116525$ — reproducido
\item Paso 6: $2.07414\times 10^{6}$ — reproducido
\end{itemize}

\noindent\textbf{Comprobación de la respuesta} la aritmética del último paso reproduce la respuesta enunciada a la precisión con que está escrita. Es una reproducción aritmética; no valida el modelo físico ni sus hipótesis.

\noindent\textbf{Limitaciones}
\begin{itemize}
\item Modelo educativo; la aplicabilidad a una máquina real requiere revisión experta.
\end{itemize}

\noindent\textbf{Errores frecuentes}
\begin{itemize}
\item Mantener simultáneamente Cx constante y escala k de todo el vector.
\end{itemize}
% studium:end-record eyJpZCI6ICJQUkItMDAyOSIsICJzZWN0aW9uIjogImMwNCJ9
\end{studiumworked}

\begin{studiumworked}{Problema resuelto C04-03-P}


\noindent\textbf{Enunciado}
\(C04-03-P \cdot\) Rotor con reducción de velocidad relativa. Ejercicio de transferencia. Un rotor recibe W tangencial de entrada positiva y descarga con dirección relativa simétrica y módulo reducido por k. Halle la potencia extraída. Antes de consultar la solución, prediga el efecto de cambiar el primer dato respecto del caso resuelto de esta familia. Justifique la tendencia y explique qué restricción podría convertirse en limitante. Datos: \(C_x=187.44\, \mathrm{m}\,\mathrm{s}^{-1}\); \(C_\theta=493\, \mathrm{m}\,\mathrm{s}^{-1}\); \(U=238\, \mathrm{m}\,\mathrm{s}^{-1}\); \(k_W=0.92\); \(\dot m=17.8\, \mathrm{kg}\,\mathrm{s}^{-1}\).

\noindent\textbf{Datos}
\begin{itemize}
\item $C_x = 187.44\,\mathrm{m}\,\mathrm{s}^{-1}$ — componente axial entrada
\item $C_\theta = 493\,\mathrm{m}\,\mathrm{s}^{-1}$ — remolino entrada
\item $U = 238\,\mathrm{m}\,\mathrm{s}^{-1}$ — velocidad periférica
\item $k_W = 0.92$ — cociente Wsal/Went
\item $\dot m = 17.8\,\mathrm{kg}\,\mathrm{s}^{-1}$ — caudal
\end{itemize}

\noindent\textbf{Incógnita}
$\dot{W}$, en $\mathrm{W}$; también las magnitudes intermedias desarrolladas.

\noindent\textbf{Modelo físico}

\noindent\textit{Rotor con reducción de velocidad relativa}

\[
\fitmath{\dot{W}=\dot m U \left(C_\theta + k_W \left(C_\theta - U\right) - U\right)}
\]

\noindent\textbf{Hipótesis}
\begin{itemize}
\item Radio constante; k y dirección relativa son datos del modelo.
\end{itemize}

\noindent\textbf{Desarrollo paso a paso}

\begin{enumerate}

\item La salida tangencial absoluta combina U y la componente relativa reducida.

\[\fitmath{\mathrm{cts}=\mathrm{uu}-\mathrm{kk}\cdot \left(\mathrm{ct}-\mathrm{uu}\right)}\]

\[\fitmath{\mathrm{cts}=(238)-(0.92)\cdot \left((493)-(238)\right)\approx 3.4\,\mathrm{m/s}}\]

\item Se calcula el trabajo de Euler.

\[\fitmath{\mathrm{wt}=\mathrm{uu}\cdot \left(\mathrm{ct}-\left(\mathrm{uu}-\mathrm{kk}\cdot \left(\mathrm{ct}-\mathrm{uu}\right)\right)\right)}\]

\[\fitmath{\mathrm{wt}=(238)\cdot \left((493)-\left((238)-(0.92)\cdot \left((493)-(238)\right)\right)\right)\approx 116524.8\,\mathrm{J/kg}}\]

\item El caudal convierte trabajo en potencia.

\[\fitmath{\mathrm{Pow}=\mathrm{md}\cdot \mathrm{uu}\cdot \left(\mathrm{ct}-\left(\mathrm{uu}-\mathrm{kk}\cdot \left(\mathrm{ct}-\mathrm{uu}\right)\right)\right)}\]

\[\fitmath{\mathrm{Pow}=(17.8)\cdot (238)\cdot \left((493)-\left((238)-(0.92)\cdot \left((493)-(238)\right)\right)\right)\approx 2074141\,\mathrm{W}}\]

\end{enumerate}


\noindent\textbf{Sustitución numérica}

\[
\fitmath{\left(17.8\,\mathrm{kg}\,\mathrm{s}^{-1}\right) \left(238\,\mathrm{m}\,\mathrm{s}^{-1}\right) \left(\left(493\,\mathrm{m}\,\mathrm{s}^{-1}\right) + \left(0.92\right) \left(\left(493\,\mathrm{m}\,\mathrm{s}^{-1}\right) - \left(238\,\mathrm{m}\,\mathrm{s}^{-1}\right)\right) - \left(238\,\mathrm{m}\,\mathrm{s}^{-1}\right)\right)}
\]

\noindent\textbf{Resultado}

\[
\fitmath{\dot{W}=2.074141440\times10^{6}\,\mathrm{W}}
\]

\noindent\textbf{Interpretación}
La componente axial de salida es k veces la de entrada al mantener el ángulo relativo simétrico. No se impone además acción estática perfecta.

\noindent\textbf{Comprobación de los pasos}
\begin{itemize}
\item Paso 2: $3.4$ — reproducido
\item Paso 4: $116525$ — reproducido
\item Paso 6: $2.07414\times 10^{6}$ — reproducido
\end{itemize}

\noindent\textbf{Comprobación de la respuesta} la aritmética del último paso reproduce la respuesta enunciada a la precisión con que está escrita. Es una reproducción aritmética; no valida el modelo físico ni sus hipótesis.

\noindent\textbf{Limitaciones}
\begin{itemize}
\item Modelo educativo; la aplicabilidad a una máquina real requiere revisión experta.
\end{itemize}

\noindent\textbf{Errores frecuentes}
\begin{itemize}
\item Mantener simultáneamente Cx constante y escala k de todo el vector.
\end{itemize}
% studium:end-record eyJpZCI6ICJQUkItMDAzMCIsICJzZWN0aW9uIjogImMwNCJ9
\end{studiumworked}

\begin{studiumworked}{Problema resuelto C04-04-R}


\noindent\textbf{Enunciado}
\(C04-04-R \cdot\) Carga de cascada de turbina. Estime el parámetro de carga Z de una fila con ángulos algebraicos dados, en el modelo de baja compresibilidad definido en el capítulo. Datos: \(\alpha_e=0.8203047\, \mathrm{rad}\); \(\alpha_s=-1.012291\, \mathrm{rad}\); \(s=0.024\, \mathrm{m}\); \(c_x=0.049\, \mathrm{m}\).

\noindent\textbf{Datos}
\begin{itemize}
\item $\alpha_e = 0.82030474843733\,\mathrm{rad}$ — ángulo entrada en referencia de fila
\item $\alpha_s = -1.0122909661567\,\mathrm{rad}$ — ángulo salida
\item $s = 0.024\,\mathrm{m}$ — paso
\item $c_x = 0.049\,\mathrm{m}$ — cuerda axial
\end{itemize}

\noindent\textbf{Incógnita}
$Z$ (adimensional); también las magnitudes intermedias desarrolladas.

\noindent\textbf{Modelo físico}

\noindent\textit{Carga de cascada de turbina}

\[
\fitmath{Z=\frac{2 s \left(\tan{\left(\alpha_e \right)} - \tan{\left(\alpha_s \right)}\right) \cos^{2}{\left(\alpha_s \right)}}{c_x}}
\]

\noindent\textbf{Hipótesis}
\begin{itemize}
\item Densidad y velocidad axial aproximadamente constantes; ángulos algebraicos.
\end{itemize}

\noindent\textbf{Desarrollo paso a paso}

\begin{enumerate}

\item Se calcula la razón de paso sobre cuerda axial.

\[\fitmath{\mathrm{sc}=\frac{\mathrm{ss}}{\mathrm{cc}}}\]

\[\fitmath{\mathrm{sc}=\frac{(0.024)}{(0.049)}\approx 0.4897959\,}\]

\item Se evalúa la carga a partir del cambio tangencial normalizado.

\[\fitmath{Z=2\cdot \frac{\mathrm{ss}}{\mathrm{cc}}\cdot {\cos\left(\mathrm{ps}\right)}^{2}\cdot \left(\tan\left(\mathrm{pe}\right)-\tan\left(\mathrm{ps}\right)\right)}\]

\[\fitmath{Z=2\cdot \frac{(0.024)}{(0.049)}\cdot {\cos\left((-1.012290966)\right)}^{2}\cdot \left(\tan\left((0.8203047484)\right)-\tan\left((-1.012290966)\right)\right)\approx 0.7352166\,}\]

\end{enumerate}


\noindent\textbf{Sustitución numérica}

\[
\fitmath{\frac{2 \left(0.024\,\mathrm{m}\right) \left(\tan{\left(\left(0.82030474843733\,\mathrm{rad}\right) \right)} - \tan{\left(\left(-1.0122909661567\,\mathrm{rad}\right) \right)}\right) \cos^{2}{\left(\left(-1.0122909661567\,\mathrm{rad}\right) \right)}}{\left(0.049\,\mathrm{m}\right)}}
\]

\noindent\textbf{Resultado}

\[
\fitmath{Z=7.352166148\times10^{-1}}
\]

\noindent\textbf{Interpretación}
El resultado es un indicador de carga del modelo. No permite asignar una pérdida ni una eficiencia sin correlación contrastada.

\noindent\textbf{Comprobación de los pasos}
\begin{itemize}
\item Paso 2: $0.489796$ — reproducido
\item Paso 4: $0.735217$ — reproducido
\end{itemize}

\noindent\textbf{Comprobación de la respuesta} la aritmética del último paso reproduce la respuesta enunciada a la precisión con que está escrita. Es una reproducción aritmética; no valida el modelo físico ni sus hipótesis.

\noindent\textbf{Limitaciones}
\begin{itemize}
\item Modelo educativo; la aplicabilidad a una máquina real requiere revisión experta.
\end{itemize}

\noindent\textbf{Errores frecuentes}
\begin{itemize}
\item Usar cuerda real donde el parámetro pide cuerda axial.
\end{itemize}
% studium:end-record eyJpZCI6ICJQUkItMDAzMSIsICJzZWN0aW9uIjogImMwNCJ9
\end{studiumworked}

\begin{studiumworked}{Problema resuelto C04-04-P}


\noindent\textbf{Enunciado}
\(C04-04-P \cdot\) Carga de cascada de turbina. Ejercicio de transferencia. Estime el parámetro de carga Z de una fila con ángulos algebraicos dados, en el modelo de baja compresibilidad definido en el capítulo. Antes de consultar la solución, prediga el efecto de cambiar el primer dato respecto del caso resuelto de esta familia. Justifique la tendencia y explique qué restricción podría convertirse en limitante. Datos: \(\alpha_e=0.8736246\, \mathrm{rad}\); \(\alpha_s=-1.012291\, \mathrm{rad}\); \(s=0.024\, \mathrm{m}\); \(c_x=0.049\, \mathrm{m}\).

\noindent\textbf{Datos}
\begin{itemize}
\item $\alpha_e = 0.87362455708576\,\mathrm{rad}$ — ángulo entrada en referencia de fila
\item $\alpha_s = -1.0122909661567\,\mathrm{rad}$ — ángulo salida
\item $s = 0.024\,\mathrm{m}$ — paso
\item $c_x = 0.049\,\mathrm{m}$ — cuerda axial
\end{itemize}

\noindent\textbf{Incógnita}
$Z$ (adimensional); también las magnitudes intermedias desarrolladas.

\noindent\textbf{Modelo físico}

\noindent\textit{Carga de cascada de turbina}

\[
\fitmath{Z=\frac{2 s \left(\tan{\left(\alpha_e \right)} - \tan{\left(\alpha_s \right)}\right) \cos^{2}{\left(\alpha_s \right)}}{c_x}}
\]

\noindent\textbf{Hipótesis}
\begin{itemize}
\item Densidad y velocidad axial aproximadamente constantes; ángulos algebraicos.
\end{itemize}

\noindent\textbf{Desarrollo paso a paso}

\begin{enumerate}

\item Se calcula la razón de paso sobre cuerda axial.

\[\fitmath{\mathrm{sc}=\frac{\mathrm{ss}}{\mathrm{cc}}}\]

\[\fitmath{\mathrm{sc}=\frac{(0.024)}{(0.049)}\approx 0.4897959\,}\]

\item Se evalúa la carga a partir del cambio tangencial normalizado.

\[\fitmath{Z=2\cdot \frac{\mathrm{ss}}{\mathrm{cc}}\cdot {\cos\left(\mathrm{ps}\right)}^{2}\cdot \left(\tan\left(\mathrm{pe}\right)-\tan\left(\mathrm{ps}\right)\right)}\]

\[\fitmath{Z=2\cdot \frac{(0.024)}{(0.049)}\cdot {\cos\left((-1.012290966)\right)}^{2}\cdot \left(\tan\left((0.8736245571)\right)-\tan\left((-1.012290966)\right)\right)\approx 0.7686973\,}\]

\end{enumerate}


\noindent\textbf{Sustitución numérica}

\[
\fitmath{\frac{2 \left(0.024\,\mathrm{m}\right) \left(\tan{\left(\left(0.87362455708576\,\mathrm{rad}\right) \right)} - \tan{\left(\left(-1.0122909661567\,\mathrm{rad}\right) \right)}\right) \cos^{2}{\left(\left(-1.0122909661567\,\mathrm{rad}\right) \right)}}{\left(0.049\,\mathrm{m}\right)}}
\]

\noindent\textbf{Resultado}

\[
\fitmath{Z=7.686972619\times10^{-1}}
\]

\noindent\textbf{Interpretación}
El resultado es un indicador de carga del modelo. No permite asignar una pérdida ni una eficiencia sin correlación contrastada.

\noindent\textbf{Comprobación de los pasos}
\begin{itemize}
\item Paso 2: $0.489796$ — reproducido
\item Paso 4: $0.768697$ — reproducido
\end{itemize}

\noindent\textbf{Comprobación de la respuesta} la aritmética del último paso reproduce la respuesta enunciada a la precisión con que está escrita. Es una reproducción aritmética; no valida el modelo físico ni sus hipótesis.

\noindent\textbf{Limitaciones}
\begin{itemize}
\item Modelo educativo; la aplicabilidad a una máquina real requiere revisión experta.
\end{itemize}

\noindent\textbf{Errores frecuentes}
\begin{itemize}
\item Usar cuerda real donde el parámetro pide cuerda axial.
\end{itemize}
% studium:end-record eyJpZCI6ICJQUkItMDAzMiIsICJzZWN0aW9uIjogImMwNCJ9
\end{studiumworked}

\chapter{Turbinas de reacción y rendimiento del escalón}
\begin{studiumworked}{Problema resuelto C05-01-R}


\noindent\textbf{Enunciado}
\(C05-01-R \cdot\) Triángulos desde carga y reacción. Diseñe las componentes del escalón repetitivo con carga, reacción y velocidad periférica dadas. Obtenga la componente relativa tangencial a la salida del rotor. Datos: \(U=305\, \mathrm{m}\,\mathrm{s}^{-1}\); \(\psi=1.42\); \(\mathcal R=0.46\).

\noindent\textbf{Datos}
\begin{itemize}
\item $U = 305\,\mathrm{m}\,\mathrm{s}^{-1}$ — velocidad periférica
\item $\psi = 1.42$ — carga turbina
\item $\mathcal R = 0.46$ — reacción
\end{itemize}

\noindent\textbf{Incógnita}
$W_{\theta 3}$, en $\mathrm{m\,s^{-1}}$; también las magnitudes intermedias desarrolladas.

\noindent\textbf{Modelo físico}

\noindent\textit{Triángulos desde carga y reacción}

\[
\fitmath{W_{\theta 3}=U \left(- \frac{\psi}{2} - \mathcal R\right)}
\]

\noindent\textbf{Hipótesis}
\begin{itemize}
\item Radio y componente axial constantes; igual energía cinética en extremos del escalón.
\end{itemize}

\noindent\textbf{Desarrollo paso a paso}

\begin{enumerate}

\item La suma y diferencia de remolinos dan la entrada.

\[\fitmath{ct2=\mathrm{uu}\cdot \left(1-\mathrm{rea}+\frac{\mathrm{psi}}{2}\right)}\]

\[\fitmath{ct2=(305)\cdot \left(1-(0.46)+\frac{(1.42)}{2}\right)\approx 381.25\,\mathrm{m/s}}\]

\item Se obtiene la descarga absoluta.

\[\fitmath{ct3=\mathrm{uu}\cdot \left(1-\mathrm{rea}-\frac{\mathrm{psi}}{2}\right)}\]

\[\fitmath{ct3=(305)\cdot \left(1-(0.46)-\frac{(1.42)}{2}\right)\approx -51.85\,\mathrm{m/s}}\]

\item La referencia relativa resta U a la descarga.

\[\fitmath{wt3=\mathrm{uu}\cdot \left(-\mathrm{rea}-\frac{\mathrm{psi}}{2}\right)}\]

\[\fitmath{wt3=(305)\cdot \left(-(0.46)-\frac{(1.42)}{2}\right)\approx -356.85\,\mathrm{m/s}}\]

\end{enumerate}


\noindent\textbf{Sustitución numérica}

\[
\fitmath{\left(305\,\mathrm{m}\,\mathrm{s}^{-1}\right) \left(- \frac{\left(1.42\right)}{2} - \left(0.46\right)\right)}
\]

\noindent\textbf{Resultado}

\[
\fitmath{W_{\theta 3}=-3.568500000\times10^{2}\,\mathrm{m}\,\mathrm{s}^{-1}}
\]

\noindent\textbf{Interpretación}
La componente relativa de salida es negativa; su módulo no debe sustituir su signo en la suma vectorial.

\noindent\textbf{Comprobación de los pasos}
\begin{itemize}
\item Paso 2: $381.25$ — reproducido
\item Paso 4: $-51.85$ — reproducido
\item Paso 6: $-356.85$ — reproducido
\end{itemize}

\noindent\textbf{Comprobación de la respuesta} la aritmética del último paso reproduce la respuesta enunciada a la precisión con que está escrita. Es una reproducción aritmética; no valida el modelo físico ni sus hipótesis.

\noindent\textbf{Limitaciones}
\begin{itemize}
\item Modelo educativo; la aplicabilidad a una máquina real requiere revisión experta.
\end{itemize}

\noindent\textbf{Errores frecuentes}
\begin{itemize}
\item Elegir carga, reacción y descarga axial como tres datos independientes.
\end{itemize}
% studium:end-record eyJpZCI6ICJQUkItMDAzMyIsICJzZWN0aW9uIjogImMwNSJ9
\end{studiumworked}

\begin{studiumworked}{Problema resuelto C05-01-P}


\noindent\textbf{Enunciado}
\(C05-01-P \cdot\) Triángulos desde carga y reacción. Ejercicio de transferencia. Diseñe las componentes del escalón repetitivo con carga, reacción y velocidad periférica dadas. Obtenga la componente relativa tangencial a la salida del rotor. Antes de consultar la solución, prediga el efecto de cambiar el primer dato respecto del caso resuelto de esta familia. Justifique la tendencia y explique qué restricción podría convertirse en limitante. Datos: \(U=324.825\, \mathrm{m}\,\mathrm{s}^{-1}\); \(\psi=1.42\); \(\mathcal R=0.46\).

\noindent\textbf{Datos}
\begin{itemize}
\item $U = 324.825\,\mathrm{m}\,\mathrm{s}^{-1}$ — velocidad periférica
\item $\psi = 1.42$ — carga turbina
\item $\mathcal R = 0.46$ — reacción
\end{itemize}

\noindent\textbf{Incógnita}
$W_{\theta 3}$, en $\mathrm{m\,s^{-1}}$; también las magnitudes intermedias desarrolladas.

\noindent\textbf{Modelo físico}

\noindent\textit{Triángulos desde carga y reacción}

\[
\fitmath{W_{\theta 3}=U \left(- \frac{\psi}{2} - \mathcal R\right)}
\]

\noindent\textbf{Hipótesis}
\begin{itemize}
\item Radio y componente axial constantes; igual energía cinética en extremos del escalón.
\end{itemize}

\noindent\textbf{Desarrollo paso a paso}

\begin{enumerate}

\item La suma y diferencia de remolinos dan la entrada.

\[\fitmath{\mathrm{ct}_{2}=\mathrm{uu}\cdot \left(1-\mathrm{rea}+\frac{\psi}{2}\right)}\]

\[\fitmath{\mathrm{ct}_{2}=(324.825)\cdot \left(1-(0.46)+\frac{(1.42)}{2}\right)\approx 406.0312\,\mathrm{m/s}}\]

\item Se obtiene la descarga absoluta.

\[\fitmath{\mathrm{ct}_{3}=\mathrm{uu}\cdot \left(1-\mathrm{rea}-\frac{\psi}{2}\right)}\]

\[\fitmath{\mathrm{ct}_{3}=(324.825)\cdot \left(1-(0.46)-\frac{(1.42)}{2}\right)\approx -55.22025\,\mathrm{m/s}}\]

\item La referencia relativa resta U a la descarga.

\[\fitmath{\mathrm{wt}_{3}=\mathrm{uu}\cdot \left(-\mathrm{rea}-\frac{\psi}{2}\right)}\]

\[\fitmath{\mathrm{wt}_{3}=(324.825)\cdot \left(-(0.46)-\frac{(1.42)}{2}\right)\approx -380.0453\,\mathrm{m/s}}\]

\end{enumerate}


\noindent\textbf{Sustitución numérica}

\[
\fitmath{\left(324.825\,\mathrm{m}\,\mathrm{s}^{-1}\right) \left(- \frac{\left(1.42\right)}{2} - \left(0.46\right)\right)}
\]

\noindent\textbf{Resultado}

\[
\fitmath{W_{\theta 3}=-3.800452500\times10^{2}\,\mathrm{m}\,\mathrm{s}^{-1}}
\]

\noindent\textbf{Interpretación}
La componente relativa de salida es negativa; su módulo no debe sustituir su signo en la suma vectorial.

\noindent\textbf{Comprobación de los pasos}
\begin{itemize}
\item Paso 2: $406.031$ — reproducido
\item Paso 4: $-55.2202$ — reproducido
\item Paso 6: $-380.045$ — reproducido
\end{itemize}

\noindent\textbf{Comprobación de la respuesta} la aritmética del último paso reproduce la respuesta enunciada a la precisión con que está escrita. Es una reproducción aritmética; no valida el modelo físico ni sus hipótesis.

\noindent\textbf{Limitaciones}
\begin{itemize}
\item Modelo educativo; la aplicabilidad a una máquina real requiere revisión experta.
\end{itemize}

\noindent\textbf{Errores frecuentes}
\begin{itemize}
\item Elegir carga, reacción y descarga axial como tres datos independientes.
\end{itemize}
% studium:end-record eyJpZCI6ICJQUkItMDAzNCIsICJzZWN0aW9uIjogImMwNSJ9
\end{studiumworked}

\begin{studiumworked}{Problema resuelto C05-02-R}


\noindent\textbf{Enunciado}
\(C05-02-R \cdot\) Reacción general no repetitiva. Calcule la reacción estática de una turbina cuando las velocidades absolutas de entrada y salida del escalón difieren. Datos: \(U=287\, \mathrm{m}\,\mathrm{s}^{-1}\); \(C_{\theta2}=421\, \mathrm{m}\,\mathrm{s}^{-1}\); \(C_{\theta3}=-38\, \mathrm{m}\,\mathrm{s}^{-1}\); \(C_1=127\, \mathrm{m}\,\mathrm{s}^{-1}\); \(C_2=467\, \mathrm{m}\,\mathrm{s}^{-1}\); \(C_3=203\, \mathrm{m}\,\mathrm{s}^{-1}\).

\noindent\textbf{Datos}
\begin{itemize}
\item $U = 287\,\mathrm{m}\,\mathrm{s}^{-1}$ — velocidad periférica
\item $C_{\theta 2} = 421\,\mathrm{m}\,\mathrm{s}^{-1}$ — remolino entrada rotor
\item $C_{\theta 3} = -38\,\mathrm{m}\,\mathrm{s}^{-1}$ — remolino salida rotor
\item $c_{1} = 127\,\mathrm{m}\,\mathrm{s}^{-1}$ — módulo entrada escalón
\item $c_{2} = 467\,\mathrm{m}\,\mathrm{s}^{-1}$ — módulo entrada rotor
\item $c_{3} = 203\,\mathrm{m}\,\mathrm{s}^{-1}$ — módulo salida escalón
\end{itemize}

\noindent\textbf{Incógnita}
$\mathcal{R}$ (adimensional); también las magnitudes intermedias desarrolladas.

\noindent\textbf{Modelo físico}

\noindent\textit{Reacción general no repetitiva}

\[
\fitmath{\mathcal R=\frac{- \frac{c_{2}^{2}}{2} + \frac{c_{3}^{2}}{2} + U \left(C_{\theta 2} - C_{\theta 3}\right)}{- \frac{c_{1}^{2}}{2} + \frac{c_{3}^{2}}{2} + U \left(C_{\theta 2} - C_{\theta 3}\right)}}
\]

\noindent\textbf{Hipótesis}
\begin{itemize}
\item Rotor y estator adiabáticos; radios constantes para Euler.
\end{itemize}

\noindent\textbf{Desarrollo paso a paso}

\begin{enumerate}

\item Euler da el trabajo extraído.

\[\fitmath{\mathrm{wt}=\mathrm{uu}\cdot \left(\mathrm{ct}_{2}-\mathrm{ct}_{3}\right)}\]

\[\fitmath{\mathrm{wt}=(287)\cdot \left((421)-(-38)\right)\approx 131733\,\mathrm{J/kg}}\]

\item La caída estática del rotor incluye el cambio cinético.

\[\fitmath{\mathrm{dhr}=\mathrm{uu}\cdot \left(\mathrm{ct}_{2}-\mathrm{ct}_{3}\right)+\frac{{c_{3}}^{2}-{c_{2}}^{2}}{2}}\]

\[\fitmath{\mathrm{dhr}=(287)\cdot \left((421)-(-38)\right)+\frac{{(203)}^{2}-{(467)}^{2}}{2}\approx 43293\,\mathrm{J/kg}}\]

\item La caída estática del escalón tiene extremos 1 y 3.

\[\fitmath{\mathrm{rea}=\frac{\mathrm{uu}\cdot \left(\mathrm{ct}_{2}-\mathrm{ct}_{3}\right)+\frac{{c_{3}}^{2}-{c_{2}}^{2}}{2}}{\mathrm{uu}\cdot \left(\mathrm{ct}_{2}-\mathrm{ct}_{3}\right)+\frac{{c_{3}}^{2}-{c_{1}}^{2}}{2}}}\]

\[\fitmath{\mathrm{rea}=\frac{(287)\cdot \left((421)-(-38)\right)+\frac{{(203)}^{2}-{(467)}^{2}}{2}}{(287)\cdot \left((421)-(-38)\right)+\frac{{(203)}^{2}-{(127)}^{2}}{2}}\approx 0.3000769\,}\]

\end{enumerate}


\noindent\textbf{Sustitución numérica}

\[
\fitmath{\frac{- \frac{\left(\left(467\,\mathrm{m}\,\mathrm{s}^{-1}\right)\right)^{2}}{2} + \frac{\left(\left(203\,\mathrm{m}\,\mathrm{s}^{-1}\right)\right)^{2}}{2} + \left(287\,\mathrm{m}\,\mathrm{s}^{-1}\right) \left(\left(421\,\mathrm{m}\,\mathrm{s}^{-1}\right) - \left(-38\,\mathrm{m}\,\mathrm{s}^{-1}\right)\right)}{- \frac{\left(\left(127\,\mathrm{m}\,\mathrm{s}^{-1}\right)\right)^{2}}{2} + \frac{\left(\left(203\,\mathrm{m}\,\mathrm{s}^{-1}\right)\right)^{2}}{2} + \left(287\,\mathrm{m}\,\mathrm{s}^{-1}\right) \left(\left(421\,\mathrm{m}\,\mathrm{s}^{-1}\right) - \left(-38\,\mathrm{m}\,\mathrm{s}^{-1}\right)\right)}}
\]

\noindent\textbf{Resultado}

\[
\fitmath{\mathcal R=3.000769375\times10^{-1}}
\]

\noindent\textbf{Interpretación}
La reacción general difiere de la fórmula reducida si cambia la energía cinética entre extremos.

\noindent\textbf{Comprobación de los pasos}
\begin{itemize}
\item Paso 2: $131733$ — reproducido
\item Paso 4: $43293$ — reproducido
\item Paso 6: $0.300077$ — reproducido
\end{itemize}

\noindent\textbf{Comprobación de la respuesta} la aritmética del último paso reproduce la respuesta enunciada a la precisión con que está escrita. Es una reproducción aritmética; no valida el modelo físico ni sus hipótesis.

\noindent\textbf{Limitaciones}
\begin{itemize}
\item Modelo educativo; la aplicabilidad a una máquina real requiere revisión experta.
\end{itemize}

\noindent\textbf{Errores frecuentes}
\begin{itemize}
\item Usar caída total en el denominador de una reacción definida como estática.
\end{itemize}
% studium:end-record eyJpZCI6ICJQUkItMDAzNSIsICJzZWN0aW9uIjogImMwNSJ9
\end{studiumworked}

\begin{studiumworked}{Problema resuelto C05-02-P}


\noindent\textbf{Enunciado}
\(C05-02-P \cdot\) Reacción general no repetitiva. Ejercicio de transferencia. Calcule la reacción estática de una turbina cuando las velocidades absolutas de entrada y salida del escalón difieren. Antes de consultar la solución, prediga el efecto de cambiar el primer dato respecto del caso resuelto de esta familia. Justifique la tendencia y explique qué restricción podría convertirse en limitante. Datos: \(U=305.655\, \mathrm{m}\,\mathrm{s}^{-1}\); \(C_{\theta2}=421\, \mathrm{m}\,\mathrm{s}^{-1}\); \(C_{\theta3}=-38\, \mathrm{m}\,\mathrm{s}^{-1}\); \(C_1=127\, \mathrm{m}\,\mathrm{s}^{-1}\); \(C_2=467\, \mathrm{m}\,\mathrm{s}^{-1}\); \(C_3=203\, \mathrm{m}\,\mathrm{s}^{-1}\).

\noindent\textbf{Datos}
\begin{itemize}
\item $U = 305.655\,\mathrm{m}\,\mathrm{s}^{-1}$ — velocidad periférica
\item $C_{\theta 2} = 421\,\mathrm{m}\,\mathrm{s}^{-1}$ — remolino entrada rotor
\item $C_{\theta 3} = -38\,\mathrm{m}\,\mathrm{s}^{-1}$ — remolino salida rotor
\item $c_{1} = 127\,\mathrm{m}\,\mathrm{s}^{-1}$ — módulo entrada escalón
\item $c_{2} = 467\,\mathrm{m}\,\mathrm{s}^{-1}$ — módulo entrada rotor
\item $c_{3} = 203\,\mathrm{m}\,\mathrm{s}^{-1}$ — módulo salida escalón
\end{itemize}

\noindent\textbf{Incógnita}
$\mathcal{R}$ (adimensional); también las magnitudes intermedias desarrolladas.

\noindent\textbf{Modelo físico}

\noindent\textit{Reacción general no repetitiva}

\[
\fitmath{\mathcal R=\frac{- \frac{c_{2}^{2}}{2} + \frac{c_{3}^{2}}{2} + U \left(C_{\theta 2} - C_{\theta 3}\right)}{- \frac{c_{1}^{2}}{2} + \frac{c_{3}^{2}}{2} + U \left(C_{\theta 2} - C_{\theta 3}\right)}}
\]

\noindent\textbf{Hipótesis}
\begin{itemize}
\item Rotor y estator adiabáticos; radios constantes para Euler.
\end{itemize}

\noindent\textbf{Desarrollo paso a paso}

\begin{enumerate}

\item Euler da el trabajo extraído.

\[\fitmath{\mathrm{wt}=\mathrm{uu}\cdot \left(\mathrm{ct}_{2}-\mathrm{ct}_{3}\right)}\]

\[\fitmath{\mathrm{wt}=(305.655)\cdot \left((421)-(-38)\right)\approx 140295.6\,\mathrm{J/kg}}\]

\item La caída estática del rotor incluye el cambio cinético.

\[\fitmath{\mathrm{dhr}=\mathrm{uu}\cdot \left(\mathrm{ct}_{2}-\mathrm{ct}_{3}\right)+\frac{{c_{3}}^{2}-{c_{2}}^{2}}{2}}\]

\[\fitmath{\mathrm{dhr}=(305.655)\cdot \left((421)-(-38)\right)+\frac{{(203)}^{2}-{(467)}^{2}}{2}\approx 51855.64\,\mathrm{J/kg}}\]

\item La caída estática del escalón tiene extremos 1 y 3.

\[\fitmath{\mathrm{rea}=\frac{\mathrm{uu}\cdot \left(\mathrm{ct}_{2}-\mathrm{ct}_{3}\right)+\frac{{c_{3}}^{2}-{c_{2}}^{2}}{2}}{\mathrm{uu}\cdot \left(\mathrm{ct}_{2}-\mathrm{ct}_{3}\right)+\frac{{c_{3}}^{2}-{c_{1}}^{2}}{2}}}\]

\[\fitmath{\mathrm{rea}=\frac{(305.655)\cdot \left((421)-(-38)\right)+\frac{{(203)}^{2}-{(467)}^{2}}{2}}{(305.655)\cdot \left((421)-(-38)\right)+\frac{{(203)}^{2}-{(127)}^{2}}{2}}\approx 0.3392903\,}\]

\end{enumerate}


\noindent\textbf{Sustitución numérica}

\[
\fitmath{\frac{- \frac{\left(\left(467\,\mathrm{m}\,\mathrm{s}^{-1}\right)\right)^{2}}{2} + \frac{\left(\left(203\,\mathrm{m}\,\mathrm{s}^{-1}\right)\right)^{2}}{2} + \left(305.655\,\mathrm{m}\,\mathrm{s}^{-1}\right) \left(\left(421\,\mathrm{m}\,\mathrm{s}^{-1}\right) - \left(-38\,\mathrm{m}\,\mathrm{s}^{-1}\right)\right)}{- \frac{\left(\left(127\,\mathrm{m}\,\mathrm{s}^{-1}\right)\right)^{2}}{2} + \frac{\left(\left(203\,\mathrm{m}\,\mathrm{s}^{-1}\right)\right)^{2}}{2} + \left(305.655\,\mathrm{m}\,\mathrm{s}^{-1}\right) \left(\left(421\,\mathrm{m}\,\mathrm{s}^{-1}\right) - \left(-38\,\mathrm{m}\,\mathrm{s}^{-1}\right)\right)}}
\]

\noindent\textbf{Resultado}

\[
\fitmath{\mathcal R=3.392902552\times10^{-1}}
\]

\noindent\textbf{Interpretación}
La reacción general difiere de la fórmula reducida si cambia la energía cinética entre extremos.

\noindent\textbf{Comprobación de los pasos}
\begin{itemize}
\item Paso 2: $140296$ — reproducido
\item Paso 4: $51855.6$ — reproducido
\item Paso 6: $0.33929$ — reproducido
\end{itemize}

\noindent\textbf{Comprobación de la respuesta} la aritmética del último paso reproduce la respuesta enunciada a la precisión con que está escrita. Es una reproducción aritmética; no valida el modelo físico ni sus hipótesis.

\noindent\textbf{Limitaciones}
\begin{itemize}
\item Modelo educativo; la aplicabilidad a una máquina real requiere revisión experta.
\end{itemize}

\noindent\textbf{Errores frecuentes}
\begin{itemize}
\item Usar caída total en el denominador de una reacción definida como estática.
\end{itemize}
% studium:end-record eyJpZCI6ICJQUkItMDAzNiIsICJzZWN0aW9uIjogImMwNSJ9
\end{studiumworked}

\begin{studiumworked}{Problema resuelto C05-03-R}


\noindent\textbf{Enunciado}
\(C05-03-R \cdot\) Presión de descarga de turbina. Dados trabajo y rendimiento total-total de un escalón, halle la presión total de salida. Datos: \(T_{0e}=1510\, \mathrm{K}\); \(p_{0e}=1240000\, \mathrm{Pa}\); \(w_t=156000\, \mathrm{J}\,\mathrm{kg}^{-1}\); \(\eta=0.885\); \(c_p=1156.697\, \mathrm{J}\,\mathrm{kg}^{-1}\,\mathrm{K}^{-1}\); \(\gamma=1.33\).

\noindent\textbf{Datos}
\begin{itemize}
\item $T_{0e} = 1510\,\mathrm{K}$ — temperatura total entrada
\item $p_{0e} = 1240000\,\mathrm{Pa}$ — presión total entrada
\item $w_t = 156000\,\mathrm{J}\,\mathrm{kg}^{-1}$ — trabajo extraído
\item $\eta = 0.885$ — rendimiento total-total
\item $c_p = 1156.696969697\,\mathrm{J}\,\mathrm{kg}^{-1}\,\mathrm{K}^{-1}$ — calor específico
\item $\gamma = 1.33$ — razón de calores
\end{itemize}

\noindent\textbf{Incógnita}
$p_{0s}$, en $\mathrm{Pa}$; también las magnitudes intermedias desarrolladas.

\noindent\textbf{Modelo físico}

\noindent\textit{Presión de descarga de turbina}

\[
\fitmath{p_{0s}=p_{0e} \left(1 - \frac{w_t}{T_{0e} c_p \eta}\right)^{\frac{\gamma}{\gamma - 1}}}
\]

\noindent\textbf{Hipótesis}
\begin{itemize}
\item Sin refrigeración ni sangrado; gas perfecto constante.
\end{itemize}

\noindent\textbf{Desarrollo paso a paso}

\begin{enumerate}

\item Se determina la caída de temperatura ideal equivalente.

\[\fitmath{\mathrm{dTs}=\frac{\mathrm{wt}}{\eta\cdot \mathrm{cp}}}\]

\[\fitmath{\mathrm{dTs}=\frac{(156000)}{(0.885)\cdot (1156.69697)}\approx 152.3918\,\mathrm{K}}\]

\item Se calcula la razón de temperaturas isentrópicas.

\[\fitmath{\mathrm{rr}=1-\frac{\mathrm{wt}}{\eta\cdot \mathrm{cp}\cdot \mathrm{Tin}}}\]

\[\fitmath{\mathrm{rr}=1-\frac{(156000)}{(0.885)\cdot (1156.69697)\cdot (1510)}\approx 0.8990782\,}\]

\item La razón de presiones corresponde al estado ideal equivalente.

\[\fitmath{\mathrm{pout}=\mathrm{pin}\cdot {\left(1-\frac{\mathrm{wt}}{\eta\cdot \mathrm{cp}\cdot \mathrm{Tin}}\right)}^{\frac{\gamma}{\gamma-1}}}\]

\[\fitmath{\mathrm{pout}=(1240000)\cdot {\left(1-\frac{(156000)}{(0.885)\cdot (1156.69697)\cdot (1510)}\right)}^{\frac{(1.33)}{(1.33)-1}}\approx 807628.4\,\mathrm{Pa}}\]

\end{enumerate}


\noindent\textbf{Sustitución numérica}

\[
\fitmath{\left(1240000\,\mathrm{Pa}\right) \left(1 - \frac{\left(156000\,\mathrm{J}\,\mathrm{kg}^{-1}\right)}{\left(1510\,\mathrm{K}\right) \left(1156.696969697\,\mathrm{J}\,\mathrm{kg}^{-1}\,\mathrm{K}^{-1}\right) \left(0.885\right)}\right)^{\frac{\left(1.33\right)}{\left(1.33\right) - 1}}}
\]

\noindent\textbf{Resultado}

\[
\fitmath{p_{0s}=8.076283789\times10^{5}\,\mathrm{Pa}}
\]

\noindent\textbf{Interpretación}
La base debe ser positiva. La presión resultante es total absoluta, no estática.

\noindent\textbf{Comprobación de los pasos}
\begin{itemize}
\item Paso 2: $152.392$ — reproducido
\item Paso 4: $0.899078$ — reproducido
\item Paso 6: $807628$ — reproducido
\end{itemize}

\noindent\textbf{Comprobación de la respuesta} la aritmética del último paso reproduce la respuesta enunciada a la precisión con que está escrita. Es una reproducción aritmética; no valida el modelo físico ni sus hipótesis.

\noindent\textbf{Limitaciones}
\begin{itemize}
\item Modelo educativo; la aplicabilidad a una máquina real requiere revisión experta.
\end{itemize}

\noindent\textbf{Errores frecuentes}
\begin{itemize}
\item Introducir eficiencia total-estática en una expresión total-total.
\end{itemize}
% studium:end-record eyJpZCI6ICJQUkItMDAzNyIsICJzZWN0aW9uIjogImMwNSJ9
\end{studiumworked}

\begin{studiumworked}{Problema resuelto C05-03-P}


\noindent\textbf{Enunciado}
\(C05-03-P \cdot\) Presión de descarga de turbina. Ejercicio de transferencia. Dados trabajo y rendimiento total-total de un escalón, halle la presión total de salida. Antes de consultar la solución, prediga el efecto de cambiar el primer dato respecto del caso resuelto de esta familia. Justifique la tendencia y explique qué restricción podría convertirse en limitante. Datos: \(T_{0e}=1608.15\, \mathrm{K}\); \(p_{0e}=1240000\, \mathrm{Pa}\); \(w_t=156000\, \mathrm{J}\,\mathrm{kg}^{-1}\); \(\eta=0.885\); \(c_p=1156.697\, \mathrm{J}\,\mathrm{kg}^{-1}\,\mathrm{K}^{-1}\); \(\gamma=1.33\).

\noindent\textbf{Datos}
\begin{itemize}
\item $T_{0e} = 1608.15\,\mathrm{K}$ — temperatura total entrada
\item $p_{0e} = 1240000\,\mathrm{Pa}$ — presión total entrada
\item $w_t = 156000\,\mathrm{J}\,\mathrm{kg}^{-1}$ — trabajo extraído
\item $\eta = 0.885$ — rendimiento total-total
\item $c_p = 1156.696969697\,\mathrm{J}\,\mathrm{kg}^{-1}\,\mathrm{K}^{-1}$ — calor específico
\item $\gamma = 1.33$ — razón de calores
\end{itemize}

\noindent\textbf{Incógnita}
$p_{0s}$, en $\mathrm{Pa}$; también las magnitudes intermedias desarrolladas.

\noindent\textbf{Modelo físico}

\noindent\textit{Presión de descarga de turbina}

\[
\fitmath{p_{0s}=p_{0e} \left(1 - \frac{w_t}{T_{0e} c_p \eta}\right)^{\frac{\gamma}{\gamma - 1}}}
\]

\noindent\textbf{Hipótesis}
\begin{itemize}
\item Sin refrigeración ni sangrado; gas perfecto constante.
\end{itemize}

\noindent\textbf{Desarrollo paso a paso}

\begin{enumerate}

\item Se determina la caída de temperatura ideal equivalente.

\[\fitmath{\mathrm{dTs}=\frac{\mathrm{wt}}{\eta\cdot \mathrm{cp}}}\]

\[\fitmath{\mathrm{dTs}=\frac{(156000)}{(0.885)\cdot (1156.69697)}\approx 152.3918\,\mathrm{K}}\]

\item Se calcula la razón de temperaturas isentrópicas.

\[\fitmath{\mathrm{rr}=1-\frac{\mathrm{wt}}{\eta\cdot \mathrm{cp}\cdot \mathrm{Tin}}}\]

\[\fitmath{\mathrm{rr}=1-\frac{(156000)}{(0.885)\cdot (1156.69697)\cdot (1608.15)}\approx 0.9052378\,}\]

\item La razón de presiones corresponde al estado ideal equivalente.

\[\fitmath{\mathrm{pout}=\mathrm{pin}\cdot {\left(1-\frac{\mathrm{wt}}{\eta\cdot \mathrm{cp}\cdot \mathrm{Tin}}\right)}^{\frac{\gamma}{\gamma-1}}}\]

\[\fitmath{\mathrm{pout}=(1240000)\cdot {\left(1-\frac{(156000)}{(0.885)\cdot (1156.69697)\cdot (1608.15)}\right)}^{\frac{(1.33)}{(1.33)-1}}\approx 830160.7\,\mathrm{Pa}}\]

\end{enumerate}


\noindent\textbf{Sustitución numérica}

\[
\fitmath{\left(1240000\,\mathrm{Pa}\right) \left(1 - \frac{\left(156000\,\mathrm{J}\,\mathrm{kg}^{-1}\right)}{\left(1608.15\,\mathrm{K}\right) \left(1156.696969697\,\mathrm{J}\,\mathrm{kg}^{-1}\,\mathrm{K}^{-1}\right) \left(0.885\right)}\right)^{\frac{\left(1.33\right)}{\left(1.33\right) - 1}}}
\]

\noindent\textbf{Resultado}

\[
\fitmath{p_{0s}=8.301606971\times10^{5}\,\mathrm{Pa}}
\]

\noindent\textbf{Interpretación}
La base debe ser positiva. La presión resultante es total absoluta, no estática.

\noindent\textbf{Comprobación de los pasos}
\begin{itemize}
\item Paso 2: $152.392$ — reproducido
\item Paso 4: $0.905238$ — reproducido
\item Paso 6: $830161$ — reproducido
\end{itemize}

\noindent\textbf{Comprobación de la respuesta} la aritmética del último paso reproduce la respuesta enunciada a la precisión con que está escrita. Es una reproducción aritmética; no valida el modelo físico ni sus hipótesis.

\noindent\textbf{Limitaciones}
\begin{itemize}
\item Modelo educativo; la aplicabilidad a una máquina real requiere revisión experta.
\end{itemize}

\noindent\textbf{Errores frecuentes}
\begin{itemize}
\item Introducir eficiencia total-estática en una expresión total-total.
\end{itemize}
% studium:end-record eyJpZCI6ICJQUkItMDAzOCIsICJzZWN0aW9uIjogImMwNSJ9
\end{studiumworked}

\begin{studiumworked}{Problema resuelto C05-04-R}


\noindent\textbf{Enunciado}
\(C05-04-R \cdot\) Descarga axial y potencia. Una turbina repetitiva de reacción dada descarga sin remolino. Calcule su carga y potencia útil. Datos: \(\mathcal R=0.52\); \(U=322\, \mathrm{m}\,\mathrm{s}^{-1}\); \(\dot m=26\, \mathrm{kg}\,\mathrm{s}^{-1}\); \(\eta_m=0.975\).

\noindent\textbf{Datos}
\begin{itemize}
\item $\mathcal R = 0.52$ — reacción
\item $U = 322\,\mathrm{m}\,\mathrm{s}^{-1}$ — velocidad periférica
\item $\dot m = 26\,\mathrm{kg}\,\mathrm{s}^{-1}$ — caudal
\item $\eta_m = 0.975$ — rendimiento mecánico
\end{itemize}

\noindent\textbf{Incógnita}
$\dot{W}$, en $\mathrm{W}$; también las magnitudes intermedias desarrolladas.

\noindent\textbf{Modelo físico}

\noindent\textit{Descarga axial y potencia}

\[
\fitmath{\dot{W}=2 \eta_m \dot m U^{2} \left(1 - \mathcal R\right)}
\]

\noindent\textbf{Hipótesis}
\begin{itemize}
\item Escalón repetitivo, radio y Cx constantes, eficiencia mecánica separada.
\end{itemize}

\noindent\textbf{Desarrollo paso a paso}

\begin{enumerate}

\item La descarga axial liga la carga a la reacción.

\[\fitmath{\psi=2\cdot \left(1-\mathrm{rea}\right)}\]

\[\fitmath{\psi=2\cdot \left(1-(0.52)\right)\approx 0.96\,}\]

\item El trabajo se obtiene de la carga.

\[\fitmath{\mathrm{wt}=2\cdot \left(1-\mathrm{rea}\right)\cdot {\mathrm{uu}}^{2}}\]

\[\fitmath{\mathrm{wt}=2\cdot \left(1-(0.52)\right)\cdot {(322)}^{2}\approx 99536.64\,\mathrm{J/kg}}\]

\item Se aplica una sola vez el rendimiento mecánico.

\[\fitmath{\mathrm{Pow}=\mathrm{em}\cdot \mathrm{md}\cdot 2\cdot \left(1-\mathrm{rea}\right)\cdot {\mathrm{uu}}^{2}}\]

\[\fitmath{\mathrm{Pow}=(0.975)\cdot (26)\cdot 2\cdot \left(1-(0.52)\right)\cdot {(322)}^{2}\approx 2523254\,\mathrm{W}}\]

\end{enumerate}


\noindent\textbf{Sustitución numérica}

\[
\fitmath{2 \left(0.975\right) \left(26\,\mathrm{kg}\,\mathrm{s}^{-1}\right) \left(\left(322\,\mathrm{m}\,\mathrm{s}^{-1}\right)\right)^{2} \left(1 - \left(0.52\right)\right)}
\]

\noindent\textbf{Resultado}

\[
\fitmath{\dot{W}=2.523253824\times10^{6}\,\mathrm{W}}
\]

\noindent\textbf{Interpretación}
La condición de descarga impide variar carga y reacción independientemente.

\noindent\textbf{Comprobación de los pasos}
\begin{itemize}
\item Paso 2: $0.96$ — reproducido
\item Paso 4: $99536.6$ — reproducido
\item Paso 6: $2.52325\times 10^{6}$ — reproducido
\end{itemize}

\noindent\textbf{Comprobación de la respuesta} la aritmética del último paso reproduce la respuesta enunciada a la precisión con que está escrita. Es una reproducción aritmética; no valida el modelo físico ni sus hipótesis.

\noindent\textbf{Limitaciones}
\begin{itemize}
\item Modelo educativo; la aplicabilidad a una máquina real requiere revisión experta.
\end{itemize}

\noindent\textbf{Errores frecuentes}
\begin{itemize}
\item Multiplicar por eficiencia mecánica dos veces.
\end{itemize}
% studium:end-record eyJpZCI6ICJQUkItMDAzOSIsICJzZWN0aW9uIjogImMwNSJ9
\end{studiumworked}

\begin{studiumworked}{Problema resuelto C05-04-P}


\noindent\textbf{Enunciado}
\(C05-04-P \cdot\) Descarga axial y potencia. Ejercicio de transferencia. Una turbina repetitiva de reacción dada descarga sin remolino. Calcule su carga y potencia útil. Antes de consultar la solución, prediga el efecto de cambiar el primer dato respecto del caso resuelto de esta familia. Justifique la tendencia y explique qué restricción podría convertirse en limitante. Datos: \(\mathcal R=0.565\); \(U=322\, \mathrm{m}\,\mathrm{s}^{-1}\); \(\dot m=26\, \mathrm{kg}\,\mathrm{s}^{-1}\); \(\eta_m=0.975\).

\noindent\textbf{Datos}
\begin{itemize}
\item $\mathcal R = 0.565$ — reacción
\item $U = 322\,\mathrm{m}\,\mathrm{s}^{-1}$ — velocidad periférica
\item $\dot m = 26\,\mathrm{kg}\,\mathrm{s}^{-1}$ — caudal
\item $\eta_m = 0.975$ — rendimiento mecánico
\end{itemize}

\noindent\textbf{Incógnita}
$\dot{W}$, en $\mathrm{W}$; también las magnitudes intermedias desarrolladas.

\noindent\textbf{Modelo físico}

\noindent\textit{Descarga axial y potencia}

\[
\fitmath{\dot{W}=2 \eta_m \dot m U^{2} \left(1 - \mathcal R\right)}
\]

\noindent\textbf{Hipótesis}
\begin{itemize}
\item Escalón repetitivo, radio y Cx constantes, eficiencia mecánica separada.
\end{itemize}

\noindent\textbf{Desarrollo paso a paso}

\begin{enumerate}

\item La descarga axial liga la carga a la reacción.

\[\fitmath{\mathrm{psi}=2\cdot \left(1-\mathrm{rea}\right)}\]

\[\fitmath{\mathrm{psi}=2\cdot \left(1-(0.565)\right)\approx 0.87\,}\]

\item El trabajo se obtiene de la carga.

\[\fitmath{\mathrm{wt}=2\cdot \left(1-\mathrm{rea}\right)\cdot {\mathrm{uu}}^{2}}\]

\[\fitmath{\mathrm{wt}=2\cdot \left(1-(0.565)\right)\cdot {(322)}^{2}\approx 90205.08\,\mathrm{J/kg}}\]

\item Se aplica una sola vez el rendimiento mecánico.

\[\fitmath{\mathrm{Pow}=\mathrm{em}\cdot \mathrm{md}\cdot 2\cdot \left(1-\mathrm{rea}\right)\cdot {\mathrm{uu}}^{2}}\]

\[\fitmath{\mathrm{Pow}=(0.975)\cdot (26)\cdot 2\cdot \left(1-(0.565)\right)\cdot {(322)}^{2}\approx 2286699\,\mathrm{W}}\]

\end{enumerate}


\noindent\textbf{Sustitución numérica}

\[
\fitmath{2 \left(0.975\right) \left(26\,\mathrm{kg}\,\mathrm{s}^{-1}\right) \left(\left(322\,\mathrm{m}\,\mathrm{s}^{-1}\right)\right)^{2} \left(1 - \left(0.565\right)\right)}
\]

\noindent\textbf{Resultado}

\[
\fitmath{\dot{W}=2.286698778\times10^{6}\,\mathrm{W}}
\]

\noindent\textbf{Interpretación}
La condición de descarga impide variar carga y reacción independientemente.

\noindent\textbf{Comprobación de los pasos}
\begin{itemize}
\item Paso 2: $0.87$ — reproducido
\item Paso 4: $90205.1$ — reproducido
\item Paso 6: $2.2867\times 10^{6}$ — reproducido
\end{itemize}

\noindent\textbf{Comprobación de la respuesta} la aritmética del último paso reproduce la respuesta enunciada a la precisión con que está escrita. Es una reproducción aritmética; no valida el modelo físico ni sus hipótesis.

\noindent\textbf{Limitaciones}
\begin{itemize}
\item Modelo educativo; la aplicabilidad a una máquina real requiere revisión experta.
\end{itemize}

\noindent\textbf{Errores frecuentes}
\begin{itemize}
\item Multiplicar por eficiencia mecánica dos veces.
\end{itemize}
% studium:end-record eyJpZCI6ICJQUkItMDA0MCIsICJzZWN0aW9uIjogImMwNSJ9
\end{studiumworked}

\chapter{Caudal, bloqueo sónico y diseño de turbinas}
\begin{studiumworked}{Problema resuelto C06-01-R}


\noindent\textbf{Enunciado}
\(C06-01-R \cdot\) Área crítica. Determine el gasto máximo por área y el área crítica necesaria para el caudal dado. Datos: \(p_{0e}=680000\, \mathrm{Pa}\); \(T_{0e}=1320\, \mathrm{K}\); \(\gamma=1.33\); \(R_g=287\, \mathrm{J}\,\mathrm{kg}^{-1}\,\mathrm{K}^{-1}\); \(\dot m=24\, \mathrm{kg}\,\mathrm{s}^{-1}\).

\noindent\textbf{Datos}
\begin{itemize}
\item $p_{0e} = 680000\,\mathrm{Pa}$ — presión total garganta de referencia
\item $T_{0e} = 1320\,\mathrm{K}$ — temperatura total
\item $\gamma = 1.33$ — razón de calores
\item $R_g = 287\,\mathrm{J}\,\mathrm{kg}^{-1}\,\mathrm{K}^{-1}$ — constante del gas
\item $\dot m = 24\,\mathrm{kg}\,\mathrm{s}^{-1}$ — caudal
\end{itemize}

\noindent\textbf{Incógnita}
$A^{*}$, en $\mathrm{m^{2}}$; también las magnitudes intermedias desarrolladas.

\noindent\textbf{Modelo físico}

\noindent\textit{Área crítica}

\[
\fitmath{A^{*}=\frac{\sqrt{T_{0e}} \dot m \left(\frac{2}{\gamma + 1}\right)^{- \frac{\gamma + 1}{2 \gamma - 2}}}{p_{0e} \sqrt{\frac{\gamma}{R_g}}}}
\]

\noindent\textbf{Hipótesis}
\begin{itemize}
\item Garganta ideal unidimensional isentrópica en la referencia indicada.
\end{itemize}

\noindent\textbf{Desarrollo paso a paso}

\begin{enumerate}

\item Se evalúa el cuadrado del factor crítico a M=1.

\[\fitmath{F=\frac{\gamma}{R_g}\cdot {\left(\frac{2}{\gamma+1}\right)}^{\frac{\gamma+1}{\gamma-1}}}\]

\[\fitmath{F=\frac{(1.33)}{(287)}\cdot {\left(\frac{2}{(1.33)+1}\right)}^{\frac{(1.33)+1}{(1.33)-1}}\approx 0.001576408\,\mathrm{s^{2}\,K/m^{2}}}\]

\item Se calcula el flujo másico máximo por área.

\[\fitmath{\mathrm{flux}=\frac{\mathrm{pin}}{\sqrt{\mathrm{Tin}}}\cdot \sqrt{\frac{\gamma}{R_g}}\cdot {\left(\frac{2}{\gamma+1}\right)}^{\frac{\gamma+1}{2\cdot \left(\gamma-1\right)}}}\]

\[\fitmath{\mathrm{flux}=\frac{(680000)}{\sqrt{(1320)}}\cdot \sqrt{\frac{(1.33)}{(287)}}\cdot {\left(\frac{2}{(1.33)+1}\right)}^{\frac{(1.33)+1}{2\cdot \left((1.33)-1\right)}}\approx 743.1153\,\mathrm{kg/(m^{2}\,s)}}\]

\item Se dimensiona el área normal efectiva.

\[\fitmath{\mathrm{Ast}=\frac{\mathrm{md}}{\frac{\mathrm{pin}}{\sqrt{\mathrm{Tin}}}\cdot \sqrt{\frac{\gamma}{R_g}}\cdot {\left(\frac{2}{\gamma+1}\right)}^{\frac{\gamma+1}{2\cdot \left(\gamma-1\right)}}}}\]

\[\fitmath{\mathrm{Ast}=\frac{(24)}{\frac{(680000)}{\sqrt{(1320)}}\cdot \sqrt{\frac{(1.33)}{(287)}}\cdot {\left(\frac{2}{(1.33)+1}\right)}^{\frac{(1.33)+1}{2\cdot \left((1.33)-1\right)}}}\approx 0.03229647\,\mathrm{m^{2}}}\]

\end{enumerate}


\noindent\textbf{Sustitución numérica}

\[
\fitmath{\frac{\sqrt{\left(1320\,\mathrm{K}\right)} \left(24\,\mathrm{kg}\,\mathrm{s}^{-1}\right) \left(\frac{2}{\left(1.33\right) + 1}\right)^{- \frac{\left(1.33\right) + 1}{2 \left(1.33\right) - 2}}}{\left(680000\,\mathrm{Pa}\right) \sqrt{\frac{\left(1.33\right)}{\left(287\,\mathrm{J}\,\mathrm{kg}^{-1}\,\mathrm{K}^{-1}\right)}}}}
\]

\noindent\textbf{Resultado}

\[
\fitmath{A^{*}=3.229646860\times10^{-2}\,\mathrm{m}^{2}}
\]

\noindent\textbf{Interpretación}
El área es normal al flujo y efectiva. Un coeficiente de descarga inferior a uno exigiría mayor área geométrica.

\noindent\textbf{Comprobación de los pasos}
\begin{itemize}
\item Paso 2: $0.00157641$ — reproducido
\item Paso 4: $743.115$ — reproducido
\item Paso 6: $0.0322965$ — reproducido
\end{itemize}

\noindent\textbf{Comprobación de la respuesta} la aritmética del último paso reproduce la respuesta enunciada a la precisión con que está escrita. Es una reproducción aritmética; no valida el modelo físico ni sus hipótesis.

\noindent\textbf{Limitaciones}
\begin{itemize}
\item Modelo educativo; la aplicabilidad a una máquina real requiere revisión experta.
\end{itemize}

\noindent\textbf{Errores frecuentes}
\begin{itemize}
\item Usar área frontal anular como garganta sin proyección.
\end{itemize}
% studium:end-record eyJpZCI6ICJQUkItMDA0MSIsICJzZWN0aW9uIjogImMwNiJ9
\end{studiumworked}

\begin{studiumworked}{Problema resuelto C06-01-P}


\noindent\textbf{Enunciado}
\(C06-01-P \cdot\) Área crítica. Ejercicio de transferencia. Determine el gasto máximo por área y el área crítica necesaria para el caudal dado. Antes de consultar la solución, prediga el efecto de cambiar el primer dato respecto del caso resuelto de esta familia. Justifique la tendencia y explique qué restricción podría convertirse en limitante. Datos: \(p_{0e}=724200\, \mathrm{Pa}\); \(T_{0e}=1320\, \mathrm{K}\); \(\gamma=1.33\); \(R_g=287\, \mathrm{J}\,\mathrm{kg}^{-1}\,\mathrm{K}^{-1}\); \(\dot m=24\, \mathrm{kg}\,\mathrm{s}^{-1}\).

\noindent\textbf{Datos}
\begin{itemize}
\item $p_{0e} = 724200\,\mathrm{Pa}$ — presión total garganta de referencia
\item $T_{0e} = 1320\,\mathrm{K}$ — temperatura total
\item $\gamma = 1.33$ — razón de calores
\item $R_g = 287\,\mathrm{J}\,\mathrm{kg}^{-1}\,\mathrm{K}^{-1}$ — constante del gas
\item $\dot m = 24\,\mathrm{kg}\,\mathrm{s}^{-1}$ — caudal
\end{itemize}

\noindent\textbf{Incógnita}
$A^{*}$, en $\mathrm{m^{2}}$; también las magnitudes intermedias desarrolladas.

\noindent\textbf{Modelo físico}

\noindent\textit{Área crítica}

\[
\fitmath{A^{*}=\frac{\sqrt{T_{0e}} \dot m \left(\frac{2}{\gamma + 1}\right)^{- \frac{\gamma + 1}{2 \gamma - 2}}}{p_{0e} \sqrt{\frac{\gamma}{R_g}}}}
\]

\noindent\textbf{Hipótesis}
\begin{itemize}
\item Garganta ideal unidimensional isentrópica en la referencia indicada.
\end{itemize}

\noindent\textbf{Desarrollo paso a paso}

\begin{enumerate}

\item Se evalúa el cuadrado del factor crítico a M=1.

\[\fitmath{F=\frac{\gamma}{R_g}\cdot {\left(\frac{2}{\gamma+1}\right)}^{\frac{\gamma+1}{\gamma-1}}}\]

\[\fitmath{F=\frac{(1.33)}{(287)}\cdot {\left(\frac{2}{(1.33)+1}\right)}^{\frac{(1.33)+1}{(1.33)-1}}\approx 0.001576408\,\mathrm{s^{2}\,K/m^{2}}}\]

\item Se calcula el flujo másico máximo por área.

\[\fitmath{\mathrm{flux}=\frac{\mathrm{pin}}{\sqrt{\mathrm{Tin}}}\cdot \sqrt{\frac{\gamma}{R_g}}\cdot {\left(\frac{2}{\gamma+1}\right)}^{\frac{\gamma+1}{2\cdot \left(\gamma-1\right)}}}\]

\[\fitmath{\mathrm{flux}=\frac{(724200)}{\sqrt{(1320)}}\cdot \sqrt{\frac{(1.33)}{(287)}}\cdot {\left(\frac{2}{(1.33)+1}\right)}^{\frac{(1.33)+1}{2\cdot \left((1.33)-1\right)}}\approx 791.4178\,\mathrm{kg/(m^{2}\,s)}}\]

\item Se dimensiona el área normal efectiva.

\[\fitmath{\mathrm{Ast}=\frac{\mathrm{md}}{\frac{\mathrm{pin}}{\sqrt{\mathrm{Tin}}}\cdot \sqrt{\frac{\gamma}{R_g}}\cdot {\left(\frac{2}{\gamma+1}\right)}^{\frac{\gamma+1}{2\cdot \left(\gamma-1\right)}}}}\]

\[\fitmath{\mathrm{Ast}=\frac{(24)}{\frac{(724200)}{\sqrt{(1320)}}\cdot \sqrt{\frac{(1.33)}{(287)}}\cdot {\left(\frac{2}{(1.33)+1}\right)}^{\frac{(1.33)+1}{2\cdot \left((1.33)-1\right)}}}\approx 0.03032532\,\mathrm{m^{2}}}\]

\end{enumerate}


\noindent\textbf{Sustitución numérica}

\[
\fitmath{\frac{\sqrt{\left(1320\,\mathrm{K}\right)} \left(24\,\mathrm{kg}\,\mathrm{s}^{-1}\right) \left(\frac{2}{\left(1.33\right) + 1}\right)^{- \frac{\left(1.33\right) + 1}{2 \left(1.33\right) - 2}}}{\left(724200\,\mathrm{Pa}\right) \sqrt{\frac{\left(1.33\right)}{\left(287\,\mathrm{J}\,\mathrm{kg}^{-1}\,\mathrm{K}^{-1}\right)}}}}
\]

\noindent\textbf{Resultado}

\[
\fitmath{A^{*}=3.032532263\times10^{-2}\,\mathrm{m}^{2}}
\]

\noindent\textbf{Interpretación}
El área es normal al flujo y efectiva. Un coeficiente de descarga inferior a uno exigiría mayor área geométrica.

\noindent\textbf{Comprobación de los pasos}
\begin{itemize}
\item Paso 2: $0.00157641$ — reproducido
\item Paso 4: $791.418$ — reproducido
\item Paso 6: $0.0303253$ — reproducido
\end{itemize}

\noindent\textbf{Comprobación de la respuesta} la aritmética del último paso reproduce la respuesta enunciada a la precisión con que está escrita. Es una reproducción aritmética; no valida el modelo físico ni sus hipótesis.

\noindent\textbf{Limitaciones}
\begin{itemize}
\item Modelo educativo; la aplicabilidad a una máquina real requiere revisión experta.
\end{itemize}

\noindent\textbf{Errores frecuentes}
\begin{itemize}
\item Usar área frontal anular como garganta sin proyección.
\end{itemize}
% studium:end-record eyJpZCI6ICJQUkItMDA0MiIsICJzZWN0aW9uIjogImMwNiJ9
\end{studiumworked}

\begin{studiumworked}{Problema resuelto C06-02-R}


\noindent\textbf{Enunciado}
\(C06-02-R \cdot\) Cociente crítico. Calcule la presión crítica estática de una tobera convergente ideal. Datos: \(p_{0e}=915000\, \mathrm{Pa}\); \(\gamma=1.32\).

\noindent\textbf{Datos}
\begin{itemize}
\item $p_{0e} = 915000\,\mathrm{Pa}$ — presión total depósito
\item $\gamma = 1.32$ — razón de calores
\end{itemize}

\noindent\textbf{Incógnita}
$p^{*}$, en $\mathrm{Pa}$; también las magnitudes intermedias desarrolladas.

\noindent\textbf{Modelo físico}

\noindent\textit{Cociente crítico}

\[
\fitmath{p^{*}=p_{0e} \left(\frac{2}{\gamma + 1}\right)^{\frac{\gamma}{\gamma - 1}}}
\]

\noindent\textbf{Hipótesis}
\begin{itemize}
\item Depósito y tobera convergente ideal; gas perfecto.
\end{itemize}

\noindent\textbf{Desarrollo paso a paso}

\begin{enumerate}

\item Se aplica la razón isentrópica a Mach uno.

\[\fitmath{\mathrm{ratio}={\left(\frac{2}{\gamma+1}\right)}^{\frac{\gamma}{\gamma-1}}}\]

\[\fitmath{\mathrm{ratio}={\left(\frac{2}{(1.32)+1}\right)}^{\frac{(1.32)}{(1.32)-1}}\approx 0.5421392\,}\]

\item Se obtiene la presión de garganta crítica.

\[\fitmath{\mathrm{pstar}=\mathrm{pin}\cdot {\left(\frac{2}{\gamma+1}\right)}^{\frac{\gamma}{\gamma-1}}}\]

\[\fitmath{\mathrm{pstar}=(915000)\cdot {\left(\frac{2}{(1.32)+1}\right)}^{\frac{(1.32)}{(1.32)-1}}\approx 496057.3\,\mathrm{Pa}}\]

\end{enumerate}


\noindent\textbf{Sustitución numérica}

\[
\fitmath{\left(915000\,\mathrm{Pa}\right) \left(\frac{2}{\left(1.32\right) + 1}\right)^{\frac{\left(1.32\right)}{\left(1.32\right) - 1}}}
\]

\noindent\textbf{Resultado}

\[
\fitmath{p^{*}=4.960573496\times10^{5}\,\mathrm{Pa}}
\]

\noindent\textbf{Interpretación}
Una contrapresión inferior permite bloqueo en la tobera ideal indicada; no define toda la expansión de una cascada real.

\noindent\textbf{Comprobación de los pasos}
\begin{itemize}
\item Paso 2: $0.542139$ — reproducido
\item Paso 4: $496057$ — reproducido
\end{itemize}

\noindent\textbf{Comprobación de la respuesta} la aritmética del último paso reproduce la respuesta enunciada a la precisión con que está escrita. Es una reproducción aritmética; no valida el modelo físico ni sus hipótesis.

\noindent\textbf{Limitaciones}
\begin{itemize}
\item Modelo educativo; la aplicabilidad a una máquina real requiere revisión experta.
\end{itemize}

\noindent\textbf{Errores frecuentes}
\begin{itemize}
\item Interpretar la razón crítica como pérdida de presión total.
\end{itemize}
% studium:end-record eyJpZCI6ICJQUkItMDA0MyIsICJzZWN0aW9uIjogImMwNiJ9
\end{studiumworked}

\begin{studiumworked}{Problema resuelto C06-02-P}


\noindent\textbf{Enunciado}
\(C06-02-P \cdot\) Cociente crítico. Ejercicio de transferencia. Calcule la presión crítica estática de una tobera convergente ideal. Antes de consultar la solución, prediga el efecto de cambiar el primer dato respecto del caso resuelto de esta familia. Justifique la tendencia y explique qué restricción podría convertirse en limitante. Datos: \(p_{0e}=974475\, \mathrm{Pa}\); \(\gamma=1.32\).

\noindent\textbf{Datos}
\begin{itemize}
\item $p_{0e} = 974475\,\mathrm{Pa}$ — presión total depósito
\item $\gamma = 1.32$ — razón de calores
\end{itemize}

\noindent\textbf{Incógnita}
$p^{*}$, en $\mathrm{Pa}$; también las magnitudes intermedias desarrolladas.

\noindent\textbf{Modelo físico}

\noindent\textit{Cociente crítico}

\[
\fitmath{p^{*}=p_{0e} \left(\frac{2}{\gamma + 1}\right)^{\frac{\gamma}{\gamma - 1}}}
\]

\noindent\textbf{Hipótesis}
\begin{itemize}
\item Depósito y tobera convergente ideal; gas perfecto.
\end{itemize}

\noindent\textbf{Desarrollo paso a paso}

\begin{enumerate}

\item Se aplica la razón isentrópica a Mach uno.

\[\fitmath{\mathrm{ratio}={\left(\frac{2}{\gamma+1}\right)}^{\frac{\gamma}{\gamma-1}}}\]

\[\fitmath{\mathrm{ratio}={\left(\frac{2}{(1.32)+1}\right)}^{\frac{(1.32)}{(1.32)-1}}\approx 0.5421392\,}\]

\item Se obtiene la presión de garganta crítica.

\[\fitmath{\mathrm{pstar}=\mathrm{pin}\cdot {\left(\frac{2}{\gamma+1}\right)}^{\frac{\gamma}{\gamma-1}}}\]

\[\fitmath{\mathrm{pstar}=(974475)\cdot {\left(\frac{2}{(1.32)+1}\right)}^{\frac{(1.32)}{(1.32)-1}}\approx 528301.1\,\mathrm{Pa}}\]

\end{enumerate}


\noindent\textbf{Sustitución numérica}

\[
\fitmath{\left(974475\,\mathrm{Pa}\right) \left(\frac{2}{\left(1.32\right) + 1}\right)^{\frac{\left(1.32\right)}{\left(1.32\right) - 1}}}
\]

\noindent\textbf{Resultado}

\[
\fitmath{p^{*}=5.283010773\times10^{5}\,\mathrm{Pa}}
\]

\noindent\textbf{Interpretación}
Una contrapresión inferior permite bloqueo en la tobera ideal indicada; no define toda la expansión de una cascada real.

\noindent\textbf{Comprobación de los pasos}
\begin{itemize}
\item Paso 2: $0.542139$ — reproducido
\item Paso 4: $528301$ — reproducido
\end{itemize}

\noindent\textbf{Comprobación de la respuesta} la aritmética del último paso reproduce la respuesta enunciada a la precisión con que está escrita. Es una reproducción aritmética; no valida el modelo físico ni sus hipótesis.

\noindent\textbf{Limitaciones}
\begin{itemize}
\item Modelo educativo; la aplicabilidad a una máquina real requiere revisión experta.
\end{itemize}

\noindent\textbf{Errores frecuentes}
\begin{itemize}
\item Interpretar la razón crítica como pérdida de presión total.
\end{itemize}
% studium:end-record eyJpZCI6ICJQUkItMDA0NCIsICJzZWN0aW9uIjogImMwNiJ9
\end{studiumworked}

\begin{studiumworked}{Problema resuelto C06-03-R}


\noindent\textbf{Enunciado}
\(C06-03-R \cdot\) Reparto entero de escalones. Se ha elegido un número entero de escalones de igual trabajo. Calcule la carga real que cierra exactamente el trabajo global. Datos: \(w_t=384000\, \mathrm{J}\,\mathrm{kg}^{-1}\); \(N_{\mathrm{esc}}=3\); \(U=338\, \mathrm{m}\,\mathrm{s}^{-1}\).

\noindent\textbf{Datos}
\begin{itemize}
\item $w_t = 384000\,\mathrm{J}\,\mathrm{kg}^{-1}$ — trabajo total
\item $N_{\mathrm{esc}} = 3$ — número de escalones
\item $U = 338\,\mathrm{m}\,\mathrm{s}^{-1}$ — velocidad periférica común de referencia
\end{itemize}

\noindent\textbf{Incógnita}
$\psi$ (adimensional); también las magnitudes intermedias desarrolladas.

\noindent\textbf{Modelo físico}

\noindent\textit{Reparto entero de escalones}

\[
\fitmath{\psi=\frac{w_t}{N_{\mathrm{esc}} U^{2}}}
\]

\noindent\textbf{Hipótesis}
\begin{itemize}
\item Trabajo igual por escalón y misma U de referencia.
\end{itemize}

\noindent\textbf{Desarrollo paso a paso}

\begin{enumerate}

\item El reparto conserva la suma del trabajo.

\[\fitmath{\mathrm{we}=\frac{\mathrm{wt}}{\mathrm{nn}}}\]

\[\fitmath{\mathrm{we}=\frac{(384000)}{(3)}\approx 128000\,\mathrm{J/kg}}\]

\item Se calcula la carga efectiva tras elegir el entero.

\[\fitmath{\mathrm{psi}=\frac{\mathrm{wt}}{\mathrm{nn}\cdot {\mathrm{uu}}^{2}}}\]

\[\fitmath{\mathrm{psi}=\frac{(384000)}{(3)\cdot {(338)}^{2}}\approx 1.120409\,}\]

\end{enumerate}


\noindent\textbf{Sustitución numérica}

\[
\fitmath{\frac{\left(384000\,\mathrm{J}\,\mathrm{kg}^{-1}\right)}{\left(3\right) \left(\left(338\,\mathrm{m}\,\mathrm{s}^{-1}\right)\right)^{2}}}
\]

\noindent\textbf{Resultado}

\[
\fitmath{\psi=1.120408949\times10^{0}}
\]

\noindent\textbf{Interpretación}
El entero elegido obliga a recalcular psi. La factibilidad aerodinámica se revisa después.

\noindent\textbf{Comprobación de los pasos}
\begin{itemize}
\item Paso 2: $128000$ — reproducido
\item Paso 4: $1.12041$ — reproducido
\end{itemize}

\noindent\textbf{Comprobación de la respuesta} la aritmética del último paso reproduce la respuesta enunciada a la precisión con que está escrita. Es una reproducción aritmética; no valida el modelo físico ni sus hipótesis.

\noindent\textbf{Limitaciones}
\begin{itemize}
\item Modelo educativo; la aplicabilidad a una máquina real requiere revisión experta.
\end{itemize}

\noindent\textbf{Errores frecuentes}
\begin{itemize}
\item Redondear n y conservar una carga previa incompatible.
\end{itemize}
% studium:end-record eyJpZCI6ICJQUkItMDA0NSIsICJzZWN0aW9uIjogImMwNiJ9
\end{studiumworked}

\begin{studiumworked}{Problema resuelto C06-03-P}


\noindent\textbf{Enunciado}
\(C06-03-P \cdot\) Reparto entero de escalones. Ejercicio de transferencia. Se ha elegido un número entero de escalones de igual trabajo. Calcule la carga real que cierra exactamente el trabajo global. Antes de consultar la solución, prediga el efecto de cambiar el primer dato respecto del caso resuelto de esta familia. Justifique la tendencia y explique qué restricción podría convertirse en limitante. Datos: \(w_t=408960\, \mathrm{J}\,\mathrm{kg}^{-1}\); \(N_{\mathrm{esc}}=3\); \(U=338\, \mathrm{m}\,\mathrm{s}^{-1}\).

\noindent\textbf{Datos}
\begin{itemize}
\item $w_t = 408960\,\mathrm{J}\,\mathrm{kg}^{-1}$ — trabajo total
\item $N_{\mathrm{esc}} = 3$ — número de escalones
\item $U = 338\,\mathrm{m}\,\mathrm{s}^{-1}$ — velocidad periférica común de referencia
\end{itemize}

\noindent\textbf{Incógnita}
$\psi$ (adimensional); también las magnitudes intermedias desarrolladas.

\noindent\textbf{Modelo físico}

\noindent\textit{Reparto entero de escalones}

\[
\fitmath{\psi=\frac{w_t}{N_{\mathrm{esc}} U^{2}}}
\]

\noindent\textbf{Hipótesis}
\begin{itemize}
\item Trabajo igual por escalón y misma U de referencia.
\end{itemize}

\noindent\textbf{Desarrollo paso a paso}

\begin{enumerate}

\item El reparto conserva la suma del trabajo.

\[\fitmath{\mathrm{we}=\frac{\mathrm{wt}}{\mathrm{nn}}}\]

\[\fitmath{\mathrm{we}=\frac{(408960)}{(3)}\approx 136320\,\mathrm{J/kg}}\]

\item Se calcula la carga efectiva tras elegir el entero.

\[\fitmath{\psi=\frac{\mathrm{wt}}{\mathrm{nn}\cdot {\mathrm{uu}}^{2}}}\]

\[\fitmath{\psi=\frac{(408960)}{(3)\cdot {(338)}^{2}}\approx 1.193236\,}\]

\end{enumerate}


\noindent\textbf{Sustitución numérica}

\[
\fitmath{\frac{\left(408960\,\mathrm{J}\,\mathrm{kg}^{-1}\right)}{\left(3\right) \left(\left(338\,\mathrm{m}\,\mathrm{s}^{-1}\right)\right)^{2}}}
\]

\noindent\textbf{Resultado}

\[
\fitmath{\psi=1.193235531\times10^{0}}
\]

\noindent\textbf{Interpretación}
El entero elegido obliga a recalcular psi. La factibilidad aerodinámica se revisa después.

\noindent\textbf{Comprobación de los pasos}
\begin{itemize}
\item Paso 2: $136320$ — reproducido
\item Paso 4: $1.19324$ — reproducido
\end{itemize}

\noindent\textbf{Comprobación de la respuesta} la aritmética del último paso reproduce la respuesta enunciada a la precisión con que está escrita. Es una reproducción aritmética; no valida el modelo físico ni sus hipótesis.

\noindent\textbf{Limitaciones}
\begin{itemize}
\item Modelo educativo; la aplicabilidad a una máquina real requiere revisión experta.
\end{itemize}

\noindent\textbf{Errores frecuentes}
\begin{itemize}
\item Redondear n y conservar una carga previa incompatible.
\end{itemize}
% studium:end-record eyJpZCI6ICJQUkItMDA0NiIsICJzZWN0aW9uIjogImMwNiJ9
\end{studiumworked}

\begin{studiumworked}{Problema resuelto C06-04-R}


\noindent\textbf{Enunciado}
\(C06-04-R \cdot\) Caudal de garganta oblicua. Una sección axial opera con módulo de velocidad sónico y ángulo dado. Halle el caudal usando su proyección axial. Datos: \(\rho=1.92\, \mathrm{kg}\,\mathrm{m}^{-3}\); \(C=626\, \mathrm{m}\,\mathrm{s}^{-1}\); \(\alpha_e=1.064651\, \mathrm{rad}\); \(A=0.045\, \mathrm{m}^{2}\).

\noindent\textbf{Datos}
\begin{itemize}
\item $\rho = 1.92\,\mathrm{kg}\,\mathrm{m}^{-3}$ — densidad estática
\item $C = 626\,\mathrm{m}\,\mathrm{s}^{-1}$ — módulo sónico local
\item $\alpha_e = 1.0646508437165\,\mathrm{rad}$ — ángulo con el eje
\item $A = 0.045\,\mathrm{m}^{2}$ — área anular frontal
\end{itemize}

\noindent\textbf{Incógnita}
$\dot{m}$, en $\mathrm{kg\,s^{-1}}$; también las magnitudes intermedias desarrolladas.

\noindent\textbf{Modelo físico}

\noindent\textit{Caudal de garganta oblicua}

\[
\fitmath{\dot m=A \rho C \cos{\left(\alpha_e \right)}}
\]

\noindent\textbf{Hipótesis}
\begin{itemize}
\item Corte axial con densidad y velocidad uniformes.
\end{itemize}

\noindent\textbf{Desarrollo paso a paso}

\begin{enumerate}

\item Se proyecta el módulo sobre la normal axial.

\[\fitmath{\mathrm{cx}=\mathrm{vel}\cdot \cos\left(\mathrm{ang}\right)}\]

\[\fitmath{\mathrm{cx}=(626)\cdot \cos\left((1.064650844)\right)\approx 303.4908\,\mathrm{m/s}}\]

\item La continuidad usa esa componente.

\[\fitmath{\mathrm{md}=\rho\cdot \mathrm{vel}\cdot \cos\left(\mathrm{ang}\right)\cdot \mathrm{area}}\]

\[\fitmath{\mathrm{md}=(1.92)\cdot (626)\cdot \cos\left((1.064650844)\right)\cdot (0.045)\approx 26.22161\,\mathrm{kg/s}}\]

\end{enumerate}


\noindent\textbf{Sustitución numérica}

\[
\fitmath{\left(0.045\,\mathrm{m}^{2}\right) \left(1.92\,\mathrm{kg}\,\mathrm{m}^{-3}\right) \left(626\,\mathrm{m}\,\mathrm{s}^{-1}\right) \cos{\left(\left(1.0646508437165\,\mathrm{rad}\right) \right)}}
\]

\noindent\textbf{Resultado}

\[
\fitmath{\dot m=2.622160704\times10^{1}\,\mathrm{kg}\,\mathrm{s}^{-1}}
\]

\noindent\textbf{Interpretación}
Mach axial y Mach del módulo son distintos. La garganta real puede tener otra orientación y área.

\noindent\textbf{Comprobación de los pasos}
\begin{itemize}
\item Paso 2: $303.491$ — reproducido
\item Paso 4: $26.2216$ — reproducido
\end{itemize}

\noindent\textbf{Comprobación de la respuesta} la aritmética del último paso reproduce la respuesta enunciada a la precisión con que está escrita. Es una reproducción aritmética; no valida el modelo físico ni sus hipótesis.

\noindent\textbf{Limitaciones}
\begin{itemize}
\item Modelo educativo; la aplicabilidad a una máquina real requiere revisión experta.
\end{itemize}

\noindent\textbf{Errores frecuentes}
\begin{itemize}
\item Omitir el coseno por estar el módulo a Mach uno.
\end{itemize}
% studium:end-record eyJpZCI6ICJQUkItMDA0NyIsICJzZWN0aW9uIjogImMwNiJ9
\end{studiumworked}

\begin{studiumworked}{Problema resuelto C06-04-P}


\noindent\textbf{Enunciado}
\(C06-04-P \cdot\) Caudal de garganta oblicua. Ejercicio de transferencia. Una sección axial opera con módulo de velocidad sónico y ángulo dado. Halle el caudal usando su proyección axial. Antes de consultar la solución, prediga el efecto de cambiar el primer dato respecto del caso resuelto de esta familia. Justifique la tendencia y explique qué restricción podría convertirse en limitante. Datos: \(\rho=2.0448\, \mathrm{kg}\,\mathrm{m}^{-3}\); \(C=626\, \mathrm{m}\,\mathrm{s}^{-1}\); \(\alpha_e=1.064651\, \mathrm{rad}\); \(A=0.045\, \mathrm{m}^{2}\).

\noindent\textbf{Datos}
\begin{itemize}
\item $\rho = 2.0448\,\mathrm{kg}\,\mathrm{m}^{-3}$ — densidad estática
\item $C = 626\,\mathrm{m}\,\mathrm{s}^{-1}$ — módulo sónico local
\item $\alpha_e = 1.0646508437165\,\mathrm{rad}$ — ángulo con el eje
\item $A = 0.045\,\mathrm{m}^{2}$ — área anular frontal
\end{itemize}

\noindent\textbf{Incógnita}
$\dot{m}$, en $\mathrm{kg\,s^{-1}}$; también las magnitudes intermedias desarrolladas.

\noindent\textbf{Modelo físico}

\noindent\textit{Caudal de garganta oblicua}

\[
\fitmath{\dot m=A \rho C \cos{\left(\alpha_e \right)}}
\]

\noindent\textbf{Hipótesis}
\begin{itemize}
\item Corte axial con densidad y velocidad uniformes.
\end{itemize}

\noindent\textbf{Desarrollo paso a paso}

\begin{enumerate}

\item Se proyecta el módulo sobre la normal axial.

\[\fitmath{\mathrm{cx}=\mathrm{vel}\cdot \cos\left(\mathrm{ang}\right)}\]

\[\fitmath{\mathrm{cx}=(626)\cdot \cos\left((1.064650844)\right)\approx 303.4908\,\mathrm{m/s}}\]

\item La continuidad usa esa componente.

\[\fitmath{\mathrm{md}=\rho\cdot \mathrm{vel}\cdot \cos\left(\mathrm{ang}\right)\cdot \mathrm{area}}\]

\[\fitmath{\mathrm{md}=(2.0448)\cdot (626)\cdot \cos\left((1.064650844)\right)\cdot (0.045)\approx 27.92601\,\mathrm{kg/s}}\]

\end{enumerate}


\noindent\textbf{Sustitución numérica}

\[
\fitmath{\left(0.045\,\mathrm{m}^{2}\right) \left(2.0448\,\mathrm{kg}\,\mathrm{m}^{-3}\right) \left(626\,\mathrm{m}\,\mathrm{s}^{-1}\right) \cos{\left(\left(1.0646508437165\,\mathrm{rad}\right) \right)}}
\]

\noindent\textbf{Resultado}

\[
\fitmath{\dot m=2.792601150\times10^{1}\,\mathrm{kg}\,\mathrm{s}^{-1}}
\]

\noindent\textbf{Interpretación}
Mach axial y Mach del módulo son distintos. La garganta real puede tener otra orientación y área.

\noindent\textbf{Comprobación de los pasos}
\begin{itemize}
\item Paso 2: $303.491$ — reproducido
\item Paso 4: $27.926$ — reproducido
\end{itemize}

\noindent\textbf{Comprobación de la respuesta} la aritmética del último paso reproduce la respuesta enunciada a la precisión con que está escrita. Es una reproducción aritmética; no valida el modelo físico ni sus hipótesis.

\noindent\textbf{Limitaciones}
\begin{itemize}
\item Modelo educativo; la aplicabilidad a una máquina real requiere revisión experta.
\end{itemize}

\noindent\textbf{Errores frecuentes}
\begin{itemize}
\item Omitir el coseno por estar el módulo a Mach uno.
\end{itemize}
% studium:end-record eyJpZCI6ICJQUkItMDA0OCIsICJzZWN0aW9uIjogImMwNiJ9
\end{studiumworked}

\chapter{Escalón de compresor axial al radio medio}
\begin{studiumworked}{Problema resuelto C07-01-R}


\noindent\textbf{Enunciado}
\(C07-01-R \cdot\) Trabajo y presión del escalón. Dados los parámetros del escalón compresor, calcule trabajo, incremento total de temperatura y relación de presión. Datos: \(U=273\, \mathrm{m}\,\mathrm{s}^{-1}\); \(\psi=0.39\); \(T_{0e}=301\, \mathrm{K}\); \(\eta=0.88\); \(c_p=1004.5\, \mathrm{J}\,\mathrm{kg}^{-1}\,\mathrm{K}^{-1}\); \(\gamma=1.4\).

\noindent\textbf{Datos}
\begin{itemize}
\item $U = 273\,\mathrm{m}\,\mathrm{s}^{-1}$ — velocidad periférica
\item $\psi = 0.39$ — carga
\item $T_{0e} = 301\,\mathrm{K}$ — temperatura total entrada
\item $\eta = 0.88$ — eficiencia total-total
\item $c_p = 1004.5\,\mathrm{J}\,\mathrm{kg}^{-1}\,\mathrm{K}^{-1}$ — calor específico
\item $\gamma = 1.4$ — razón de calores
\end{itemize}

\noindent\textbf{Incógnita}
$\pi_c$ (adimensional); también las magnitudes intermedias desarrolladas.

\noindent\textbf{Modelo físico}

\noindent\textit{Trabajo y presión del escalón}

\[
\fitmath{\pi_c=\left(1 + \frac{\eta \psi U^{2}}{T_{0e} c_p}\right)^{\frac{\gamma}{\gamma - 1}}}
\]

\noindent\textbf{Hipótesis}
\begin{itemize}
\item Escalón adiabático; eficiencia global y gas perfecto constantes.
\end{itemize}

\noindent\textbf{Desarrollo paso a paso}

\begin{enumerate}

\item La carga define el trabajo comunicado.

\[\fitmath{\mathrm{wc}=\psi\cdot {\mathrm{uu}}^{2}}\]

\[\fitmath{\mathrm{wc}=(0.39)\cdot {(273)}^{2}\approx 29066.31\,\mathrm{J/kg}}\]

\item El balance energético fija el incremento térmico.

\[\fitmath{\mathrm{dT}=\frac{\psi\cdot {\mathrm{uu}}^{2}}{\mathrm{cp}}}\]

\[\fitmath{\mathrm{dT}=\frac{(0.39)\cdot {(273)}^{2}}{(1004.5)}\approx 28.9361\,\mathrm{K}}\]

\item El trabajo ideal equivalente fija la relación de presión.

\[\fitmath{\mathrm{pr}={\left(1+\frac{\eta\cdot \psi\cdot {\mathrm{uu}}^{2}}{\mathrm{cp}\cdot \mathrm{Tin}}\right)}^{\frac{\gamma}{\gamma-1}}}\]

\[\fitmath{\mathrm{pr}={\left(1+\frac{(0.88)\cdot (0.39)\cdot {(273)}^{2}}{(1004.5)\cdot (301)}\right)}^{\frac{(1.4)}{(1.4)-1}}\approx 1.328739\,}\]

\end{enumerate}


\noindent\textbf{Sustitución numérica}

\[
\fitmath{\left(1 + \frac{\left(0.88\right) \left(0.39\right) \left(\left(273\,\mathrm{m}\,\mathrm{s}^{-1}\right)\right)^{2}}{\left(301\,\mathrm{K}\right) \left(1004.5\,\mathrm{J}\,\mathrm{kg}^{-1}\,\mathrm{K}^{-1}\right)}\right)^{\frac{\left(1.4\right)}{\left(1.4\right) - 1}}}
\]

\noindent\textbf{Resultado}

\[
\fitmath{\pi_c=1.328739105\times10^{0}}
\]

\noindent\textbf{Interpretación}
La relación debe ser mayor que uno. Compare el trabajo reconstruido desde la temperatura con Euler.

\noindent\textbf{Comprobación de los pasos}
\begin{itemize}
\item Paso 2: $29066.3$ — reproducido
\item Paso 4: $28.9361$ — reproducido
\item Paso 6: $1.32874$ — reproducido
\end{itemize}

\noindent\textbf{Comprobación de la respuesta} la aritmética del último paso reproduce la respuesta enunciada a la precisión con que está escrita. Es una reproducción aritmética; no valida el modelo físico ni sus hipótesis.

\noindent\textbf{Limitaciones}
\begin{itemize}
\item Modelo educativo; la aplicabilidad a una máquina real requiere revisión experta.
\end{itemize}

\noindent\textbf{Errores frecuentes}
\begin{itemize}
\item Aplicar la razón de temperaturas reales como si fuera isentrópica.
\end{itemize}
% studium:end-record eyJpZCI6ICJQUkItMDA0OSIsICJzZWN0aW9uIjogImMwNyJ9
\end{studiumworked}

\begin{studiumworked}{Problema resuelto C07-01-P}


\noindent\textbf{Enunciado}
\(C07-01-P \cdot\) Trabajo y presión del escalón. Ejercicio de transferencia. Dados los parámetros del escalón compresor, calcule trabajo, incremento total de temperatura y relación de presión. Antes de consultar la solución, prediga el efecto de cambiar el primer dato respecto del caso resuelto de esta familia. Justifique la tendencia y explique qué restricción podría convertirse en limitante. Datos: \(U=290.745\, \mathrm{m}\,\mathrm{s}^{-1}\); \(\psi=0.39\); \(T_{0e}=301\, \mathrm{K}\); \(\eta=0.88\); \(c_p=1004.5\, \mathrm{J}\,\mathrm{kg}^{-1}\,\mathrm{K}^{-1}\); \(\gamma=1.4\).

\noindent\textbf{Datos}
\begin{itemize}
\item $U = 290.745\,\mathrm{m}\,\mathrm{s}^{-1}$ — velocidad periférica
\item $\psi = 0.39$ — carga
\item $T_{0e} = 301\,\mathrm{K}$ — temperatura total entrada
\item $\eta = 0.88$ — eficiencia total-total
\item $c_p = 1004.5\,\mathrm{J}\,\mathrm{kg}^{-1}\,\mathrm{K}^{-1}$ — calor específico
\item $\gamma = 1.4$ — razón de calores
\end{itemize}

\noindent\textbf{Incógnita}
$\pi_c$ (adimensional); también las magnitudes intermedias desarrolladas.

\noindent\textbf{Modelo físico}

\noindent\textit{Trabajo y presión del escalón}

\[
\fitmath{\pi_c=\left(1 + \frac{\eta \psi U^{2}}{T_{0e} c_p}\right)^{\frac{\gamma}{\gamma - 1}}}
\]

\noindent\textbf{Hipótesis}
\begin{itemize}
\item Escalón adiabático; eficiencia global y gas perfecto constantes.
\end{itemize}

\noindent\textbf{Desarrollo paso a paso}

\begin{enumerate}

\item La carga define el trabajo comunicado.

\[\fitmath{\mathrm{wc}=\mathrm{psi}\cdot {\mathrm{uu}}^{2}}\]

\[\fitmath{\mathrm{wc}=(0.39)\cdot {(290.745)}^{2}\approx 32967.74\,\mathrm{J/kg}}\]

\item El balance energético fija el incremento térmico.

\[\fitmath{\mathrm{dT}=\frac{\mathrm{psi}\cdot {\mathrm{uu}}^{2}}{\mathrm{cp}}}\]

\[\fitmath{\mathrm{dT}=\frac{(0.39)\cdot {(290.745)}^{2}}{(1004.5)}\approx 32.82005\,\mathrm{K}}\]

\item El trabajo ideal equivalente fija la relación de presión.

\[\fitmath{\mathrm{pr}={\left(1+\frac{\mathrm{eta}\cdot \mathrm{psi}\cdot {\mathrm{uu}}^{2}}{\mathrm{cp}\cdot \mathrm{Tin}}\right)}^{\frac{\gamma}{\gamma-1}}}\]

\[\fitmath{\mathrm{pr}={\left(1+\frac{(0.88)\cdot (0.39)\cdot {(290.745)}^{2}}{(1004.5)\cdot (301)}\right)}^{\frac{(1.4)}{(1.4)-1}}\approx 1.378068\,}\]

\end{enumerate}


\noindent\textbf{Sustitución numérica}

\[
\fitmath{\left(1 + \frac{\left(0.88\right) \left(0.39\right) \left(\left(290.745\,\mathrm{m}\,\mathrm{s}^{-1}\right)\right)^{2}}{\left(301\,\mathrm{K}\right) \left(1004.5\,\mathrm{J}\,\mathrm{kg}^{-1}\,\mathrm{K}^{-1}\right)}\right)^{\frac{\left(1.4\right)}{\left(1.4\right) - 1}}}
\]

\noindent\textbf{Resultado}

\[
\fitmath{\pi_c=1.378068393\times10^{0}}
\]

\noindent\textbf{Interpretación}
La relación debe ser mayor que uno. Compare el trabajo reconstruido desde la temperatura con Euler.

\noindent\textbf{Comprobación de los pasos}
\begin{itemize}
\item Paso 2: $32967.7$ — reproducido
\item Paso 4: $32.82$ — reproducido
\item Paso 6: $1.37807$ — reproducido
\end{itemize}

\noindent\textbf{Comprobación de la respuesta} la aritmética del último paso reproduce la respuesta enunciada a la precisión con que está escrita. Es una reproducción aritmética; no valida el modelo físico ni sus hipótesis.

\noindent\textbf{Limitaciones}
\begin{itemize}
\item Modelo educativo; la aplicabilidad a una máquina real requiere revisión experta.
\end{itemize}

\noindent\textbf{Errores frecuentes}
\begin{itemize}
\item Aplicar la razón de temperaturas reales como si fuera isentrópica.
\end{itemize}
% studium:end-record eyJpZCI6ICJQUkItMDA1MCIsICJzZWN0aW9uIjogImMwNyJ9
\end{studiumworked}

\begin{studiumworked}{Problema resuelto C07-02-R}


\noindent\textbf{Enunciado}
\(C07-02-R \cdot\) Reacción y remolinos. Construya un rotor compresor repetitivo a partir de reacción y carga. La incógnita final es el ángulo relativo de salida, algebraico en radianes. Datos: \(U=296\, \mathrm{m}\,\mathrm{s}^{-1}\); \(\psi=0.43\); \(\mathcal R=0.55\); \(C_x=164\, \mathrm{m}\,\mathrm{s}^{-1}\).

\noindent\textbf{Datos}
\begin{itemize}
\item $U = 296\,\mathrm{m}\,\mathrm{s}^{-1}$ — velocidad periférica
\item $\psi = 0.43$ — carga
\item $\mathcal R = 0.55$ — reacción
\item $C_x = 164\,\mathrm{m}\,\mathrm{s}^{-1}$ — velocidad axial
\end{itemize}

\noindent\textbf{Incógnita}
$\beta_2$, en $\mathrm{rad}$; también las magnitudes intermedias desarrolladas.

\noindent\textbf{Modelo físico}

\noindent\textit{Reacción y remolinos}

\[
\fitmath{\beta_2=\operatorname{atan}{\left(\frac{U \left(\frac{\psi}{2} - \mathcal R\right)}{C_x} \right)}}
\]

\noindent\textbf{Hipótesis}
\begin{itemize}
\item Radio y velocidad axial constantes; C3 al cuadrado igual a C1 al cuadrado.
\end{itemize}

\noindent\textbf{Desarrollo paso a paso}

\begin{enumerate}

\item La suma y diferencia fijan el remolino de entrada.

\[\fitmath{ct1=\mathrm{uu}\cdot \left(1-\mathrm{rea}-\frac{\mathrm{psi}}{2}\right)}\]

\[\fitmath{ct1=(296)\cdot \left(1-(0.55)-\frac{(0.43)}{2}\right)\approx 69.56\,\mathrm{m/s}}\]

\item Se calcula la componente absoluta de salida.

\[\fitmath{ct2=\mathrm{uu}\cdot \left(1-\mathrm{rea}+\frac{\mathrm{psi}}{2}\right)}\]

\[\fitmath{ct2=(296)\cdot \left(1-(0.55)+\frac{(0.43)}{2}\right)\approx 196.84\,\mathrm{m/s}}\]

\item Se resta U y se obtiene el ángulo desde el eje.

\[\fitmath{beta2=\arctan\left(\frac{\mathrm{uu}\cdot \left(-\mathrm{rea}+\frac{\mathrm{psi}}{2}\right)}{\mathrm{cx}}\right)}\]

\[\fitmath{beta2=\arctan\left(\frac{(296)\cdot \left(-(0.55)+\frac{(0.43)}{2}\right)}{(164)}\right)\approx -0.54382\,\mathrm{rad}}\]

\end{enumerate}


\noindent\textbf{Sustitución numérica}

\[
\fitmath{\operatorname{atan}{\left(\frac{\left(296\,\mathrm{m}\,\mathrm{s}^{-1}\right) \left(\frac{\left(0.43\right)}{2} - \left(0.55\right)\right)}{\left(164\,\mathrm{m}\,\mathrm{s}^{-1}\right)} \right)}}
\]

\noindent\textbf{Resultado}

\[
\fitmath{\beta_2=-5.438199954\times10^{-1}\,\mathrm{rad}}
\]

\noindent\textbf{Interpretación}
La beta negativa se conserva en la construcción. La reacción dada solo coincide con la definición bajo la hipótesis repetitiva.

\noindent\textbf{Comprobación de los pasos}
\begin{itemize}
\item Paso 2: $69.56$ — reproducido
\item Paso 4: $196.84$ — reproducido
\item Paso 6: $-0.54382$ — reproducido
\end{itemize}

\noindent\textbf{Comprobación de la respuesta} la aritmética del último paso reproduce la respuesta enunciada a la precisión con que está escrita. Es una reproducción aritmética; no valida el modelo físico ni sus hipótesis.

\noindent\textbf{Limitaciones}
\begin{itemize}
\item Modelo educativo; la aplicabilidad a una máquina real requiere revisión experta.
\end{itemize}

\noindent\textbf{Errores frecuentes}
\begin{itemize}
\item Usar el módulo positivo de beta de otro convenio.
\end{itemize}
% studium:end-record eyJpZCI6ICJQUkItMDA1MSIsICJzZWN0aW9uIjogImMwNyJ9
\end{studiumworked}

\begin{studiumworked}{Problema resuelto C07-02-P}


\noindent\textbf{Enunciado}
\(C07-02-P \cdot\) Reacción y remolinos. Ejercicio de transferencia. Construya un rotor compresor repetitivo a partir de reacción y carga. La incógnita final es el ángulo relativo de salida, algebraico en radianes. Antes de consultar la solución, prediga el efecto de cambiar el primer dato respecto del caso resuelto de esta familia. Justifique la tendencia y explique qué restricción podría convertirse en limitante. Datos: \(U=315.24\, \mathrm{m}\,\mathrm{s}^{-1}\); \(\psi=0.43\); \(\mathcal R=0.55\); \(C_x=164\, \mathrm{m}\,\mathrm{s}^{-1}\).

\noindent\textbf{Datos}
\begin{itemize}
\item $U = 315.24\,\mathrm{m}\,\mathrm{s}^{-1}$ — velocidad periférica
\item $\psi = 0.43$ — carga
\item $\mathcal R = 0.55$ — reacción
\item $C_x = 164\,\mathrm{m}\,\mathrm{s}^{-1}$ — velocidad axial
\end{itemize}

\noindent\textbf{Incógnita}
$\beta_2$, en $\mathrm{rad}$; también las magnitudes intermedias desarrolladas.

\noindent\textbf{Modelo físico}

\noindent\textit{Reacción y remolinos}

\[
\fitmath{\beta_2=\operatorname{atan}{\left(\frac{U \left(\frac{\psi}{2} - \mathcal R\right)}{C_x} \right)}}
\]

\noindent\textbf{Hipótesis}
\begin{itemize}
\item Radio y velocidad axial constantes; C3 al cuadrado igual a C1 al cuadrado.
\end{itemize}

\noindent\textbf{Desarrollo paso a paso}

\begin{enumerate}

\item La suma y diferencia fijan el remolino de entrada.

\[\fitmath{ct1=\mathrm{uu}\cdot \left(1-\mathrm{rea}-\frac{\mathrm{psi}}{2}\right)}\]

\[\fitmath{ct1=(315.24)\cdot \left(1-(0.55)-\frac{(0.43)}{2}\right)\approx 74.0814\,\mathrm{m/s}}\]

\item Se calcula la componente absoluta de salida.

\[\fitmath{ct2=\mathrm{uu}\cdot \left(1-\mathrm{rea}+\frac{\mathrm{psi}}{2}\right)}\]

\[\fitmath{ct2=(315.24)\cdot \left(1-(0.55)+\frac{(0.43)}{2}\right)\approx 209.6346\,\mathrm{m/s}}\]

\item Se resta U y se obtiene el ángulo desde el eje.

\[\fitmath{beta2=\arctan\left(\frac{\mathrm{uu}\cdot \left(-\mathrm{rea}+\frac{\mathrm{psi}}{2}\right)}{\mathrm{cx}}\right)}\]

\[\fitmath{beta2=\arctan\left(\frac{(315.24)\cdot \left(-(0.55)+\frac{(0.43)}{2}\right)}{(164)}\right)\approx -0.5721\,\mathrm{rad}}\]

\end{enumerate}


\noindent\textbf{Sustitución numérica}

\[
\fitmath{\operatorname{atan}{\left(\frac{\left(315.24\,\mathrm{m}\,\mathrm{s}^{-1}\right) \left(\frac{\left(0.43\right)}{2} - \left(0.55\right)\right)}{\left(164\,\mathrm{m}\,\mathrm{s}^{-1}\right)} \right)}}
\]

\noindent\textbf{Resultado}

\[
\fitmath{\beta_2=-5.721000360\times10^{-1}\,\mathrm{rad}}
\]

\noindent\textbf{Interpretación}
La beta negativa se conserva en la construcción. La reacción dada solo coincide con la definición bajo la hipótesis repetitiva.

\noindent\textbf{Comprobación de los pasos}
\begin{itemize}
\item Paso 2: $74.0814$ — reproducido
\item Paso 4: $209.635$ — reproducido
\item Paso 6: $-0.5721$ — reproducido
\end{itemize}

\noindent\textbf{Comprobación de la respuesta} la aritmética del último paso reproduce la respuesta enunciada a la precisión con que está escrita. Es una reproducción aritmética; no valida el modelo físico ni sus hipótesis.

\noindent\textbf{Limitaciones}
\begin{itemize}
\item Modelo educativo; la aplicabilidad a una máquina real requiere revisión experta.
\end{itemize}

\noindent\textbf{Errores frecuentes}
\begin{itemize}
\item Usar el módulo positivo de beta de otro convenio.
\end{itemize}
% studium:end-record eyJpZCI6ICJQUkItMDA1MiIsICJzZWN0aW9uIjogImMwNyJ9
\end{studiumworked}

\begin{studiumworked}{Problema resuelto C07-03-R}


\noindent\textbf{Enunciado}
\(C07-03-R \cdot\) Mach relativo. Calcule temperatura estática de entrada y Mach relativo del rotor a partir de remanso y componentes absolutas. Datos: \(T_0=309\, \mathrm{K}\); \(C_x=172\, \mathrm{m}\,\mathrm{s}^{-1}\); \(C_\theta=43\, \mathrm{m}\,\mathrm{s}^{-1}\); \(U=319\, \mathrm{m}\,\mathrm{s}^{-1}\); \(c_p=1004.5\, \mathrm{J}\,\mathrm{kg}^{-1}\,\mathrm{K}^{-1}\); \(\gamma=1.4\); \(R_g=287\, \mathrm{J}\,\mathrm{kg}^{-1}\,\mathrm{K}^{-1}\).

\noindent\textbf{Datos}
\begin{itemize}
\item $T_0 = 309\,\mathrm{K}$ — temperatura total absoluta
\item $C_x = 172\,\mathrm{m}\,\mathrm{s}^{-1}$ — componente axial
\item $C_\theta = 43\,\mathrm{m}\,\mathrm{s}^{-1}$ — remolino absoluto
\item $U = 319\,\mathrm{m}\,\mathrm{s}^{-1}$ — velocidad periférica
\item $c_p = 1004.5\,\mathrm{J}\,\mathrm{kg}^{-1}\,\mathrm{K}^{-1}$ — calor específico
\item $\gamma = 1.4$ — razón de calores
\item $R_g = 287\,\mathrm{J}\,\mathrm{kg}^{-1}\,\mathrm{K}^{-1}$ — constante del gas
\end{itemize}

\noindent\textbf{Incógnita}
$\mathrm{Ma}_r$ (adimensional); también las magnitudes intermedias desarrolladas.

\noindent\textbf{Modelo físico}

\noindent\textit{Mach relativo}

\[
\fitmath{\mathrm{Ma}_r=\frac{\sqrt{C_x^{2} + \left(C_\theta - U\right)^{2}}}{\sqrt{R_g \gamma \left(T_0 - \frac{C_\theta^{2} + C_x^{2}}{2 c_p}\right)}}}
\]

\noindent\textbf{Hipótesis}
\begin{itemize}
\item Gas perfecto y referencia local de rotor.
\end{itemize}

\noindent\textbf{Desarrollo paso a paso}

\begin{enumerate}

\item La energía cinética absoluta determina la temperatura estática.

\[\fitmath{T=T_0-\frac{{\mathrm{cx}}^{2}+{\mathrm{ct}}^{2}}{2\cdot \mathrm{cp}}}\]

\[\fitmath{T=(309)-\frac{{(172)}^{2}+{(43)}^{2}}{2\cdot (1004.5)}\approx 293.3539\,\mathrm{K}}\]

\item El módulo relativo se construye vectorialmente.

\[\fitmath{W=\sqrt{{\mathrm{cx}}^{2}+{\left(\mathrm{ct}-\mathrm{uu}\right)}^{2}}}\]

\[\fitmath{W=\sqrt{{(172)}^{2}+{\left((43)-(319)\right)}^{2}}\approx 325.2076\,\mathrm{m/s}}\]

\item El sonido depende de temperatura estática.

\[\fitmath{\mathrm{Mr}=\frac{\sqrt{{\mathrm{cx}}^{2}+{\left(\mathrm{ct}-\mathrm{uu}\right)}^{2}}}{\sqrt{\gamma\cdot R_g\cdot \left(T_0-\frac{{\mathrm{cx}}^{2}+{\mathrm{ct}}^{2}}{2\cdot \mathrm{cp}}\right)}}}\]

\[\fitmath{\mathrm{Mr}=\frac{\sqrt{{(172)}^{2}+{\left((43)-(319)\right)}^{2}}}{\sqrt{(1.4)\cdot (287)\cdot \left((309)-\frac{{(172)}^{2}+{(43)}^{2}}{2\cdot (1004.5)}\right)}}\approx 0.9472395\,}\]

\end{enumerate}


\noindent\textbf{Sustitución numérica}

\[
\fitmath{\frac{\sqrt{\left(\left(172\,\mathrm{m}\,\mathrm{s}^{-1}\right)\right)^{2} + \left(\left(43\,\mathrm{m}\,\mathrm{s}^{-1}\right) - \left(319\,\mathrm{m}\,\mathrm{s}^{-1}\right)\right)^{2}}}{\sqrt{\left(287\,\mathrm{J}\,\mathrm{kg}^{-1}\,\mathrm{K}^{-1}\right) \left(1.4\right) \left(\left(309\,\mathrm{K}\right) - \frac{\left(\left(43\,\mathrm{m}\,\mathrm{s}^{-1}\right)\right)^{2} + \left(\left(172\,\mathrm{m}\,\mathrm{s}^{-1}\right)\right)^{2}}{2 \left(1004.5\,\mathrm{J}\,\mathrm{kg}^{-1}\,\mathrm{K}^{-1}\right)}\right)}}}
\]

\noindent\textbf{Resultado}

\[
\fitmath{\mathrm{Ma}_r=9.472395356\times10^{-1}}
\]

\noindent\textbf{Interpretación}
El Mach relativo puede superar al absoluto. La propiedad acústica utiliza la misma temperatura estática en ambas referencias.

\noindent\textbf{Comprobación de los pasos}
\begin{itemize}
\item Paso 2: $293.354$ — reproducido
\item Paso 4: $325.208$ — reproducido
\item Paso 6: $0.94724$ — reproducido
\end{itemize}

\noindent\textbf{Comprobación de la respuesta} la aritmética del último paso reproduce la respuesta enunciada a la precisión con que está escrita. Es una reproducción aritmética; no valida el modelo físico ni sus hipótesis.

\noindent\textbf{Limitaciones}
\begin{itemize}
\item Modelo educativo; la aplicabilidad a una máquina real requiere revisión experta.
\end{itemize}

\noindent\textbf{Errores frecuentes}
\begin{itemize}
\item Usar temperatura total en la velocidad del sonido.
\end{itemize}
% studium:end-record eyJpZCI6ICJQUkItMDA1MyIsICJzZWN0aW9uIjogImMwNyJ9
\end{studiumworked}

\begin{studiumworked}{Problema resuelto C07-03-P}


\noindent\textbf{Enunciado}
\(C07-03-P \cdot\) Mach relativo. Ejercicio de transferencia. Calcule temperatura estática de entrada y Mach relativo del rotor a partir de remanso y componentes absolutas. Antes de consultar la solución, prediga el efecto de cambiar el primer dato respecto del caso resuelto de esta familia. Justifique la tendencia y explique qué restricción podría convertirse en limitante. Datos: \(T_0=329.085\, \mathrm{K}\); \(C_x=172\, \mathrm{m}\,\mathrm{s}^{-1}\); \(C_\theta=43\, \mathrm{m}\,\mathrm{s}^{-1}\); \(U=319\, \mathrm{m}\,\mathrm{s}^{-1}\); \(c_p=1004.5\, \mathrm{J}\,\mathrm{kg}^{-1}\,\mathrm{K}^{-1}\); \(\gamma=1.4\); \(R_g=287\, \mathrm{J}\,\mathrm{kg}^{-1}\,\mathrm{K}^{-1}\).

\noindent\textbf{Datos}
\begin{itemize}
\item $T_0 = 329.085\,\mathrm{K}$ — temperatura total absoluta
\item $C_x = 172\,\mathrm{m}\,\mathrm{s}^{-1}$ — componente axial
\item $C_\theta = 43\,\mathrm{m}\,\mathrm{s}^{-1}$ — remolino absoluto
\item $U = 319\,\mathrm{m}\,\mathrm{s}^{-1}$ — velocidad periférica
\item $c_p = 1004.5\,\mathrm{J}\,\mathrm{kg}^{-1}\,\mathrm{K}^{-1}$ — calor específico
\item $\gamma = 1.4$ — razón de calores
\item $R_g = 287\,\mathrm{J}\,\mathrm{kg}^{-1}\,\mathrm{K}^{-1}$ — constante del gas
\end{itemize}

\noindent\textbf{Incógnita}
$\mathrm{Ma}_r$ (adimensional); también las magnitudes intermedias desarrolladas.

\noindent\textbf{Modelo físico}

\noindent\textit{Mach relativo}

\[
\fitmath{\mathrm{Ma}_r=\frac{\sqrt{C_x^{2} + \left(C_\theta - U\right)^{2}}}{\sqrt{R_g \gamma \left(T_0 - \frac{C_\theta^{2} + C_x^{2}}{2 c_p}\right)}}}
\]

\noindent\textbf{Hipótesis}
\begin{itemize}
\item Gas perfecto y referencia local de rotor.
\end{itemize}

\noindent\textbf{Desarrollo paso a paso}

\begin{enumerate}

\item La energía cinética absoluta determina la temperatura estática.

\[\fitmath{T=T_0-\frac{{\mathrm{cx}}^{2}+{\mathrm{ct}}^{2}}{2\cdot \mathrm{cp}}}\]

\[\fitmath{T=(329.085)-\frac{{(172)}^{2}+{(43)}^{2}}{2\cdot (1004.5)}\approx 313.4389\,\mathrm{K}}\]

\item El módulo relativo se construye vectorialmente.

\[\fitmath{W=\sqrt{{\mathrm{cx}}^{2}+{\left(\mathrm{ct}-\mathrm{uu}\right)}^{2}}}\]

\[\fitmath{W=\sqrt{{(172)}^{2}+{\left((43)-(319)\right)}^{2}}\approx 325.2076\,\mathrm{m/s}}\]

\item El sonido depende de temperatura estática.

\[\fitmath{\mathrm{Mr}=\frac{\sqrt{{\mathrm{cx}}^{2}+{\left(\mathrm{ct}-\mathrm{uu}\right)}^{2}}}{\sqrt{\gamma\cdot R_g\cdot \left(T_0-\frac{{\mathrm{cx}}^{2}+{\mathrm{ct}}^{2}}{2\cdot \mathrm{cp}}\right)}}}\]

\[\fitmath{\mathrm{Mr}=\frac{\sqrt{{(172)}^{2}+{\left((43)-(319)\right)}^{2}}}{\sqrt{(1.4)\cdot (287)\cdot \left((329.085)-\frac{{(172)}^{2}+{(43)}^{2}}{2\cdot (1004.5)}\right)}}\approx 0.9163878\,}\]

\end{enumerate}


\noindent\textbf{Sustitución numérica}

\[
\fitmath{\frac{\sqrt{\left(\left(172\,\mathrm{m}\,\mathrm{s}^{-1}\right)\right)^{2} + \left(\left(43\,\mathrm{m}\,\mathrm{s}^{-1}\right) - \left(319\,\mathrm{m}\,\mathrm{s}^{-1}\right)\right)^{2}}}{\sqrt{\left(287\,\mathrm{J}\,\mathrm{kg}^{-1}\,\mathrm{K}^{-1}\right) \left(1.4\right) \left(\left(329.085\,\mathrm{K}\right) - \frac{\left(\left(43\,\mathrm{m}\,\mathrm{s}^{-1}\right)\right)^{2} + \left(\left(172\,\mathrm{m}\,\mathrm{s}^{-1}\right)\right)^{2}}{2 \left(1004.5\,\mathrm{J}\,\mathrm{kg}^{-1}\,\mathrm{K}^{-1}\right)}\right)}}}
\]

\noindent\textbf{Resultado}

\[
\fitmath{\mathrm{Ma}_r=9.163878079\times10^{-1}}
\]

\noindent\textbf{Interpretación}
El Mach relativo puede superar al absoluto. La propiedad acústica utiliza la misma temperatura estática en ambas referencias.

\noindent\textbf{Comprobación de los pasos}
\begin{itemize}
\item Paso 2: $313.439$ — reproducido
\item Paso 4: $325.208$ — reproducido
\item Paso 6: $0.916388$ — reproducido
\end{itemize}

\noindent\textbf{Comprobación de la respuesta} la aritmética del último paso reproduce la respuesta enunciada a la precisión con que está escrita. Es una reproducción aritmética; no valida el modelo físico ni sus hipótesis.

\noindent\textbf{Limitaciones}
\begin{itemize}
\item Modelo educativo; la aplicabilidad a una máquina real requiere revisión experta.
\end{itemize}

\noindent\textbf{Errores frecuentes}
\begin{itemize}
\item Usar temperatura total en la velocidad del sonido.
\end{itemize}
% studium:end-record eyJpZCI6ICJQUkItMDA1NCIsICJzZWN0aW9uIjogImMwNyJ9
\end{studiumworked}

\begin{studiumworked}{Problema resuelto C07-04-R}


\noindent\textbf{Enunciado}
\(C07-04-R \cdot\) Cierre de reacción por entalpías. Un rotor tiene los remolinos indicados y velocidad axial constante; el escalón es repetitivo. Calcule su reacción desde incrementos estáticos. Datos: \(U=284\, \mathrm{m}\,\mathrm{s}^{-1}\); \(C_{\theta e}=61\, \mathrm{m}\,\mathrm{s}^{-1}\); \(C_{\theta s}=172\, \mathrm{m}\,\mathrm{s}^{-1}\).

\noindent\textbf{Datos}
\begin{itemize}
\item $U = 284\,\mathrm{m}\,\mathrm{s}^{-1}$ — velocidad periférica
\item $C_{\theta e} = 61\,\mathrm{m}\,\mathrm{s}^{-1}$ — remolino entrada rotor
\item $C_{\theta s} = 172\,\mathrm{m}\,\mathrm{s}^{-1}$ — remolino salida rotor
\end{itemize}

\noindent\textbf{Incógnita}
$\mathcal{R}$ (adimensional); también las magnitudes intermedias desarrolladas.

\noindent\textbf{Modelo físico}

\noindent\textit{Cierre de reacción por entalpías}

\[
\fitmath{\mathcal R=\frac{\frac{C_{\theta e}^{2}}{2} - \frac{C_{\theta s}^{2}}{2} + U \left(- C_{\theta e} + C_{\theta s}\right)}{U \left(- C_{\theta e} + C_{\theta s}\right)}}
\]

\noindent\textbf{Hipótesis}
\begin{itemize}
\item Radio y Cx constantes, igual velocidad absoluta entre extremos del escalón.
\end{itemize}

\noindent\textbf{Desarrollo paso a paso}

\begin{enumerate}

\item Euler da el incremento total de entalpía.

\[\fitmath{\mathrm{dhz}=\mathrm{uu}\cdot \left(\mathrm{cs}-\mathrm{ce}\right)}\]

\[\fitmath{\mathrm{dhz}=(284)\cdot \left((172)-(61)\right)\approx 31524\,\mathrm{J/kg}}\]

\item Se resta el incremento de energía cinética absoluta.

\[\fitmath{\mathrm{dh}=\mathrm{uu}\cdot \left(\mathrm{cs}-\mathrm{ce}\right)-\frac{{\mathrm{cs}}^{2}-{\mathrm{ce}}^{2}}{2}}\]

\[\fitmath{\mathrm{dh}=(284)\cdot \left((172)-(61)\right)-\frac{{(172)}^{2}-{(61)}^{2}}{2}\approx 18592.5\,\mathrm{J/kg}}\]

\item Se forma la razón de incremento estático del rotor y del escalón.

\[\fitmath{\mathrm{rea}=\frac{\mathrm{uu}\cdot \left(\mathrm{cs}-\mathrm{ce}\right)-\frac{{\mathrm{cs}}^{2}-{\mathrm{ce}}^{2}}{2}}{\mathrm{uu}\cdot \left(\mathrm{cs}-\mathrm{ce}\right)}}\]

\[\fitmath{\mathrm{rea}=\frac{(284)\cdot \left((172)-(61)\right)-\frac{{(172)}^{2}-{(61)}^{2}}{2}}{(284)\cdot \left((172)-(61)\right)}\approx 0.5897887\,}\]

\end{enumerate}


\noindent\textbf{Sustitución numérica}

\[
\fitmath{\frac{\frac{\left(\left(61\,\mathrm{m}\,\mathrm{s}^{-1}\right)\right)^{2}}{2} - \frac{\left(\left(172\,\mathrm{m}\,\mathrm{s}^{-1}\right)\right)^{2}}{2} + \left(284\,\mathrm{m}\,\mathrm{s}^{-1}\right) \left(- \left(61\,\mathrm{m}\,\mathrm{s}^{-1}\right) + \left(172\,\mathrm{m}\,\mathrm{s}^{-1}\right)\right)}{\left(284\,\mathrm{m}\,\mathrm{s}^{-1}\right) \left(- \left(61\,\mathrm{m}\,\mathrm{s}^{-1}\right) + \left(172\,\mathrm{m}\,\mathrm{s}^{-1}\right)\right)}}
\]

\noindent\textbf{Resultado}

\[
\fitmath{\mathcal R=5.897887324\times10^{-1}}
\]

\noindent\textbf{Interpretación}
El resultado debe coincidir con 1-(ce+cs)/(2U). Esa identidad es una comprobación algebraica, no experimental.

\noindent\textbf{Comprobación de los pasos}
\begin{itemize}
\item Paso 2: $31524$ — reproducido
\item Paso 4: $18592.5$ — reproducido
\item Paso 6: $0.589789$ — reproducido
\end{itemize}

\noindent\textbf{Comprobación de la respuesta} la aritmética del último paso reproduce la respuesta enunciada a la precisión con que está escrita. Es una reproducción aritmética; no valida el modelo físico ni sus hipótesis.

\noindent\textbf{Limitaciones}
\begin{itemize}
\item Modelo educativo; la aplicabilidad a una máquina real requiere revisión experta.
\end{itemize}

\noindent\textbf{Errores frecuentes}
\begin{itemize}
\item Confundir el incremento estático de rotor y el total.
\end{itemize}
% studium:end-record eyJpZCI6ICJQUkItMDA1NSIsICJzZWN0aW9uIjogImMwNyJ9
\end{studiumworked}

\begin{studiumworked}{Problema resuelto C07-04-P}


\noindent\textbf{Enunciado}
\(C07-04-P \cdot\) Cierre de reacción por entalpías. Ejercicio de transferencia. Un rotor tiene los remolinos indicados y velocidad axial constante; el escalón es repetitivo. Calcule su reacción desde incrementos estáticos. Antes de consultar la solución, prediga el efecto de cambiar el primer dato respecto del caso resuelto de esta familia. Justifique la tendencia y explique qué restricción podría convertirse en limitante. Datos: \(U=302.46\, \mathrm{m}\,\mathrm{s}^{-1}\); \(C_{\theta e}=61\, \mathrm{m}\,\mathrm{s}^{-1}\); \(C_{\theta s}=172\, \mathrm{m}\,\mathrm{s}^{-1}\).

\noindent\textbf{Datos}
\begin{itemize}
\item $U = 302.46\,\mathrm{m}\,\mathrm{s}^{-1}$ — velocidad periférica
\item $C_{\theta e} = 61\,\mathrm{m}\,\mathrm{s}^{-1}$ — remolino entrada rotor
\item $C_{\theta s} = 172\,\mathrm{m}\,\mathrm{s}^{-1}$ — remolino salida rotor
\end{itemize}

\noindent\textbf{Incógnita}
$\mathcal{R}$ (adimensional); también las magnitudes intermedias desarrolladas.

\noindent\textbf{Modelo físico}

\noindent\textit{Cierre de reacción por entalpías}

\[
\fitmath{\mathcal R=\frac{\frac{C_{\theta e}^{2}}{2} - \frac{C_{\theta s}^{2}}{2} + U \left(- C_{\theta e} + C_{\theta s}\right)}{U \left(- C_{\theta e} + C_{\theta s}\right)}}
\]

\noindent\textbf{Hipótesis}
\begin{itemize}
\item Radio y Cx constantes, igual velocidad absoluta entre extremos del escalón.
\end{itemize}

\noindent\textbf{Desarrollo paso a paso}

\begin{enumerate}

\item Euler da el incremento total de entalpía.

\[\fitmath{\mathrm{dhz}=\mathrm{uu}\cdot \left(\mathrm{cs}-\mathrm{ce}\right)}\]

\[\fitmath{\mathrm{dhz}=(302.46)\cdot \left((172)-(61)\right)\approx 33573.06\,\mathrm{J/kg}}\]

\item Se resta el incremento de energía cinética absoluta.

\[\fitmath{\mathrm{dh}=\mathrm{uu}\cdot \left(\mathrm{cs}-\mathrm{ce}\right)-\frac{{\mathrm{cs}}^{2}-{\mathrm{ce}}^{2}}{2}}\]

\[\fitmath{\mathrm{dh}=(302.46)\cdot \left((172)-(61)\right)-\frac{{(172)}^{2}-{(61)}^{2}}{2}\approx 20641.56\,\mathrm{J/kg}}\]

\item Se forma la razón de incremento estático del rotor y del escalón.

\[\fitmath{\mathrm{rea}=\frac{\mathrm{uu}\cdot \left(\mathrm{cs}-\mathrm{ce}\right)-\frac{{\mathrm{cs}}^{2}-{\mathrm{ce}}^{2}}{2}}{\mathrm{uu}\cdot \left(\mathrm{cs}-\mathrm{ce}\right)}}\]

\[\fitmath{\mathrm{rea}=\frac{(302.46)\cdot \left((172)-(61)\right)-\frac{{(172)}^{2}-{(61)}^{2}}{2}}{(302.46)\cdot \left((172)-(61)\right)}\approx 0.6148251\,}\]

\end{enumerate}


\noindent\textbf{Sustitución numérica}

\[
\fitmath{\frac{\frac{\left(\left(61\,\mathrm{m}\,\mathrm{s}^{-1}\right)\right)^{2}}{2} - \frac{\left(\left(172\,\mathrm{m}\,\mathrm{s}^{-1}\right)\right)^{2}}{2} + \left(302.46\,\mathrm{m}\,\mathrm{s}^{-1}\right) \left(- \left(61\,\mathrm{m}\,\mathrm{s}^{-1}\right) + \left(172\,\mathrm{m}\,\mathrm{s}^{-1}\right)\right)}{\left(302.46\,\mathrm{m}\,\mathrm{s}^{-1}\right) \left(- \left(61\,\mathrm{m}\,\mathrm{s}^{-1}\right) + \left(172\,\mathrm{m}\,\mathrm{s}^{-1}\right)\right)}}
\]

\noindent\textbf{Resultado}

\[
\fitmath{\mathcal R=6.148251008\times10^{-1}}
\]

\noindent\textbf{Interpretación}
El resultado debe coincidir con 1-(ce+cs)/(2U). Esa identidad es una comprobación algebraica, no experimental.

\noindent\textbf{Comprobación de los pasos}
\begin{itemize}
\item Paso 2: $33573.1$ — reproducido
\item Paso 4: $20641.6$ — reproducido
\item Paso 6: $0.614825$ — reproducido
\end{itemize}

\noindent\textbf{Comprobación de la respuesta} la aritmética del último paso reproduce la respuesta enunciada a la precisión con que está escrita. Es una reproducción aritmética; no valida el modelo físico ni sus hipótesis.

\noindent\textbf{Limitaciones}
\begin{itemize}
\item Modelo educativo; la aplicabilidad a una máquina real requiere revisión experta.
\end{itemize}

\noindent\textbf{Errores frecuentes}
\begin{itemize}
\item Confundir el incremento estático de rotor y el total.
\end{itemize}
% studium:end-record eyJpZCI6ICJQUkItMDA1NiIsICJzZWN0aW9uIjogImMwNyJ9
\end{studiumworked}

\chapter{Cascadas: difusión, incidencia, desviación y pérdidas}
\begin{studiumworked}{Problema resuelto C08-01-R}


\noindent\textbf{Enunciado}
\(C08-01-R \cdot\) Indicador de difusión. Estime D para una fila móvil con módulos relativos y cambio tangencial dados. Datos: \(W_e=282\, \mathrm{m}\,\mathrm{s}^{-1}\); \(W_s=216\, \mathrm{m}\,\mathrm{s}^{-1}\); \(|\Delta W_\theta|=97\, \mathrm{m}\,\mathrm{s}^{-1}\); \(\sigma=1.24\).

\noindent\textbf{Datos}
\begin{itemize}
\item $W_e = 282\,\mathrm{m}\,\mathrm{s}^{-1}$ — módulo relativo entrada
\item $W_s = 216\,\mathrm{m}\,\mathrm{s}^{-1}$ — módulo relativo salida
\item $|\Delta W_\theta| = 97\,\mathrm{m}\,\mathrm{s}^{-1}$ — magnitud de cambio tangencial relativo
\item $\sigma = 1.24$ — solidez
\end{itemize}

\noindent\textbf{Incógnita}
$D$ (adimensional); también las magnitudes intermedias desarrolladas.

\noindent\textbf{Modelo físico}

\noindent\textit{Indicador de difusión}

\[
\fitmath{D=\frac{|\Delta W_\theta|}{2 \sigma W_e} + 1 - \frac{W_s}{W_e}}
\]

\noindent\textbf{Hipótesis}
\begin{itemize}
\item Factor de difusión definido en el capítulo, referencia relativa.
\end{itemize}

\noindent\textbf{Desarrollo paso a paso}

\begin{enumerate}

\item Se calcula la contribución de deceleración.

\[\fitmath{\mathrm{dec}=1-\frac{\mathrm{vs}}{\mathrm{ve}}}\]

\[\fitmath{\mathrm{dec}=1-\frac{(216)}{(282)}\approx 0.2340426\,}\]

\item La carga tangencial añade la segunda contribución.

\[\fitmath{D=1-\frac{\mathrm{vs}}{\mathrm{ve}}+\frac{\mathrm{dv}}{2\cdot \mathrm{sol}\cdot \mathrm{ve}}}\]

\[\fitmath{D=1-\frac{(216)}{(282)}+\frac{(97)}{2\cdot (1.24)\cdot (282)}\approx 0.3727408\,}\]

\end{enumerate}


\noindent\textbf{Sustitución numérica}

\[
\fitmath{\frac{\left(97\,\mathrm{m}\,\mathrm{s}^{-1}\right)}{2 \left(1.24\right) \left(282\,\mathrm{m}\,\mathrm{s}^{-1}\right)} + 1 - \frac{\left(216\,\mathrm{m}\,\mathrm{s}^{-1}\right)}{\left(282\,\mathrm{m}\,\mathrm{s}^{-1}\right)}}
\]

\noindent\textbf{Resultado}

\[
\fitmath{D=3.727407916\times10^{-1}}
\]

\noindent\textbf{Interpretación}
D es un indicador preliminar. La aceptabilidad exige criterio de perfil y régimen declarado.

\noindent\textbf{Comprobación de los pasos}
\begin{itemize}
\item Paso 2: $0.234043$ — reproducido
\item Paso 4: $0.372741$ — reproducido
\end{itemize}

\noindent\textbf{Comprobación de la respuesta} la aritmética del último paso reproduce la respuesta enunciada a la precisión con que está escrita. Es una reproducción aritmética; no valida el modelo físico ni sus hipótesis.

\noindent\textbf{Limitaciones}
\begin{itemize}
\item Modelo educativo; la aplicabilidad a una máquina real requiere revisión experta.
\end{itemize}

\noindent\textbf{Errores frecuentes}
\begin{itemize}
\item Usar velocidades absolutas para un rotor.
\end{itemize}
% studium:end-record eyJpZCI6ICJQUkItMDA1NyIsICJzZWN0aW9uIjogImMwOCJ9
\end{studiumworked}

\begin{studiumworked}{Problema resuelto C08-01-P}


\noindent\textbf{Enunciado}
\(C08-01-P \cdot\) Indicador de difusión. Ejercicio de transferencia. Estime D para una fila móvil con módulos relativos y cambio tangencial dados. Antes de consultar la solución, prediga el efecto de cambiar el primer dato respecto del caso resuelto de esta familia. Justifique la tendencia y explique qué restricción podría convertirse en limitante. Datos: \(W_e=300.33\, \mathrm{m}\,\mathrm{s}^{-1}\); \(W_s=216\, \mathrm{m}\,\mathrm{s}^{-1}\); \(|\Delta W_\theta|=97\, \mathrm{m}\,\mathrm{s}^{-1}\); \(\sigma=1.24\).

\noindent\textbf{Datos}
\begin{itemize}
\item $W_e = 300.33\,\mathrm{m}\,\mathrm{s}^{-1}$ — módulo relativo entrada
\item $W_s = 216\,\mathrm{m}\,\mathrm{s}^{-1}$ — módulo relativo salida
\item $|\Delta W_\theta| = 97\,\mathrm{m}\,\mathrm{s}^{-1}$ — magnitud de cambio tangencial relativo
\item $\sigma = 1.24$ — solidez
\end{itemize}

\noindent\textbf{Incógnita}
$D$ (adimensional); también las magnitudes intermedias desarrolladas.

\noindent\textbf{Modelo físico}

\noindent\textit{Indicador de difusión}

\[
\fitmath{D=\frac{|\Delta W_\theta|}{2 \sigma W_e} + 1 - \frac{W_s}{W_e}}
\]

\noindent\textbf{Hipótesis}
\begin{itemize}
\item Factor de difusión definido en el capítulo, referencia relativa.
\end{itemize}

\noindent\textbf{Desarrollo paso a paso}

\begin{enumerate}

\item Se calcula la contribución de deceleración.

\[\fitmath{\mathrm{dec}=1-\frac{\mathrm{vs}}{\mathrm{ve}}}\]

\[\fitmath{\mathrm{dec}=1-\frac{(216)}{(300.33)}\approx 0.2807911\,}\]

\item La carga tangencial añade la segunda contribución.

\[\fitmath{D=1-\frac{\mathrm{vs}}{\mathrm{ve}}+\frac{\mathrm{dv}}{2\cdot \mathrm{sol}\cdot \mathrm{ve}}}\]

\[\fitmath{D=1-\frac{(216)}{(300.33)}+\frac{(97)}{2\cdot (1.24)\cdot (300.33)}\approx 0.4110242\,}\]

\end{enumerate}


\noindent\textbf{Sustitución numérica}

\[
\fitmath{\frac{\left(97\,\mathrm{m}\,\mathrm{s}^{-1}\right)}{2 \left(1.24\right) \left(300.33\,\mathrm{m}\,\mathrm{s}^{-1}\right)} + 1 - \frac{\left(216\,\mathrm{m}\,\mathrm{s}^{-1}\right)}{\left(300.33\,\mathrm{m}\,\mathrm{s}^{-1}\right)}}
\]

\noindent\textbf{Resultado}

\[
\fitmath{D=4.110242174\times10^{-1}}
\]

\noindent\textbf{Interpretación}
D es un indicador preliminar. La aceptabilidad exige criterio de perfil y régimen declarado.

\noindent\textbf{Comprobación de los pasos}
\begin{itemize}
\item Paso 2: $0.280791$ — reproducido
\item Paso 4: $0.411024$ — reproducido
\end{itemize}

\noindent\textbf{Comprobación de la respuesta} la aritmética del último paso reproduce la respuesta enunciada a la precisión con que está escrita. Es una reproducción aritmética; no valida el modelo físico ni sus hipótesis.

\noindent\textbf{Limitaciones}
\begin{itemize}
\item Modelo educativo; la aplicabilidad a una máquina real requiere revisión experta.
\end{itemize}

\noindent\textbf{Errores frecuentes}
\begin{itemize}
\item Usar velocidades absolutas para un rotor.
\end{itemize}
% studium:end-record eyJpZCI6ICJQUkItMDA1OCIsICJzZWN0aW9uIjogImMwOCJ9
\end{studiumworked}

\begin{studiumworked}{Problema resuelto C08-02-R}


\noindent\textbf{Enunciado}
\(C08-02-R \cdot\) Solidez mínima del modelo. Un diseño impone D no superior al valor indicado. Calcule la solidez que satisface exactamente el límite y compruebe la positividad del denominador. Datos: \(W_e=301\, \mathrm{m}\,\mathrm{s}^{-1}\); \(W_s=232\, \mathrm{m}\,\mathrm{s}^{-1}\); \(|\Delta W_\theta|=112\, \mathrm{m}\,\mathrm{s}^{-1}\); \(D_{\mathrm{lim}}=0.47\).

\noindent\textbf{Datos}
\begin{itemize}
\item $W_e = 301\,\mathrm{m}\,\mathrm{s}^{-1}$ — módulo relativo entrada
\item $W_s = 232\,\mathrm{m}\,\mathrm{s}^{-1}$ — módulo relativo salida
\item $|\Delta W_\theta| = 112\,\mathrm{m}\,\mathrm{s}^{-1}$ — cambio tangencial de magnitud positiva
\item $D_{\mathrm{lim}} = 0.47$ — límite educativo de difusión
\end{itemize}

\noindent\textbf{Incógnita}
$\sigma$ (adimensional); también las magnitudes intermedias desarrolladas.

\noindent\textbf{Modelo físico}

\noindent\textit{Solidez mínima del modelo}

\[
\fitmath{\sigma=\frac{|\Delta W_\theta|}{2 W_e \left(D_{\mathrm{lim}} - 1 + \frac{W_s}{W_e}\right)}}
\]

\noindent\textbf{Hipótesis}
\begin{itemize}
\item Criterio de difusión prescrito como dato; sin garantía de pérdida real.
\end{itemize}

\noindent\textbf{Desarrollo paso a paso}

\begin{enumerate}

\item Se determina el margen restante después de la deceleración.

\[\fitmath{\mathrm{den}=\mathrm{lim}-1+\frac{\mathrm{vs}}{\mathrm{ve}}}\]

\[\fitmath{\mathrm{den}=(0.47)-1+\frac{(232)}{(301)}\approx 0.2407641\,}\]

\item Se despeja la solidez.

\[\fitmath{\mathrm{sol}=\frac{\mathrm{dv}}{2\cdot \mathrm{ve}\cdot \left(\mathrm{lim}-1+\frac{\mathrm{vs}}{\mathrm{ve}}\right)}}\]

\[\fitmath{\mathrm{sol}=\frac{(112)}{2\cdot (301)\cdot \left((0.47)-1+\frac{(232)}{(301)}\right)}\approx 0.7727335\,}\]

\end{enumerate}


\noindent\textbf{Sustitución numérica}

\[
\fitmath{\frac{\left(112\,\mathrm{m}\,\mathrm{s}^{-1}\right)}{2 \left(301\,\mathrm{m}\,\mathrm{s}^{-1}\right) \left(\left(0.47\right) - 1 + \frac{\left(232\,\mathrm{m}\,\mathrm{s}^{-1}\right)}{\left(301\,\mathrm{m}\,\mathrm{s}^{-1}\right)}\right)}}
\]

\noindent\textbf{Resultado}

\[
\fitmath{\sigma=7.727335449\times10^{-1}}
\]

\noindent\textbf{Interpretación}
Si el denominador fuera no positivo, este giro y cociente de módulos no podrían cumplir el límite por aumentar solidez.

\noindent\textbf{Comprobación de los pasos}
\begin{itemize}
\item Paso 2: $0.240764$ — reproducido
\item Paso 4: $0.772734$ — reproducido
\end{itemize}

\noindent\textbf{Comprobación de la respuesta} la aritmética del último paso reproduce la respuesta enunciada a la precisión con que está escrita. Es una reproducción aritmética; no valida el modelo físico ni sus hipótesis.

\noindent\textbf{Limitaciones}
\begin{itemize}
\item Modelo educativo; la aplicabilidad a una máquina real requiere revisión experta.
\end{itemize}

\noindent\textbf{Errores frecuentes}
\begin{itemize}
\item Aceptar una solidez negativa.
\end{itemize}
% studium:end-record eyJpZCI6ICJQUkItMDA1OSIsICJzZWN0aW9uIjogImMwOCJ9
\end{studiumworked}

\begin{studiumworked}{Problema resuelto C08-02-P}


\noindent\textbf{Enunciado}
\(C08-02-P \cdot\) Solidez mínima del modelo. Ejercicio de transferencia. Un diseño impone D no superior al valor indicado. Calcule la solidez que satisface exactamente el límite y compruebe la positividad del denominador. Antes de consultar la solución, prediga el efecto de cambiar el primer dato respecto del caso resuelto de esta familia. Justifique la tendencia y explique qué restricción podría convertirse en limitante. Datos: \(W_e=320.565\, \mathrm{m}\,\mathrm{s}^{-1}\); \(W_s=232\, \mathrm{m}\,\mathrm{s}^{-1}\); \(|\Delta W_\theta|=112\, \mathrm{m}\,\mathrm{s}^{-1}\); \(D_{\mathrm{lim}}=0.47\).

\noindent\textbf{Datos}
\begin{itemize}
\item $W_e = 320.565\,\mathrm{m}\,\mathrm{s}^{-1}$ — módulo relativo entrada
\item $W_s = 232\,\mathrm{m}\,\mathrm{s}^{-1}$ — módulo relativo salida
\item $|\Delta W_\theta| = 112\,\mathrm{m}\,\mathrm{s}^{-1}$ — cambio tangencial de magnitud positiva
\item $D_{\mathrm{lim}} = 0.47$ — límite educativo de difusión
\end{itemize}

\noindent\textbf{Incógnita}
$\sigma$ (adimensional); también las magnitudes intermedias desarrolladas.

\noindent\textbf{Modelo físico}

\noindent\textit{Solidez mínima del modelo}

\[
\fitmath{\sigma=\frac{|\Delta W_\theta|}{2 W_e \left(D_{\mathrm{lim}} - 1 + \frac{W_s}{W_e}\right)}}
\]

\noindent\textbf{Hipótesis}
\begin{itemize}
\item Criterio de difusión prescrito como dato; sin garantía de pérdida real.
\end{itemize}

\noindent\textbf{Desarrollo paso a paso}

\begin{enumerate}

\item Se determina el margen restante después de la deceleración.

\[\fitmath{\mathrm{den}=\mathrm{lim}-1+\frac{\mathrm{vs}}{\mathrm{ve}}}\]

\[\fitmath{\mathrm{den}=(0.47)-1+\frac{(232)}{(320.565)}\approx 0.1937222\,}\]

\item Se despeja la solidez.

\[\fitmath{\mathrm{sol}=\frac{\mathrm{dv}}{2\cdot \mathrm{ve}\cdot \left(\mathrm{lim}-1+\frac{\mathrm{vs}}{\mathrm{ve}}\right)}}\]

\[\fitmath{\mathrm{sol}=\frac{(112)}{2\cdot (320.565)\cdot \left((0.47)-1+\frac{(232)}{(320.565)}\right)}\approx 0.9017633\,}\]

\end{enumerate}


\noindent\textbf{Sustitución numérica}

\[
\fitmath{\frac{\left(112\,\mathrm{m}\,\mathrm{s}^{-1}\right)}{2 \left(320.565\,\mathrm{m}\,\mathrm{s}^{-1}\right) \left(\left(0.47\right) - 1 + \frac{\left(232\,\mathrm{m}\,\mathrm{s}^{-1}\right)}{\left(320.565\,\mathrm{m}\,\mathrm{s}^{-1}\right)}\right)}}
\]

\noindent\textbf{Resultado}

\[
\fitmath{\sigma=9.017633499\times10^{-1}}
\]

\noindent\textbf{Interpretación}
Si el denominador fuera no positivo, este giro y cociente de módulos no podrían cumplir el límite por aumentar solidez.

\noindent\textbf{Comprobación de los pasos}
\begin{itemize}
\item Paso 2: $0.193722$ — reproducido
\item Paso 4: $0.901763$ — reproducido
\end{itemize}

\noindent\textbf{Comprobación de la respuesta} la aritmética del último paso reproduce la respuesta enunciada a la precisión con que está escrita. Es una reproducción aritmética; no valida el modelo físico ni sus hipótesis.

\noindent\textbf{Limitaciones}
\begin{itemize}
\item Modelo educativo; la aplicabilidad a una máquina real requiere revisión experta.
\end{itemize}

\noindent\textbf{Errores frecuentes}
\begin{itemize}
\item Aceptar una solidez negativa.
\end{itemize}
% studium:end-record eyJpZCI6ICJQUkItMDA2MCIsICJzZWN0aW9uIjogImMwOCJ9
\end{studiumworked}

\begin{studiumworked}{Problema resuelto C08-03-R}


\noindent\textbf{Enunciado}
\(C08-03-R \cdot\) Número entero y solidez real. Se eligen Z álabes enteros para un radio y una cuerda dados. Calcule paso y solidez real tras la elección. Datos: \(r=0.263\, \mathrm{m}\); \(c=0.052\, \mathrm{m}\); \(Z=41\).

\noindent\textbf{Datos}
\begin{itemize}
\item $r = 0.263\,\mathrm{m}$ — radio
\item $c = 0.052\,\mathrm{m}$ — cuerda
\item $Z = 41$ — número entero de álabes
\end{itemize}

\noindent\textbf{Incógnita}
$\sigma$ (adimensional); también las magnitudes intermedias desarrolladas.

\noindent\textbf{Modelo físico}

\noindent\textit{Número entero y solidez real}

\[
\fitmath{\sigma=\frac{c Z}{2 \pi r}}
\]

\noindent\textbf{Hipótesis}
\begin{itemize}
\item Fila uniforme al radio de referencia.
\end{itemize}

\noindent\textbf{Desarrollo paso a paso}

\begin{enumerate}

\item La circunferencia se reparte en pasos iguales.

\[\fitmath{\mathrm{ss}=\frac{2\cdot \pi\cdot \mathrm{rad}}{\mathrm{zz}}}\]

\[\fitmath{\mathrm{ss}=\frac{2\cdot \pi\cdot (0.263)}{(41)}\approx 0.04030434\,\mathrm{m}}\]

\item La solidez usa el paso realmente obtenido.

\[\fitmath{\mathrm{sol}=\frac{\mathrm{cc}\cdot \mathrm{zz}}{2\cdot \pi\cdot \mathrm{rad}}}\]

\[\fitmath{\mathrm{sol}=\frac{(0.052)\cdot (41)}{2\cdot \pi\cdot (0.263)}\approx 1.290184\,}\]

\end{enumerate}


\noindent\textbf{Sustitución numérica}

\[
\fitmath{\frac{\left(0.052\,\mathrm{m}\right) \left(41\right)}{2 \pi \left(0.263\,\mathrm{m}\right)}}
\]

\noindent\textbf{Resultado}

\[
\fitmath{\sigma=1.290183797\times10^{0}}
\]

\noindent\textbf{Interpretación}
El redondeo del número de álabes modifica la solidez; debe evaluarse después el parámetro de carga.

\noindent\textbf{Comprobación de los pasos}
\begin{itemize}
\item Paso 2: $0.0403043$ — reproducido
\item Paso 4: $1.29018$ — reproducido
\end{itemize}

\noindent\textbf{Comprobación de la respuesta} la aritmética del último paso reproduce la respuesta enunciada a la precisión con que está escrita. Es una reproducción aritmética; no valida el modelo físico ni sus hipótesis.

\noindent\textbf{Limitaciones}
\begin{itemize}
\item Modelo educativo; la aplicabilidad a una máquina real requiere revisión experta.
\end{itemize}

\noindent\textbf{Errores frecuentes}
\begin{itemize}
\item Conservar exactamente la solidez objetivo después del redondeo.
\end{itemize}
% studium:end-record eyJpZCI6ICJQUkItMDA2MSIsICJzZWN0aW9uIjogImMwOCJ9
\end{studiumworked}

\begin{studiumworked}{Problema resuelto C08-03-P}


\noindent\textbf{Enunciado}
\(C08-03-P \cdot\) Número entero y solidez real. Ejercicio de transferencia. Se eligen Z álabes enteros para un radio y una cuerda dados. Calcule paso y solidez real tras la elección. Antes de consultar la solución, prediga el efecto de cambiar el primer dato respecto del caso resuelto de esta familia. Justifique la tendencia y explique qué restricción podría convertirse en limitante. Datos: \(r=0.280095\, \mathrm{m}\); \(c=0.052\, \mathrm{m}\); \(Z=41\).

\noindent\textbf{Datos}
\begin{itemize}
\item $r = 0.280095\,\mathrm{m}$ — radio
\item $c = 0.052\,\mathrm{m}$ — cuerda
\item $Z = 41$ — número entero de álabes
\end{itemize}

\noindent\textbf{Incógnita}
$\sigma$ (adimensional); también las magnitudes intermedias desarrolladas.

\noindent\textbf{Modelo físico}

\noindent\textit{Número entero y solidez real}

\[
\fitmath{\sigma=\frac{c Z}{2 \pi r}}
\]

\noindent\textbf{Hipótesis}
\begin{itemize}
\item Fila uniforme al radio de referencia.
\end{itemize}

\noindent\textbf{Desarrollo paso a paso}

\begin{enumerate}

\item La circunferencia se reparte en pasos iguales.

\[\fitmath{\mathrm{ss}=\frac{2\cdot \pi\cdot \mathrm{rad}}{\mathrm{zz}}}\]

\[\fitmath{\mathrm{ss}=\frac{2\cdot \pi\cdot (0.280095)}{(41)}\approx 0.04292412\,\mathrm{m}}\]

\item La solidez usa el paso realmente obtenido.

\[\fitmath{\mathrm{sol}=\frac{\mathrm{cc}\cdot \mathrm{zz}}{2\cdot \pi\cdot \mathrm{rad}}}\]

\[\fitmath{\mathrm{sol}=\frac{(0.052)\cdot (41)}{2\cdot \pi\cdot (0.280095)}\approx 1.21144\,}\]

\end{enumerate}


\noindent\textbf{Sustitución numérica}

\[
\fitmath{\frac{\left(0.052\,\mathrm{m}\right) \left(41\right)}{2 \pi \left(0.280095\,\mathrm{m}\right)}}
\]

\noindent\textbf{Resultado}

\[
\fitmath{\sigma=1.211440185\times10^{0}}
\]

\noindent\textbf{Interpretación}
El redondeo del número de álabes modifica la solidez; debe evaluarse después el parámetro de carga.

\noindent\textbf{Comprobación de los pasos}
\begin{itemize}
\item Paso 2: $0.0429241$ — reproducido
\item Paso 4: $1.21144$ — reproducido
\end{itemize}

\noindent\textbf{Comprobación de la respuesta} la aritmética del último paso reproduce la respuesta enunciada a la precisión con que está escrita. Es una reproducción aritmética; no valida el modelo físico ni sus hipótesis.

\noindent\textbf{Limitaciones}
\begin{itemize}
\item Modelo educativo; la aplicabilidad a una máquina real requiere revisión experta.
\end{itemize}

\noindent\textbf{Errores frecuentes}
\begin{itemize}
\item Conservar exactamente la solidez objetivo después del redondeo.
\end{itemize}
% studium:end-record eyJpZCI6ICJQUkItMDA2MiIsICJzZWN0aW9uIjogImMwOCJ9
\end{studiumworked}

\begin{studiumworked}{Problema resuelto C08-04-R}


\noindent\textbf{Enunciado}
\(C08-04-R \cdot\) Incidencia y desviación algebraicas. En una fila fija se miden ángulos de gas y se conocen los metálicos. Determine incidencia y desviación; la incógnita final es delta. Datos: \(\alpha_e=0.8552113\, \mathrm{rad}\); \(\alpha_s=0.3839724\, \mathrm{rad}\); \(\kappa_e=0.8028515\, \mathrm{rad}\); \(\kappa_s=0.3141593\, \mathrm{rad}\).

\noindent\textbf{Datos}
\begin{itemize}
\item $\alpha_e = 0.85521133347722\,\mathrm{rad}$ — ángulo gas entrada
\item $\alpha_s = 0.38397243543875\,\mathrm{rad}$ — ángulo gas salida
\item $\kappa_e = 0.80285145591739\,\mathrm{rad}$ — ángulo metal entrada
\item $\kappa_s = 0.31415926535898\,\mathrm{rad}$ — ángulo metal salida
\end{itemize}

\noindent\textbf{Incógnita}
$\delta$, en $\mathrm{rad}$; también las magnitudes intermedias desarrolladas.

\noindent\textbf{Modelo físico}

\noindent\textit{Incidencia y desviación algebraicas}

\[
\fitmath{\delta=\alpha_s - \kappa_s}
\]

\noindent\textbf{Hipótesis}
\begin{itemize}
\item Ángulos desde el eje y mismo sentido positivo.
\end{itemize}

\noindent\textbf{Desarrollo paso a paso}

\begin{enumerate}

\item La incidencia compara gas y metal en entrada.

\[\fitmath{\mathrm{inc}=\mathrm{ae}-\mathrm{ke}}\]

\[\fitmath{\mathrm{inc}=(0.8552113335)-(0.8028514559)\approx 0.05235988\,\mathrm{rad}}\]

\item La desviación compara gas y metal en salida según nuestro convenio.

\[\fitmath{\mathrm{delta}=\mathrm{asv}-\mathrm{ks}}\]

\[\fitmath{\mathrm{delta}=(0.3839724354)-(0.3141592654)\approx 0.06981317\,\mathrm{rad}}\]

\end{enumerate}


\noindent\textbf{Sustitución numérica}

\[
\fitmath{\left(0.38397243543875\,\mathrm{rad}\right) - \left(0.31415926535898\,\mathrm{rad}\right)}
\]

\noindent\textbf{Resultado}

\[
\fitmath{\delta=6.981317008\times10^{-2}\,\mathrm{rad}}
\]

\noindent\textbf{Interpretación}
Los signos dependen de una definición explícita. Para expresar en grados multiplique los radianes por 180/pi.

\noindent\textbf{Comprobación de los pasos}
\begin{itemize}
\item Paso 2: $0.0523599$ — reproducido
\item Paso 4: $0.0698132$ — reproducido
\end{itemize}

\noindent\textbf{Comprobación de la respuesta} la aritmética del último paso reproduce la respuesta enunciada a la precisión con que está escrita. Es una reproducción aritmética; no valida el modelo físico ni sus hipótesis.

\noindent\textbf{Limitaciones}
\begin{itemize}
\item Modelo educativo; la aplicabilidad a una máquina real requiere revisión experta.
\end{itemize}

\noindent\textbf{Errores frecuentes}
\begin{itemize}
\item Comparar ángulo absoluto con beta relativa.
\end{itemize}
% studium:end-record eyJpZCI6ICJQUkItMDA2MyIsICJzZWN0aW9uIjogImMwOCJ9
\end{studiumworked}

\begin{studiumworked}{Problema resuelto C08-04-P}


\noindent\textbf{Enunciado}
\(C08-04-P \cdot\) Incidencia y desviación algebraicas. Ejercicio de transferencia. En una fila fija se miden ángulos de gas y se conocen los metálicos. Determine incidencia y desviación; la incógnita final es delta. Antes de consultar la solución, prediga el efecto de cambiar el primer dato respecto del caso resuelto de esta familia. Justifique la tendencia y explique qué restricción podría convertirse en limitante. Datos: \(\alpha_e=0.9108001\, \mathrm{rad}\); \(\alpha_s=0.3839724\, \mathrm{rad}\); \(\kappa_e=0.8028515\, \mathrm{rad}\); \(\kappa_s=0.3141593\, \mathrm{rad}\).

\noindent\textbf{Datos}
\begin{itemize}
\item $\alpha_e = 0.91080007015324\,\mathrm{rad}$ — ángulo gas entrada
\item $\alpha_s = 0.38397243543875\,\mathrm{rad}$ — ángulo gas salida
\item $\kappa_e = 0.80285145591739\,\mathrm{rad}$ — ángulo metal entrada
\item $\kappa_s = 0.31415926535898\,\mathrm{rad}$ — ángulo metal salida
\end{itemize}

\noindent\textbf{Incógnita}
$\delta$, en $\mathrm{rad}$; también las magnitudes intermedias desarrolladas.

\noindent\textbf{Modelo físico}

\noindent\textit{Incidencia y desviación algebraicas}

\[
\fitmath{\delta=\alpha_s - \kappa_s}
\]

\noindent\textbf{Hipótesis}
\begin{itemize}
\item Ángulos desde el eje y mismo sentido positivo.
\end{itemize}

\noindent\textbf{Desarrollo paso a paso}

\begin{enumerate}

\item La incidencia compara gas y metal en entrada.

\[\fitmath{\mathrm{inc}=\mathrm{ae}-\mathrm{ke}}\]

\[\fitmath{\mathrm{inc}=(0.9108000702)-(0.8028514559)\approx 0.1079486\,\mathrm{rad}}\]

\item La desviación compara gas y metal en salida según nuestro convenio.

\[\fitmath{\mathrm{delta}=\mathrm{asv}-\mathrm{ks}}\]

\[\fitmath{\mathrm{delta}=(0.3839724354)-(0.3141592654)\approx 0.06981317\,\mathrm{rad}}\]

\end{enumerate}


\noindent\textbf{Sustitución numérica}

\[
\fitmath{\left(0.38397243543875\,\mathrm{rad}\right) - \left(0.31415926535898\,\mathrm{rad}\right)}
\]

\noindent\textbf{Resultado}

\[
\fitmath{\delta=6.981317008\times10^{-2}\,\mathrm{rad}}
\]

\noindent\textbf{Interpretación}
Los signos dependen de una definición explícita. Para expresar en grados multiplique los radianes por 180/pi.

\noindent\textbf{Comprobación de los pasos}
\begin{itemize}
\item Paso 2: $0.107949$ — reproducido
\item Paso 4: $0.0698132$ — reproducido
\end{itemize}

\noindent\textbf{Comprobación de la respuesta} la aritmética del último paso reproduce la respuesta enunciada a la precisión con que está escrita. Es una reproducción aritmética; no valida el modelo físico ni sus hipótesis.

\noindent\textbf{Limitaciones}
\begin{itemize}
\item Modelo educativo; la aplicabilidad a una máquina real requiere revisión experta.
\end{itemize}

\noindent\textbf{Errores frecuentes}
\begin{itemize}
\item Comparar ángulo absoluto con beta relativa.
\end{itemize}
% studium:end-record eyJpZCI6ICJQUkItMDA2NCIsICJzZWN0aW9uIjogImMwOCJ9
\end{studiumworked}

\chapter{Compresión multiescalón, anillo y acoplamiento}
\begin{studiumworked}{Problema resuelto C09-01-R}


\noindent\textbf{Enunciado}
\(C09-01-R \cdot\) Compresión de varios escalones. Una compresión politrópica global alimenta un eje. Halle temperatura final y potencia requerida en el eje. Datos: \(T_{0e}=289\, \mathrm{K}\); \(\pi_c=9.3\); \(\eta_p=0.9\); \(\gamma=1.4\); \(c_p=1004.5\, \mathrm{J}\,\mathrm{kg}^{-1}\,\mathrm{K}^{-1}\); \(\dot m=23\, \mathrm{kg}\,\mathrm{s}^{-1}\); \(\eta_m=0.98\).

\noindent\textbf{Datos}
\begin{itemize}
\item $T_{0e} = 289\,\mathrm{K}$ — temperatura total entrada
\item $\pi_c = 9.3$ — relación global de presión
\item $\eta_p = 0.9$ — eficiencia politrópica
\item $\gamma = 1.4$ — razón de calores
\item $c_p = 1004.5\,\mathrm{J}\,\mathrm{kg}^{-1}\,\mathrm{K}^{-1}$ — calor específico
\item $\dot m = 23\,\mathrm{kg}\,\mathrm{s}^{-1}$ — caudal
\item $\eta_m = 0.98$ — rendimiento mecánico compresor
\end{itemize}

\noindent\textbf{Incógnita}
$\dot{W}$, en $\mathrm{W}$; también las magnitudes intermedias desarrolladas.

\noindent\textbf{Modelo físico}

\noindent\textit{Compresión de varios escalones}

\[
\fitmath{\dot{W}=\frac{T_{0e} c_p \dot m \left(\pi_c^{\frac{\gamma - 1}{\eta_p \gamma}} - 1\right)}{\eta_m}}
\]

\noindent\textbf{Hipótesis}
\begin{itemize}
\item Propiedades y eficiencia politrópica constantes, sin sangrado.
\end{itemize}

\noindent\textbf{Desarrollo paso a paso}

\begin{enumerate}

\item Se integra el modelo politrópico global.

\[\fitmath{\mathrm{Tout}=\mathrm{Tin}\cdot {\mathrm{pr}}^{\frac{\gamma-1}{\gamma\cdot \mathrm{ep}}}}\]

\[\fitmath{\mathrm{Tout}=(289)\cdot {(9.3)}^{\frac{(1.4)-1}{(1.4)\cdot (0.9)}}\approx 586.6144\,\mathrm{K}}\]

\item La potencia de eje supera la potencia comunicada al fluido.

\[\fitmath{\mathrm{Pow}=\frac{\mathrm{md}\cdot \mathrm{cp}\cdot \mathrm{Tin}\cdot \left({\mathrm{pr}}^{\frac{\gamma-1}{\gamma\cdot \mathrm{ep}}}-1\right)}{\mathrm{em}}}\]

\[\fitmath{\mathrm{Pow}=\frac{(23)\cdot (1004.5)\cdot (289)\cdot \left({(9.3)}^{\frac{(1.4)-1}{(1.4)\cdot (0.9)}}-1\right)}{(0.98)}\approx 7016260\,\mathrm{W}}\]

\end{enumerate}


\noindent\textbf{Sustitución numérica}

\[
\fitmath{\frac{\left(289\,\mathrm{K}\right) \left(1004.5\,\mathrm{J}\,\mathrm{kg}^{-1}\,\mathrm{K}^{-1}\right) \left(23\,\mathrm{kg}\,\mathrm{s}^{-1}\right) \left(\left(9.3\right)^{\frac{\left(1.4\right) - 1}{\left(0.9\right) \left(1.4\right)}} - 1\right)}{\left(0.98\right)}}
\]

\noindent\textbf{Resultado}

\[
\fitmath{\dot{W}=7.016260352\times10^{6}\,\mathrm{W}}
\]

\noindent\textbf{Interpretación}
No se multiplica la eficiencia de cada escalón. El caudal y el incremento total de entalpía fijan la potencia.

\noindent\textbf{Comprobación de los pasos}
\begin{itemize}
\item Paso 2: $586.614$ — reproducido
\item Paso 4: $7.01626\times 10^{6}$ — reproducido
\end{itemize}

\noindent\textbf{Comprobación de la respuesta} la aritmética del último paso reproduce la respuesta enunciada a la precisión con que está escrita. Es una reproducción aritmética; no valida el modelo físico ni sus hipótesis.

\noindent\textbf{Limitaciones}
\begin{itemize}
\item Modelo educativo; la aplicabilidad a una máquina real requiere revisión experta.
\end{itemize}

\noindent\textbf{Errores frecuentes}
\begin{itemize}
\item Usar el producto de eficiencias como eficiencia global.
\end{itemize}
% studium:end-record eyJpZCI6ICJQUkItMDA2NSIsICJzZWN0aW9uIjogImMwOSJ9
\end{studiumworked}

\begin{studiumworked}{Problema resuelto C09-01-P}


\noindent\textbf{Enunciado}
\(C09-01-P \cdot\) Compresión de varios escalones. Ejercicio de transferencia. Una compresión politrópica global alimenta un eje. Halle temperatura final y potencia requerida en el eje. Antes de consultar la solución, prediga el efecto de cambiar el primer dato respecto del caso resuelto de esta familia. Justifique la tendencia y explique qué restricción podría convertirse en limitante. Datos: \(T_{0e}=307.785\, \mathrm{K}\); \(\pi_c=9.3\); \(\eta_p=0.9\); \(\gamma=1.4\); \(c_p=1004.5\, \mathrm{J}\,\mathrm{kg}^{-1}\,\mathrm{K}^{-1}\); \(\dot m=23\, \mathrm{kg}\,\mathrm{s}^{-1}\); \(\eta_m=0.98\).

\noindent\textbf{Datos}
\begin{itemize}
\item $T_{0e} = 307.785\,\mathrm{K}$ — temperatura total entrada
\item $\pi_c = 9.3$ — relación global de presión
\item $\eta_p = 0.9$ — eficiencia politrópica
\item $\gamma = 1.4$ — razón de calores
\item $c_p = 1004.5\,\mathrm{J}\,\mathrm{kg}^{-1}\,\mathrm{K}^{-1}$ — calor específico
\item $\dot m = 23\,\mathrm{kg}\,\mathrm{s}^{-1}$ — caudal
\item $\eta_m = 0.98$ — rendimiento mecánico compresor
\end{itemize}

\noindent\textbf{Incógnita}
$\dot{W}$, en $\mathrm{W}$; también las magnitudes intermedias desarrolladas.

\noindent\textbf{Modelo físico}

\noindent\textit{Compresión de varios escalones}

\[
\fitmath{\dot{W}=\frac{T_{0e} c_p \dot m \left(\pi_c^{\frac{\gamma - 1}{\eta_p \gamma}} - 1\right)}{\eta_m}}
\]

\noindent\textbf{Hipótesis}
\begin{itemize}
\item Propiedades y eficiencia politrópica constantes, sin sangrado.
\end{itemize}

\noindent\textbf{Desarrollo paso a paso}

\begin{enumerate}

\item Se integra el modelo politrópico global.

\[\fitmath{\mathrm{Tout}=\mathrm{Tin}\cdot {\mathrm{pr}}^{\frac{\gamma-1}{\gamma\cdot \mathrm{ep}}}}\]

\[\fitmath{\mathrm{Tout}=(307.785)\cdot {(9.3)}^{\frac{(1.4)-1}{(1.4)\cdot (0.9)}}\approx 624.7444\,\mathrm{K}}\]

\item La potencia de eje supera la potencia comunicada al fluido.

\[\fitmath{\mathrm{Pow}=\frac{\mathrm{md}\cdot \mathrm{cp}\cdot \mathrm{Tin}\cdot \left({\mathrm{pr}}^{\frac{\gamma-1}{\gamma\cdot \mathrm{ep}}}-1\right)}{\mathrm{em}}}\]

\[\fitmath{\mathrm{Pow}=\frac{(23)\cdot (1004.5)\cdot (307.785)\cdot \left({(9.3)}^{\frac{(1.4)-1}{(1.4)\cdot (0.9)}}-1\right)}{(0.98)}\approx 7472317\,\mathrm{W}}\]

\end{enumerate}


\noindent\textbf{Sustitución numérica}

\[
\fitmath{\frac{\left(307.785\,\mathrm{K}\right) \left(1004.5\,\mathrm{J}\,\mathrm{kg}^{-1}\,\mathrm{K}^{-1}\right) \left(23\,\mathrm{kg}\,\mathrm{s}^{-1}\right) \left(\left(9.3\right)^{\frac{\left(1.4\right) - 1}{\left(0.9\right) \left(1.4\right)}} - 1\right)}{\left(0.98\right)}}
\]

\noindent\textbf{Resultado}

\[
\fitmath{\dot{W}=7.472317275\times10^{6}\,\mathrm{W}}
\]

\noindent\textbf{Interpretación}
No se multiplica la eficiencia de cada escalón. El caudal y el incremento total de entalpía fijan la potencia.

\noindent\textbf{Comprobación de los pasos}
\begin{itemize}
\item Paso 2: $624.744$ — reproducido
\item Paso 4: $7.47232\times 10^{6}$ — reproducido
\end{itemize}

\noindent\textbf{Comprobación de la respuesta} la aritmética del último paso reproduce la respuesta enunciada a la precisión con que está escrita. Es una reproducción aritmética; no valida el modelo físico ni sus hipótesis.

\noindent\textbf{Limitaciones}
\begin{itemize}
\item Modelo educativo; la aplicabilidad a una máquina real requiere revisión experta.
\end{itemize}

\noindent\textbf{Errores frecuentes}
\begin{itemize}
\item Usar el producto de eficiencias como eficiencia global.
\end{itemize}
% studium:end-record eyJpZCI6ICJQUkItMDA2NiIsICJzZWN0aW9uIjogImMwOSJ9
\end{studiumworked}

\begin{studiumworked}{Problema resuelto C09-02-R}


\noindent\textbf{Enunciado}
\(C09-02-R \cdot\) Altura de anillo con bloqueo. Dimensione el área geométrica y la altura del anillo con un factor B de área disponible. Datos: \(\dot m=27\, \mathrm{kg}\,\mathrm{s}^{-1}\); \(\rho=3.6\, \mathrm{kg}\,\mathrm{m}^{-3}\); \(C_x=151\, \mathrm{m}\,\mathrm{s}^{-1}\); \(B=0.955\); \(r_m=0.285\, \mathrm{m}\).

\noindent\textbf{Datos}
\begin{itemize}
\item $\dot m = 27\,\mathrm{kg}\,\mathrm{s}^{-1}$ — caudal
\item $\rho = 3.6\,\mathrm{kg}\,\mathrm{m}^{-3}$ — densidad estática
\item $C_x = 151\,\mathrm{m}\,\mathrm{s}^{-1}$ — velocidad axial efectiva
\item $B = 0.955$ — fracción de área disponible
\item $r_m = 0.285\,\mathrm{m}$ — radio medio aritmético
\end{itemize}

\noindent\textbf{Incógnita}
$b$, en $\mathrm{m}$; también las magnitudes intermedias desarrolladas.

\noindent\textbf{Modelo físico}

\noindent\textit{Altura de anillo con bloqueo}

\[
\fitmath{b=\frac{\dot m}{2 \pi B C_x \rho r_m}}
\]

\noindent\textbf{Hipótesis}
\begin{itemize}
\item Flujo uniforme en el área disponible.
\end{itemize}

\noindent\textbf{Desarrollo paso a paso}

\begin{enumerate}

\item La continuidad se aplica al área efectiva B A.

\[\fitmath{\mathrm{area}=\frac{\mathrm{md}}{\rho\cdot \mathrm{cx}\cdot \mathrm{bb}}}\]

\[\fitmath{\mathrm{area}=\frac{(27)}{(3.6)\cdot (151)\cdot (0.955)}\approx 0.05200929\,\mathrm{m^{2}}}\]

\item Se convierte el área geométrica en altura.

\[\fitmath{b=\frac{\mathrm{md}}{\rho\cdot \mathrm{cx}\cdot \mathrm{bb}\cdot 2\cdot \pi\cdot \mathrm{rm}}}\]

\[\fitmath{b=\frac{(27)}{(3.6)\cdot (151)\cdot (0.955)\cdot 2\cdot \pi\cdot (0.285)}\approx 0.02904399\,\mathrm{m}}\]

\end{enumerate}


\noindent\textbf{Sustitución numérica}

\[
\fitmath{\frac{\left(27\,\mathrm{kg}\,\mathrm{s}^{-1}\right)}{2 \pi \left(0.955\right) \left(151\,\mathrm{m}\,\mathrm{s}^{-1}\right) \left(3.6\,\mathrm{kg}\,\mathrm{m}^{-3}\right) \left(0.285\,\mathrm{m}\right)}}
\]

\noindent\textbf{Resultado}

\[
\fitmath{b=2.904398583\times10^{-2}\,\mathrm{m}}
\]

\noindent\textbf{Interpretación}
B inferior a uno exige mayor área geométrica. La raíz debe permanecer positiva al restar b/2 al radio medio.

\noindent\textbf{Comprobación de los pasos}
\begin{itemize}
\item Paso 2: $0.0520093$ — reproducido
\item Paso 4: $0.029044$ — reproducido
\end{itemize}

\noindent\textbf{Comprobación de la respuesta} la aritmética del último paso reproduce la respuesta enunciada a la precisión con que está escrita. Es una reproducción aritmética; no valida el modelo físico ni sus hipótesis.

\noindent\textbf{Limitaciones}
\begin{itemize}
\item Modelo educativo; la aplicabilidad a una máquina real requiere revisión experta.
\end{itemize}

\noindent\textbf{Errores frecuentes}
\begin{itemize}
\item Confundir B con la fracción bloqueada 1-B.
\end{itemize}
% studium:end-record eyJpZCI6ICJQUkItMDA2NyIsICJzZWN0aW9uIjogImMwOSJ9
\end{studiumworked}

\begin{studiumworked}{Problema resuelto C09-02-P}


\noindent\textbf{Enunciado}
\(C09-02-P \cdot\) Altura de anillo con bloqueo. Ejercicio de transferencia. Dimensione el área geométrica y la altura del anillo con un factor B de área disponible. Antes de consultar la solución, prediga el efecto de cambiar el primer dato respecto del caso resuelto de esta familia. Justifique la tendencia y explique qué restricción podría convertirse en limitante. Datos: \(\dot m=28.755\, \mathrm{kg}\,\mathrm{s}^{-1}\); \(\rho=3.6\, \mathrm{kg}\,\mathrm{m}^{-3}\); \(C_x=151\, \mathrm{m}\,\mathrm{s}^{-1}\); \(B=0.955\); \(r_m=0.285\, \mathrm{m}\).

\noindent\textbf{Datos}
\begin{itemize}
\item $\dot m = 28.755\,\mathrm{kg}\,\mathrm{s}^{-1}$ — caudal
\item $\rho = 3.6\,\mathrm{kg}\,\mathrm{m}^{-3}$ — densidad estática
\item $C_x = 151\,\mathrm{m}\,\mathrm{s}^{-1}$ — velocidad axial efectiva
\item $B = 0.955$ — fracción de área disponible
\item $r_m = 0.285\,\mathrm{m}$ — radio medio aritmético
\end{itemize}

\noindent\textbf{Incógnita}
$b$, en $\mathrm{m}$; también las magnitudes intermedias desarrolladas.

\noindent\textbf{Modelo físico}

\noindent\textit{Altura de anillo con bloqueo}

\[
\fitmath{b=\frac{\dot m}{2 \pi B C_x \rho r_m}}
\]

\noindent\textbf{Hipótesis}
\begin{itemize}
\item Flujo uniforme en el área disponible.
\end{itemize}

\noindent\textbf{Desarrollo paso a paso}

\begin{enumerate}

\item La continuidad se aplica al área efectiva B A.

\[\fitmath{\mathrm{area}=\frac{\mathrm{md}}{\rho\cdot \mathrm{cx}\cdot \mathrm{bb}}}\]

\[\fitmath{\mathrm{area}=\frac{(28.755)}{(3.6)\cdot (151)\cdot (0.955)}\approx 0.0553899\,\mathrm{m^{2}}}\]

\item Se convierte el área geométrica en altura.

\[\fitmath{b=\frac{\mathrm{md}}{\rho\cdot \mathrm{cx}\cdot \mathrm{bb}\cdot 2\cdot \pi\cdot \mathrm{rm}}}\]

\[\fitmath{b=\frac{(28.755)}{(3.6)\cdot (151)\cdot (0.955)\cdot 2\cdot \pi\cdot (0.285)}\approx 0.03093184\,\mathrm{m}}\]

\end{enumerate}


\noindent\textbf{Sustitución numérica}

\[
\fitmath{\frac{\left(28.755\,\mathrm{kg}\,\mathrm{s}^{-1}\right)}{2 \pi \left(0.955\right) \left(151\,\mathrm{m}\,\mathrm{s}^{-1}\right) \left(3.6\,\mathrm{kg}\,\mathrm{m}^{-3}\right) \left(0.285\,\mathrm{m}\right)}}
\]

\noindent\textbf{Resultado}

\[
\fitmath{b=3.093184491\times10^{-2}\,\mathrm{m}}
\]

\noindent\textbf{Interpretación}
B inferior a uno exige mayor área geométrica. La raíz debe permanecer positiva al restar b/2 al radio medio.

\noindent\textbf{Comprobación de los pasos}
\begin{itemize}
\item Paso 2: $0.0553899$ — reproducido
\item Paso 4: $0.0309318$ — reproducido
\end{itemize}

\noindent\textbf{Comprobación de la respuesta} la aritmética del último paso reproduce la respuesta enunciada a la precisión con que está escrita. Es una reproducción aritmética; no valida el modelo físico ni sus hipótesis.

\noindent\textbf{Limitaciones}
\begin{itemize}
\item Modelo educativo; la aplicabilidad a una máquina real requiere revisión experta.
\end{itemize}

\noindent\textbf{Errores frecuentes}
\begin{itemize}
\item Confundir B con la fracción bloqueada 1-B.
\end{itemize}
% studium:end-record eyJpZCI6ICJQUkItMDA2OCIsICJzZWN0aW9uIjogImMwOSJ9
\end{studiumworked}

\begin{studiumworked}{Problema resuelto C09-03-R}


\noindent\textbf{Enunciado}
\(C09-03-R \cdot\) Trabajo por escalón y razón de presión. Un escalón tardío recibe gas más caliente. Calcule su relación de presión para un trabajo y eficiencia fijos. Datos: \(T_{0e}=487\, \mathrm{K}\); \(w_c=29800\, \mathrm{J}\,\mathrm{kg}^{-1}\); \(\eta=0.87\); \(\gamma=1.4\); \(c_p=1004.5\, \mathrm{J}\,\mathrm{kg}^{-1}\,\mathrm{K}^{-1}\).

\noindent\textbf{Datos}
\begin{itemize}
\item $T_{0e} = 487\,\mathrm{K}$ — temperatura total entrada
\item $w_c = 29800\,\mathrm{J}\,\mathrm{kg}^{-1}$ — trabajo específico
\item $\eta = 0.87$ — eficiencia total-total
\item $\gamma = 1.4$ — razón de calores
\item $c_p = 1004.5\,\mathrm{J}\,\mathrm{kg}^{-1}\,\mathrm{K}^{-1}$ — calor específico
\end{itemize}

\noindent\textbf{Incógnita}
$\pi_c$ (adimensional); también las magnitudes intermedias desarrolladas.

\noindent\textbf{Modelo físico}

\noindent\textit{Trabajo por escalón y razón de presión}

\[
\fitmath{\pi_c=\left(1 + \frac{\eta w_c}{T_{0e} c_p}\right)^{\frac{\gamma}{\gamma - 1}}}
\]

\noindent\textbf{Hipótesis}
\begin{itemize}
\item Gas perfecto constante y escalón adiabático.
\end{itemize}

\noindent\textbf{Desarrollo paso a paso}

\begin{enumerate}

\item Se calcula el aumento relativo de temperatura ideal.

\[\fitmath{\mathrm{rr}=\frac{\mathrm{eta}\cdot \mathrm{wc}}{\mathrm{cp}\cdot \mathrm{Tin}}}\]

\[\fitmath{\mathrm{rr}=\frac{(0.87)\cdot (29800)}{(1004.5)\cdot (487)}\approx 0.05299765\,}\]

\item Se recupera la relación de presión.

\[\fitmath{\mathrm{pr}={\left(1+\frac{\mathrm{eta}\cdot \mathrm{wc}}{\mathrm{cp}\cdot \mathrm{Tin}}\right)}^{\frac{\gamma}{\gamma-1}}}\]

\[\fitmath{\mathrm{pr}={\left(1+\frac{(0.87)\cdot (29800)}{(1004.5)\cdot (487)}\right)}^{\frac{(1.4)}{(1.4)-1}}\approx 1.198108\,}\]

\end{enumerate}


\noindent\textbf{Sustitución numérica}

\[
\fitmath{\left(1 + \frac{\left(0.87\right) \left(29800\,\mathrm{J}\,\mathrm{kg}^{-1}\right)}{\left(487\,\mathrm{K}\right) \left(1004.5\,\mathrm{J}\,\mathrm{kg}^{-1}\,\mathrm{K}^{-1}\right)}\right)^{\frac{\left(1.4\right)}{\left(1.4\right) - 1}}}
\]

\noindent\textbf{Resultado}

\[
\fitmath{\pi_c=1.198107832\times10^{0}}
\]

\noindent\textbf{Interpretación}
A mayor temperatura de entrada y mismo trabajo, la razón de presión es menor.

\noindent\textbf{Comprobación de los pasos}
\begin{itemize}
\item Paso 2: $0.0529977$ — reproducido
\item Paso 4: $1.19811$ — reproducido
\end{itemize}

\noindent\textbf{Comprobación de la respuesta} la aritmética del último paso reproduce la respuesta enunciada a la precisión con que está escrita. Es una reproducción aritmética; no valida el modelo físico ni sus hipótesis.

\noindent\textbf{Limitaciones}
\begin{itemize}
\item Modelo educativo; la aplicabilidad a una máquina real requiere revisión experta.
\end{itemize}

\noindent\textbf{Errores frecuentes}
\begin{itemize}
\item Asignar igual razón de presión a todos los escalones de igual trabajo.
\end{itemize}
% studium:end-record eyJpZCI6ICJQUkItMDA2OSIsICJzZWN0aW9uIjogImMwOSJ9
\end{studiumworked}

\begin{studiumworked}{Problema resuelto C09-03-P}


\noindent\textbf{Enunciado}
\(C09-03-P \cdot\) Trabajo por escalón y razón de presión. Ejercicio de transferencia. Un escalón tardío recibe gas más caliente. Calcule su relación de presión para un trabajo y eficiencia fijos. Antes de consultar la solución, prediga el efecto de cambiar el primer dato respecto del caso resuelto de esta familia. Justifique la tendencia y explique qué restricción podría convertirse en limitante. Datos: \(T_{0e}=518.655\, \mathrm{K}\); \(w_c=29800\, \mathrm{J}\,\mathrm{kg}^{-1}\); \(\eta=0.87\); \(\gamma=1.4\); \(c_p=1004.5\, \mathrm{J}\,\mathrm{kg}^{-1}\,\mathrm{K}^{-1}\).

\noindent\textbf{Datos}
\begin{itemize}
\item $T_{0e} = 518.655\,\mathrm{K}$ — temperatura total entrada
\item $w_c = 29800\,\mathrm{J}\,\mathrm{kg}^{-1}$ — trabajo específico
\item $\eta = 0.87$ — eficiencia total-total
\item $\gamma = 1.4$ — razón de calores
\item $c_p = 1004.5\,\mathrm{J}\,\mathrm{kg}^{-1}\,\mathrm{K}^{-1}$ — calor específico
\end{itemize}

\noindent\textbf{Incógnita}
$\pi_c$ (adimensional); también las magnitudes intermedias desarrolladas.

\noindent\textbf{Modelo físico}

\noindent\textit{Trabajo por escalón y razón de presión}

\[
\fitmath{\pi_c=\left(1 + \frac{\eta w_c}{T_{0e} c_p}\right)^{\frac{\gamma}{\gamma - 1}}}
\]

\noindent\textbf{Hipótesis}
\begin{itemize}
\item Gas perfecto constante y escalón adiabático.
\end{itemize}

\noindent\textbf{Desarrollo paso a paso}

\begin{enumerate}

\item Se calcula el aumento relativo de temperatura ideal.

\[\fitmath{\mathrm{rr}=\frac{\mathrm{eta}\cdot \mathrm{wc}}{\mathrm{cp}\cdot \mathrm{Tin}}}\]

\[\fitmath{\mathrm{rr}=\frac{(0.87)\cdot (29800)}{(1004.5)\cdot (518.655)}\approx 0.04976305\,}\]

\item Se recupera la relación de presión.

\[\fitmath{\mathrm{pr}={\left(1+\frac{\mathrm{eta}\cdot \mathrm{wc}}{\mathrm{cp}\cdot \mathrm{Tin}}\right)}^{\frac{\gamma}{\gamma-1}}}\]

\[\fitmath{\mathrm{pr}={\left(1+\frac{(0.87)\cdot (29800)}{(1004.5)\cdot (518.655)}\right)}^{\frac{(1.4)}{(1.4)-1}}\approx 1.185276\,}\]

\end{enumerate}


\noindent\textbf{Sustitución numérica}

\[
\fitmath{\left(1 + \frac{\left(0.87\right) \left(29800\,\mathrm{J}\,\mathrm{kg}^{-1}\right)}{\left(518.655\,\mathrm{K}\right) \left(1004.5\,\mathrm{J}\,\mathrm{kg}^{-1}\,\mathrm{K}^{-1}\right)}\right)^{\frac{\left(1.4\right)}{\left(1.4\right) - 1}}}
\]

\noindent\textbf{Resultado}

\[
\fitmath{\pi_c=1.185275999\times10^{0}}
\]

\noindent\textbf{Interpretación}
A mayor temperatura de entrada y mismo trabajo, la razón de presión es menor.

\noindent\textbf{Comprobación de los pasos}
\begin{itemize}
\item Paso 2: $0.0497631$ — reproducido
\item Paso 4: $1.18528$ — reproducido
\end{itemize}

\noindent\textbf{Comprobación de la respuesta} la aritmética del último paso reproduce la respuesta enunciada a la precisión con que está escrita. Es una reproducción aritmética; no valida el modelo físico ni sus hipótesis.

\noindent\textbf{Limitaciones}
\begin{itemize}
\item Modelo educativo; la aplicabilidad a una máquina real requiere revisión experta.
\end{itemize}

\noindent\textbf{Errores frecuentes}
\begin{itemize}
\item Asignar igual razón de presión a todos los escalones de igual trabajo.
\end{itemize}
% studium:end-record eyJpZCI6ICJQUkItMDA3MCIsICJzZWN0aW9uIjogImMwOSJ9
\end{studiumworked}

\begin{studiumworked}{Problema resuelto C09-04-R}


\noindent\textbf{Enunciado}
\(C09-04-R \cdot\) Acoplamiento de eje con caudales distintos. El compresor y la turbina comparten eje con rendimiento mecánico dado. Calcule el trabajo que debe extraer cada kilogramo de gas en la turbina. Datos: \(\dot{m}_c=25\, \mathrm{kg}\,\mathrm{s}^{-1}\); \(\dot{m}_t=25.6\, \mathrm{kg}\,\mathrm{s}^{-1}\); \(w_c=214000\, \mathrm{J}\,\mathrm{kg}^{-1}\); \(\dot{W}_{\mathrm{acc}}=85000\, \mathrm{W}\); \(\eta_m=0.975\).

\noindent\textbf{Datos}
\begin{itemize}
\item $\dot{m}_c = 25\,\mathrm{kg}\,\mathrm{s}^{-1}$ — caudal compresor
\item $\dot{m}_t = 25.6\,\mathrm{kg}\,\mathrm{s}^{-1}$ — caudal turbina
\item $w_c = 214000\,\mathrm{J}\,\mathrm{kg}^{-1}$ — trabajo comunicado al aire
\item $\dot{W}_{\mathrm{acc}} = 85000\,\mathrm{W}$ — potencia de accesorios
\item $\eta_m = 0.975$ — eficiencia de transmisión desde turbina
\end{itemize}

\noindent\textbf{Incógnita}
$w_t$, en $\mathrm{J\,kg^{-1}}$; también las magnitudes intermedias desarrolladas.

\noindent\textbf{Modelo físico}

\noindent\textit{Acoplamiento de eje con caudales distintos}

\[
\fitmath{w_t=\frac{\dot{W}_{\mathrm{acc}} + \dot{m}_c w_c}{\eta_m \dot{m}_t}}
\]

\noindent\textbf{Hipótesis}
\begin{itemize}
\item wc ya excluye pérdidas de transmisión; em se aplica una vez al lado turbina.
\end{itemize}

\noindent\textbf{Desarrollo paso a paso}

\begin{enumerate}

\item Se suma la demanda de potencia de aire y accesorios.

\[\fitmath{\mathrm{need}=\mathrm{mc}\cdot \mathrm{wc}+\mathrm{acc}}\]

\[\fitmath{\mathrm{need}=(25)\cdot (214000)+(85000)\approx 5435000\,\mathrm{W}}\]

\item El caudal de turbina y la transmisión fijan su trabajo.

\[\fitmath{\mathrm{wt}=\frac{\mathrm{mc}\cdot \mathrm{wc}+\mathrm{acc}}{\mathrm{em}\cdot \mathrm{mt}}}\]

\[\fitmath{\mathrm{wt}=\frac{(25)\cdot (214000)+(85000)}{(0.975)\cdot (25.6)}\approx 217748.4\,\mathrm{J/kg}}\]

\end{enumerate}


\noindent\textbf{Sustitución numérica}

\[
\fitmath{\frac{\left(85000\,\mathrm{W}\right) + \left(25\,\mathrm{kg}\,\mathrm{s}^{-1}\right) \left(214000\,\mathrm{J}\,\mathrm{kg}^{-1}\right)}{\left(0.975\right) \left(25.6\,\mathrm{kg}\,\mathrm{s}^{-1}\right)}}
\]

\noindent\textbf{Resultado}

\[
\fitmath{w_t=2.177483974\times10^{5}\,\mathrm{J}\,\mathrm{kg}^{-1}}
\]

\noindent\textbf{Interpretación}
Usar el caudal del compresor en la turbina omite la diferencia de masas. El reparto de accesorios está explicitado en este modelo.

\noindent\textbf{Comprobación de los pasos}
\begin{itemize}
\item Paso 2: $5435000$ — reproducido
\item Paso 4: $217748$ — reproducido
\end{itemize}

\noindent\textbf{Comprobación de la respuesta} la aritmética del último paso reproduce la respuesta enunciada a la precisión con que está escrita. Es una reproducción aritmética; no valida el modelo físico ni sus hipótesis.

\noindent\textbf{Limitaciones}
\begin{itemize}
\item Modelo educativo; la aplicabilidad a una máquina real requiere revisión experta.
\end{itemize}

\noindent\textbf{Errores frecuentes}
\begin{itemize}
\item Suponer igualdad de caudales sin revisar combustible y sangrados.
\end{itemize}
% studium:end-record eyJpZCI6ICJQUkItMDA3MSIsICJzZWN0aW9uIjogImMwOSJ9
\end{studiumworked}

\begin{studiumworked}{Problema resuelto C09-04-P}


\noindent\textbf{Enunciado}
\(C09-04-P \cdot\) Acoplamiento de eje con caudales distintos. Ejercicio de transferencia. El compresor y la turbina comparten eje con rendimiento mecánico dado. Calcule el trabajo que debe extraer cada kilogramo de gas en la turbina. Antes de consultar la solución, prediga el efecto de cambiar el primer dato respecto del caso resuelto de esta familia. Justifique la tendencia y explique qué restricción podría convertirse en limitante. Datos: \(\dot{m}_c=26.625\, \mathrm{kg}\,\mathrm{s}^{-1}\); \(\dot{m}_t=25.6\, \mathrm{kg}\,\mathrm{s}^{-1}\); \(w_c=214000\, \mathrm{J}\,\mathrm{kg}^{-1}\); \(\dot{W}_{\mathrm{acc}}=85000\, \mathrm{W}\); \(\eta_m=0.975\).

\noindent\textbf{Datos}
\begin{itemize}
\item $\dot{m}_c = 26.625\,\mathrm{kg}\,\mathrm{s}^{-1}$ — caudal compresor
\item $\dot{m}_t = 25.6\,\mathrm{kg}\,\mathrm{s}^{-1}$ — caudal turbina
\item $w_c = 214000\,\mathrm{J}\,\mathrm{kg}^{-1}$ — trabajo comunicado al aire
\item $\dot{W}_{\mathrm{acc}} = 85000\,\mathrm{W}$ — potencia de accesorios
\item $\eta_m = 0.975$ — eficiencia de transmisión desde turbina
\end{itemize}

\noindent\textbf{Incógnita}
$w_t$, en $\mathrm{J\,kg^{-1}}$; también las magnitudes intermedias desarrolladas.

\noindent\textbf{Modelo físico}

\noindent\textit{Acoplamiento de eje con caudales distintos}

\[
\fitmath{w_t=\frac{\dot{W}_{\mathrm{acc}} + \dot{m}_c w_c}{\eta_m \dot{m}_t}}
\]

\noindent\textbf{Hipótesis}
\begin{itemize}
\item wc ya excluye pérdidas de transmisión; em se aplica una vez al lado turbina.
\end{itemize}

\noindent\textbf{Desarrollo paso a paso}

\begin{enumerate}

\item Se suma la demanda de potencia de aire y accesorios.

\[\fitmath{\mathrm{need}=\mathrm{mc}\cdot \mathrm{wc}+\mathrm{acc}}\]

\[\fitmath{\mathrm{need}=(26.625)\cdot (214000)+(85000)\approx 5782750\,\mathrm{W}}\]

\item El caudal de turbina y la transmisión fijan su trabajo.

\[\fitmath{\mathrm{wt}=\frac{\mathrm{mc}\cdot \mathrm{wc}+\mathrm{acc}}{\mathrm{em}\cdot \mathrm{mt}}}\]

\[\fitmath{\mathrm{wt}=\frac{(26.625)\cdot (214000)+(85000)}{(0.975)\cdot (25.6)}\approx 231680.7\,\mathrm{J/kg}}\]

\end{enumerate}


\noindent\textbf{Sustitución numérica}

\[
\fitmath{\frac{\left(85000\,\mathrm{W}\right) + \left(26.625\,\mathrm{kg}\,\mathrm{s}^{-1}\right) \left(214000\,\mathrm{J}\,\mathrm{kg}^{-1}\right)}{\left(0.975\right) \left(25.6\,\mathrm{kg}\,\mathrm{s}^{-1}\right)}}
\]

\noindent\textbf{Resultado}

\[
\fitmath{w_t=2.316806891\times10^{5}\,\mathrm{J}\,\mathrm{kg}^{-1}}
\]

\noindent\textbf{Interpretación}
Usar el caudal del compresor en la turbina omite la diferencia de masas. El reparto de accesorios está explicitado en este modelo.

\noindent\textbf{Comprobación de los pasos}
\begin{itemize}
\item Paso 2: $5782750$ — reproducido
\item Paso 4: $231681$ — reproducido
\end{itemize}

\noindent\textbf{Comprobación de la respuesta} la aritmética del último paso reproduce la respuesta enunciada a la precisión con que está escrita. Es una reproducción aritmética; no valida el modelo físico ni sus hipótesis.

\noindent\textbf{Limitaciones}
\begin{itemize}
\item Modelo educativo; la aplicabilidad a una máquina real requiere revisión experta.
\end{itemize}

\noindent\textbf{Errores frecuentes}
\begin{itemize}
\item Suponer igualdad de caudales sin revisar combustible y sangrados.
\end{itemize}
% studium:end-record eyJpZCI6ICJQUkItMDA3MiIsICJzZWN0aW9uIjogImMwOSJ9
\end{studiumworked}

\chapter{Arranque, pérdida rotativa, bombeo y mapas}
\begin{studiumworked}{Problema resuelto C10-01-R}


\noindent\textbf{Enunciado}
\(C10-01-R \cdot\) Incidencia durante variación de caudal. Un rotor sin prerrotación conserva velocidad periférica y ángulo metálico. Calcule su beta de entrada y la incidencia a la nueva velocidad axial. Datos: \(U=247\, \mathrm{m}\,\mathrm{s}^{-1}\); \(C_x=119\, \mathrm{m}\,\mathrm{s}^{-1}\); \(\kappa_e=-1.064651\, \mathrm{rad}\).

\noindent\textbf{Datos}
\begin{itemize}
\item $U = 247\,\mathrm{m}\,\mathrm{s}^{-1}$ — velocidad periférica
\item $C_x = 119\,\mathrm{m}\,\mathrm{s}^{-1}$ — velocidad axial nueva
\item $\kappa_e = -1.0646508437165\,\mathrm{rad}$ — ángulo metálico de entrada
\end{itemize}

\noindent\textbf{Incógnita}
$i$, en $\mathrm{rad}$; también las magnitudes intermedias desarrolladas.

\noindent\textbf{Modelo físico}

\noindent\textit{Incidencia durante variación de caudal}

\[
\fitmath{i=- \kappa_e - \operatorname{atan}{\left(\frac{U}{C_x} \right)}}
\]

\noindent\textbf{Hipótesis}
\begin{itemize}
\item Velocidad axial positiva, sin prerrotación, metal fijo.
\end{itemize}

\noindent\textbf{Desarrollo paso a paso}

\begin{enumerate}

\item El triángulo relativo utiliza tangencial -U.

\[\fitmath{\beta=\arctan\left(\frac{-\mathrm{uu}}{\mathrm{cx}}\right)}\]

\[\fitmath{\beta=\arctan\left(\frac{-(247)}{(119)}\right)\approx -1.12183\,\mathrm{rad}}\]

\item Se compara con la dirección del metal.

\[\fitmath{\mathrm{inc}=\arctan\left(\frac{-\mathrm{uu}}{\mathrm{cx}}\right)-\mathrm{kap}}\]

\[\fitmath{\mathrm{inc}=\arctan\left(\frac{-(247)}{(119)}\right)-(-1.064650844)\approx -0.05717871\,\mathrm{rad}}\]

\end{enumerate}


\noindent\textbf{Sustitución numérica}

\[
\fitmath{- \left(-1.0646508437165\,\mathrm{rad}\right) - \operatorname{atan}{\left(\frac{\left(247\,\mathrm{m}\,\mathrm{s}^{-1}\right)}{\left(119\,\mathrm{m}\,\mathrm{s}^{-1}\right)} \right)}}
\]

\noindent\textbf{Resultado}

\[
\fitmath{i=-5.717871161\times10^{-2}\,\mathrm{rad}}
\]

\noindent\textbf{Interpretación}
El signo corresponde exclusivamente al convenio algebraico de este manual. Debe compararse con una incidencia nominal calculada igual.

\noindent\textbf{Comprobación de los pasos}
\begin{itemize}
\item Paso 2: $-1.12183$ — reproducido
\item Paso 4: $-0.0571787$ — reproducido
\end{itemize}

\noindent\textbf{Comprobación de la respuesta} la aritmética del último paso reproduce la respuesta enunciada a la precisión con que está escrita. Es una reproducción aritmética; no valida el modelo físico ni sus hipótesis.

\noindent\textbf{Limitaciones}
\begin{itemize}
\item Modelo educativo; la aplicabilidad a una máquina real requiere revisión experta.
\end{itemize}

\noindent\textbf{Errores frecuentes}
\begin{itemize}
\item Conservar beta al cambiar Cx/U.
\end{itemize}
% studium:end-record eyJpZCI6ICJQUkItMDA3MyIsICJzZWN0aW9uIjogImMxMCJ9
\end{studiumworked}

\begin{studiumworked}{Problema resuelto C10-01-P}


\noindent\textbf{Enunciado}
\(C10-01-P \cdot\) Incidencia durante variación de caudal. Ejercicio de transferencia. Un rotor sin prerrotación conserva velocidad periférica y ángulo metálico. Calcule su beta de entrada y la incidencia a la nueva velocidad axial. Antes de consultar la solución, prediga el efecto de cambiar el primer dato respecto del caso resuelto de esta familia. Justifique la tendencia y explique qué restricción podría convertirse en limitante. Datos: \(U=263.055\, \mathrm{m}\,\mathrm{s}^{-1}\); \(C_x=119\, \mathrm{m}\,\mathrm{s}^{-1}\); \(\kappa_e=-1.064651\, \mathrm{rad}\).

\noindent\textbf{Datos}
\begin{itemize}
\item $U = 263.055\,\mathrm{m}\,\mathrm{s}^{-1}$ — velocidad periférica
\item $C_x = 119\,\mathrm{m}\,\mathrm{s}^{-1}$ — velocidad axial nueva
\item $\kappa_e = -1.0646508437165\,\mathrm{rad}$ — ángulo metálico de entrada
\end{itemize}

\noindent\textbf{Incógnita}
$i$, en $\mathrm{rad}$; también las magnitudes intermedias desarrolladas.

\noindent\textbf{Modelo físico}

\noindent\textit{Incidencia durante variación de caudal}

\[
\fitmath{i=- \kappa_e - \operatorname{atan}{\left(\frac{U}{C_x} \right)}}
\]

\noindent\textbf{Hipótesis}
\begin{itemize}
\item Velocidad axial positiva, sin prerrotación, metal fijo.
\end{itemize}

\noindent\textbf{Desarrollo paso a paso}

\begin{enumerate}

\item El triángulo relativo utiliza tangencial -U.

\[\fitmath{\beta=\arctan\left(\frac{-\mathrm{uu}}{\mathrm{cx}}\right)}\]

\[\fitmath{\beta=\arctan\left(\frac{-(263.055)}{(119)}\right)\approx -1.145968\,\mathrm{rad}}\]

\item Se compara con la dirección del metal.

\[\fitmath{\mathrm{inc}=\arctan\left(\frac{-\mathrm{uu}}{\mathrm{cx}}\right)-\mathrm{kap}}\]

\[\fitmath{\mathrm{inc}=\arctan\left(\frac{-(263.055)}{(119)}\right)-(-1.064650844)\approx -0.0813167\,\mathrm{rad}}\]

\end{enumerate}


\noindent\textbf{Sustitución numérica}

\[
\fitmath{- \left(-1.0646508437165\,\mathrm{rad}\right) - \operatorname{atan}{\left(\frac{\left(263.055\,\mathrm{m}\,\mathrm{s}^{-1}\right)}{\left(119\,\mathrm{m}\,\mathrm{s}^{-1}\right)} \right)}}
\]

\noindent\textbf{Resultado}

\[
\fitmath{i=-8.131670125\times10^{-2}\,\mathrm{rad}}
\]

\noindent\textbf{Interpretación}
El signo corresponde exclusivamente al convenio algebraico de este manual. Debe compararse con una incidencia nominal calculada igual.

\noindent\textbf{Comprobación de los pasos}
\begin{itemize}
\item Paso 2: $-1.14597$ — reproducido
\item Paso 4: $-0.0813167$ — reproducido
\end{itemize}

\noindent\textbf{Comprobación de la respuesta} la aritmética del último paso reproduce la respuesta enunciada a la precisión con que está escrita. Es una reproducción aritmética; no valida el modelo físico ni sus hipótesis.

\noindent\textbf{Limitaciones}
\begin{itemize}
\item Modelo educativo; la aplicabilidad a una máquina real requiere revisión experta.
\end{itemize}

\noindent\textbf{Errores frecuentes}
\begin{itemize}
\item Conservar beta al cambiar Cx/U.
\end{itemize}
% studium:end-record eyJpZCI6ICJQUkItMDA3NCIsICJzZWN0aW9uIjogImMxMCJ9
\end{studiumworked}

\begin{studiumworked}{Problema resuelto C10-02-R}


\noindent\textbf{Enunciado}
\(C10-02-R \cdot\) Margen de caudal definido. A igual velocidad corregida, un punto operativo y un límite prescrito tienen los caudales indicados. Calcule el margen según la definición del capítulo. Datos: \(\dot{m}_{c,\mathrm{op}}=28.4\, \mathrm{kg}\,\mathrm{s}^{-1}\); \(\dot{m}_{c,\mathrm{lim}}=23.9\, \mathrm{kg}\,\mathrm{s}^{-1}\).

\noindent\textbf{Datos}
\begin{itemize}
\item $\dot{m}_{c,\mathrm{op}} = 28.4\,\mathrm{kg}\,\mathrm{s}^{-1}$ — caudal corregido operativo
\item $\dot{m}_{c,\mathrm{lim}} = 23.9\,\mathrm{kg}\,\mathrm{s}^{-1}$ — caudal corregido límite
\end{itemize}

\noindent\textbf{Incógnita}
$\mathrm{SM}_m$ (adimensional); también las magnitudes intermedias desarrolladas.

\noindent\textbf{Modelo físico}

\noindent\textit{Margen de caudal definido}

\[
\fitmath{\mathrm{SM}_m=\frac{- 100 \dot{m}_{c,\mathrm{lim}} + 100 \dot{m}_{c,\mathrm{op}}}{\dot{m}_{c,\mathrm{op}}}}
\]

\noindent\textbf{Hipótesis}
\begin{itemize}
\item Caudales corregidos a las mismas referencias.
\end{itemize}

\noindent\textbf{Desarrollo paso a paso}

\begin{enumerate}

\item Se calcula la separación de caudal.

\[\fitmath{\mathrm{dm}=\mathrm{op}-\mathrm{lim}}\]

\[\fitmath{\mathrm{dm}=(28.4)-(23.9)\approx 4.5\,\mathrm{kg/s}}\]

\item Se normaliza por el caudal operativo y se expresa en porcentaje.

\[\fitmath{\mathrm{SM}=\frac{100\cdot \left(\mathrm{op}-\mathrm{lim}\right)}{\mathrm{op}}}\]

\[\fitmath{\mathrm{SM}=\frac{100\cdot \left((28.4)-(23.9)\right)}{(28.4)}\approx 15.84507\,}\]

\end{enumerate}


\noindent\textbf{Sustitución numérica}

\[
\fitmath{\frac{- 100 \left(23.9\,\mathrm{kg}\,\mathrm{s}^{-1}\right) + 100 \left(28.4\,\mathrm{kg}\,\mathrm{s}^{-1}\right)}{\left(28.4\,\mathrm{kg}\,\mathrm{s}^{-1}\right)}}
\]

\noindent\textbf{Resultado}

\[
\fitmath{\mathrm{SM}_m=1.584507042\times10^{1}}
\]

\noindent\textbf{Interpretación}
El porcentaje solo es comparable con la misma definición y velocidad corregida. El límite es un dato, no una predicción del modelo.

\noindent\textbf{Comprobación de los pasos}
\begin{itemize}
\item Paso 2: $4.5$ — reproducido
\item Paso 4: $15.8451$ — reproducido
\end{itemize}

\noindent\textbf{Comprobación de la respuesta} la aritmética del último paso reproduce la respuesta enunciada a la precisión con que está escrita. Es una reproducción aritmética; no valida el modelo físico ni sus hipótesis.

\noindent\textbf{Limitaciones}
\begin{itemize}
\item Modelo educativo; la aplicabilidad a una máquina real requiere revisión experta.
\end{itemize}

\noindent\textbf{Errores frecuentes}
\begin{itemize}
\item Normalizar con el límite y llamar igual al margen resultante.
\end{itemize}
% studium:end-record eyJpZCI6ICJQUkItMDA3NSIsICJzZWN0aW9uIjogImMxMCJ9
\end{studiumworked}

\begin{studiumworked}{Problema resuelto C10-02-P}


\noindent\textbf{Enunciado}
\(C10-02-P \cdot\) Margen de caudal definido. Ejercicio de transferencia. A igual velocidad corregida, un punto operativo y un límite prescrito tienen los caudales indicados. Calcule el margen según la definición del capítulo. Antes de consultar la solución, prediga el efecto de cambiar el primer dato respecto del caso resuelto de esta familia. Justifique la tendencia y explique qué restricción podría convertirse en limitante. Datos: \(\dot{m}_{c,\mathrm{op}}=30.246\, \mathrm{kg}\,\mathrm{s}^{-1}\); \(\dot{m}_{c,\mathrm{lim}}=23.9\, \mathrm{kg}\,\mathrm{s}^{-1}\).

\noindent\textbf{Datos}
\begin{itemize}
\item $\dot{m}_{c,\mathrm{op}} = 30.246\,\mathrm{kg}\,\mathrm{s}^{-1}$ — caudal corregido operativo
\item $\dot{m}_{c,\mathrm{lim}} = 23.9\,\mathrm{kg}\,\mathrm{s}^{-1}$ — caudal corregido límite
\end{itemize}

\noindent\textbf{Incógnita}
$\mathrm{SM}_m$ (adimensional); también las magnitudes intermedias desarrolladas.

\noindent\textbf{Modelo físico}

\noindent\textit{Margen de caudal definido}

\[
\fitmath{\mathrm{SM}_m=\frac{- 100 \dot{m}_{c,\mathrm{lim}} + 100 \dot{m}_{c,\mathrm{op}}}{\dot{m}_{c,\mathrm{op}}}}
\]

\noindent\textbf{Hipótesis}
\begin{itemize}
\item Caudales corregidos a las mismas referencias.
\end{itemize}

\noindent\textbf{Desarrollo paso a paso}

\begin{enumerate}

\item Se calcula la separación de caudal.

\[\fitmath{\mathrm{dm}=\mathrm{op}-\mathrm{lim}}\]

\[\fitmath{\mathrm{dm}=(30.246)-(23.9)\approx 6.346\,\mathrm{kg/s}}\]

\item Se normaliza por el caudal operativo y se expresa en porcentaje.

\[\fitmath{\mathrm{SM}=\frac{100\cdot \left(\mathrm{op}-\mathrm{lim}\right)}{\mathrm{op}}}\]

\[\fitmath{\mathrm{SM}=\frac{100\cdot \left((30.246)-(23.9)\right)}{(30.246)}\approx 20.98129\,}\]

\end{enumerate}


\noindent\textbf{Sustitución numérica}

\[
\fitmath{\frac{- 100 \left(23.9\,\mathrm{kg}\,\mathrm{s}^{-1}\right) + 100 \left(30.246\,\mathrm{kg}\,\mathrm{s}^{-1}\right)}{\left(30.246\,\mathrm{kg}\,\mathrm{s}^{-1}\right)}}
\]

\noindent\textbf{Resultado}

\[
\fitmath{\mathrm{SM}_m=2.098128678\times10^{1}}
\]

\noindent\textbf{Interpretación}
El porcentaje solo es comparable con la misma definición y velocidad corregida. El límite es un dato, no una predicción del modelo.

\noindent\textbf{Comprobación de los pasos}
\begin{itemize}
\item Paso 2: $6.346$ — reproducido
\item Paso 4: $20.9813$ — reproducido
\end{itemize}

\noindent\textbf{Comprobación de la respuesta} la aritmética del último paso reproduce la respuesta enunciada a la precisión con que está escrita. Es una reproducción aritmética; no valida el modelo físico ni sus hipótesis.

\noindent\textbf{Limitaciones}
\begin{itemize}
\item Modelo educativo; la aplicabilidad a una máquina real requiere revisión experta.
\end{itemize}

\noindent\textbf{Errores frecuentes}
\begin{itemize}
\item Normalizar con el límite y llamar igual al margen resultante.
\end{itemize}
% studium:end-record eyJpZCI6ICJQUkItMDA3NiIsICJzZWN0aW9uIjogImMxMCJ9
\end{studiumworked}

\begin{studiumworked}{Problema resuelto C10-03-R}


\noindent\textbf{Enunciado}
\(C10-03-R \cdot\) Corrección de velocidad. Halle la velocidad corregida y la velocidad periférica física al radio dado. La incógnita final es U real. Datos: \(N=196.6667\, \mathrm{s}^{-1}\); \(T_{0e}=259\, \mathrm{K}\); \(T_{\mathrm{ref}}=288.15\, \mathrm{K}\); \(r=0.276\, \mathrm{m}\).

\noindent\textbf{Datos}
\begin{itemize}
\item $N = 196.66666666667\,\mathrm{s}^{-1}$ — giro físico en revoluciones por segundo; equivale a 11800 rpm
\item $T_{0e} = 259\,\mathrm{K}$ — temperatura total entrada
\item $T_{\mathrm{ref}} = 288.15\,\mathrm{K}$ — temperatura referencia
\item $r = 0.276\,\mathrm{m}$ — radio
\end{itemize}

\noindent\textbf{Incógnita}
$U$, en $\mathrm{m\,s^{-1}}$; también las magnitudes intermedias desarrolladas.

\noindent\textbf{Modelo físico}

\noindent\textit{Corrección de velocidad}

\[
\fitmath{U=2 \pi N r}
\]

\noindent\textbf{Hipótesis}
\begin{itemize}
\item Giro en revoluciones por segundo y radio fijo.
\end{itemize}

\noindent\textbf{Desarrollo paso a paso}

\begin{enumerate}

\item La corrección compara velocidades a la referencia térmica.

\[\fitmath{\mathrm{Nc}=\frac{\mathrm{nn}}{\sqrt{\frac{\mathrm{Tin}}{\mathrm{Tr}}}}}\]

\[\fitmath{\mathrm{Nc}=\frac{(196.6666667)}{\sqrt{\frac{(259)}{(288.15)}}}\approx 207.4389\,\mathrm{s^{-1}}}\]

\item El giro en revoluciones por segundo se convierte a velocidad angular para U.

\[\fitmath{U=2\cdot \pi\cdot \mathrm{nn}\cdot \mathrm{rad}}\]

\[\fitmath{U=2\cdot \pi\cdot (196.6666667)\cdot (0.276)\approx 341.0513\,\mathrm{m/s}}\]

\end{enumerate}


\noindent\textbf{Sustitución numérica}

\[
\fitmath{2 \pi \left(196.66666666667\,\mathrm{s}^{-1}\right) \left(0.276\,\mathrm{m}\right)}
\]

\noindent\textbf{Resultado}

\[
\fitmath{U=3.410512985\times10^{2}\,\mathrm{m}\,\mathrm{s}^{-1}}
\]

\noindent\textbf{Interpretación}
El dato se expresa en revoluciones por segundo; U se calcula con N físico, no con N corregido.

\noindent\textbf{Comprobación de los pasos}
\begin{itemize}
\item Paso 2: $207.439$ — reproducido
\item Paso 4: $341.051$ — reproducido
\end{itemize}

\noindent\textbf{Comprobación de la respuesta} la aritmética del último paso reproduce la respuesta enunciada a la precisión con que está escrita. Es una reproducción aritmética; no valida el modelo físico ni sus hipótesis.

\noindent\textbf{Limitaciones}
\begin{itemize}
\item Modelo educativo; la aplicabilidad a una máquina real requiere revisión experta.
\end{itemize}

\noindent\textbf{Errores frecuentes}
\begin{itemize}
\item Usar N corregido en una velocidad mecánica real.
\end{itemize}
% studium:end-record eyJpZCI6ICJQUkItMDA3NyIsICJzZWN0aW9uIjogImMxMCJ9
\end{studiumworked}

\begin{studiumworked}{Problema resuelto C10-03-P}


\noindent\textbf{Enunciado}
\(C10-03-P \cdot\) Corrección de velocidad. Ejercicio de transferencia. Halle la velocidad corregida y la velocidad periférica física al radio dado. La incógnita final es U real. Antes de consultar la solución, prediga el efecto de cambiar el primer dato respecto del caso resuelto de esta familia. Justifique la tendencia y explique qué restricción podría convertirse en limitante. Datos: \(N=209.45\, \mathrm{s}^{-1}\); \(T_{0e}=259\, \mathrm{K}\); \(T_{\mathrm{ref}}=288.15\, \mathrm{K}\); \(r=0.276\, \mathrm{m}\).

\noindent\textbf{Datos}
\begin{itemize}
\item $N = 209.45\,\mathrm{s}^{-1}$ — giro físico en revoluciones por segundo; equivale a 12567 rpm
\item $T_{0e} = 259\,\mathrm{K}$ — temperatura total entrada
\item $T_{\mathrm{ref}} = 288.15\,\mathrm{K}$ — temperatura referencia
\item $r = 0.276\,\mathrm{m}$ — radio
\end{itemize}

\noindent\textbf{Incógnita}
$U$, en $\mathrm{m\,s^{-1}}$; también las magnitudes intermedias desarrolladas.

\noindent\textbf{Modelo físico}

\noindent\textit{Corrección de velocidad}

\[
\fitmath{U=2 \pi N r}
\]

\noindent\textbf{Hipótesis}
\begin{itemize}
\item Giro en revoluciones por segundo y radio fijo.
\end{itemize}

\noindent\textbf{Desarrollo paso a paso}

\begin{enumerate}

\item La corrección compara velocidades a la referencia térmica.

\[\fitmath{\mathrm{Nc}=\frac{\mathrm{nn}}{\sqrt{\frac{\mathrm{Tin}}{\mathrm{Tr}}}}}\]

\[\fitmath{\mathrm{Nc}=\frac{(209.45)}{\sqrt{\frac{(259)}{(288.15)}}}\approx 220.9224\,\mathrm{s^{-1}}}\]

\item El giro en revoluciones por segundo se convierte a velocidad angular para U.

\[\fitmath{U=2\cdot \pi\cdot \mathrm{nn}\cdot \mathrm{rad}}\]

\[\fitmath{U=2\cdot \pi\cdot (209.45)\cdot (0.276)\approx 363.2196\,\mathrm{m/s}}\]

\end{enumerate}


\noindent\textbf{Sustitución numérica}

\[
\fitmath{2 \pi \left(209.45\,\mathrm{s}^{-1}\right) \left(0.276\,\mathrm{m}\right)}
\]

\noindent\textbf{Resultado}

\[
\fitmath{U=3.632196329\times10^{2}\,\mathrm{m}\,\mathrm{s}^{-1}}
\]

\noindent\textbf{Interpretación}
El dato se expresa en revoluciones por segundo; U se calcula con N físico, no con N corregido.

\noindent\textbf{Comprobación de los pasos}
\begin{itemize}
\item Paso 2: $220.922$ — reproducido
\item Paso 4: $363.22$ — reproducido
\end{itemize}

\noindent\textbf{Comprobación de la respuesta} la aritmética del último paso reproduce la respuesta enunciada a la precisión con que está escrita. Es una reproducción aritmética; no valida el modelo físico ni sus hipótesis.

\noindent\textbf{Limitaciones}
\begin{itemize}
\item Modelo educativo; la aplicabilidad a una máquina real requiere revisión experta.
\end{itemize}

\noindent\textbf{Errores frecuentes}
\begin{itemize}
\item Usar N corregido en una velocidad mecánica real.
\end{itemize}
% studium:end-record eyJpZCI6ICJQUkItMDIwOCIsICJzZWN0aW9uIjogImMxMCJ9
\end{studiumworked}

\begin{studiumworked}{Problema resuelto C10-04-R}


\noindent\textbf{Enunciado}
\(C10-04-R \cdot\) Estabilidad lineal del modelo ampliado. Para el modelo de dos estados del capítulo, calcule traza y determinante de la matriz lineal. El resultado final es el determinante; interprete estabilidad. Datos: \(a_c=-1800\, \mathrm{Pa}\,\mathrm{s}\,\mathrm{kg}^{-1}\); \(b_v=0.00016\, \mathrm{kg}\,\mathrm{s}^{-1}\,\mathrm{Pa}^{-1}\); \(L_c=260\, \mathrm{Pa}\,\mathrm{s}^{2}\,\mathrm{kg}^{-1}\); \(K_p=0.0018\, \mathrm{kg}\,\mathrm{Pa}^{-1}\).

\noindent\textbf{Datos}
\begin{itemize}
\item $a_c = -1800\,\mathrm{Pa}\,\mathrm{s}\,\mathrm{kg}^{-1}$ — pendiente de compresor
\item $b_v = 0.00016\,\mathrm{kg}\,\mathrm{s}^{-1}\,\mathrm{Pa}^{-1}$ — pendiente de descarga
\item $L_c = 260\,\mathrm{Pa}\,\mathrm{s}^{2}\,\mathrm{kg}^{-1}$ — inercia del modelo
\item $K_p = 0.0018\,\mathrm{kg}\,\mathrm{Pa}^{-1}$ — capacidad del modelo
\end{itemize}

\noindent\textbf{Incógnita}
$\det(J)$, en $\mathrm{s^{-2}}$; también las magnitudes intermedias desarrolladas.

\noindent\textbf{Modelo físico}

\noindent\textit{Estabilidad lineal del modelo ampliado}

\[
\fitmath{\det(J)=\frac{- a_c b_v + 1}{K_p L_c}}
\]

\noindent\textbf{Hipótesis}
\begin{itemize}
\item Linealización local; parámetros positivos L,K; modelo pedagógico sin retardo.
\end{itemize}

\noindent\textbf{Desarrollo paso a paso}

\begin{enumerate}

\item La traza suma los términos diagonales.

\[\fitmath{\mathrm{tr}=\frac{\mathrm{aa}}{\mathrm{ll}}-\frac{\mathrm{bb}}{\mathrm{kk}}}\]

\[\fitmath{\mathrm{tr}=\frac{(-1800)}{(260)}-\frac{(0.00016)}{(0.0018)}\approx -7.011966\,\mathrm{s^{-1}}}\]

\item El determinante incluye el acoplamiento cruzado.

\[\fitmath{\mathrm{det}=\frac{1-\mathrm{aa}\cdot \mathrm{bb}}{\mathrm{ll}\cdot \mathrm{kk}}}\]

\[\fitmath{\mathrm{det}=\frac{1-(-1800)\cdot (0.00016)}{(260)\cdot (0.0018)}\approx 2.752137\,\mathrm{s^{-2}}}\]

\end{enumerate}


\noindent\textbf{Sustitución numérica}

\[
\fitmath{\frac{- \left(-1800\,\mathrm{Pa}\,\mathrm{s}\,\mathrm{kg}^{-1}\right) \left(0.00016\,\mathrm{kg}\,\mathrm{s}^{-1}\,\mathrm{Pa}^{-1}\right) + 1}{\left(0.0018\,\mathrm{kg}\,\mathrm{Pa}^{-1}\right) \left(260\,\mathrm{Pa}\,\mathrm{s}^{2}\,\mathrm{kg}^{-1}\right)}}
\]

\noindent\textbf{Resultado}

\[
\fitmath{\det(J)=2.752136752\times10^{0}\,\mathrm{s}^{-2}}
\]

\noindent\textbf{Interpretación}
Traza negativa y determinante positivo dan estabilidad lineal del modelo de dos estados. Esto no predice pérdida rotativa ni estabilidad de un motor real.

\noindent\textbf{Comprobación de los pasos}
\begin{itemize}
\item Paso 2: $-7.01197$ — reproducido
\item Paso 4: $2.75214$ — reproducido
\end{itemize}

\noindent\textbf{Comprobación de la respuesta} la aritmética del último paso reproduce la respuesta enunciada a la precisión con que está escrita. Es una reproducción aritmética; no valida el modelo físico ni sus hipótesis.

\noindent\textbf{Limitaciones}
\begin{itemize}
\item Modelo educativo; la aplicabilidad a una máquina real requiere revisión experta.
\end{itemize}

\noindent\textbf{Errores frecuentes}
\begin{itemize}
\item Tomar L como longitud sin unidades equivalentes.
\end{itemize}
% studium:end-record eyJpZCI6ICJQUkItMDA3OSIsICJzZWN0aW9uIjogImMxMCJ9
\end{studiumworked}

\begin{studiumworked}{Problema resuelto C10-04-P}


\noindent\textbf{Enunciado}
\(C10-04-P \cdot\) Estabilidad lineal del modelo ampliado. Ejercicio de transferencia. Para el modelo de dos estados del capítulo, calcule traza y determinante de la matriz lineal. El resultado final es el determinante; interprete estabilidad. Antes de consultar la solución, prediga el efecto de cambiar el primer dato respecto del caso resuelto de esta familia. Justifique la tendencia y explique qué restricción podría convertirse en limitante. Datos: \(a_c=-1917\, \mathrm{Pa}\,\mathrm{s}\,\mathrm{kg}^{-1}\); \(b_v=0.00016\, \mathrm{kg}\,\mathrm{s}^{-1}\,\mathrm{Pa}^{-1}\); \(L_c=260\, \mathrm{Pa}\,\mathrm{s}^{2}\,\mathrm{kg}^{-1}\); \(K_p=0.0018\, \mathrm{kg}\,\mathrm{Pa}^{-1}\).

\noindent\textbf{Datos}
\begin{itemize}
\item $a_c = -1917\,\mathrm{Pa}\,\mathrm{s}\,\mathrm{kg}^{-1}$ — pendiente de compresor
\item $b_v = 0.00016\,\mathrm{kg}\,\mathrm{s}^{-1}\,\mathrm{Pa}^{-1}$ — pendiente de descarga
\item $L_c = 260\,\mathrm{Pa}\,\mathrm{s}^{2}\,\mathrm{kg}^{-1}$ — inercia del modelo
\item $K_p = 0.0018\,\mathrm{kg}\,\mathrm{Pa}^{-1}$ — capacidad del modelo
\end{itemize}

\noindent\textbf{Incógnita}
$\det(J)$, en $\mathrm{s^{-2}}$; también las magnitudes intermedias desarrolladas.

\noindent\textbf{Modelo físico}

\noindent\textit{Estabilidad lineal del modelo ampliado}

\[
\fitmath{\det(J)=\frac{- a_c b_v + 1}{K_p L_c}}
\]

\noindent\textbf{Hipótesis}
\begin{itemize}
\item Linealización local; parámetros positivos L,K; modelo pedagógico sin retardo.
\end{itemize}

\noindent\textbf{Desarrollo paso a paso}

\begin{enumerate}

\item La traza suma los términos diagonales.

\[\fitmath{\mathrm{tr}=\frac{\mathrm{aa}}{\mathrm{ll}}-\frac{\mathrm{bb}}{\mathrm{kk}}}\]

\[\fitmath{\mathrm{tr}=\frac{(-1917)}{(260)}-\frac{(0.00016)}{(0.0018)}\approx -7.461966\,\mathrm{s^{-1}}}\]

\item El determinante incluye el acoplamiento cruzado.

\[\fitmath{\mathrm{det}=\frac{1-\mathrm{aa}\cdot \mathrm{bb}}{\mathrm{ll}\cdot \mathrm{kk}}}\]

\[\fitmath{\mathrm{det}=\frac{1-(-1917)\cdot (0.00016)}{(260)\cdot (0.0018)}\approx 2.792137\,\mathrm{s^{-2}}}\]

\end{enumerate}


\noindent\textbf{Sustitución numérica}

\[
\fitmath{\frac{- \left(-1917\,\mathrm{Pa}\,\mathrm{s}\,\mathrm{kg}^{-1}\right) \left(0.00016\,\mathrm{kg}\,\mathrm{s}^{-1}\,\mathrm{Pa}^{-1}\right) + 1}{\left(0.0018\,\mathrm{kg}\,\mathrm{Pa}^{-1}\right) \left(260\,\mathrm{Pa}\,\mathrm{s}^{2}\,\mathrm{kg}^{-1}\right)}}
\]

\noindent\textbf{Resultado}

\[
\fitmath{\det(J)=2.792136752\times10^{0}\,\mathrm{s}^{-2}}
\]

\noindent\textbf{Interpretación}
Traza negativa y determinante positivo dan estabilidad lineal del modelo de dos estados. Esto no predice pérdida rotativa ni estabilidad de un motor real.

\noindent\textbf{Comprobación de los pasos}
\begin{itemize}
\item Paso 2: $-7.46197$ — reproducido
\item Paso 4: $2.79214$ — reproducido
\end{itemize}

\noindent\textbf{Comprobación de la respuesta} la aritmética del último paso reproduce la respuesta enunciada a la precisión con que está escrita. Es una reproducción aritmética; no valida el modelo físico ni sus hipótesis.

\noindent\textbf{Limitaciones}
\begin{itemize}
\item Modelo educativo; la aplicabilidad a una máquina real requiere revisión experta.
\end{itemize}

\noindent\textbf{Errores frecuentes}
\begin{itemize}
\item Tomar L como longitud sin unidades equivalentes.
\end{itemize}
% studium:end-record eyJpZCI6ICJQUkItMDA4MCIsICJzZWN0aW9uIjogImMxMCJ9
\end{studiumworked}

\chapter{Equilibrio radial y leyes de remolino}
\begin{studiumworked}{Problema resuelto C11-01-R}


\noindent\textbf{Enunciado}
\(C11-01-R \cdot\) Presión de vórtice libre. Integre el equilibrio radial de un fluido de densidad constante y remolino K/r entre radio medio y punta. Datos: \(\rho=2.3\, \mathrm{kg}\,\mathrm{m}^{-3}\); \(K_v=36\, \mathrm{m}^{2}\,\mathrm{s}^{-1}\); \(r_m=0.29\, \mathrm{m}\); \(r_t=0.36\, \mathrm{m}\).

\noindent\textbf{Datos}
\begin{itemize}
\item $\rho = 2.3\,\mathrm{kg}\,\mathrm{m}^{-3}$ — densidad uniforme
\item $K_v = 36\,\mathrm{m}^{2}\,\mathrm{s}^{-1}$ — constante de remolino
\item $r_m = 0.29\,\mathrm{m}$ — radio referencia
\item $r_t = 0.36\,\mathrm{m}$ — radio exterior
\end{itemize}

\noindent\textbf{Incógnita}
$\Delta p$, en $\mathrm{Pa}$; también las magnitudes intermedias desarrolladas.

\noindent\textbf{Modelo físico}

\noindent\textit{Presión de vórtice libre}

\[
\fitmath{\Delta p=\frac{K_v^{2} \rho \left(- \frac{1}{r_t^{2}} + \frac{1}{r_m^{2}}\right)}{2}}
\]

\noindent\textbf{Hipótesis}
\begin{itemize}
\item Equilibrio radial simple, axisimetría, densidad constante.
\end{itemize}

\noindent\textbf{Desarrollo paso a paso}

\begin{enumerate}

\item El momento angular constante fija Ctheta en la punta.

\[\fitmath{\mathrm{ct}=\frac{\mathrm{kv}}{\mathrm{rt}}}\]

\[\fitmath{\mathrm{ct}=\frac{(36)}{(0.36)}\approx 100\,\mathrm{m/s}}\]

\item Se integra el gradiente positivo de presión.

\[\fitmath{\mathrm{dp}=\frac{\rho\cdot {\mathrm{kv}}^{2}\cdot \left(\frac{1}{{\mathrm{rm}}^{2}}-\frac{1}{{\mathrm{rt}}^{2}}\right)}{2}}\]

\[\fitmath{\mathrm{dp}=\frac{(2.3)\cdot {(36)}^{2}\cdot \left(\frac{1}{{(0.29)}^{2}}-\frac{1}{{(0.36)}^{2}}\right)}{2}\approx 6221.76\,\mathrm{Pa}}\]

\end{enumerate}


\noindent\textbf{Sustitución numérica}

\[
\fitmath{\frac{\left(\left(36\,\mathrm{m}^{2}\,\mathrm{s}^{-1}\right)\right)^{2} \left(2.3\,\mathrm{kg}\,\mathrm{m}^{-3}\right) \left(- \frac{1}{\left(0.36\,\mathrm{m}\right)^{2}} + \frac{1}{\left(0.29\,\mathrm{m}\right)^{2}}\right)}{2}}
\]

\noindent\textbf{Resultado}

\[
\fitmath{\Delta p=6.221759810\times10^{3}\,\mathrm{Pa}}
\]

\noindent\textbf{Interpretación}
La presión exterior aumenta. La aproximación de densidad constante debe justificarse antes de aplicarla a un gas real.

\noindent\textbf{Comprobación de los pasos}
\begin{itemize}
\item Paso 2: $100$ — reproducido
\item Paso 4: $6221.76$ — reproducido
\end{itemize}

\noindent\textbf{Comprobación de la respuesta} la aritmética del último paso reproduce la respuesta enunciada a la precisión con que está escrita. Es una reproducción aritmética; no valida el modelo físico ni sus hipótesis.

\noindent\textbf{Limitaciones}
\begin{itemize}
\item Modelo educativo; la aplicabilidad a una máquina real requiere revisión experta.
\end{itemize}

\noindent\textbf{Errores frecuentes}
\begin{itemize}
\item Invertir los radios y obtener caída de presión hacia fuera.
\end{itemize}
% studium:end-record eyJpZCI6ICJQUkItMDA4MSIsICJzZWN0aW9uIjogImMxMSJ9
\end{studiumworked}

\begin{studiumworked}{Problema resuelto C11-01-P}


\noindent\textbf{Enunciado}
\(C11-01-P \cdot\) Presión de vórtice libre. Ejercicio de transferencia. Integre el equilibrio radial de un fluido de densidad constante y remolino K/r entre radio medio y punta. Antes de consultar la solución, prediga el efecto de cambiar el primer dato respecto del caso resuelto de esta familia. Justifique la tendencia y explique qué restricción podría convertirse en limitante. Datos: \(\rho=2.4495\, \mathrm{kg}\,\mathrm{m}^{-3}\); \(K_v=36\, \mathrm{m}^{2}\,\mathrm{s}^{-1}\); \(r_m=0.29\, \mathrm{m}\); \(r_t=0.36\, \mathrm{m}\).

\noindent\textbf{Datos}
\begin{itemize}
\item $\rho = 2.4495\,\mathrm{kg}\,\mathrm{m}^{-3}$ — densidad uniforme
\item $K_v = 36\,\mathrm{m}^{2}\,\mathrm{s}^{-1}$ — constante de remolino
\item $r_m = 0.29\,\mathrm{m}$ — radio referencia
\item $r_t = 0.36\,\mathrm{m}$ — radio exterior
\end{itemize}

\noindent\textbf{Incógnita}
$\Delta p$, en $\mathrm{Pa}$; también las magnitudes intermedias desarrolladas.

\noindent\textbf{Modelo físico}

\noindent\textit{Presión de vórtice libre}

\[
\fitmath{\Delta p=\frac{K_v^{2} \rho \left(- \frac{1}{r_t^{2}} + \frac{1}{r_m^{2}}\right)}{2}}
\]

\noindent\textbf{Hipótesis}
\begin{itemize}
\item Equilibrio radial simple, axisimetría, densidad constante.
\end{itemize}

\noindent\textbf{Desarrollo paso a paso}

\begin{enumerate}

\item El momento angular constante fija Ctheta en la punta.

\[\fitmath{\mathrm{ct}=\frac{\mathrm{kv}}{\mathrm{rt}}}\]

\[\fitmath{\mathrm{ct}=\frac{(36)}{(0.36)}\approx 100\,\mathrm{m/s}}\]

\item Se integra el gradiente positivo de presión.

\[\fitmath{\mathrm{dp}=\frac{\rho\cdot {\mathrm{kv}}^{2}\cdot \left(\frac{1}{{\mathrm{rm}}^{2}}-\frac{1}{{\mathrm{rt}}^{2}}\right)}{2}}\]

\[\fitmath{\mathrm{dp}=\frac{(2.4495)\cdot {(36)}^{2}\cdot \left(\frac{1}{{(0.29)}^{2}}-\frac{1}{{(0.36)}^{2}}\right)}{2}\approx 6626.174\,\mathrm{Pa}}\]

\end{enumerate}


\noindent\textbf{Sustitución numérica}

\[
\fitmath{\frac{\left(\left(36\,\mathrm{m}^{2}\,\mathrm{s}^{-1}\right)\right)^{2} \left(2.4495\,\mathrm{kg}\,\mathrm{m}^{-3}\right) \left(- \frac{1}{\left(0.36\,\mathrm{m}\right)^{2}} + \frac{1}{\left(0.29\,\mathrm{m}\right)^{2}}\right)}{2}}
\]

\noindent\textbf{Resultado}

\[
\fitmath{\Delta p=6.626174197\times10^{3}\,\mathrm{Pa}}
\]

\noindent\textbf{Interpretación}
La presión exterior aumenta. La aproximación de densidad constante debe justificarse antes de aplicarla a un gas real.

\noindent\textbf{Comprobación de los pasos}
\begin{itemize}
\item Paso 2: $100$ — reproducido
\item Paso 4: $6626.17$ — reproducido
\end{itemize}

\noindent\textbf{Comprobación de la respuesta} la aritmética del último paso reproduce la respuesta enunciada a la precisión con que está escrita. Es una reproducción aritmética; no valida el modelo físico ni sus hipótesis.

\noindent\textbf{Limitaciones}
\begin{itemize}
\item Modelo educativo; la aplicabilidad a una máquina real requiere revisión experta.
\end{itemize}

\noindent\textbf{Errores frecuentes}
\begin{itemize}
\item Invertir los radios y obtener caída de presión hacia fuera.
\end{itemize}
% studium:end-record eyJpZCI6ICJQUkItMDA4MiIsICJzZWN0aW9uIjogImMxMSJ9
\end{studiumworked}

\begin{studiumworked}{Problema resuelto C11-02-R}


\noindent\textbf{Enunciado}
\(C11-02-R \cdot\) Trabajo uniforme con vórtice libre. Una fila móvil tiene constantes de vórtice libre distintas en entrada y salida. Calcule el trabajo y el incremento de temperatura total. Datos: \(\Omega=1080\, \mathrm{s}^{-1}\); \(K_e=11\, \mathrm{m}^{2}\,\mathrm{s}^{-1}\); \(K_s=39\, \mathrm{m}^{2}\,\mathrm{s}^{-1}\); \(c_p=1004.5\, \mathrm{J}\,\mathrm{kg}^{-1}\,\mathrm{K}^{-1}\).

\noindent\textbf{Datos}
\begin{itemize}
\item $\Omega = 1080\,\mathrm{s}^{-1}$ — velocidad angular
\item $K_e = 11\,\mathrm{m}^{2}\,\mathrm{s}^{-1}$ — constante entrada
\item $K_s = 39\,\mathrm{m}^{2}\,\mathrm{s}^{-1}$ — constante salida
\item $c_p = 1004.5\,\mathrm{J}\,\mathrm{kg}^{-1}\,\mathrm{K}^{-1}$ — calor específico
\end{itemize}

\noindent\textbf{Incógnita}
$\Delta T_0$, en $\mathrm{K}$; también las magnitudes intermedias desarrolladas.

\noindent\textbf{Modelo físico}

\noindent\textit{Trabajo uniforme con vórtice libre}

\[
\fitmath{\Delta T_0=\frac{\Omega \left(- K_e + K_s\right)}{c_p}}
\]

\noindent\textbf{Hipótesis}
\begin{itemize}
\item Misma velocidad angular y líneas de corriente del modelo ideal.
\end{itemize}

\noindent\textbf{Desarrollo paso a paso}

\begin{enumerate}

\item Euler cancela la dependencia radial r.

\[\fitmath{\mathrm{wc}=\Omega\cdot \left(\mathrm{ks}-\mathrm{ke}\right)}\]

\[\fitmath{\mathrm{wc}=(1080)\cdot \left((39)-(11)\right)\approx 30240\,\mathrm{J/kg}}\]

\item Se convierte el trabajo en incremento total.

\[\fitmath{\mathrm{dT}=\frac{\Omega\cdot \left(\mathrm{ks}-\mathrm{ke}\right)}{\mathrm{cp}}}\]

\[\fitmath{\mathrm{dT}=\frac{(1080)\cdot \left((39)-(11)\right)}{(1004.5)}\approx 30.10453\,\mathrm{K}}\]

\end{enumerate}


\noindent\textbf{Sustitución numérica}

\[
\fitmath{\frac{\left(1080\,\mathrm{s}^{-1}\right) \left(- \left(11\,\mathrm{m}^{2}\,\mathrm{s}^{-1}\right) + \left(39\,\mathrm{m}^{2}\,\mathrm{s}^{-1}\right)\right)}{\left(1004.5\,\mathrm{J}\,\mathrm{kg}^{-1}\,\mathrm{K}^{-1}\right)}}
\]

\noindent\textbf{Resultado}

\[
\fitmath{\Delta T_0=3.010452962\times10^{1}\,\mathrm{K}}
\]

\noindent\textbf{Interpretación}
El trabajo es uniforme con el radio, aunque la reacción no lo sea.

\noindent\textbf{Comprobación de los pasos}
\begin{itemize}
\item Paso 2: $30240$ — reproducido
\item Paso 4: $30.1045$ — reproducido
\end{itemize}

\noindent\textbf{Comprobación de la respuesta} la aritmética del último paso reproduce la respuesta enunciada a la precisión con que está escrita. Es una reproducción aritmética; no valida el modelo físico ni sus hipótesis.

\noindent\textbf{Limitaciones}
\begin{itemize}
\item Modelo educativo; la aplicabilidad a una máquina real requiere revisión experta.
\end{itemize}

\noindent\textbf{Errores frecuentes}
\begin{itemize}
\item Usar igual constante de vórtice a ambos lados de un rotor con trabajo.
\end{itemize}
% studium:end-record eyJpZCI6ICJQUkItMDA4MyIsICJzZWN0aW9uIjogImMxMSJ9
\end{studiumworked}

\begin{studiumworked}{Problema resuelto C11-02-P}


\noindent\textbf{Enunciado}
\(C11-02-P \cdot\) Trabajo uniforme con vórtice libre. Ejercicio de transferencia. Una fila móvil tiene constantes de vórtice libre distintas en entrada y salida. Calcule el trabajo y el incremento de temperatura total. Antes de consultar la solución, prediga el efecto de cambiar el primer dato respecto del caso resuelto de esta familia. Justifique la tendencia y explique qué restricción podría convertirse en limitante. Datos: \(\Omega=1150.2\, \mathrm{s}^{-1}\); \(K_e=11\, \mathrm{m}^{2}\,\mathrm{s}^{-1}\); \(K_s=39\, \mathrm{m}^{2}\,\mathrm{s}^{-1}\); \(c_p=1004.5\, \mathrm{J}\,\mathrm{kg}^{-1}\,\mathrm{K}^{-1}\).

\noindent\textbf{Datos}
\begin{itemize}
\item $\Omega = 1150.2\,\mathrm{s}^{-1}$ — velocidad angular
\item $K_e = 11\,\mathrm{m}^{2}\,\mathrm{s}^{-1}$ — constante entrada
\item $K_s = 39\,\mathrm{m}^{2}\,\mathrm{s}^{-1}$ — constante salida
\item $c_p = 1004.5\,\mathrm{J}\,\mathrm{kg}^{-1}\,\mathrm{K}^{-1}$ — calor específico
\end{itemize}

\noindent\textbf{Incógnita}
$\Delta T_0$, en $\mathrm{K}$; también las magnitudes intermedias desarrolladas.

\noindent\textbf{Modelo físico}

\noindent\textit{Trabajo uniforme con vórtice libre}

\[
\fitmath{\Delta T_0=\frac{\Omega \left(- K_e + K_s\right)}{c_p}}
\]

\noindent\textbf{Hipótesis}
\begin{itemize}
\item Misma velocidad angular y líneas de corriente del modelo ideal.
\end{itemize}

\noindent\textbf{Desarrollo paso a paso}

\begin{enumerate}

\item Euler cancela la dependencia radial r.

\[\fitmath{\mathrm{wc}=\Omega\cdot \left(\mathrm{ks}-\mathrm{ke}\right)}\]

\[\fitmath{\mathrm{wc}=(1150.2)\cdot \left((39)-(11)\right)\approx 32205.6\,\mathrm{J/kg}}\]

\item Se convierte el trabajo en incremento total.

\[\fitmath{\mathrm{dT}=\frac{\Omega\cdot \left(\mathrm{ks}-\mathrm{ke}\right)}{\mathrm{cp}}}\]

\[\fitmath{\mathrm{dT}=\frac{(1150.2)\cdot \left((39)-(11)\right)}{(1004.5)}\approx 32.06132\,\mathrm{K}}\]

\end{enumerate}


\noindent\textbf{Sustitución numérica}

\[
\fitmath{\frac{\left(1150.2\,\mathrm{s}^{-1}\right) \left(- \left(11\,\mathrm{m}^{2}\,\mathrm{s}^{-1}\right) + \left(39\,\mathrm{m}^{2}\,\mathrm{s}^{-1}\right)\right)}{\left(1004.5\,\mathrm{J}\,\mathrm{kg}^{-1}\,\mathrm{K}^{-1}\right)}}
\]

\noindent\textbf{Resultado}

\[
\fitmath{\Delta T_0=3.206132404\times10^{1}\,\mathrm{K}}
\]

\noindent\textbf{Interpretación}
El trabajo es uniforme con el radio, aunque la reacción no lo sea.

\noindent\textbf{Comprobación de los pasos}
\begin{itemize}
\item Paso 2: $32205.6$ — reproducido
\item Paso 4: $32.0613$ — reproducido
\end{itemize}

\noindent\textbf{Comprobación de la respuesta} la aritmética del último paso reproduce la respuesta enunciada a la precisión con que está escrita. Es una reproducción aritmética; no valida el modelo físico ni sus hipótesis.

\noindent\textbf{Limitaciones}
\begin{itemize}
\item Modelo educativo; la aplicabilidad a una máquina real requiere revisión experta.
\end{itemize}

\noindent\textbf{Errores frecuentes}
\begin{itemize}
\item Usar igual constante de vórtice a ambos lados de un rotor con trabajo.
\end{itemize}
% studium:end-record eyJpZCI6ICJQUkItMDA4NCIsICJzZWN0aW9uIjogImMxMSJ9
\end{studiumworked}

\begin{studiumworked}{Problema resuelto C11-03-R}


\noindent\textbf{Enunciado}
\(C11-03-R \cdot\) Reacción hacia la raíz. Calcule la reacción reducida en la raíz para un escalón repetitivo de vórtice libre. Datos: \(K_e=16\, \mathrm{m}^{2}\,\mathrm{s}^{-1}\); \(K_s=49\, \mathrm{m}^{2}\,\mathrm{s}^{-1}\); \(\Omega=1120\, \mathrm{s}^{-1}\); \(r_h=0.205\, \mathrm{m}\).

\noindent\textbf{Datos}
\begin{itemize}
\item $K_e = 16\,\mathrm{m}^{2}\,\mathrm{s}^{-1}$ — constante tangencial entrada
\item $K_s = 49\,\mathrm{m}^{2}\,\mathrm{s}^{-1}$ — constante tangencial salida
\item $\Omega = 1120\,\mathrm{s}^{-1}$ — velocidad angular
\item $r_h = 0.205\,\mathrm{m}$ — radio raíz
\end{itemize}

\noindent\textbf{Incógnita}
$\mathcal{R}$ (adimensional); también las magnitudes intermedias desarrolladas.

\noindent\textbf{Modelo físico}

\noindent\textit{Reacción hacia la raíz}

\[
\fitmath{\mathcal R=1 - \frac{K_e + K_s}{2 \Omega r_h^{2}}}
\]

\noindent\textbf{Hipótesis}
\begin{itemize}
\item Criterio reducido válido para radio y Cx constantes en cada escalón repetitivo.
\end{itemize}

\noindent\textbf{Desarrollo paso a paso}

\begin{enumerate}

\item Se calcula la suma de componentes tangenciales en la raíz.

\[\fitmath{\mathrm{sum}=\frac{\mathrm{ke}+\mathrm{ks}}{\mathrm{rh}}}\]

\[\fitmath{\mathrm{sum}=\frac{(16)+(49)}{(0.205)}\approx 317.0732\,\mathrm{m/s}}\]

\item Se normaliza con dos veces U local.

\[\fitmath{\mathrm{rea}=1-\frac{\mathrm{ke}+\mathrm{ks}}{2\cdot \Omega\cdot {\mathrm{rh}}^{2}}}\]

\[\fitmath{\mathrm{rea}=1-\frac{(16)+(49)}{2\cdot (1120)\cdot {(0.205)}^{2}}\approx 0.3095096\,}\]

\end{enumerate}


\noindent\textbf{Sustitución numérica}

\[
\fitmath{1 - \frac{\left(16\,\mathrm{m}^{2}\,\mathrm{s}^{-1}\right) + \left(49\,\mathrm{m}^{2}\,\mathrm{s}^{-1}\right)}{2 \left(1120\,\mathrm{s}^{-1}\right) \left(0.205\,\mathrm{m}\right)^{2}}}
\]

\noindent\textbf{Resultado}

\[
\fitmath{\mathcal R=3.095096456\times10^{-1}}
\]

\noindent\textbf{Interpretación}
La raíz puede tener reacción mucho menor que el radio medio. Si resulta negativa se debe revisar el diseño, no cambiar el signo.

\noindent\textbf{Comprobación de los pasos}
\begin{itemize}
\item Paso 2: $317.073$ — reproducido
\item Paso 4: $0.30951$ — reproducido
\end{itemize}

\noindent\textbf{Comprobación de la respuesta} la aritmética del último paso reproduce la respuesta enunciada a la precisión con que está escrita. Es una reproducción aritmética; no valida el modelo físico ni sus hipótesis.

\noindent\textbf{Limitaciones}
\begin{itemize}
\item Modelo educativo; la aplicabilidad a una máquina real requiere revisión experta.
\end{itemize}

\noindent\textbf{Errores frecuentes}
\begin{itemize}
\item Asignar a toda la pala la reacción del radio medio.
\end{itemize}
% studium:end-record eyJpZCI6ICJQUkItMDA4NSIsICJzZWN0aW9uIjogImMxMSJ9
\end{studiumworked}

\begin{studiumworked}{Problema resuelto C11-03-P}


\noindent\textbf{Enunciado}
\(C11-03-P \cdot\) Reacción hacia la raíz. Ejercicio de transferencia. Calcule la reacción reducida en la raíz para un escalón repetitivo de vórtice libre. Antes de consultar la solución, prediga el efecto de cambiar el primer dato respecto del caso resuelto de esta familia. Justifique la tendencia y explique qué restricción podría convertirse en limitante. Datos: \(K_e=17.04\, \mathrm{m}^{2}\,\mathrm{s}^{-1}\); \(K_s=49\, \mathrm{m}^{2}\,\mathrm{s}^{-1}\); \(\Omega=1120\, \mathrm{s}^{-1}\); \(r_h=0.205\, \mathrm{m}\).

\noindent\textbf{Datos}
\begin{itemize}
\item $K_e = 17.04\,\mathrm{m}^{2}\,\mathrm{s}^{-1}$ — constante tangencial entrada
\item $K_s = 49\,\mathrm{m}^{2}\,\mathrm{s}^{-1}$ — constante tangencial salida
\item $\Omega = 1120\,\mathrm{s}^{-1}$ — velocidad angular
\item $r_h = 0.205\,\mathrm{m}$ — radio raíz
\end{itemize}

\noindent\textbf{Incógnita}
$\mathcal{R}$ (adimensional); también las magnitudes intermedias desarrolladas.

\noindent\textbf{Modelo físico}

\noindent\textit{Reacción hacia la raíz}

\[
\fitmath{\mathcal R=1 - \frac{K_e + K_s}{2 \Omega r_h^{2}}}
\]

\noindent\textbf{Hipótesis}
\begin{itemize}
\item Criterio reducido válido para radio y Cx constantes en cada escalón repetitivo.
\end{itemize}

\noindent\textbf{Desarrollo paso a paso}

\begin{enumerate}

\item Se calcula la suma de componentes tangenciales en la raíz.

\[\fitmath{\mathrm{sum}=\frac{\mathrm{ke}+\mathrm{ks}}{\mathrm{rh}}}\]

\[\fitmath{\mathrm{sum}=\frac{(17.04)+(49)}{(0.205)}\approx 322.1463\,\mathrm{m/s}}\]

\item Se normaliza con dos veces U local.

\[\fitmath{\mathrm{rea}=1-\frac{\mathrm{ke}+\mathrm{ks}}{2\cdot \Omega\cdot {\mathrm{rh}}^{2}}}\]

\[\fitmath{\mathrm{rea}=1-\frac{(17.04)+(49)}{2\cdot (1120)\cdot {(0.205)}^{2}}\approx 0.2984618\,}\]

\end{enumerate}


\noindent\textbf{Sustitución numérica}

\[
\fitmath{1 - \frac{\left(17.04\,\mathrm{m}^{2}\,\mathrm{s}^{-1}\right) + \left(49\,\mathrm{m}^{2}\,\mathrm{s}^{-1}\right)}{2 \left(1120\,\mathrm{s}^{-1}\right) \left(0.205\,\mathrm{m}\right)^{2}}}
\]

\noindent\textbf{Resultado}

\[
\fitmath{\mathcal R=2.984617999\times10^{-1}}
\]

\noindent\textbf{Interpretación}
La raíz puede tener reacción mucho menor que el radio medio. Si resulta negativa se debe revisar el diseño, no cambiar el signo.

\noindent\textbf{Comprobación de los pasos}
\begin{itemize}
\item Paso 2: $322.146$ — reproducido
\item Paso 4: $0.298462$ — reproducido
\end{itemize}

\noindent\textbf{Comprobación de la respuesta} la aritmética del último paso reproduce la respuesta enunciada a la precisión con que está escrita. Es una reproducción aritmética; no valida el modelo físico ni sus hipótesis.

\noindent\textbf{Limitaciones}
\begin{itemize}
\item Modelo educativo; la aplicabilidad a una máquina real requiere revisión experta.
\end{itemize}

\noindent\textbf{Errores frecuentes}
\begin{itemize}
\item Asignar a toda la pala la reacción del radio medio.
\end{itemize}
% studium:end-record eyJpZCI6ICJQUkItMDA4NiIsICJzZWN0aW9uIjogImMxMSJ9
\end{studiumworked}

\begin{studiumworked}{Problema resuelto C11-04-R}


\noindent\textbf{Enunciado}
\(C11-04-R \cdot\) Vórtice forzado admisible. Con entalpía total y entropía uniformes, calcule la velocidad axial exterior exigida por un remolino forzado. Datos: \(C_{x,m}=188\, \mathrm{m}\,\mathrm{s}^{-1}\); \(k_f=320\, \mathrm{s}^{-1}\); \(r_m=0.25\, \mathrm{m}\); \(r_t=0.33\, \mathrm{m}\).

\noindent\textbf{Datos}
\begin{itemize}
\item $C_{x,m} = 188\,\mathrm{m}\,\mathrm{s}^{-1}$ — velocidad axial al radio medio
\item $k_f = 320\,\mathrm{s}^{-1}$ — coeficiente $C_\theta=k_f r$
\item $r_m = 0.25\,\mathrm{m}$ — radio medio
\item $r_t = 0.33\,\mathrm{m}$ — radio exterior
\end{itemize}

\noindent\textbf{Incógnita}
$C_x$, en $\mathrm{m\,s^{-1}}$; también las magnitudes intermedias desarrolladas.

\noindent\textbf{Modelo físico}

\noindent\textit{Vórtice forzado admisible}

\[
\fitmath{C_x=\sqrt{C_{x,m}^{2} - 2 k_f^{2} \left(- r_m^{2} + r_t^{2}\right)}}
\]

\noindent\textbf{Hipótesis}
\begin{itemize}
\item Equilibrio radial simple, h0 y s uniformes.
\end{itemize}

\noindent\textbf{Desarrollo paso a paso}

\begin{enumerate}

\item El equilibrio energético fija el cuadrado axial exterior.

\[\fitmath{\mathrm{cx}_{2}={\mathrm{cxm}}^{2}-2\cdot {\mathrm{kv}}^{2}\cdot \left({\mathrm{rt}}^{2}-{\mathrm{rm}}^{2}\right)}\]

\[\fitmath{\mathrm{cx}_{2}={(188)}^{2}-2\cdot {(320)}^{2}\cdot \left({(0.33)}^{2}-{(0.25)}^{2}\right)\approx 25841.28\,\mathrm{m^{2}/s^{2}}}\]

\item Se selecciona la rama axial positiva.

\[\fitmath{\mathrm{cx}=\sqrt{{\mathrm{cxm}}^{2}-2\cdot {\mathrm{kv}}^{2}\cdot \left({\mathrm{rt}}^{2}-{\mathrm{rm}}^{2}\right)}}\]

\[\fitmath{\mathrm{cx}=\sqrt{{(188)}^{2}-2\cdot {(320)}^{2}\cdot \left({(0.33)}^{2}-{(0.25)}^{2}\right)}\approx 160.7522\,\mathrm{m/s}}\]

\end{enumerate}


\noindent\textbf{Sustitución numérica}

\[
\fitmath{\sqrt{\left(\left(188\,\mathrm{m}\,\mathrm{s}^{-1}\right)\right)^{2} - 2 \left(\left(320\,\mathrm{s}^{-1}\right)\right)^{2} \left(- \left(0.25\,\mathrm{m}\right)^{2} + \left(0.33\,\mathrm{m}\right)^{2}\right)}}
\]

\noindent\textbf{Resultado}

\[
\fitmath{C_x=1.607522317\times10^{2}\,\mathrm{m}\,\mathrm{s}^{-1}}
\]

\noindent\textbf{Interpretación}
El radicando debe ser positivo. Si no lo es, la ley propuesta no admite esta solución axial en todo el anillo.

\noindent\textbf{Comprobación de los pasos}
\begin{itemize}
\item Paso 2: $25841.3$ — reproducido
\item Paso 4: $160.752$ — reproducido
\end{itemize}

\noindent\textbf{Comprobación de la respuesta} la aritmética del último paso reproduce la respuesta enunciada a la precisión con que está escrita. Es una reproducción aritmética; no valida el modelo físico ni sus hipótesis.

\noindent\textbf{Limitaciones}
\begin{itemize}
\item Modelo educativo; la aplicabilidad a una máquina real requiere revisión experta.
\end{itemize}

\noindent\textbf{Errores frecuentes}
\begin{itemize}
\item Imponer Cx constante también para vórtice forzado.
\end{itemize}
% studium:end-record eyJpZCI6ICJQUkItMDA4NyIsICJzZWN0aW9uIjogImMxMSJ9
\end{studiumworked}

\begin{studiumworked}{Problema resuelto C11-04-P}


\noindent\textbf{Enunciado}
\(C11-04-P \cdot\) Vórtice forzado admisible. Ejercicio de transferencia. Con entalpía total y entropía uniformes, calcule la velocidad axial exterior exigida por un remolino forzado. Antes de consultar la solución, prediga el efecto de cambiar el primer dato respecto del caso resuelto de esta familia. Justifique la tendencia y explique qué restricción podría convertirse en limitante. Datos: \(C_{x,m}=200.22\, \mathrm{m}\,\mathrm{s}^{-1}\); \(k_f=320\, \mathrm{s}^{-1}\); \(r_m=0.25\, \mathrm{m}\); \(r_t=0.33\, \mathrm{m}\).

\noindent\textbf{Datos}
\begin{itemize}
\item $C_{x,m} = 200.22\,\mathrm{m}\,\mathrm{s}^{-1}$ — velocidad axial al radio medio
\item $k_f = 320\,\mathrm{s}^{-1}$ — coeficiente $C_\theta=k_f r$
\item $r_m = 0.25\,\mathrm{m}$ — radio medio
\item $r_t = 0.33\,\mathrm{m}$ — radio exterior
\end{itemize}

\noindent\textbf{Incógnita}
$C_x$, en $\mathrm{m\,s^{-1}}$; también las magnitudes intermedias desarrolladas.

\noindent\textbf{Modelo físico}

\noindent\textit{Vórtice forzado admisible}

\[
\fitmath{C_x=\sqrt{C_{x,m}^{2} - 2 k_f^{2} \left(- r_m^{2} + r_t^{2}\right)}}
\]

\noindent\textbf{Hipótesis}
\begin{itemize}
\item Equilibrio radial simple, h0 y s uniformes.
\end{itemize}

\noindent\textbf{Desarrollo paso a paso}

\begin{enumerate}

\item El equilibrio energético fija el cuadrado axial exterior.

\[\fitmath{\mathrm{cx}_{2}={\mathrm{cxm}}^{2}-2\cdot {\mathrm{kv}}^{2}\cdot \left({\mathrm{rt}}^{2}-{\mathrm{rm}}^{2}\right)}\]

\[\fitmath{\mathrm{cx}_{2}={(200.22)}^{2}-2\cdot {(320)}^{2}\cdot \left({(0.33)}^{2}-{(0.25)}^{2}\right)\approx 30585.33\,\mathrm{m^{2}/s^{2}}}\]

\item Se selecciona la rama axial positiva.

\[\fitmath{\mathrm{cx}=\sqrt{{\mathrm{cxm}}^{2}-2\cdot {\mathrm{kv}}^{2}\cdot \left({\mathrm{rt}}^{2}-{\mathrm{rm}}^{2}\right)}}\]

\[\fitmath{\mathrm{cx}=\sqrt{{(200.22)}^{2}-2\cdot {(320)}^{2}\cdot \left({(0.33)}^{2}-{(0.25)}^{2}\right)}\approx 174.8866\,\mathrm{m/s}}\]

\end{enumerate}


\noindent\textbf{Sustitución numérica}

\[
\fitmath{\sqrt{\left(\left(200.22\,\mathrm{m}\,\mathrm{s}^{-1}\right)\right)^{2} - 2 \left(\left(320\,\mathrm{s}^{-1}\right)\right)^{2} \left(- \left(0.25\,\mathrm{m}\right)^{2} + \left(0.33\,\mathrm{m}\right)^{2}\right)}}
\]

\noindent\textbf{Resultado}

\[
\fitmath{C_x=1.748866158\times10^{2}\,\mathrm{m}\,\mathrm{s}^{-1}}
\]

\noindent\textbf{Interpretación}
El radicando debe ser positivo. Si no lo es, la ley propuesta no admite esta solución axial en todo el anillo.

\noindent\textbf{Comprobación de los pasos}
\begin{itemize}
\item Paso 2: $30585.3$ — reproducido
\item Paso 4: $174.887$ — reproducido
\end{itemize}

\noindent\textbf{Comprobación de la respuesta} la aritmética del último paso reproduce la respuesta enunciada a la precisión con que está escrita. Es una reproducción aritmética; no valida el modelo físico ni sus hipótesis.

\noindent\textbf{Limitaciones}
\begin{itemize}
\item Modelo educativo; la aplicabilidad a una máquina real requiere revisión experta.
\end{itemize}

\noindent\textbf{Errores frecuentes}
\begin{itemize}
\item Imponer Cx constante también para vórtice forzado.
\end{itemize}
% studium:end-record eyJpZCI6ICJQUkItMDA4OCIsICJzZWN0aW9uIjogImMxMSJ9
\end{studiumworked}

\chapter{Torsión, secciones y construcción geométrica de álabes}
\begin{studiumworked}{Problema resuelto C12-01-R}


\noindent\textbf{Enunciado}
\(C12-01-R \cdot\) Torsión de entrada de rotor. Un rotor sin prerrotación recibe velocidad axial uniforme. Calcule beta a raíz y punta y la diferencia de torsión de entrada. Datos: \(\Omega=1060\, \mathrm{s}^{-1}\); \(C_x=157\, \mathrm{m}\,\mathrm{s}^{-1}\); \(r_h=0.198\, \mathrm{m}\); \(r_t=0.314\, \mathrm{m}\).

\noindent\textbf{Datos}
\begin{itemize}
\item $\Omega = 1060\,\mathrm{s}^{-1}$ — velocidad angular
\item $C_x = 157\,\mathrm{m}\,\mathrm{s}^{-1}$ — velocidad axial
\item $r_h = 0.198\,\mathrm{m}$ — radio raíz
\item $r_t = 0.314\,\mathrm{m}$ — radio punta
\end{itemize}

\noindent\textbf{Incógnita}
$\Delta\beta$, en $\mathrm{rad}$; también las magnitudes intermedias desarrolladas.

\noindent\textbf{Modelo físico}

\noindent\textit{Torsión de entrada de rotor}

\[
\fitmath{\Delta\beta=\operatorname{atan}{\left(\frac{\Omega r_h}{C_x} \right)} - \operatorname{atan}{\left(\frac{\Omega r_t}{C_x} \right)}}
\]

\noindent\textbf{Hipótesis}
\begin{itemize}
\item Sin prerrotación, Cx uniforme, incidencia cero.
\end{itemize}

\noindent\textbf{Desarrollo paso a paso}

\begin{enumerate}

\item Se obtiene beta de raíz con componente tangencial -U.

\[\fitmath{\mathrm{bh}=\arctan\left(\frac{-\Omega\cdot \mathrm{rh}}{\mathrm{cx}}\right)}\]

\[\fitmath{\mathrm{bh}=\arctan\left(\frac{-(1060)\cdot (0.198)}{(157)}\right)\approx -0.9285466\,\mathrm{rad}}\]

\item Se obtiene beta de punta.

\[\fitmath{\mathrm{bt}=\arctan\left(\frac{-\Omega\cdot \mathrm{rt}}{\mathrm{cx}}\right)}\]

\[\fitmath{\mathrm{bt}=\arctan\left(\frac{-(1060)\cdot (0.314)}{(157)}\right)\approx -1.130045\,\mathrm{rad}}\]

\item La torsión es una diferencia algebraica de orientaciones.

\[\fitmath{\mathrm{twist}=\arctan\left(\frac{-\Omega\cdot \mathrm{rt}}{\mathrm{cx}}\right)-\arctan\left(\frac{-\Omega\cdot \mathrm{rh}}{\mathrm{cx}}\right)}\]

\[\fitmath{\mathrm{twist}=\arctan\left(\frac{-(1060)\cdot (0.314)}{(157)}\right)-\arctan\left(\frac{-(1060)\cdot (0.198)}{(157)}\right)\approx -0.2014989\,\mathrm{rad}}\]

\end{enumerate}


\noindent\textbf{Sustitución numérica}

\[
\fitmath{\operatorname{atan}{\left(\frac{\left(1060\,\mathrm{s}^{-1}\right) \left(0.198\,\mathrm{m}\right)}{\left(157\,\mathrm{m}\,\mathrm{s}^{-1}\right)} \right)} - \operatorname{atan}{\left(\frac{\left(1060\,\mathrm{s}^{-1}\right) \left(0.314\,\mathrm{m}\right)}{\left(157\,\mathrm{m}\,\mathrm{s}^{-1}\right)} \right)}}
\]

\noindent\textbf{Resultado}

\[
\fitmath{\Delta\beta=-2.014988500\times10^{-1}\,\mathrm{rad}}
\]

\noindent\textbf{Interpretación}
La orientación de punta es más negativa. La tangente del metal coincide solo bajo la hipótesis de incidencia prescrita igual a cero.

\noindent\textbf{Comprobación de los pasos}
\begin{itemize}
\item Paso 2: $-0.928547$ — reproducido
\item Paso 4: $-1.13005$ — reproducido
\item Paso 6: $-0.201499$ — reproducido
\end{itemize}

\noindent\textbf{Comprobación de la respuesta} la aritmética del último paso reproduce la respuesta enunciada a la precisión con que está escrita. Es una reproducción aritmética; no valida el modelo físico ni sus hipótesis.

\noindent\textbf{Limitaciones}
\begin{itemize}
\item Modelo educativo; la aplicabilidad a una máquina real requiere revisión experta.
\end{itemize}

\noindent\textbf{Errores frecuentes}
\begin{itemize}
\item Atribuir el mismo ángulo a todos los radios.
\end{itemize}
% studium:end-record eyJpZCI6ICJQUkItMDA4OSIsICJzZWN0aW9uIjogImMxMiJ9
\end{studiumworked}

\begin{studiumworked}{Problema resuelto C12-01-P}


\noindent\textbf{Enunciado}
\(C12-01-P \cdot\) Torsión de entrada de rotor. Ejercicio de transferencia. Un rotor sin prerrotación recibe velocidad axial uniforme. Calcule beta a raíz y punta y la diferencia de torsión de entrada. Antes de consultar la solución, prediga el efecto de cambiar el primer dato respecto del caso resuelto de esta familia. Justifique la tendencia y explique qué restricción podría convertirse en limitante. Datos: \(\Omega=1128.9\, \mathrm{s}^{-1}\); \(C_x=157\, \mathrm{m}\,\mathrm{s}^{-1}\); \(r_h=0.198\, \mathrm{m}\); \(r_t=0.314\, \mathrm{m}\).

\noindent\textbf{Datos}
\begin{itemize}
\item $\Omega = 1128.9\,\mathrm{s}^{-1}$ — velocidad angular
\item $C_x = 157\,\mathrm{m}\,\mathrm{s}^{-1}$ — velocidad axial
\item $r_h = 0.198\,\mathrm{m}$ — radio raíz
\item $r_t = 0.314\,\mathrm{m}$ — radio punta
\end{itemize}

\noindent\textbf{Incógnita}
$\Delta\beta$, en $\mathrm{rad}$; también las magnitudes intermedias desarrolladas.

\noindent\textbf{Modelo físico}

\noindent\textit{Torsión de entrada de rotor}

\[
\fitmath{\Delta\beta=\operatorname{atan}{\left(\frac{\Omega r_h}{C_x} \right)} - \operatorname{atan}{\left(\frac{\Omega r_t}{C_x} \right)}}
\]

\noindent\textbf{Hipótesis}
\begin{itemize}
\item Sin prerrotación, Cx uniforme, incidencia cero.
\end{itemize}

\noindent\textbf{Desarrollo paso a paso}

\begin{enumerate}

\item Se obtiene beta de raíz con componente tangencial -U.

\[\fitmath{\mathrm{bh}=\arctan\left(\frac{-\Omega\cdot \mathrm{rh}}{\mathrm{cx}}\right)}\]

\[\fitmath{\mathrm{bh}=\arctan\left(\frac{-(1128.9)\cdot (0.198)}{(157)}\right)\approx -0.9584674\,\mathrm{rad}}\]

\item Se obtiene beta de punta.

\[\fitmath{\mathrm{bt}=\arctan\left(\frac{-\Omega\cdot \mathrm{rt}}{\mathrm{cx}}\right)}\]

\[\fitmath{\mathrm{bt}=\arctan\left(\frac{-(1128.9)\cdot (0.314)}{(157)}\right)\approx -1.153855\,\mathrm{rad}}\]

\item La torsión es una diferencia algebraica de orientaciones.

\[\fitmath{\mathrm{twist}=\arctan\left(\frac{-\Omega\cdot \mathrm{rt}}{\mathrm{cx}}\right)-\arctan\left(\frac{-\Omega\cdot \mathrm{rh}}{\mathrm{cx}}\right)}\]

\[\fitmath{\mathrm{twist}=\arctan\left(\frac{-(1128.9)\cdot (0.314)}{(157)}\right)-\arctan\left(\frac{-(1128.9)\cdot (0.198)}{(157)}\right)\approx -0.1953875\,\mathrm{rad}}\]

\end{enumerate}


\noindent\textbf{Sustitución numérica}

\[
\fitmath{\operatorname{atan}{\left(\frac{\left(1128.9\,\mathrm{s}^{-1}\right) \left(0.198\,\mathrm{m}\right)}{\left(157\,\mathrm{m}\,\mathrm{s}^{-1}\right)} \right)} - \operatorname{atan}{\left(\frac{\left(1128.9\,\mathrm{s}^{-1}\right) \left(0.314\,\mathrm{m}\right)}{\left(157\,\mathrm{m}\,\mathrm{s}^{-1}\right)} \right)}}
\]

\noindent\textbf{Resultado}

\[
\fitmath{\Delta\beta=-1.953874696\times10^{-1}\,\mathrm{rad}}
\]

\noindent\textbf{Interpretación}
La orientación de punta es más negativa. La tangente del metal coincide solo bajo la hipótesis de incidencia prescrita igual a cero.

\noindent\textbf{Comprobación de los pasos}
\begin{itemize}
\item Paso 2: $-0.958467$ — reproducido
\item Paso 4: $-1.15385$ — reproducido
\item Paso 6: $-0.195387$ — reproducido
\end{itemize}

\noindent\textbf{Comprobación de la respuesta} la aritmética del último paso reproduce la respuesta enunciada a la precisión con que está escrita. Es una reproducción aritmética; no valida el modelo físico ni sus hipótesis.

\noindent\textbf{Limitaciones}
\begin{itemize}
\item Modelo educativo; la aplicabilidad a una máquina real requiere revisión experta.
\end{itemize}

\noindent\textbf{Errores frecuentes}
\begin{itemize}
\item Atribuir el mismo ángulo a todos los radios.
\end{itemize}
% studium:end-record eyJpZCI6ICJQUkItMDA5MCIsICJzZWN0aW9uIjogImMxMiJ9
\end{studiumworked}

\begin{studiumworked}{Problema resuelto C12-02-R}


\noindent\textbf{Enunciado}
\(C12-02-R \cdot\) Cuerda y proyección. Una sección tiene cuerda axial y calado dados. Calcule cuerda real y solidez a paso fijado. Datos: \(c_x=0.044\, \mathrm{m}\); \(\chi=0.6457718\, \mathrm{rad}\); \(s=0.039\, \mathrm{m}\).

\noindent\textbf{Datos}
\begin{itemize}
\item $c_x = 0.044\,\mathrm{m}$ — cuerda axial
\item $\chi = 0.6457718232379\,\mathrm{rad}$ — calado
\item $s = 0.039\,\mathrm{m}$ — paso
\end{itemize}

\noindent\textbf{Incógnita}
$\sigma$ (adimensional); también las magnitudes intermedias desarrolladas.

\noindent\textbf{Modelo físico}

\noindent\textit{Cuerda y proyección}

\[
\fitmath{\sigma=\frac{c_x}{s \cos{\left(\chi \right)}}}
\]

\noindent\textbf{Hipótesis}
\begin{itemize}
\item Calado menor de noventa grados en módulo.
\end{itemize}

\noindent\textbf{Desarrollo paso a paso}

\begin{enumerate}

\item La cuerda axial es la proyección de la cuerda real.

\[\fitmath{\mathrm{chord}=\frac{\mathrm{cxg}}{\cos\left(\chi\right)}}\]

\[\fitmath{\mathrm{chord}=\frac{(0.044)}{\cos\left((0.6457718232)\right)}\approx 0.05509397\,\mathrm{m}}\]

\item La solidez usa cuerda real.

\[\fitmath{\mathrm{sol}=\frac{\mathrm{cxg}}{\cos\left(\chi\right)\cdot \mathrm{ss}}}\]

\[\fitmath{\mathrm{sol}=\frac{(0.044)}{\cos\left((0.6457718232)\right)\cdot (0.039)}\approx 1.412666\,}\]

\end{enumerate}


\noindent\textbf{Sustitución numérica}

\[
\fitmath{\frac{\left(0.044\,\mathrm{m}\right)}{\left(0.039\,\mathrm{m}\right) \cos{\left(\left(0.6457718232379\,\mathrm{rad}\right) \right)}}}
\]

\noindent\textbf{Resultado}

\[
\fitmath{\sigma=1.412665871\times10^{0}}
\]

\noindent\textbf{Interpretación}
Cuerda axial y cuerda real no pueden intercambiarse en la solidez sin cambiar su definición.

\noindent\textbf{Comprobación de los pasos}
\begin{itemize}
\item Paso 2: $0.055094$ — reproducido
\item Paso 4: $1.41267$ — reproducido
\end{itemize}

\noindent\textbf{Comprobación de la respuesta} la aritmética del último paso reproduce la respuesta enunciada a la precisión con que está escrita. Es una reproducción aritmética; no valida el modelo físico ni sus hipótesis.

\noindent\textbf{Limitaciones}
\begin{itemize}
\item Modelo educativo; la aplicabilidad a una máquina real requiere revisión experta.
\end{itemize}

\noindent\textbf{Errores frecuentes}
\begin{itemize}
\item Multiplicar por coseno al despejar la cuerda real.
\end{itemize}
% studium:end-record eyJpZCI6ICJQUkItMDA5MSIsICJzZWN0aW9uIjogImMxMiJ9
\end{studiumworked}

\begin{studiumworked}{Problema resuelto C12-02-P}


\noindent\textbf{Enunciado}
\(C12-02-P \cdot\) Cuerda y proyección. Ejercicio de transferencia. Una sección tiene cuerda axial y calado dados. Calcule cuerda real y solidez a paso fijado. Antes de consultar la solución, prediga el efecto de cambiar el primer dato respecto del caso resuelto de esta familia. Justifique la tendencia y explique qué restricción podría convertirse en limitante. Datos: \(c_x=0.04686\, \mathrm{m}\); \(\chi=0.6457718\, \mathrm{rad}\); \(s=0.039\, \mathrm{m}\).

\noindent\textbf{Datos}
\begin{itemize}
\item $c_x = 0.04686\,\mathrm{m}$ — cuerda axial
\item $\chi = 0.6457718232379\,\mathrm{rad}$ — calado
\item $s = 0.039\,\mathrm{m}$ — paso
\end{itemize}

\noindent\textbf{Incógnita}
$\sigma$ (adimensional); también las magnitudes intermedias desarrolladas.

\noindent\textbf{Modelo físico}

\noindent\textit{Cuerda y proyección}

\[
\fitmath{\sigma=\frac{c_x}{s \cos{\left(\chi \right)}}}
\]

\noindent\textbf{Hipótesis}
\begin{itemize}
\item Calado menor de noventa grados en módulo.
\end{itemize}

\noindent\textbf{Desarrollo paso a paso}

\begin{enumerate}

\item La cuerda axial es la proyección de la cuerda real.

\[\fitmath{\mathrm{chord}=\frac{\mathrm{cxg}}{\cos\left(\chi\right)}}\]

\[\fitmath{\mathrm{chord}=\frac{(0.04686)}{\cos\left((0.6457718232)\right)}\approx 0.05867508\,\mathrm{m}}\]

\item La solidez usa cuerda real.

\[\fitmath{\mathrm{sol}=\frac{\mathrm{cxg}}{\cos\left(\chi\right)\cdot \mathrm{ss}}}\]

\[\fitmath{\mathrm{sol}=\frac{(0.04686)}{\cos\left((0.6457718232)\right)\cdot (0.039)}\approx 1.504489\,}\]

\end{enumerate}


\noindent\textbf{Sustitución numérica}

\[
\fitmath{\frac{\left(0.04686\,\mathrm{m}\right)}{\left(0.039\,\mathrm{m}\right) \cos{\left(\left(0.6457718232379\,\mathrm{rad}\right) \right)}}}
\]

\noindent\textbf{Resultado}

\[
\fitmath{\sigma=1.504489152\times10^{0}}
\]

\noindent\textbf{Interpretación}
Cuerda axial y cuerda real no pueden intercambiarse en la solidez sin cambiar su definición.

\noindent\textbf{Comprobación de los pasos}
\begin{itemize}
\item Paso 2: $0.0586751$ — reproducido
\item Paso 4: $1.50449$ — reproducido
\end{itemize}

\noindent\textbf{Comprobación de la respuesta} la aritmética del último paso reproduce la respuesta enunciada a la precisión con que está escrita. Es una reproducción aritmética; no valida el modelo físico ni sus hipótesis.

\noindent\textbf{Limitaciones}
\begin{itemize}
\item Modelo educativo; la aplicabilidad a una máquina real requiere revisión experta.
\end{itemize}

\noindent\textbf{Errores frecuentes}
\begin{itemize}
\item Multiplicar por coseno al despejar la cuerda real.
\end{itemize}
% studium:end-record eyJpZCI6ICJQUkItMDA5MiIsICJzZWN0aW9uIjogImMxMiJ9
\end{studiumworked}

\begin{studiumworked}{Problema resuelto C12-03-R}


\noindent\textbf{Enunciado}
\(C12-03-R \cdot\) Línea de Hermite. Para pendientes de entrada y salida dadas y desplazamiento d, calcule la ordenada de línea media al parámetro xi. Datos: \(c_x=0.047\, \mathrm{m}\); \(m_e=0.86\); \(m_s=-0.24\); \(\delta_y=0.31\); \(\xi=0.4\).

\noindent\textbf{Datos}
\begin{itemize}
\item $c_x = 0.047\,\mathrm{m}$ — cuerda axial
\item $m_e = 0.86$ — pendiente entrada
\item $m_s = -0.24$ — pendiente salida
\item $\delta_y = 0.31$ — desplazamiento final normalizado
\item $\xi = 0.4$ — parámetro xi
\end{itemize}

\noindent\textbf{Incógnita}
$y$, en $\mathrm{m}$; también las magnitudes intermedias desarrolladas.

\noindent\textbf{Modelo físico}

\noindent\textit{Línea de Hermite}

\[
\fitmath{y=c_x \left(\delta_y \left(- 2 \xi^{3} + 3 \xi^{2}\right) + m_e \left(\xi^{3} - 2 \xi^{2} + \xi\right) + m_s \left(\xi^{3} - \xi^{2}\right)\right)}
\]

\noindent\textbf{Hipótesis}
\begin{itemize}
\item Perfil educativo definido como función de x.
\end{itemize}

\noindent\textbf{Desarrollo paso a paso}

\begin{enumerate}

\item Se evalúa la línea media normalizada.

\[\fitmath{\mathrm{fy}=\left({\mathrm{xx}}^{3}-2\cdot {\mathrm{xx}}^{2}+\mathrm{xx}\right)\cdot \mathrm{me}+\left(-2\cdot {\mathrm{xx}}^{3}+3\cdot {\mathrm{xx}}^{2}\right)\cdot \mathrm{dd}+\left({\mathrm{xx}}^{3}-{\mathrm{xx}}^{2}\right)\cdot \mathrm{ms}}\]

\[\fitmath{\mathrm{fy}=\left({(0.4)}^{3}-2\cdot {(0.4)}^{2}+(0.4)\right)\cdot (0.86)+\left(-2\cdot {(0.4)}^{3}+3\cdot {(0.4)}^{2}\right)\cdot (0.31)+\left({(0.4)}^{3}-{(0.4)}^{2}\right)\cdot (-0.24)\approx 0.256\,}\]

\item Se devuelve la longitud física transversal.

\[\fitmath{\mathrm{yy}=\mathrm{cxg}\cdot \left(\left({\mathrm{xx}}^{3}-2\cdot {\mathrm{xx}}^{2}+\mathrm{xx}\right)\cdot \mathrm{me}+\left(-2\cdot {\mathrm{xx}}^{3}+3\cdot {\mathrm{xx}}^{2}\right)\cdot \mathrm{dd}+\left({\mathrm{xx}}^{3}-{\mathrm{xx}}^{2}\right)\cdot \mathrm{ms}\right)}\]

\[\fitmath{\mathrm{yy}=(0.047)\cdot \left(\left({(0.4)}^{3}-2\cdot {(0.4)}^{2}+(0.4)\right)\cdot (0.86)+\left(-2\cdot {(0.4)}^{3}+3\cdot {(0.4)}^{2}\right)\cdot (0.31)+\left({(0.4)}^{3}-{(0.4)}^{2}\right)\cdot (-0.24)\right)\approx 0.012032\,\mathrm{m}}\]

\end{enumerate}


\noindent\textbf{Sustitución numérica}

\[
\fitmath{\left(0.047\,\mathrm{m}\right) \left(\left(0.31\right) \left(- 2 \left(0.4\right)^{3} + 3 \left(0.4\right)^{2}\right) + \left(0.86\right) \left(\left(0.4\right)^{3} - 2 \left(0.4\right)^{2} + \left(0.4\right)\right) + \left(-0.24\right) \left(\left(0.4\right)^{3} - \left(0.4\right)^{2}\right)\right)}
\]

\noindent\textbf{Resultado}

\[
\fitmath{y=1.203200000\times10^{-2}\,\mathrm{m}}
\]

\noindent\textbf{Interpretación}
El resultado es una coordenada geométrica, no una predicción aerodinámica. Compruebe los cuatro extremos y pendientes del polinomio.

\noindent\textbf{Comprobación de los pasos}
\begin{itemize}
\item Paso 2: $0.256$ — reproducido
\item Paso 4: $0.012032$ — reproducido
\end{itemize}

\noindent\textbf{Comprobación de la respuesta} la aritmética del último paso reproduce la respuesta enunciada a la precisión con que está escrita. Es una reproducción aritmética; no valida el modelo físico ni sus hipótesis.

\noindent\textbf{Limitaciones}
\begin{itemize}
\item Modelo educativo; la aplicabilidad a una máquina real requiere revisión experta.
\end{itemize}

\noindent\textbf{Errores frecuentes}
\begin{itemize}
\item Confundir pendiente con ángulo en grados.
\end{itemize}
% studium:end-record eyJpZCI6ICJQUkItMDA5MyIsICJzZWN0aW9uIjogImMxMiJ9
\end{studiumworked}

\begin{studiumworked}{Problema resuelto C12-03-P}


\noindent\textbf{Enunciado}
\(C12-03-P \cdot\) Línea de Hermite. Ejercicio de transferencia. Para pendientes de entrada y salida dadas y desplazamiento d, calcule la ordenada de línea media al parámetro xi. Antes de consultar la solución, prediga el efecto de cambiar el primer dato respecto del caso resuelto de esta familia. Justifique la tendencia y explique qué restricción podría convertirse en limitante. Datos: \(c_x=0.050055\, \mathrm{m}\); \(m_e=0.86\); \(m_s=-0.24\); \(\delta_y=0.31\); \(\xi=0.4\).

\noindent\textbf{Datos}
\begin{itemize}
\item $c_x = 0.050055\,\mathrm{m}$ — cuerda axial
\item $m_e = 0.86$ — pendiente entrada
\item $m_s = -0.24$ — pendiente salida
\item $\delta_y = 0.31$ — desplazamiento final normalizado
\item $\xi = 0.4$ — parámetro xi
\end{itemize}

\noindent\textbf{Incógnita}
$y$, en $\mathrm{m}$; también las magnitudes intermedias desarrolladas.

\noindent\textbf{Modelo físico}

\noindent\textit{Línea de Hermite}

\[
\fitmath{y=c_x \left(\delta_y \left(- 2 \xi^{3} + 3 \xi^{2}\right) + m_e \left(\xi^{3} - 2 \xi^{2} + \xi\right) + m_s \left(\xi^{3} - \xi^{2}\right)\right)}
\]

\noindent\textbf{Hipótesis}
\begin{itemize}
\item Perfil educativo definido como función de x.
\end{itemize}

\noindent\textbf{Desarrollo paso a paso}

\begin{enumerate}

\item Se evalúa la línea media normalizada.

\[\fitmath{\mathrm{fy}=\left({\mathrm{xx}}^{3}-2\cdot {\mathrm{xx}}^{2}+\mathrm{xx}\right)\cdot \mathrm{me}+\left(-2\cdot {\mathrm{xx}}^{3}+3\cdot {\mathrm{xx}}^{2}\right)\cdot \mathrm{dd}+\left({\mathrm{xx}}^{3}-{\mathrm{xx}}^{2}\right)\cdot \mathrm{ms}}\]

\[\fitmath{\mathrm{fy}=\left({(0.4)}^{3}-2\cdot {(0.4)}^{2}+(0.4)\right)\cdot (0.86)+\left(-2\cdot {(0.4)}^{3}+3\cdot {(0.4)}^{2}\right)\cdot (0.31)+\left({(0.4)}^{3}-{(0.4)}^{2}\right)\cdot (-0.24)\approx 0.256\,}\]

\item Se devuelve la longitud física transversal.

\[\fitmath{\mathrm{yy}=\mathrm{cxg}\cdot \left(\left({\mathrm{xx}}^{3}-2\cdot {\mathrm{xx}}^{2}+\mathrm{xx}\right)\cdot \mathrm{me}+\left(-2\cdot {\mathrm{xx}}^{3}+3\cdot {\mathrm{xx}}^{2}\right)\cdot \mathrm{dd}+\left({\mathrm{xx}}^{3}-{\mathrm{xx}}^{2}\right)\cdot \mathrm{ms}\right)}\]

\[\fitmath{\mathrm{yy}=(0.050055)\cdot \left(\left({(0.4)}^{3}-2\cdot {(0.4)}^{2}+(0.4)\right)\cdot (0.86)+\left(-2\cdot {(0.4)}^{3}+3\cdot {(0.4)}^{2}\right)\cdot (0.31)+\left({(0.4)}^{3}-{(0.4)}^{2}\right)\cdot (-0.24)\right)\approx 0.01281408\,\mathrm{m}}\]

\end{enumerate}


\noindent\textbf{Sustitución numérica}

\[
\fitmath{\left(0.050055\,\mathrm{m}\right) \left(\left(0.31\right) \left(- 2 \left(0.4\right)^{3} + 3 \left(0.4\right)^{2}\right) + \left(0.86\right) \left(\left(0.4\right)^{3} - 2 \left(0.4\right)^{2} + \left(0.4\right)\right) + \left(-0.24\right) \left(\left(0.4\right)^{3} - \left(0.4\right)^{2}\right)\right)}
\]

\noindent\textbf{Resultado}

\[
\fitmath{y=1.281408000\times10^{-2}\,\mathrm{m}}
\]

\noindent\textbf{Interpretación}
El resultado es una coordenada geométrica, no una predicción aerodinámica. Compruebe los cuatro extremos y pendientes del polinomio.

\noindent\textbf{Comprobación de los pasos}
\begin{itemize}
\item Paso 2: $0.256$ — reproducido
\item Paso 4: $0.0128141$ — reproducido
\end{itemize}

\noindent\textbf{Comprobación de la respuesta} la aritmética del último paso reproduce la respuesta enunciada a la precisión con que está escrita. Es una reproducción aritmética; no valida el modelo físico ni sus hipótesis.

\noindent\textbf{Limitaciones}
\begin{itemize}
\item Modelo educativo; la aplicabilidad a una máquina real requiere revisión experta.
\end{itemize}

\noindent\textbf{Errores frecuentes}
\begin{itemize}
\item Confundir pendiente con ángulo en grados.
\end{itemize}
% studium:end-record eyJpZCI6ICJQUkItMDA5NCIsICJzZWN0aW9uIjogImMxMiJ9
\end{studiumworked}

\begin{studiumworked}{Problema resuelto C12-04-R}


\noindent\textbf{Enunciado}
\(C12-04-R \cdot\) Desplazamiento normal por espesor. En un punto de una línea media se conocen pendiente y semiespesor normal. Halle las dos componentes del desplazamiento normal y su componente axial como resultado final. Datos: \(t_n=0.0024\, \mathrm{m}\); \(m_y=0.73\).

\noindent\textbf{Datos}
\begin{itemize}
\item $t_n = 0.0024\,\mathrm{m}$ — semiespesor normal local
\item $m_y = 0.73$ — pendiente local
\end{itemize}

\noindent\textbf{Incógnita}
$\Delta x$, en $\mathrm{m}$; también las magnitudes intermedias desarrolladas.

\noindent\textbf{Modelo físico}

\noindent\textit{Desplazamiento normal por espesor}

\[
\fitmath{\Delta x=- \frac{m_y t_n}{\sqrt{m_y^{2} + 1}}}
\]

\noindent\textbf{Hipótesis}
\begin{itemize}
\item Superficie obtenida por desplazamiento normal local.
\end{itemize}

\noindent\textbf{Desarrollo paso a paso}

\begin{enumerate}

\item La normal unitaria tiene componente transversal positiva.

\[\fitmath{\mathrm{dy}=\frac{\mathrm{tt}}{\sqrt{1+{\mathrm{slope}}^{2}}}}\]

\[\fitmath{\mathrm{dy}=\frac{(0.0024)}{\sqrt{1+{(0.73)}^{2}}}\approx 0.001938449\,\mathrm{m}}\]

\item La componente axial compensa la inclinación local.

\[\fitmath{\mathrm{dx}=\frac{-\mathrm{tt}\cdot \mathrm{slope}}{\sqrt{1+{\mathrm{slope}}^{2}}}}\]

\[\fitmath{\mathrm{dx}=\frac{-(0.0024)\cdot (0.73)}{\sqrt{1+{(0.73)}^{2}}}\approx -0.001415068\,\mathrm{m}}\]

\end{enumerate}


\noindent\textbf{Sustitución numérica}

\[
\fitmath{- \frac{\left(0.73\right) \left(0.0024\,\mathrm{m}\right)}{\sqrt{\left(0.73\right)^{2} + 1}}}
\]

\noindent\textbf{Resultado}

\[
\fitmath{\Delta x=-1.415067608\times10^{-3}\,\mathrm{m}}
\]

\noindent\textbf{Interpretación}
La longitud del vector desplazamiento es tt. Añadir solo dy produce un espesor vertical distinto.

\noindent\textbf{Comprobación de los pasos}
\begin{itemize}
\item Paso 2: $0.00193845$ — reproducido
\item Paso 4: $-0.00141507$ — reproducido
\end{itemize}

\noindent\textbf{Comprobación de la respuesta} la aritmética del último paso reproduce la respuesta enunciada a la precisión con que está escrita. Es una reproducción aritmética; no valida el modelo físico ni sus hipótesis.

\noindent\textbf{Limitaciones}
\begin{itemize}
\item Modelo educativo; la aplicabilidad a una máquina real requiere revisión experta.
\end{itemize}

\noindent\textbf{Errores frecuentes}
\begin{itemize}
\item Olvidar el desplazamiento axial al añadir espesor.
\end{itemize}
% studium:end-record eyJpZCI6ICJQUkItMDA5NSIsICJzZWN0aW9uIjogImMxMiJ9
\end{studiumworked}

\begin{studiumworked}{Problema resuelto C12-04-P}


\noindent\textbf{Enunciado}
\(C12-04-P \cdot\) Desplazamiento normal por espesor. Ejercicio de transferencia. En un punto de una línea media se conocen pendiente y semiespesor normal. Halle las dos componentes del desplazamiento normal y su componente axial como resultado final. Antes de consultar la solución, prediga el efecto de cambiar el primer dato respecto del caso resuelto de esta familia. Justifique la tendencia y explique qué restricción podría convertirse en limitante. Datos: \(t_n=0.002556\, \mathrm{m}\); \(m_y=0.73\).

\noindent\textbf{Datos}
\begin{itemize}
\item $t_n = 0.002556\,\mathrm{m}$ — semiespesor normal local
\item $m_y = 0.73$ — pendiente local
\end{itemize}

\noindent\textbf{Incógnita}
$\Delta x$, en $\mathrm{m}$; también las magnitudes intermedias desarrolladas.

\noindent\textbf{Modelo físico}

\noindent\textit{Desplazamiento normal por espesor}

\[
\fitmath{\Delta x=- \frac{m_y t_n}{\sqrt{m_y^{2} + 1}}}
\]

\noindent\textbf{Hipótesis}
\begin{itemize}
\item Superficie obtenida por desplazamiento normal local.
\end{itemize}

\noindent\textbf{Desarrollo paso a paso}

\begin{enumerate}

\item La normal unitaria tiene componente transversal positiva.

\[\fitmath{\mathrm{dy}=\frac{\mathrm{tt}}{\sqrt{1+{\mathrm{slope}}^{2}}}}\]

\[\fitmath{\mathrm{dy}=\frac{(0.002556)}{\sqrt{1+{(0.73)}^{2}}}\approx 0.002064448\,\mathrm{m}}\]

\item La componente axial compensa la inclinación local.

\[\fitmath{\mathrm{dx}=\frac{-\mathrm{tt}\cdot \mathrm{slope}}{\sqrt{1+{\mathrm{slope}}^{2}}}}\]

\[\fitmath{\mathrm{dx}=\frac{-(0.002556)\cdot (0.73)}{\sqrt{1+{(0.73)}^{2}}}\approx -0.001507047\,\mathrm{m}}\]

\end{enumerate}


\noindent\textbf{Sustitución numérica}

\[
\fitmath{- \frac{\left(0.73\right) \left(0.002556\,\mathrm{m}\right)}{\sqrt{\left(0.73\right)^{2} + 1}}}
\]

\noindent\textbf{Resultado}

\[
\fitmath{\Delta x=-1.507047002\times10^{-3}\,\mathrm{m}}
\]

\noindent\textbf{Interpretación}
La longitud del vector desplazamiento es tt. Añadir solo dy produce un espesor vertical distinto.

\noindent\textbf{Comprobación de los pasos}
\begin{itemize}
\item Paso 2: $0.00206445$ — reproducido
\item Paso 4: $-0.00150705$ — reproducido
\end{itemize}

\noindent\textbf{Comprobación de la respuesta} la aritmética del último paso reproduce la respuesta enunciada a la precisión con que está escrita. Es una reproducción aritmética; no valida el modelo físico ni sus hipótesis.

\noindent\textbf{Limitaciones}
\begin{itemize}
\item Modelo educativo; la aplicabilidad a una máquina real requiere revisión experta.
\end{itemize}

\noindent\textbf{Errores frecuentes}
\begin{itemize}
\item Olvidar el desplazamiento axial al añadir espesor.
\end{itemize}
% studium:end-record eyJpZCI6ICJQUkItMDA5NiIsICJzZWN0aW9uIjogImMxMiJ9
\end{studiumworked}

\chapter{Esfuerzos centrífugos, flexión y vibración}
\begin{studiumworked}{Problema resuelto C13-01-R}


\noindent\textbf{Enunciado}
\(C13-01-R \cdot\) Tracción centrífuga uniforme. Calcule la tensión centrífuga nominal de raíz de una pala uniforme. Datos: \(\rho=8120\, \mathrm{kg}\,\mathrm{m}^{-3}\); \(\Omega=1170\, \mathrm{s}^{-1}\); \(r_h=0.245\, \mathrm{m}\); \(r_t=0.302\, \mathrm{m}\).

\noindent\textbf{Datos}
\begin{itemize}
\item $\rho = 8120\,\mathrm{kg}\,\mathrm{m}^{-3}$ — densidad material
\item $\Omega = 1170\,\mathrm{s}^{-1}$ — velocidad angular
\item $r_h = 0.245\,\mathrm{m}$ — radio raíz
\item $r_t = 0.302\,\mathrm{m}$ — radio punta
\end{itemize}

\noindent\textbf{Incógnita}
$\sigma$, en $\mathrm{Pa}$; también las magnitudes intermedias desarrolladas.

\noindent\textbf{Modelo físico}

\noindent\textit{Tracción centrífuga uniforme}

\[
\fitmath{\sigma=\frac{\Omega^{2} \rho \left(- r_h^{2} + r_t^{2}\right)}{2}}
\]

\noindent\textbf{Hipótesis}
\begin{itemize}
\item Barra radial uniforme, material homogéneo, régimen estacionario.
\end{itemize}

\noindent\textbf{Desarrollo paso a paso}

\begin{enumerate}

\item Se integra r entre raíz y punta.

\[\fitmath{\mathrm{integ}=\frac{{\mathrm{rt}}^{2}-{\mathrm{rh}}^{2}}{2}}\]

\[\fitmath{\mathrm{integ}=\frac{{(0.302)}^{2}-{(0.245)}^{2}}{2}\approx 0.0155895\,\mathrm{m^{2}}}\]

\item Se multiplica por densidad y velocidad angular al cuadrado.

\[\fitmath{\mathrm{stress}=\frac{\rho\cdot {\Omega}^{2}\cdot \left({\mathrm{rt}}^{2}-{\mathrm{rh}}^{2}\right)}{2}}\]

\[\fitmath{\mathrm{stress}=\frac{(8120)\cdot {(1170)}^{2}\cdot \left({(0.302)}^{2}-{(0.245)}^{2}\right)}{2}\approx 1.732846\times10^{8}\,\mathrm{Pa}}\]

\end{enumerate}


\noindent\textbf{Sustitución numérica}

\[
\fitmath{\frac{\left(\left(1170\,\mathrm{s}^{-1}\right)\right)^{2} \left(8120\,\mathrm{kg}\,\mathrm{m}^{-3}\right) \left(- \left(0.245\,\mathrm{m}\right)^{2} + \left(0.302\,\mathrm{m}\right)^{2}\right)}{2}}
\]

\noindent\textbf{Resultado}

\[
\fitmath{\sigma=1.732845884\times10^{8}\,\mathrm{Pa}}
\]

\noindent\textbf{Interpretación}
La tensión escala con el cuadrado de giro. No incluye concentraciones de raíz ni flexión.

\noindent\textbf{Comprobación de los pasos}
\begin{itemize}
\item Paso 2: $0.0155895$ — reproducido
\item Paso 4: $1.73285\times 10^{8}$ — reproducido
\end{itemize}

\noindent\textbf{Comprobación de la respuesta} la aritmética del último paso reproduce la respuesta enunciada a la precisión con que está escrita. Es una reproducción aritmética; no valida el modelo físico ni sus hipótesis.

\noindent\textbf{Limitaciones}
\begin{itemize}
\item Modelo educativo; la aplicabilidad a una máquina real requiere revisión experta.
\end{itemize}

\noindent\textbf{Errores frecuentes}
\begin{itemize}
\item Usar densidad del gas o rpm directamente.
\end{itemize}
% studium:end-record eyJpZCI6ICJQUkItMDA5NyIsICJzZWN0aW9uIjogImMxMyJ9
\end{studiumworked}

\begin{studiumworked}{Problema resuelto C13-01-P}


\noindent\textbf{Enunciado}
\(C13-01-P \cdot\) Tracción centrífuga uniforme. Ejercicio de transferencia. Calcule la tensión centrífuga nominal de raíz de una pala uniforme. Antes de consultar la solución, prediga el efecto de cambiar el primer dato respecto del caso resuelto de esta familia. Justifique la tendencia y explique qué restricción podría convertirse en limitante. Datos: \(\rho=8647.8\, \mathrm{kg}\,\mathrm{m}^{-3}\); \(\Omega=1170\, \mathrm{s}^{-1}\); \(r_h=0.245\, \mathrm{m}\); \(r_t=0.302\, \mathrm{m}\).

\noindent\textbf{Datos}
\begin{itemize}
\item $\rho = 8647.8\,\mathrm{kg}\,\mathrm{m}^{-3}$ — densidad material
\item $\Omega = 1170\,\mathrm{s}^{-1}$ — velocidad angular
\item $r_h = 0.245\,\mathrm{m}$ — radio raíz
\item $r_t = 0.302\,\mathrm{m}$ — radio punta
\end{itemize}

\noindent\textbf{Incógnita}
$\sigma$, en $\mathrm{Pa}$; también las magnitudes intermedias desarrolladas.

\noindent\textbf{Modelo físico}

\noindent\textit{Tracción centrífuga uniforme}

\[
\fitmath{\sigma=\frac{\Omega^{2} \rho \left(- r_h^{2} + r_t^{2}\right)}{2}}
\]

\noindent\textbf{Hipótesis}
\begin{itemize}
\item Barra radial uniforme, material homogéneo, régimen estacionario.
\end{itemize}

\noindent\textbf{Desarrollo paso a paso}

\begin{enumerate}

\item Se integra r entre raíz y punta.

\[\fitmath{\mathrm{integ}=\frac{{\mathrm{rt}}^{2}-{\mathrm{rh}}^{2}}{2}}\]

\[\fitmath{\mathrm{integ}=\frac{{(0.302)}^{2}-{(0.245)}^{2}}{2}\approx 0.0155895\,\mathrm{m^{2}}}\]

\item Se multiplica por densidad y velocidad angular al cuadrado.

\[\fitmath{\mathrm{stress}=\frac{\rho\cdot {\Omega}^{2}\cdot \left({\mathrm{rt}}^{2}-{\mathrm{rh}}^{2}\right)}{2}}\]

\[\fitmath{\mathrm{stress}=\frac{(8647.8)\cdot {(1170)}^{2}\cdot \left({(0.302)}^{2}-{(0.245)}^{2}\right)}{2}\approx 1.845481\times10^{8}\,\mathrm{Pa}}\]

\end{enumerate}


\noindent\textbf{Sustitución numérica}

\[
\fitmath{\frac{\left(\left(1170\,\mathrm{s}^{-1}\right)\right)^{2} \left(8647.8\,\mathrm{kg}\,\mathrm{m}^{-3}\right) \left(- \left(0.245\,\mathrm{m}\right)^{2} + \left(0.302\,\mathrm{m}\right)^{2}\right)}{2}}
\]

\noindent\textbf{Resultado}

\[
\fitmath{\sigma=1.845480866\times10^{8}\,\mathrm{Pa}}
\]

\noindent\textbf{Interpretación}
La tensión escala con el cuadrado de giro. No incluye concentraciones de raíz ni flexión.

\noindent\textbf{Comprobación de los pasos}
\begin{itemize}
\item Paso 2: $0.0155895$ — reproducido
\item Paso 4: $1.84548\times 10^{8}$ — reproducido
\end{itemize}

\noindent\textbf{Comprobación de la respuesta} la aritmética del último paso reproduce la respuesta enunciada a la precisión con que está escrita. Es una reproducción aritmética; no valida el modelo físico ni sus hipótesis.

\noindent\textbf{Limitaciones}
\begin{itemize}
\item Modelo educativo; la aplicabilidad a una máquina real requiere revisión experta.
\end{itemize}

\noindent\textbf{Errores frecuentes}
\begin{itemize}
\item Usar densidad del gas o rpm directamente.
\end{itemize}
% studium:end-record eyJpZCI6ICJQUkItMDA5OCIsICJzZWN0aW9uIjogImMxMyJ9
\end{studiumworked}

\begin{studiumworked}{Problema resuelto C13-02-R}


\noindent\textbf{Enunciado}
\(C13-02-R \cdot\) Pala estrechada. La sección disminuye linealmente desde Ah en la raíz hasta Ah(1-k) en punta. Calcule la tensión centrífuga nominal de raíz. Datos: \(\rho=8050\, \mathrm{kg}\,\mathrm{m}^{-3}\); \(\Omega=1210\, \mathrm{s}^{-1}\); \(r_h=0.238\, \mathrm{m}\); \(r_t=0.301\, \mathrm{m}\); \(\alpha_t=0.42\).

\noindent\textbf{Datos}
\begin{itemize}
\item $\rho = 8050\,\mathrm{kg}\,\mathrm{m}^{-3}$ — densidad material
\item $\Omega = 1210\,\mathrm{s}^{-1}$ — velocidad angular
\item $r_h = 0.238\,\mathrm{m}$ — radio raíz
\item $r_t = 0.301\,\mathrm{m}$ — radio punta
\item $\alpha_t = 0.42$ — fracción de estrechamiento
\end{itemize}

\noindent\textbf{Incógnita}
$\sigma$, en $\mathrm{Pa}$; también las magnitudes intermedias desarrolladas.

\noindent\textbf{Modelo físico}

\noindent\textit{Pala estrechada}

\[
\fitmath{\sigma=\Omega^{2} \rho \left(- \frac{\alpha_t \left(- \frac{r_h^{3}}{3} - \frac{r_h \left(- r_h^{2} + r_t^{2}\right)}{2} + \frac{r_t^{3}}{3}\right)}{- r_h + r_t} - \frac{r_h^{2}}{2} + \frac{r_t^{2}}{2}\right)}
\]

\noindent\textbf{Hipótesis}
\begin{itemize}
\item Área lineal positiva en todo el intervalo; sin carga aerodinámica.
\end{itemize}

\noindent\textbf{Desarrollo paso a paso}

\begin{enumerate}

\item El ancho radial define la ley de área.

\[\fitmath{b=\mathrm{rt}-\mathrm{rh}}\]

\[\fitmath{b=(0.301)-(0.238)\approx 0.063\,\mathrm{m}}\]

\item La integral se divide por el área raíz; Ah se cancela.

\[\fitmath{\mathrm{stress}=\rho\cdot {\Omega}^{2}\cdot \left(\frac{{\mathrm{rt}}^{2}-{\mathrm{rh}}^{2}}{2}-\frac{\mathrm{kk}}{\mathrm{rt}-\mathrm{rh}}\cdot \left(\frac{{\mathrm{rt}}^{3}-{\mathrm{rh}}^{3}}{3}-\frac{\mathrm{rh}\cdot \left({\mathrm{rt}}^{2}-{\mathrm{rh}}^{2}\right)}{2}\right)\right)}\]

\[\fitmath{\mathrm{stress}=(8050)\cdot {(1210)}^{2}\cdot \left(\frac{{(0.301)}^{2}-{(0.238)}^{2}}{2}-\frac{(0.42)}{(0.301)-(0.238)}\cdot \left(\frac{{(0.301)}^{3}-{(0.238)}^{3}}{3}-\frac{(0.238)\cdot \left({(0.301)}^{2}-{(0.238)}^{2}\right)}{2}\right)\right)\approx 1.564486\times10^{8}\,\mathrm{Pa}}\]

\end{enumerate}


\noindent\textbf{Sustitución numérica}

\[
\fitmath{\left(\left(1210\,\mathrm{s}^{-1}\right)\right)^{2} \left(8050\,\mathrm{kg}\,\mathrm{m}^{-3}\right) \left(- \frac{\left(0.42\right) \left(- \frac{\left(0.238\,\mathrm{m}\right)^{3}}{3} - \frac{\left(0.238\,\mathrm{m}\right) \left(- \left(0.238\,\mathrm{m}\right)^{2} + \left(0.301\,\mathrm{m}\right)^{2}\right)}{2} + \frac{\left(0.301\,\mathrm{m}\right)^{3}}{3}\right)}{- \left(0.238\,\mathrm{m}\right) + \left(0.301\,\mathrm{m}\right)} - \frac{\left(0.238\,\mathrm{m}\right)^{2}}{2} + \frac{\left(0.301\,\mathrm{m}\right)^{2}}{2}\right)}
\]

\noindent\textbf{Resultado}

\[
\fitmath{\sigma=1.564486090\times10^{8}\,\mathrm{Pa}}
\]

\noindent\textbf{Interpretación}
La reducción de masa exterior reduce la tracción raíz respecto de una pala uniforme del mismo material y radios.

\noindent\textbf{Comprobación de los pasos}
\begin{itemize}
\item Paso 2: $0.063$ — reproducido
\item Paso 4: $1.56449\times 10^{8}$ — reproducido
\end{itemize}

\noindent\textbf{Comprobación de la respuesta} la aritmética del último paso reproduce la respuesta enunciada a la precisión con que está escrita. Es una reproducción aritmética; no valida el modelo físico ni sus hipótesis.

\noindent\textbf{Limitaciones}
\begin{itemize}
\item Modelo educativo; la aplicabilidad a una máquina real requiere revisión experta.
\end{itemize}

\noindent\textbf{Errores frecuentes}
\begin{itemize}
\item Reducir tensión multiplicando por el área de punta sin integrar.
\end{itemize}
% studium:end-record eyJpZCI6ICJQUkItMDA5OSIsICJzZWN0aW9uIjogImMxMyJ9
\end{studiumworked}

\begin{studiumworked}{Problema resuelto C13-02-P}


\noindent\textbf{Enunciado}
\(C13-02-P \cdot\) Pala estrechada. Ejercicio de transferencia. La sección disminuye linealmente desde Ah en la raíz hasta Ah(1-k) en punta. Calcule la tensión centrífuga nominal de raíz. Antes de consultar la solución, prediga el efecto de cambiar el primer dato respecto del caso resuelto de esta familia. Justifique la tendencia y explique qué restricción podría convertirse en limitante. Datos: \(\rho=8573.25\, \mathrm{kg}\,\mathrm{m}^{-3}\); \(\Omega=1210\, \mathrm{s}^{-1}\); \(r_h=0.238\, \mathrm{m}\); \(r_t=0.301\, \mathrm{m}\); \(\alpha_t=0.42\).

\noindent\textbf{Datos}
\begin{itemize}
\item $\rho = 8573.25\,\mathrm{kg}\,\mathrm{m}^{-3}$ — densidad material
\item $\Omega = 1210\,\mathrm{s}^{-1}$ — velocidad angular
\item $r_h = 0.238\,\mathrm{m}$ — radio raíz
\item $r_t = 0.301\,\mathrm{m}$ — radio punta
\item $\alpha_t = 0.42$ — fracción de estrechamiento
\end{itemize}

\noindent\textbf{Incógnita}
$\sigma$, en $\mathrm{Pa}$; también las magnitudes intermedias desarrolladas.

\noindent\textbf{Modelo físico}

\noindent\textit{Pala estrechada}

\[
\fitmath{\sigma=\Omega^{2} \rho \left(- \frac{\alpha_t \left(- \frac{r_h^{3}}{3} - \frac{r_h \left(- r_h^{2} + r_t^{2}\right)}{2} + \frac{r_t^{3}}{3}\right)}{- r_h + r_t} - \frac{r_h^{2}}{2} + \frac{r_t^{2}}{2}\right)}
\]

\noindent\textbf{Hipótesis}
\begin{itemize}
\item Área lineal positiva en todo el intervalo; sin carga aerodinámica.
\end{itemize}

\noindent\textbf{Desarrollo paso a paso}

\begin{enumerate}

\item El ancho radial define la ley de área.

\[\fitmath{b=\mathrm{rt}-\mathrm{rh}}\]

\[\fitmath{b=(0.301)-(0.238)\approx 0.063\,\mathrm{m}}\]

\item La integral se divide por el área raíz; Ah se cancela.

\[\fitmath{\mathrm{stress}=\rho\cdot {\Omega}^{2}\cdot \left(\frac{{\mathrm{rt}}^{2}-{\mathrm{rh}}^{2}}{2}-\frac{\mathrm{kk}}{\mathrm{rt}-\mathrm{rh}}\cdot \left(\frac{{\mathrm{rt}}^{3}-{\mathrm{rh}}^{3}}{3}-\frac{\mathrm{rh}\cdot \left({\mathrm{rt}}^{2}-{\mathrm{rh}}^{2}\right)}{2}\right)\right)}\]

\[\fitmath{\mathrm{stress}=(8573.25)\cdot {(1210)}^{2}\cdot \left(\frac{{(0.301)}^{2}-{(0.238)}^{2}}{2}-\frac{(0.42)}{(0.301)-(0.238)}\cdot \left(\frac{{(0.301)}^{3}-{(0.238)}^{3}}{3}-\frac{(0.238)\cdot \left({(0.301)}^{2}-{(0.238)}^{2}\right)}{2}\right)\right)\approx 1.666178\times10^{8}\,\mathrm{Pa}}\]

\end{enumerate}


\noindent\textbf{Sustitución numérica}

\[
\fitmath{\left(\left(1210\,\mathrm{s}^{-1}\right)\right)^{2} \left(8573.25\,\mathrm{kg}\,\mathrm{m}^{-3}\right) \left(- \frac{\left(0.42\right) \left(- \frac{\left(0.238\,\mathrm{m}\right)^{3}}{3} - \frac{\left(0.238\,\mathrm{m}\right) \left(- \left(0.238\,\mathrm{m}\right)^{2} + \left(0.301\,\mathrm{m}\right)^{2}\right)}{2} + \frac{\left(0.301\,\mathrm{m}\right)^{3}}{3}\right)}{- \left(0.238\,\mathrm{m}\right) + \left(0.301\,\mathrm{m}\right)} - \frac{\left(0.238\,\mathrm{m}\right)^{2}}{2} + \frac{\left(0.301\,\mathrm{m}\right)^{2}}{2}\right)}
\]

\noindent\textbf{Resultado}

\[
\fitmath{\sigma=1.666177686\times10^{8}\,\mathrm{Pa}}
\]

\noindent\textbf{Interpretación}
La reducción de masa exterior reduce la tracción raíz respecto de una pala uniforme del mismo material y radios.

\noindent\textbf{Comprobación de los pasos}
\begin{itemize}
\item Paso 2: $0.063$ — reproducido
\item Paso 4: $1.66618\times 10^{8}$ — reproducido
\end{itemize}

\noindent\textbf{Comprobación de la respuesta} la aritmética del último paso reproduce la respuesta enunciada a la precisión con que está escrita. Es una reproducción aritmética; no valida el modelo físico ni sus hipótesis.

\noindent\textbf{Limitaciones}
\begin{itemize}
\item Modelo educativo; la aplicabilidad a una máquina real requiere revisión experta.
\end{itemize}

\noindent\textbf{Errores frecuentes}
\begin{itemize}
\item Reducir tensión multiplicando por el área de punta sin integrar.
\end{itemize}
% studium:end-record eyJpZCI6ICJQUkItMDEwMCIsICJzZWN0aW9uIjogImMxMyJ9
\end{studiumworked}

\begin{studiumworked}{Problema resuelto C13-03-R}


\noindent\textbf{Enunciado}
\(C13-03-R \cdot\) Flexión nominal y tracción. Una pala tiene fuerza aerodinámica equivalente aplicada al brazo dado y tracción nominal conocida. Halle la máxima tensión normal en la fibra indicada. Datos: \(F_t=186\, \mathrm{N}\); \(L_b=0.038\, \mathrm{m}\); \(y=0.0042\, \mathrm{m}\); \(I_b=8.4\times 10^{-10}\, \mathrm{m}^{4}\); \(\sigma_N=1.78\times 10^{8}\, \mathrm{Pa}\).

\noindent\textbf{Datos}
\begin{itemize}
\item $F_t = 186\,\mathrm{N}$ — fuerza transversal equivalente
\item $L_b = 0.038\,\mathrm{m}$ — brazo respecto raíz
\item $y = 0.0042\,\mathrm{m}$ — distancia de fibra al eje neutro
\item $I_b = 8.4\times 10^{-10}\,\mathrm{m}^{4}$ — segundo momento de área
\item $\sigma_N = 178000000\,\mathrm{Pa}$ — tracción centrífuga nominal
\end{itemize}

\noindent\textbf{Incógnita}
$\sigma$, en $\mathrm{Pa}$; también las magnitudes intermedias desarrolladas.

\noindent\textbf{Modelo físico}

\noindent\textit{Flexión nominal y tracción}

\[
\fitmath{\sigma=\sigma_N + \frac{L_b F_t y}{I_b}}
\]

\noindent\textbf{Hipótesis}
\begin{itemize}
\item Viga elástica lineal, cargas y sección equivalentes.
\end{itemize}

\noindent\textbf{Desarrollo paso a paso}

\begin{enumerate}

\item La fuerza equivalente genera el momento de raíz.

\[\fitmath{\mathrm{moment}=\mathrm{force}\cdot \mathrm{arm}}\]

\[\fitmath{\mathrm{moment}=(186)\cdot (0.038)\approx 7.068\,\mathrm{N\,m}}\]

\item Se añade la flexión en la fibra de máxima tracción.

\[\fitmath{\mathrm{stress}=\mathrm{sig}+\frac{\mathrm{force}\cdot \mathrm{arm}\cdot \mathrm{yy}}{\mathrm{Ib}}}\]

\[\fitmath{\mathrm{stress}=(178000000)+\frac{(186)\cdot (0.038)\cdot (0.0042)}{(8.4\times10^{-10})}\approx 2.1334\times10^{8}\,\mathrm{Pa}}\]

\end{enumerate}


\noindent\textbf{Sustitución numérica}

\[
\fitmath{\left(178000000\,\mathrm{Pa}\right) + \frac{\left(0.038\,\mathrm{m}\right) \left(186\,\mathrm{N}\right) \left(0.0042\,\mathrm{m}\right)}{\left(8.4e-10\,\mathrm{m}^{4}\right)}}
\]

\noindent\textbf{Resultado}

\[
\fitmath{\sigma=2.133400000\times10^{8}\,\mathrm{Pa}}
\]

\noindent\textbf{Interpretación}
La fibra opuesta resta la contribución de flexión. La suma nominal no evalúa fatiga ni concentraciones.

\noindent\textbf{Comprobación de los pasos}
\begin{itemize}
\item Paso 2: $7.068$ — reproducido
\item Paso 4: $213340000$ — reproducido
\end{itemize}

\noindent\textbf{Comprobación de la respuesta} la aritmética del último paso reproduce la respuesta enunciada a la precisión con que está escrita. Es una reproducción aritmética; no valida el modelo físico ni sus hipótesis.

\noindent\textbf{Limitaciones}
\begin{itemize}
\item Modelo educativo; la aplicabilidad a una máquina real requiere revisión experta.
\end{itemize}

\noindent\textbf{Errores frecuentes}
\begin{itemize}
\item Usar momento polar en una fórmula de flexión sin justificarlo.
\end{itemize}
% studium:end-record eyJpZCI6ICJQUkItMDEwMSIsICJzZWN0aW9uIjogImMxMyJ9
\end{studiumworked}

\begin{studiumworked}{Problema resuelto C13-03-P}


\noindent\textbf{Enunciado}
\(C13-03-P \cdot\) Flexión nominal y tracción. Ejercicio de transferencia. Una pala tiene fuerza aerodinámica equivalente aplicada al brazo dado y tracción nominal conocida. Halle la máxima tensión normal en la fibra indicada. Antes de consultar la solución, prediga el efecto de cambiar el primer dato respecto del caso resuelto de esta familia. Justifique la tendencia y explique qué restricción podría convertirse en limitante. Datos: \(F_t=198.09\, \mathrm{N}\); \(L_b=0.038\, \mathrm{m}\); \(y=0.0042\, \mathrm{m}\); \(I_b=8.4\times 10^{-10}\, \mathrm{m}^{4}\); \(\sigma_N=1.78\times 10^{8}\, \mathrm{Pa}\).

\noindent\textbf{Datos}
\begin{itemize}
\item $F_t = 198.09\,\mathrm{N}$ — fuerza transversal equivalente
\item $L_b = 0.038\,\mathrm{m}$ — brazo respecto raíz
\item $y = 0.0042\,\mathrm{m}$ — distancia de fibra al eje neutro
\item $I_b = 8.4\times 10^{-10}\,\mathrm{m}^{4}$ — segundo momento de área
\item $\sigma_N = 178000000\,\mathrm{Pa}$ — tracción centrífuga nominal
\end{itemize}

\noindent\textbf{Incógnita}
$\sigma$, en $\mathrm{Pa}$; también las magnitudes intermedias desarrolladas.

\noindent\textbf{Modelo físico}

\noindent\textit{Flexión nominal y tracción}

\[
\fitmath{\sigma=\sigma_N + \frac{L_b F_t y}{I_b}}
\]

\noindent\textbf{Hipótesis}
\begin{itemize}
\item Viga elástica lineal, cargas y sección equivalentes.
\end{itemize}

\noindent\textbf{Desarrollo paso a paso}

\begin{enumerate}

\item La fuerza equivalente genera el momento de raíz.

\[\fitmath{\mathrm{moment}=\mathrm{force}\cdot \mathrm{arm}}\]

\[\fitmath{\mathrm{moment}=(198.09)\cdot (0.038)\approx 7.52742\,\mathrm{N\,m}}\]

\item Se añade la flexión en la fibra de máxima tracción.

\[\fitmath{\mathrm{stress}=\mathrm{sig}+\frac{\mathrm{force}\cdot \mathrm{arm}\cdot \mathrm{yy}}{\mathrm{Ib}}}\]

\[\fitmath{\mathrm{stress}=(178000000)+\frac{(198.09)\cdot (0.038)\cdot (0.0042)}{(8.4\times10^{-10})}\approx 2.156371\times10^{8}\,\mathrm{Pa}}\]

\end{enumerate}


\noindent\textbf{Sustitución numérica}

\[
\fitmath{\left(178000000\,\mathrm{Pa}\right) + \frac{\left(0.038\,\mathrm{m}\right) \left(198.09\,\mathrm{N}\right) \left(0.0042\,\mathrm{m}\right)}{\left(8.4e-10\,\mathrm{m}^{4}\right)}}
\]

\noindent\textbf{Resultado}

\[
\fitmath{\sigma=2.156371000\times10^{8}\,\mathrm{Pa}}
\]

\noindent\textbf{Interpretación}
La fibra opuesta resta la contribución de flexión. La suma nominal no evalúa fatiga ni concentraciones.

\noindent\textbf{Comprobación de los pasos}
\begin{itemize}
\item Paso 2: $7.52742$ — reproducido
\item Paso 4: $215637100$ — reproducido
\end{itemize}

\noindent\textbf{Comprobación de la respuesta} la aritmética del último paso reproduce la respuesta enunciada a la precisión con que está escrita. Es una reproducción aritmética; no valida el modelo físico ni sus hipótesis.

\noindent\textbf{Limitaciones}
\begin{itemize}
\item Modelo educativo; la aplicabilidad a una máquina real requiere revisión experta.
\end{itemize}

\noindent\textbf{Errores frecuentes}
\begin{itemize}
\item Usar momento polar en una fórmula de flexión sin justificarlo.
\end{itemize}
% studium:end-record eyJpZCI6ICJQUkItMDEwMiIsICJzZWN0aW9uIjogImMxMyJ9
\end{studiumworked}

\begin{studiumworked}{Problema resuelto C13-04-R}


\noindent\textbf{Enunciado}
\(C13-04-R \cdot\) Frecuencia del voladizo educativo. Estime la primera frecuencia de una pala como viga uniforme no giratoria. Datos: \(E=1.85\times 10^{11}\, \mathrm{Pa}\); \(I_b=7.8\times 10^{-10}\, \mathrm{m}^{4}\); \(\rho=8030\, \mathrm{kg}\,\mathrm{m}^{-3}\); \(A_b=0.000165\, \mathrm{m}^{2}\); \(L=0.082\, \mathrm{m}\).

\noindent\textbf{Datos}
\begin{itemize}
\item $E = 185000000000\,\mathrm{Pa}$ — módulo elástico
\item $I_b = 7.8\times 10^{-10}\,\mathrm{m}^{4}$ — segundo momento
\item $\rho = 8030\,\mathrm{kg}\,\mathrm{m}^{-3}$ — densidad material
\item $A_b = 0.000165\,\mathrm{m}^{2}$ — área de sección
\item $L = 0.082\,\mathrm{m}$ — longitud
\end{itemize}

\noindent\textbf{Incógnita}
$f_1$, en $\mathrm{s^{-1}}$; también las magnitudes intermedias desarrolladas.

\noindent\textbf{Modelo físico}

\noindent\textit{Frecuencia del voladizo educativo}

\[
\fitmath{f_1=\frac{1.7578125 \sqrt{\frac{E I_b}{A_b \rho}}}{\pi L^{2}}}
\]

\noindent\textbf{Hipótesis}
\begin{itemize}
\item Voladizo uniforme Euler-Bernoulli, sin rotación ni amortiguamiento.
\end{itemize}

\noindent\textbf{Desarrollo paso a paso}

\begin{enumerate}

\item Se calcula la rigidez específica de flexión.

\[\fitmath{\mathrm{stiff}=\sqrt{\frac{\mathrm{Ee}\cdot \mathrm{Ib}}{\rho\cdot \mathrm{Ab}}}}\]

\[\fitmath{\mathrm{stiff}=\sqrt{\frac{(1.85\times10^{11})\cdot (7.8\times10^{-10})}{(8030)\cdot (0.000165)}}\approx 10.43598\,\mathrm{m^{2}/s}}\]

\item La primera raíz modal fija la frecuencia.

\[\fitmath{\mathrm{freq}=\frac{{1.875}^{2}}{2\cdot \pi\cdot {\mathrm{length}}^{2}}\cdot \sqrt{\frac{\mathrm{Ee}\cdot \mathrm{Ib}}{\rho\cdot \mathrm{Ab}}}}\]

\[\fitmath{\mathrm{freq}=\frac{{1.875}^{2}}{2\cdot \pi\cdot {(0.082)}^{2}}\cdot \sqrt{\frac{(1.85\times10^{11})\cdot (7.8\times10^{-10})}{(8030)\cdot (0.000165)}}\approx 868.4172\,\mathrm{s^{-1}}}\]

\end{enumerate}


\noindent\textbf{Sustitución numérica}

\[
\fitmath{\frac{1.7578125 \sqrt{\frac{\left(185000000000\,\mathrm{Pa}\right) \left(7.8e-10\,\mathrm{m}^{4}\right)}{\left(0.000165\,\mathrm{m}^{2}\right) \left(8030\,\mathrm{kg}\,\mathrm{m}^{-3}\right)}}}{\pi \left(0.082\,\mathrm{m}\right)^{2}}}
\]

\noindent\textbf{Resultado}

\[
\fitmath{f_1=8.684171505\times10^{2}\,\mathrm{s}^{-1}}
\]

\noindent\textbf{Interpretación}
Es una estimación no giratoria; no incluye rigidización centrífuga ni anclaje real. Compare después con órdenes de excitación definidos.

\noindent\textbf{Comprobación de los pasos}
\begin{itemize}
\item Paso 2: $10.436$ — reproducido
\item Paso 4: $868.417$ — reproducido
\end{itemize}

\noindent\textbf{Comprobación de la respuesta} la aritmética del último paso reproduce la respuesta enunciada a la precisión con que está escrita. Es una reproducción aritmética; no valida el modelo físico ni sus hipótesis.

\noindent\textbf{Limitaciones}
\begin{itemize}
\item Modelo educativo; la aplicabilidad a una máquina real requiere revisión experta.
\end{itemize}

\noindent\textbf{Errores frecuentes}
\begin{itemize}
\item Declarar ausencia de resonancia a partir de un único modelo estático.
\end{itemize}
% studium:end-record eyJpZCI6ICJQUkItMDEwMyIsICJzZWN0aW9uIjogImMxMyJ9
\end{studiumworked}

\begin{studiumworked}{Problema resuelto C13-04-P}


\noindent\textbf{Enunciado}
\(C13-04-P \cdot\) Frecuencia del voladizo educativo. Ejercicio de transferencia. Estime la primera frecuencia de una pala como viga uniforme no giratoria. Antes de consultar la solución, prediga el efecto de cambiar el primer dato respecto del caso resuelto de esta familia. Justifique la tendencia y explique qué restricción podría convertirse en limitante. Datos: \(E=1.97025\times 10^{11}\, \mathrm{Pa}\); \(I_b=7.8\times 10^{-10}\, \mathrm{m}^{4}\); \(\rho=8030\, \mathrm{kg}\,\mathrm{m}^{-3}\); \(A_b=0.000165\, \mathrm{m}^{2}\); \(L=0.082\, \mathrm{m}\).

\noindent\textbf{Datos}
\begin{itemize}
\item $E = 197025000000\,\mathrm{Pa}$ — módulo elástico
\item $I_b = 7.8\times 10^{-10}\,\mathrm{m}^{4}$ — segundo momento
\item $\rho = 8030\,\mathrm{kg}\,\mathrm{m}^{-3}$ — densidad material
\item $A_b = 0.000165\,\mathrm{m}^{2}$ — área de sección
\item $L = 0.082\,\mathrm{m}$ — longitud
\end{itemize}

\noindent\textbf{Incógnita}
$f_1$, en $\mathrm{s^{-1}}$; también las magnitudes intermedias desarrolladas.

\noindent\textbf{Modelo físico}

\noindent\textit{Frecuencia del voladizo educativo}

\[
\fitmath{f_1=\frac{1.7578125 \sqrt{\frac{E I_b}{A_b \rho}}}{\pi L^{2}}}
\]

\noindent\textbf{Hipótesis}
\begin{itemize}
\item Voladizo uniforme Euler-Bernoulli, sin rotación ni amortiguamiento.
\end{itemize}

\noindent\textbf{Desarrollo paso a paso}

\begin{enumerate}

\item Se calcula la rigidez específica de flexión.

\[\fitmath{\mathrm{stiff}=\sqrt{\frac{\mathrm{Ee}\cdot \mathrm{Ib}}{\rho\cdot \mathrm{Ab}}}}\]

\[\fitmath{\mathrm{stiff}=\sqrt{\frac{(1.97025\times10^{11})\cdot (7.8\times10^{-10})}{(8030)\cdot (0.000165)}}\approx 10.76981\,\mathrm{m^{2}/s}}\]

\item La primera raíz modal fija la frecuencia.

\[\fitmath{\mathrm{freq}=\frac{{1.875}^{2}}{2\cdot \pi\cdot {\mathrm{length}}^{2}}\cdot \sqrt{\frac{\mathrm{Ee}\cdot \mathrm{Ib}}{\rho\cdot \mathrm{Ab}}}}\]

\[\fitmath{\mathrm{freq}=\frac{{1.875}^{2}}{2\cdot \pi\cdot {(0.082)}^{2}}\cdot \sqrt{\frac{(1.97025\times10^{11})\cdot (7.8\times10^{-10})}{(8030)\cdot (0.000165)}}\approx 896.1964\,\mathrm{s^{-1}}}\]

\end{enumerate}


\noindent\textbf{Sustitución numérica}

\[
\fitmath{\frac{1.7578125 \sqrt{\frac{\left(197025000000\,\mathrm{Pa}\right) \left(7.8e-10\,\mathrm{m}^{4}\right)}{\left(0.000165\,\mathrm{m}^{2}\right) \left(8030\,\mathrm{kg}\,\mathrm{m}^{-3}\right)}}}{\pi \left(0.082\,\mathrm{m}\right)^{2}}}
\]

\noindent\textbf{Resultado}

\[
\fitmath{f_1=8.961964014\times10^{2}\,\mathrm{s}^{-1}}
\]

\noindent\textbf{Interpretación}
Es una estimación no giratoria; no incluye rigidización centrífuga ni anclaje real. Compare después con órdenes de excitación definidos.

\noindent\textbf{Comprobación de los pasos}
\begin{itemize}
\item Paso 2: $10.7698$ — reproducido
\item Paso 4: $896.196$ — reproducido
\end{itemize}

\noindent\textbf{Comprobación de la respuesta} la aritmética del último paso reproduce la respuesta enunciada a la precisión con que está escrita. Es una reproducción aritmética; no valida el modelo físico ni sus hipótesis.

\noindent\textbf{Limitaciones}
\begin{itemize}
\item Modelo educativo; la aplicabilidad a una máquina real requiere revisión experta.
\end{itemize}

\noindent\textbf{Errores frecuentes}
\begin{itemize}
\item Declarar ausencia de resonancia a partir de un único modelo estático.
\end{itemize}
% studium:end-record eyJpZCI6ICJQUkItMDEwNCIsICJzZWN0aW9uIjogImMxMyJ9
\end{studiumworked}

\chapter{Transferencia térmica y refrigeración de álabes}
\begin{studiumworked}{Problema resuelto C14-01-R}


\noindent\textbf{Enunciado}
\(C14-01-R \cdot\) Red térmica de pared. Calcule flujo térmico y temperatura exterior del metal en una pared plana con convección a ambos lados. Datos: \(T_{aw}=1620\, \mathrm{K}\); \(T_c=745\, \mathrm{K}\); \(h_g=1680\, \mathrm{W}\,\mathrm{m}^{-2}\,\mathrm{K}^{-1}\); \(h_c=2450\, \mathrm{W}\,\mathrm{m}^{-2}\,\mathrm{K}^{-1}\); \(t=0.0011\, \mathrm{m}\); \(k_m=21\, \mathrm{W}\,\mathrm{m}^{-1}\,\mathrm{K}^{-1}\).

\noindent\textbf{Datos}
\begin{itemize}
\item $T_{aw} = 1620\,\mathrm{K}$ — temperatura recuperación exterior
\item $T_c = 745\,\mathrm{K}$ — temperatura refrigerante local
\item $h_g = 1680\,\mathrm{W}\,\mathrm{m}^{-2}\,\mathrm{K}^{-1}$ — coeficiente exterior
\item $h_c = 2450\,\mathrm{W}\,\mathrm{m}^{-2}\,\mathrm{K}^{-1}$ — coeficiente interior
\item $t = 0.0011\,\mathrm{m}$ — espesor
\item $k_m = 21\,\mathrm{W}\,\mathrm{m}^{-1}\,\mathrm{K}^{-1}$ — conductividad metal
\end{itemize}

\noindent\textbf{Incógnita}
$T_w$, en $\mathrm{K}$; también las magnitudes intermedias desarrolladas.

\noindent\textbf{Modelo físico}

\noindent\textit{Red térmica de pared}

\[
\fitmath{T_w=T_{aw} - \frac{- T_c + T_{aw}}{h_g \left(\frac{1}{h_g} + \frac{1}{h_c} + \frac{t}{k_m}\right)}}
\]

\noindent\textbf{Hipótesis}
\begin{itemize}
\item Estacionario, una dimensión, sin radiación, coeficientes constantes.
\end{itemize}

\noindent\textbf{Desarrollo paso a paso}

\begin{enumerate}

\item Las resistencias se suman por unidad de área.

\[\fitmath{\mathrm{Res}=\frac{1}{\mathrm{hg}}+\frac{\mathrm{th}}{\mathrm{cond}}+\frac{1}{\mathrm{hc}}}\]

\[\fitmath{\mathrm{Res}=\frac{1}{(1680)}+\frac{(0.0011)}{(21)}+\frac{1}{(2450)}\approx 0.001055782\,\mathrm{m^{2}\,K/W}}\]

\item El salto térmico impulsa el flujo.

\[\fitmath{\mathrm{flux}=\frac{\mathrm{Tg}-\mathrm{Tc}}{\frac{1}{\mathrm{hg}}+\frac{\mathrm{th}}{\mathrm{cond}}+\frac{1}{\mathrm{hc}}}}\]

\[\fitmath{\mathrm{flux}=\frac{(1620)-(745)}{\frac{1}{(1680)}+\frac{(0.0011)}{(21)}+\frac{1}{(2450)}}\approx 828769.3\,\mathrm{W/m^{2}}}\]

\item La caída exterior determina la cara caliente.

\[\fitmath{\mathrm{Tw}=\mathrm{Tg}-\frac{\mathrm{Tg}-\mathrm{Tc}}{\mathrm{hg}\cdot \left(\frac{1}{\mathrm{hg}}+\frac{\mathrm{th}}{\mathrm{cond}}+\frac{1}{\mathrm{hc}}\right)}}\]

\[\fitmath{\mathrm{Tw}=(1620)-\frac{(1620)-(745)}{(1680)\cdot \left(\frac{1}{(1680)}+\frac{(0.0011)}{(21)}+\frac{1}{(2450)}\right)}\approx 1126.685\,\mathrm{K}}\]

\end{enumerate}


\noindent\textbf{Sustitución numérica}

\[
\fitmath{\left(1620\,\mathrm{K}\right) - \frac{- \left(745\,\mathrm{K}\right) + \left(1620\,\mathrm{K}\right)}{\left(1680\,\mathrm{W}\,\mathrm{m}^{-2}\,\mathrm{K}^{-1}\right) \left(\frac{1}{\left(1680\,\mathrm{W}\,\mathrm{m}^{-2}\,\mathrm{K}^{-1}\right)} + \frac{1}{\left(2450\,\mathrm{W}\,\mathrm{m}^{-2}\,\mathrm{K}^{-1}\right)} + \frac{\left(0.0011\,\mathrm{m}\right)}{\left(21\,\mathrm{W}\,\mathrm{m}^{-1}\,\mathrm{K}^{-1}\right)}\right)}}
\]

\noindent\textbf{Resultado}

\[
\fitmath{T_w=1.126684923\times10^{3}\,\mathrm{K}}
\]

\noindent\textbf{Interpretación}
La temperatura del metal debe quedar entre refrigerante y recuperación exterior. La diferencia de caras corresponde al flujo por conducción.

\noindent\textbf{Comprobación de los pasos}
\begin{itemize}
\item Paso 2: $0.00105578$ — reproducido
\item Paso 4: $828769$ — reproducido
\item Paso 6: $1126.68$ — reproducido
\end{itemize}

\noindent\textbf{Comprobación de la respuesta} la aritmética del último paso reproduce la respuesta enunciada a la precisión con que está escrita. Es una reproducción aritmética; no valida el modelo físico ni sus hipótesis.

\noindent\textbf{Limitaciones}
\begin{itemize}
\item Modelo educativo; la aplicabilidad a una máquina real requiere revisión experta.
\end{itemize}

\noindent\textbf{Errores frecuentes}
\begin{itemize}
\item Confundir área anular con área de intercambio.
\end{itemize}
% studium:end-record eyJpZCI6ICJQUkItMDEwNSIsICJzZWN0aW9uIjogImMxNCJ9
\end{studiumworked}

\begin{studiumworked}{Problema resuelto C14-01-P}


\noindent\textbf{Enunciado}
\(C14-01-P \cdot\) Red térmica de pared. Ejercicio de transferencia. Calcule flujo térmico y temperatura exterior del metal en una pared plana con convección a ambos lados. Antes de consultar la solución, prediga el efecto de cambiar el primer dato respecto del caso resuelto de esta familia. Justifique la tendencia y explique qué restricción podría convertirse en limitante. Datos: \(T_{aw}=1725.3\, \mathrm{K}\); \(T_c=745\, \mathrm{K}\); \(h_g=1680\, \mathrm{W}\,\mathrm{m}^{-2}\,\mathrm{K}^{-1}\); \(h_c=2450\, \mathrm{W}\,\mathrm{m}^{-2}\,\mathrm{K}^{-1}\); \(t=0.0011\, \mathrm{m}\); \(k_m=21\, \mathrm{W}\,\mathrm{m}^{-1}\,\mathrm{K}^{-1}\).

\noindent\textbf{Datos}
\begin{itemize}
\item $T_{aw} = 1725.3\,\mathrm{K}$ — temperatura recuperación exterior
\item $T_c = 745\,\mathrm{K}$ — temperatura refrigerante local
\item $h_g = 1680\,\mathrm{W}\,\mathrm{m}^{-2}\,\mathrm{K}^{-1}$ — coeficiente exterior
\item $h_c = 2450\,\mathrm{W}\,\mathrm{m}^{-2}\,\mathrm{K}^{-1}$ — coeficiente interior
\item $t = 0.0011\,\mathrm{m}$ — espesor
\item $k_m = 21\,\mathrm{W}\,\mathrm{m}^{-1}\,\mathrm{K}^{-1}$ — conductividad metal
\end{itemize}

\noindent\textbf{Incógnita}
$T_w$, en $\mathrm{K}$; también las magnitudes intermedias desarrolladas.

\noindent\textbf{Modelo físico}

\noindent\textit{Red térmica de pared}

\[
\fitmath{T_w=T_{aw} - \frac{- T_c + T_{aw}}{h_g \left(\frac{1}{h_g} + \frac{1}{h_c} + \frac{t}{k_m}\right)}}
\]

\noindent\textbf{Hipótesis}
\begin{itemize}
\item Estacionario, una dimensión, sin radiación, coeficientes constantes.
\end{itemize}

\noindent\textbf{Desarrollo paso a paso}

\begin{enumerate}

\item Las resistencias se suman por unidad de área.

\[\fitmath{\mathrm{Res}=\frac{1}{\mathrm{hg}}+\frac{\mathrm{th}}{\mathrm{cond}}+\frac{1}{\mathrm{hc}}}\]

\[\fitmath{\mathrm{Res}=\frac{1}{(1680)}+\frac{(0.0011)}{(21)}+\frac{1}{(2450)}\approx 0.001055782\,\mathrm{m^{2}\,K/W}}\]

\item El salto térmico impulsa el flujo.

\[\fitmath{\mathrm{flux}=\frac{\mathrm{Tg}-\mathrm{Tc}}{\frac{1}{\mathrm{hg}}+\frac{\mathrm{th}}{\mathrm{cond}}+\frac{1}{\mathrm{hc}}}}\]

\[\fitmath{\mathrm{flux}=\frac{(1725.3)-(745)}{\frac{1}{(1680)}+\frac{(0.0011)}{(21)}+\frac{1}{(2450)}}\approx 928505.8\,\mathrm{W/m^{2}}}\]

\item La caída exterior determina la cara caliente.

\[\fitmath{\mathrm{Tw}=\mathrm{Tg}-\frac{\mathrm{Tg}-\mathrm{Tc}}{\mathrm{hg}\cdot \left(\frac{1}{\mathrm{hg}}+\frac{\mathrm{th}}{\mathrm{cond}}+\frac{1}{\mathrm{hc}}\right)}}\]

\[\fitmath{\mathrm{Tw}=(1725.3)-\frac{(1725.3)-(745)}{(1680)\cdot \left(\frac{1}{(1680)}+\frac{(0.0011)}{(21)}+\frac{1}{(2450)}\right)}\approx 1172.618\,\mathrm{K}}\]

\end{enumerate}


\noindent\textbf{Sustitución numérica}

\[
\fitmath{\left(1725.3\,\mathrm{K}\right) - \frac{- \left(745\,\mathrm{K}\right) + \left(1725.3\,\mathrm{K}\right)}{\left(1680\,\mathrm{W}\,\mathrm{m}^{-2}\,\mathrm{K}^{-1}\right) \left(\frac{1}{\left(1680\,\mathrm{W}\,\mathrm{m}^{-2}\,\mathrm{K}^{-1}\right)} + \frac{1}{\left(2450\,\mathrm{W}\,\mathrm{m}^{-2}\,\mathrm{K}^{-1}\right)} + \frac{\left(0.0011\,\mathrm{m}\right)}{\left(21\,\mathrm{W}\,\mathrm{m}^{-1}\,\mathrm{K}^{-1}\right)}\right)}}
\]

\noindent\textbf{Resultado}

\[
\fitmath{T_w=1.172617977\times10^{3}\,\mathrm{K}}
\]

\noindent\textbf{Interpretación}
La temperatura del metal debe quedar entre refrigerante y recuperación exterior. La diferencia de caras corresponde al flujo por conducción.

\noindent\textbf{Comprobación de los pasos}
\begin{itemize}
\item Paso 2: $0.00105578$ — reproducido
\item Paso 4: $928506$ — reproducido
\item Paso 6: $1172.62$ — reproducido
\end{itemize}

\noindent\textbf{Comprobación de la respuesta} la aritmética del último paso reproduce la respuesta enunciada a la precisión con que está escrita. Es una reproducción aritmética; no valida el modelo físico ni sus hipótesis.

\noindent\textbf{Limitaciones}
\begin{itemize}
\item Modelo educativo; la aplicabilidad a una máquina real requiere revisión experta.
\end{itemize}

\noindent\textbf{Errores frecuentes}
\begin{itemize}
\item Confundir área anular con área de intercambio.
\end{itemize}
% studium:end-record eyJpZCI6ICJQUkItMDEwNiIsICJzZWN0aW9uIjogImMxNCJ9
\end{studiumworked}

\begin{studiumworked}{Problema resuelto C14-02-R}


\noindent\textbf{Enunciado}
\(C14-02-R \cdot\) Caudal energético de refrigeración. Una superficie recibe flujo térmico medio conocido. Calcule el caudal refrigerante que absorbería ese calor con el incremento permitido. Datos: \(q''=620000\, \mathrm{W}\,\mathrm{m}^{-2}\); \(A=0.034\, \mathrm{m}^{2}\); \(c_p=1090\, \mathrm{J}\,\mathrm{kg}^{-1}\,\mathrm{K}^{-1}\); \(\Delta T_c=185\, \mathrm{K}\).

\noindent\textbf{Datos}
\begin{itemize}
\item $q'' = 620000\,\mathrm{W}\,\mathrm{m}^{-2}$ — flujo térmico medio
\item $A = 0.034\,\mathrm{m}^{2}$ — área de intercambio
\item $c_p = 1090\,\mathrm{J}\,\mathrm{kg}^{-1}\,\mathrm{K}^{-1}$ — calor específico refrigerante
\item $\Delta T_c = 185\,\mathrm{K}$ — incremento de temperatura permitido
\end{itemize}

\noindent\textbf{Incógnita}
$\dot{m}_c$, en $\mathrm{kg\,s^{-1}}$; también las magnitudes intermedias desarrolladas.

\noindent\textbf{Modelo físico}

\noindent\textit{Caudal energético de refrigeración}

\[
\fitmath{\dot{m}_c=\frac{A q''}{c_p \Delta T_c}}
\]

\noindent\textbf{Hipótesis}
\begin{itemize}
\item Sin pérdidas térmicas adicionales; cp constante.
\end{itemize}

\noindent\textbf{Desarrollo paso a paso}

\begin{enumerate}

\item Se integra el flujo térmico uniforme sobre la superficie.

\[\fitmath{\mathrm{heat}=\mathrm{flux}\cdot \mathrm{area}}\]

\[\fitmath{\mathrm{heat}=(620000)\cdot (0.034)\approx 21080\,\mathrm{W}}\]

\item El balance de entalpía fija el caudal.

\[\fitmath{\mathrm{mc}=\frac{\mathrm{flux}\cdot \mathrm{area}}{\mathrm{cp}\cdot \mathrm{dt}}}\]

\[\fitmath{\mathrm{mc}=\frac{(620000)\cdot (0.034)}{(1090)\cdot (185)}\approx 0.1045376\,\mathrm{kg/s}}\]

\end{enumerate}


\noindent\textbf{Sustitución numérica}

\[
\fitmath{\frac{\left(0.034\,\mathrm{m}^{2}\right) \left(620000\,\mathrm{W}\,\mathrm{m}^{-2}\right)}{\left(1090\,\mathrm{J}\,\mathrm{kg}^{-1}\,\mathrm{K}^{-1}\right) \left(185\,\mathrm{K}\right)}}
\]

\noindent\textbf{Resultado}

\[
\fitmath{\dot{m}_c=1.045375651\times10^{-1}\,\mathrm{kg}\,\mathrm{s}^{-1}}
\]

\noindent\textbf{Interpretación}
El resultado es un mínimo del balance energético. La convección y la presión disponible deben comprobarse aparte.

\noindent\textbf{Comprobación de los pasos}
\begin{itemize}
\item Paso 2: $21080$ — reproducido
\item Paso 4: $0.104538$ — reproducido
\end{itemize}

\noindent\textbf{Comprobación de la respuesta} la aritmética del último paso reproduce la respuesta enunciada a la precisión con que está escrita. Es una reproducción aritmética; no valida el modelo físico ni sus hipótesis.

\noindent\textbf{Limitaciones}
\begin{itemize}
\item Modelo educativo; la aplicabilidad a una máquina real requiere revisión experta.
\end{itemize}

\noindent\textbf{Errores frecuentes}
\begin{itemize}
\item Tomar este mínimo como caudal seguro de una pala real.
\end{itemize}
% studium:end-record eyJpZCI6ICJQUkItMDEwNyIsICJzZWN0aW9uIjogImMxNCJ9
\end{studiumworked}

\begin{studiumworked}{Problema resuelto C14-02-P}


\noindent\textbf{Enunciado}
\(C14-02-P \cdot\) Caudal energético de refrigeración. Ejercicio de transferencia. Una superficie recibe flujo térmico medio conocido. Calcule el caudal refrigerante que absorbería ese calor con el incremento permitido. Antes de consultar la solución, prediga el efecto de cambiar el primer dato respecto del caso resuelto de esta familia. Justifique la tendencia y explique qué restricción podría convertirse en limitante. Datos: \(q''=660300\, \mathrm{W}\,\mathrm{m}^{-2}\); \(A=0.034\, \mathrm{m}^{2}\); \(c_p=1090\, \mathrm{J}\,\mathrm{kg}^{-1}\,\mathrm{K}^{-1}\); \(\Delta T_c=185\, \mathrm{K}\).

\noindent\textbf{Datos}
\begin{itemize}
\item $q'' = 660300\,\mathrm{W}\,\mathrm{m}^{-2}$ — flujo térmico medio
\item $A = 0.034\,\mathrm{m}^{2}$ — área de intercambio
\item $c_p = 1090\,\mathrm{J}\,\mathrm{kg}^{-1}\,\mathrm{K}^{-1}$ — calor específico refrigerante
\item $\Delta T_c = 185\,\mathrm{K}$ — incremento de temperatura permitido
\end{itemize}

\noindent\textbf{Incógnita}
$\dot{m}_c$, en $\mathrm{kg\,s^{-1}}$; también las magnitudes intermedias desarrolladas.

\noindent\textbf{Modelo físico}

\noindent\textit{Caudal energético de refrigeración}

\[
\fitmath{\dot{m}_c=\frac{A q''}{c_p \Delta T_c}}
\]

\noindent\textbf{Hipótesis}
\begin{itemize}
\item Sin pérdidas térmicas adicionales; cp constante.
\end{itemize}

\noindent\textbf{Desarrollo paso a paso}

\begin{enumerate}

\item Se integra el flujo térmico uniforme sobre la superficie.

\[\fitmath{\mathrm{heat}=\mathrm{flux}\cdot \mathrm{area}}\]

\[\fitmath{\mathrm{heat}=(660300)\cdot (0.034)\approx 22450.2\,\mathrm{W}}\]

\item El balance de entalpía fija el caudal.

\[\fitmath{\mathrm{mc}=\frac{\mathrm{flux}\cdot \mathrm{area}}{\mathrm{cp}\cdot \mathrm{dt}}}\]

\[\fitmath{\mathrm{mc}=\frac{(660300)\cdot (0.034)}{(1090)\cdot (185)}\approx 0.1113325\,\mathrm{kg/s}}\]

\end{enumerate}


\noindent\textbf{Sustitución numérica}

\[
\fitmath{\frac{\left(0.034\,\mathrm{m}^{2}\right) \left(660300\,\mathrm{W}\,\mathrm{m}^{-2}\right)}{\left(1090\,\mathrm{J}\,\mathrm{kg}^{-1}\,\mathrm{K}^{-1}\right) \left(185\,\mathrm{K}\right)}}
\]

\noindent\textbf{Resultado}

\[
\fitmath{\dot{m}_c=1.113325068\times10^{-1}\,\mathrm{kg}\,\mathrm{s}^{-1}}
\]

\noindent\textbf{Interpretación}
El resultado es un mínimo del balance energético. La convección y la presión disponible deben comprobarse aparte.

\noindent\textbf{Comprobación de los pasos}
\begin{itemize}
\item Paso 2: $22450.2$ — reproducido
\item Paso 4: $0.111333$ — reproducido
\end{itemize}

\noindent\textbf{Comprobación de la respuesta} la aritmética del último paso reproduce la respuesta enunciada a la precisión con que está escrita. Es una reproducción aritmética; no valida el modelo físico ni sus hipótesis.

\noindent\textbf{Limitaciones}
\begin{itemize}
\item Modelo educativo; la aplicabilidad a una máquina real requiere revisión experta.
\end{itemize}

\noindent\textbf{Errores frecuentes}
\begin{itemize}
\item Tomar este mínimo como caudal seguro de una pala real.
\end{itemize}
% studium:end-record eyJpZCI6ICJQUkItMDEwOCIsICJzZWN0aW9uIjogImMxNCJ9
\end{studiumworked}

\begin{studiumworked}{Problema resuelto C14-03-R}


\noindent\textbf{Enunciado}
\(C14-03-R \cdot\) Película y recuperación. Con efectividad de película prescrita, calcule la recuperación efectiva y el flujo térmico a una pared de temperatura dada. Datos: \(T_{aw}=1690\, \mathrm{K}\); \(T_c=775\, \mathrm{K}\); \(\eta_f=0.41\); \(T_w=1130\, \mathrm{K}\); \(h_g=1770\, \mathrm{W}\,\mathrm{m}^{-2}\,\mathrm{K}^{-1}\).

\noindent\textbf{Datos}
\begin{itemize}
\item $T_{aw} = 1690\,\mathrm{K}$ — recuperación sin película
\item $T_c = 775\,\mathrm{K}$ — temperatura refrigerante
\item $\eta_f = 0.41$ — efectividad de película
\item $T_w = 1130\,\mathrm{K}$ — temperatura pared
\item $h_g = 1770\,\mathrm{W}\,\mathrm{m}^{-2}\,\mathrm{K}^{-1}$ — coeficiente exterior compatible con inyección
\end{itemize}

\noindent\textbf{Incógnita}
$q''$, en $\mathrm{W\,m^{-2}}$; también las magnitudes intermedias desarrolladas.

\noindent\textbf{Modelo físico}

\noindent\textit{Película y recuperación}

\[
\fitmath{q''=h_g \left(T_{aw} - T_w - \eta_f \left(- T_c + T_{aw}\right)\right)}
\]

\noindent\textbf{Hipótesis}
\begin{itemize}
\item Efectividad y coeficiente dados para la misma geometría y punto.
\end{itemize}

\noindent\textbf{Desarrollo paso a paso}

\begin{enumerate}

\item Se despeja la recuperación con película.

\[\fitmath{\mathrm{Taw}=\mathrm{Tg}-\mathrm{eff}\cdot \left(\mathrm{Tg}-\mathrm{Tc}\right)}\]

\[\fitmath{\mathrm{Taw}=(1690)-(0.41)\cdot \left((1690)-(775)\right)\approx 1314.85\,\mathrm{K}}\]

\item Se aplica convección con el coeficiente especificado.

\[\fitmath{\mathrm{flux}=\mathrm{hg}\cdot \left(\mathrm{Tg}-\mathrm{eff}\cdot \left(\mathrm{Tg}-\mathrm{Tc}\right)-\mathrm{Tw}\right)}\]

\[\fitmath{\mathrm{flux}=(1770)\cdot \left((1690)-(0.41)\cdot \left((1690)-(775)\right)-(1130)\right)\approx 327184.5\,\mathrm{W/m^{2}}}\]

\end{enumerate}


\noindent\textbf{Sustitución numérica}

\[
\fitmath{\left(1770\,\mathrm{W}\,\mathrm{m}^{-2}\,\mathrm{K}^{-1}\right) \left(\left(1690\,\mathrm{K}\right) - \left(1130\,\mathrm{K}\right) - \left(0.41\right) \left(- \left(775\,\mathrm{K}\right) + \left(1690\,\mathrm{K}\right)\right)\right)}
\]

\noindent\textbf{Resultado}

\[
\fitmath{q''=3.271845000\times10^{5}\,\mathrm{W}\,\mathrm{m}^{-2}}
\]

\noindent\textbf{Interpretación}
La efectividad no sustituye al coeficiente de convección. Verifique que el flujo conserva el sentido esperado.

\noindent\textbf{Comprobación de los pasos}
\begin{itemize}
\item Paso 2: $1314.85$ — reproducido
\item Paso 4: $327185$ — reproducido
\end{itemize}

\noindent\textbf{Comprobación de la respuesta} la aritmética del último paso reproduce la respuesta enunciada a la precisión con que está escrita. Es una reproducción aritmética; no valida el modelo físico ni sus hipótesis.

\noindent\textbf{Limitaciones}
\begin{itemize}
\item Modelo educativo; la aplicabilidad a una máquina real requiere revisión experta.
\end{itemize}

\noindent\textbf{Errores frecuentes}
\begin{itemize}
\item Identificar efectividad de película y eficiencia de turbina.
\end{itemize}
% studium:end-record eyJpZCI6ICJQUkItMDEwOSIsICJzZWN0aW9uIjogImMxNCJ9
\end{studiumworked}

\begin{studiumworked}{Problema resuelto C14-03-P}


\noindent\textbf{Enunciado}
\(C14-03-P \cdot\) Película y recuperación. Ejercicio de transferencia. Con efectividad de película prescrita, calcule la recuperación efectiva y el flujo térmico a una pared de temperatura dada. Antes de consultar la solución, prediga el efecto de cambiar el primer dato respecto del caso resuelto de esta familia. Justifique la tendencia y explique qué restricción podría convertirse en limitante. Datos: \(T_{aw}=1799.85\, \mathrm{K}\); \(T_c=775\, \mathrm{K}\); \(\eta_f=0.41\); \(T_w=1130\, \mathrm{K}\); \(h_g=1770\, \mathrm{W}\,\mathrm{m}^{-2}\,\mathrm{K}^{-1}\).

\noindent\textbf{Datos}
\begin{itemize}
\item $T_{aw} = 1799.85\,\mathrm{K}$ — recuperación sin película
\item $T_c = 775\,\mathrm{K}$ — temperatura refrigerante
\item $\eta_f = 0.41$ — efectividad de película
\item $T_w = 1130\,\mathrm{K}$ — temperatura pared
\item $h_g = 1770\,\mathrm{W}\,\mathrm{m}^{-2}\,\mathrm{K}^{-1}$ — coeficiente exterior compatible con inyección
\end{itemize}

\noindent\textbf{Incógnita}
$q''$, en $\mathrm{W\,m^{-2}}$; también las magnitudes intermedias desarrolladas.

\noindent\textbf{Modelo físico}

\noindent\textit{Película y recuperación}

\[
\fitmath{q''=h_g \left(T_{aw} - T_w - \eta_f \left(- T_c + T_{aw}\right)\right)}
\]

\noindent\textbf{Hipótesis}
\begin{itemize}
\item Efectividad y coeficiente dados para la misma geometría y punto.
\end{itemize}

\noindent\textbf{Desarrollo paso a paso}

\begin{enumerate}

\item Se despeja la recuperación con película.

\[\fitmath{\mathrm{Taw}=\mathrm{Tg}-\mathrm{eff}\cdot \left(\mathrm{Tg}-\mathrm{Tc}\right)}\]

\[\fitmath{\mathrm{Taw}=(1799.85)-(0.41)\cdot \left((1799.85)-(775)\right)\approx 1379.661\,\mathrm{K}}\]

\item Se aplica convección con el coeficiente especificado.

\[\fitmath{\mathrm{flux}=\mathrm{hg}\cdot \left(\mathrm{Tg}-\mathrm{eff}\cdot \left(\mathrm{Tg}-\mathrm{Tc}\right)-\mathrm{Tw}\right)}\]

\[\fitmath{\mathrm{flux}=(1770)\cdot \left((1799.85)-(0.41)\cdot \left((1799.85)-(775)\right)-(1130)\right)\approx 441900.9\,\mathrm{W/m^{2}}}\]

\end{enumerate}


\noindent\textbf{Sustitución numérica}

\[
\fitmath{\left(1770\,\mathrm{W}\,\mathrm{m}^{-2}\,\mathrm{K}^{-1}\right) \left(\left(1799.85\,\mathrm{K}\right) - \left(1130\,\mathrm{K}\right) - \left(0.41\right) \left(- \left(775\,\mathrm{K}\right) + \left(1799.85\,\mathrm{K}\right)\right)\right)}
\]

\noindent\textbf{Resultado}

\[
\fitmath{q''=4.419008550\times10^{5}\,\mathrm{W}\,\mathrm{m}^{-2}}
\]

\noindent\textbf{Interpretación}
La efectividad no sustituye al coeficiente de convección. Verifique que el flujo conserva el sentido esperado.

\noindent\textbf{Comprobación de los pasos}
\begin{itemize}
\item Paso 2: $1379.66$ — reproducido
\item Paso 4: $441901$ — reproducido
\end{itemize}

\noindent\textbf{Comprobación de la respuesta} la aritmética del último paso reproduce la respuesta enunciada a la precisión con que está escrita. Es una reproducción aritmética; no valida el modelo físico ni sus hipótesis.

\noindent\textbf{Limitaciones}
\begin{itemize}
\item Modelo educativo; la aplicabilidad a una máquina real requiere revisión experta.
\end{itemize}

\noindent\textbf{Errores frecuentes}
\begin{itemize}
\item Identificar efectividad de película y eficiencia de turbina.
\end{itemize}
% studium:end-record eyJpZCI6ICJQUkItMDExMCIsICJzZWN0aW9uIjogImMxNCJ9
\end{studiumworked}

\begin{studiumworked}{Problema resuelto C14-04-R}


\noindent\textbf{Enunciado}
\(C14-04-R \cdot\) Restricción térmica uniaxial. Una barra homogénea tiene impedida completamente su dilatación longitudinal. Determine la deformación térmica libre y la tensión elástica de compresión. Datos: \(E=1.74\times 10^{11}\, \mathrm{Pa}\); \(\alpha_{\mathrm{th}}=1.41\times 10^{-5}\, \mathrm{K}^{-1}\); \(\Delta T=128\, \mathrm{K}\).

\noindent\textbf{Datos}
\begin{itemize}
\item $E = 174000000000\,\mathrm{Pa}$ — módulo elástico
\item $\alpha_{\mathrm{th}} = 1.41\times 10^{-5}\,\mathrm{K}^{-1}$ — coeficiente dilatación
\item $\Delta T = 128\,\mathrm{K}$ — aumento uniforme de temperatura
\end{itemize}

\noindent\textbf{Incógnita}
$\sigma$, en $\mathrm{Pa}$; también las magnitudes intermedias desarrolladas.

\noindent\textbf{Modelo físico}

\noindent\textit{Restricción térmica uniaxial}

\[
\fitmath{\sigma=- E \alpha_{\mathrm{th}} \Delta T}
\]

\noindent\textbf{Hipótesis}
\begin{itemize}
\item Restricción total uniaxial, elasticidad lineal, sin relajación.
\end{itemize}

\noindent\textbf{Desarrollo paso a paso}

\begin{enumerate}

\item La expansión libre sería positiva.

\[\fitmath{\mathrm{strain}=\mathrm{alp}\cdot \mathrm{dt}}\]

\[\fitmath{\mathrm{strain}=(1.41\times10^{-5})\cdot (128)\approx 0.0018048\,}\]

\item La restricción exige tensión elástica de signo opuesto.

\[\fitmath{\mathrm{stress}=-\mathrm{Ee}\cdot \mathrm{alp}\cdot \mathrm{dt}}\]

\[\fitmath{\mathrm{stress}=-(1.74\times10^{11})\cdot (1.41\times10^{-5})\cdot (128)\approx -3.140352\times10^{8}\,\mathrm{Pa}}\]

\end{enumerate}


\noindent\textbf{Sustitución numérica}

\[
\fitmath{- \left(174000000000\,\mathrm{Pa}\right) \left(1.41e-05\,\mathrm{K}^{-1}\right) \left(128\,\mathrm{K}\right)}
\]

\noindent\textbf{Resultado}

\[
\fitmath{\sigma=-3.140352000\times10^{8}\,\mathrm{Pa}}
\]

\noindent\textbf{Interpretación}
El signo negativo indica compresión. Una barra libre tendría tensión térmica nula con la misma temperatura uniforme.

\noindent\textbf{Comprobación de los pasos}
\begin{itemize}
\item Paso 2: $0.0018048$ — reproducido
\item Paso 4: $-314035200$ — reproducido
\end{itemize}

\noindent\textbf{Comprobación de la respuesta} la aritmética del último paso reproduce la respuesta enunciada a la precisión con que está escrita. Es una reproducción aritmética; no valida el modelo físico ni sus hipótesis.

\noindent\textbf{Limitaciones}
\begin{itemize}
\item Modelo educativo; la aplicabilidad a una máquina real requiere revisión experta.
\end{itemize}

\noindent\textbf{Errores frecuentes}
\begin{itemize}
\item Aplicar esta tensión a toda pared con gradiente y curvatura.
\end{itemize}
% studium:end-record eyJpZCI6ICJQUkItMDExMSIsICJzZWN0aW9uIjogImMxNCJ9
\end{studiumworked}

\begin{studiumworked}{Problema resuelto C14-04-P}


\noindent\textbf{Enunciado}
\(C14-04-P \cdot\) Restricción térmica uniaxial. Ejercicio de transferencia. Una barra homogénea tiene impedida completamente su dilatación longitudinal. Determine la deformación térmica libre y la tensión elástica de compresión. Antes de consultar la solución, prediga el efecto de cambiar el primer dato respecto del caso resuelto de esta familia. Justifique la tendencia y explique qué restricción podría convertirse en limitante. Datos: \(E=1.8531\times 10^{11}\, \mathrm{Pa}\); \(\alpha_{\mathrm{th}}=1.41\times 10^{-5}\, \mathrm{K}^{-1}\); \(\Delta T=128\, \mathrm{K}\).

\noindent\textbf{Datos}
\begin{itemize}
\item $E = 185310000000\,\mathrm{Pa}$ — módulo elástico
\item $\alpha_{\mathrm{th}} = 1.41\times 10^{-5}\,\mathrm{K}^{-1}$ — coeficiente dilatación
\item $\Delta T = 128\,\mathrm{K}$ — aumento uniforme de temperatura
\end{itemize}

\noindent\textbf{Incógnita}
$\sigma$, en $\mathrm{Pa}$; también las magnitudes intermedias desarrolladas.

\noindent\textbf{Modelo físico}

\noindent\textit{Restricción térmica uniaxial}

\[
\fitmath{\sigma=- E \alpha_{\mathrm{th}} \Delta T}
\]

\noindent\textbf{Hipótesis}
\begin{itemize}
\item Restricción total uniaxial, elasticidad lineal, sin relajación.
\end{itemize}

\noindent\textbf{Desarrollo paso a paso}

\begin{enumerate}

\item La expansión libre sería positiva.

\[\fitmath{\mathrm{strain}=\mathrm{alp}\cdot \mathrm{dt}}\]

\[\fitmath{\mathrm{strain}=(1.41\times10^{-5})\cdot (128)\approx 0.0018048\,}\]

\item La restricción exige tensión elástica de signo opuesto.

\[\fitmath{\mathrm{stress}=-\mathrm{Ee}\cdot \mathrm{alp}\cdot \mathrm{dt}}\]

\[\fitmath{\mathrm{stress}=-(1.8531\times10^{11})\cdot (1.41\times10^{-5})\cdot (128)\approx -3.344475\times10^{8}\,\mathrm{Pa}}\]

\end{enumerate}


\noindent\textbf{Sustitución numérica}

\[
\fitmath{- \left(185310000000\,\mathrm{Pa}\right) \left(1.41e-05\,\mathrm{K}^{-1}\right) \left(128\,\mathrm{K}\right)}
\]

\noindent\textbf{Resultado}

\[
\fitmath{\sigma=-3.344474880\times10^{8}\,\mathrm{Pa}}
\]

\noindent\textbf{Interpretación}
El signo negativo indica compresión. Una barra libre tendría tensión térmica nula con la misma temperatura uniforme.

\noindent\textbf{Comprobación de los pasos}
\begin{itemize}
\item Paso 2: $0.0018048$ — reproducido
\item Paso 4: $-334447488$ — reproducido
\end{itemize}

\noindent\textbf{Comprobación de la respuesta} la aritmética del último paso reproduce la respuesta enunciada a la precisión con que está escrita. Es una reproducción aritmética; no valida el modelo físico ni sus hipótesis.

\noindent\textbf{Limitaciones}
\begin{itemize}
\item Modelo educativo; la aplicabilidad a una máquina real requiere revisión experta.
\end{itemize}

\noindent\textbf{Errores frecuentes}
\begin{itemize}
\item Aplicar esta tensión a toda pared con gradiente y curvatura.
\end{itemize}
% studium:end-record eyJpZCI6ICJQUkItMDExMiIsICJzZWN0aW9uIjogImMxNCJ9
\end{studiumworked}

\chapter{Compresores radiales: rodete, inductor y deslizamiento}
\begin{studiumworked}{Problema resuelto C15-01-R}


\noindent\textbf{Enunciado}
\(C15-01-R \cdot\) Rodete radial con deslizamiento. Sin prerrotación y con álabes radiales ideales, calcule remolino real, trabajo y aumento total de temperatura. Datos: \(U=426\, \mathrm{m}\,\mathrm{s}^{-1}\); \(\sigma_s=0.87\); \(c_p=1004.5\, \mathrm{J}\,\mathrm{kg}^{-1}\,\mathrm{K}^{-1}\).

\noindent\textbf{Datos}
\begin{itemize}
\item $U = 426\,\mathrm{m}\,\mathrm{s}^{-1}$ — velocidad periférica salida
\item $\sigma_s = 0.87$ — factor definido como $C_\theta/C_{\theta,\mathrm{ideal}}$
\item $c_p = 1004.5\,\mathrm{J}\,\mathrm{kg}^{-1}\,\mathrm{K}^{-1}$ — calor específico
\end{itemize}

\noindent\textbf{Incógnita}
$\Delta T_0$, en $\mathrm{K}$; también las magnitudes intermedias desarrolladas.

\noindent\textbf{Modelo físico}

\noindent\textit{Rodete radial con deslizamiento}

\[
\fitmath{\Delta T_0=\frac{\sigma_s U^{2}}{c_p}}
\]

\noindent\textbf{Hipótesis}
\begin{itemize}
\item Sin prerrotación; salida ideal radial antes de aplicar slip.
\end{itemize}

\noindent\textbf{Desarrollo paso a paso}

\begin{enumerate}

\item Se reduce el remolino ideal con la definición indicada.

\[\fitmath{\mathrm{ct}=\mathrm{slip}\cdot \mathrm{uu}}\]

\[\fitmath{\mathrm{ct}=(0.87)\cdot (426)\approx 370.62\,\mathrm{m/s}}\]

\item Euler usa el remolino real.

\[\fitmath{\mathrm{wc}=\mathrm{slip}\cdot {\mathrm{uu}}^{2}}\]

\[\fitmath{\mathrm{wc}=(0.87)\cdot {(426)}^{2}\approx 157884.1\,\mathrm{J/kg}}\]

\item El balance energético fija el calentamiento total.

\[\fitmath{\mathrm{dT}=\frac{\mathrm{slip}\cdot {\mathrm{uu}}^{2}}{\mathrm{cp}}}\]

\[\fitmath{\mathrm{dT}=\frac{(0.87)\cdot {(426)}^{2}}{(1004.5)}\approx 157.1768\,\mathrm{K}}\]

\end{enumerate}


\noindent\textbf{Sustitución numérica}

\[
\fitmath{\frac{\left(0.87\right) \left(\left(426\,\mathrm{m}\,\mathrm{s}^{-1}\right)\right)^{2}}{\left(1004.5\,\mathrm{J}\,\mathrm{kg}^{-1}\,\mathrm{K}^{-1}\right)}}
\]

\noindent\textbf{Resultado}

\[
\fitmath{\Delta T_0=1.571768243\times10^{2}\,\mathrm{K}}
\]

\noindent\textbf{Interpretación}
El deslizamiento modifica trabajo. La presión necesita además una eficiencia o modelo de pérdida.

\noindent\textbf{Comprobación de los pasos}
\begin{itemize}
\item Paso 2: $370.62$ — reproducido
\item Paso 4: $157884$ — reproducido
\item Paso 6: $157.177$ — reproducido
\end{itemize}

\noindent\textbf{Comprobación de la respuesta} la aritmética del último paso reproduce la respuesta enunciada a la precisión con que está escrita. Es una reproducción aritmética; no valida el modelo físico ni sus hipótesis.

\noindent\textbf{Limitaciones}
\begin{itemize}
\item Modelo educativo; la aplicabilidad a una máquina real requiere revisión experta.
\end{itemize}

\noindent\textbf{Errores frecuentes}
\begin{itemize}
\item Multiplicar de nuevo por slip un trabajo que ya utiliza $C_\theta$ real.
\end{itemize}
% studium:end-record eyJpZCI6ICJQUkItMDExMyIsICJzZWN0aW9uIjogImMxNSJ9
\end{studiumworked}

\begin{studiumworked}{Problema resuelto C15-01-P}


\noindent\textbf{Enunciado}
\(C15-01-P \cdot\) Rodete radial con deslizamiento. Ejercicio de transferencia. Sin prerrotación y con álabes radiales ideales, calcule remolino real, trabajo y aumento total de temperatura. Antes de consultar la solución, prediga el efecto de cambiar el primer dato respecto del caso resuelto de esta familia. Justifique la tendencia y explique qué restricción podría convertirse en limitante. Datos: \(U=453.69\, \mathrm{m}\,\mathrm{s}^{-1}\); \(\sigma_s=0.87\); \(c_p=1004.5\, \mathrm{J}\,\mathrm{kg}^{-1}\,\mathrm{K}^{-1}\).

\noindent\textbf{Datos}
\begin{itemize}
\item $U = 453.69\,\mathrm{m}\,\mathrm{s}^{-1}$ — velocidad periférica salida
\item $\sigma_s = 0.87$ — factor definido como $C_\theta/C_{\theta,\mathrm{ideal}}$
\item $c_p = 1004.5\,\mathrm{J}\,\mathrm{kg}^{-1}\,\mathrm{K}^{-1}$ — calor específico
\end{itemize}

\noindent\textbf{Incógnita}
$\Delta T_0$, en $\mathrm{K}$; también las magnitudes intermedias desarrolladas.

\noindent\textbf{Modelo físico}

\noindent\textit{Rodete radial con deslizamiento}

\[
\fitmath{\Delta T_0=\frac{\sigma_s U^{2}}{c_p}}
\]

\noindent\textbf{Hipótesis}
\begin{itemize}
\item Sin prerrotación; salida ideal radial antes de aplicar slip.
\end{itemize}

\noindent\textbf{Desarrollo paso a paso}

\begin{enumerate}

\item Se reduce el remolino ideal con la definición indicada.

\[\fitmath{\mathrm{ct}=\mathrm{slip}\cdot \mathrm{uu}}\]

\[\fitmath{\mathrm{ct}=(0.87)\cdot (453.69)\approx 394.7103\,\mathrm{m/s}}\]

\item Euler usa el remolino real.

\[\fitmath{\mathrm{wc}=\mathrm{slip}\cdot {\mathrm{uu}}^{2}}\]

\[\fitmath{\mathrm{wc}=(0.87)\cdot {(453.69)}^{2}\approx 179076.1\,\mathrm{J/kg}}\]

\item El balance energético fija el calentamiento total.

\[\fitmath{\mathrm{dT}=\frac{\mathrm{slip}\cdot {\mathrm{uu}}^{2}}{\mathrm{cp}}}\]

\[\fitmath{\mathrm{dT}=\frac{(0.87)\cdot {(453.69)}^{2}}{(1004.5)}\approx 178.2739\,\mathrm{K}}\]

\end{enumerate}


\noindent\textbf{Sustitución numérica}

\[
\fitmath{\frac{\left(0.87\right) \left(\left(453.69\,\mathrm{m}\,\mathrm{s}^{-1}\right)\right)^{2}}{\left(1004.5\,\mathrm{J}\,\mathrm{kg}^{-1}\,\mathrm{K}^{-1}\right)}}
\]

\noindent\textbf{Resultado}

\[
\fitmath{\Delta T_0=1.782738835\times10^{2}\,\mathrm{K}}
\]

\noindent\textbf{Interpretación}
El deslizamiento modifica trabajo. La presión necesita además una eficiencia o modelo de pérdida.

\noindent\textbf{Comprobación de los pasos}
\begin{itemize}
\item Paso 2: $394.71$ — reproducido
\item Paso 4: $179076$ — reproducido
\item Paso 6: $178.274$ — reproducido
\end{itemize}

\noindent\textbf{Comprobación de la respuesta} la aritmética del último paso reproduce la respuesta enunciada a la precisión con que está escrita. Es una reproducción aritmética; no valida el modelo físico ni sus hipótesis.

\noindent\textbf{Limitaciones}
\begin{itemize}
\item Modelo educativo; la aplicabilidad a una máquina real requiere revisión experta.
\end{itemize}

\noindent\textbf{Errores frecuentes}
\begin{itemize}
\item Multiplicar de nuevo por slip un trabajo que ya utiliza $C_\theta$ real.
\end{itemize}
% studium:end-record eyJpZCI6ICJQUkItMDExNCIsICJzZWN0aW9uIjogImMxNSJ9
\end{studiumworked}

\begin{studiumworked}{Problema resuelto C15-02-R}


\noindent\textbf{Enunciado}
\(C15-02-R \cdot\) Curvatura hacia atrás. Un rodete tiene ángulo relativo desde la dirección meridiana negativo. Halle el trabajo con deslizamiento definido como cociente de remolinos. Datos: \(U=448\, \mathrm{m}\,\mathrm{s}^{-1}\); \(C_{m,s}=139\, \mathrm{m}\,\mathrm{s}^{-1}\); \(\lambda_s=-0.5061455\, \mathrm{rad}\); \(\sigma_s=0.89\).

\noindent\textbf{Datos}
\begin{itemize}
\item $U = 448\,\mathrm{m}\,\mathrm{s}^{-1}$ — velocidad periférica salida
\item $C_{m,s} = 139\,\mathrm{m}\,\mathrm{s}^{-1}$ — velocidad meridiana salida
\item $\lambda_s = -0.50614548307836\,\mathrm{rad}$ — ángulo relativo ideal desde meridiana
\item $\sigma_s = 0.89$ — factor de deslizamiento
\end{itemize}

\noindent\textbf{Incógnita}
$w_c$, en $\mathrm{J\,kg^{-1}}$; también las magnitudes intermedias desarrolladas.

\noindent\textbf{Modelo físico}

\noindent\textit{Curvatura hacia atrás}

\[
\fitmath{w_c=\sigma_s U \left(C_{m,s} \tan{\left(\lambda_s \right)} + U\right)}
\]

\noindent\textbf{Hipótesis}
\begin{itemize}
\item Factor prescrito positivo; salida ideal con remolino positivo.
\end{itemize}

\noindent\textbf{Desarrollo paso a paso}

\begin{enumerate}

\item El metal ideal determina la componente tangencial absoluta.

\[\fitmath{\mathrm{ctid}=\mathrm{uu}+\mathrm{cm}\cdot \tan\left(\mathrm{lam}\right)}\]

\[\fitmath{\mathrm{ctid}=(448)+(139)\cdot \tan\left((-0.5061454831)\right)\approx 370.951\,\mathrm{m/s}}\]

\item Se aplica el cociente de remolinos definido.

\[\fitmath{\mathrm{ct}=\mathrm{slip}\cdot \left(\mathrm{uu}+\mathrm{cm}\cdot \tan\left(\mathrm{lam}\right)\right)}\]

\[\fitmath{\mathrm{ct}=(0.89)\cdot \left((448)+(139)\cdot \tan\left((-0.5061454831)\right)\right)\approx 330.1464\,\mathrm{m/s}}\]

\item Sin prerrotación, Euler queda U2 Ctheta2.

\[\fitmath{\mathrm{wc}=\mathrm{uu}\cdot \mathrm{slip}\cdot \left(\mathrm{uu}+\mathrm{cm}\cdot \tan\left(\mathrm{lam}\right)\right)}\]

\[\fitmath{\mathrm{wc}=(448)\cdot (0.89)\cdot \left((448)+(139)\cdot \tan\left((-0.5061454831)\right)\right)\approx 147905.6\,\mathrm{J/kg}}\]

\end{enumerate}


\noindent\textbf{Sustitución numérica}

\[
\fitmath{\left(0.89\right) \left(448\,\mathrm{m}\,\mathrm{s}^{-1}\right) \left(\left(139\,\mathrm{m}\,\mathrm{s}^{-1}\right) \tan{\left(\left(-0.50614548307836\,\mathrm{rad}\right) \right)} + \left(448\,\mathrm{m}\,\mathrm{s}^{-1}\right)\right)}
\]

\noindent\textbf{Resultado}

\[
\fitmath{w_c=1.479055994\times10^{5}\,\mathrm{J}\,\mathrm{kg}^{-1}}
\]

\noindent\textbf{Interpretación}
La curvatura hacia atrás reduce el remolino ideal para estos datos. No intercambie esta definición con $1-dC_\theta/U$.

\noindent\textbf{Comprobación de los pasos}
\begin{itemize}
\item Paso 2: $370.951$ — reproducido
\item Paso 4: $330.146$ — reproducido
\item Paso 6: $147906$ — reproducido
\end{itemize}

\noindent\textbf{Comprobación de la respuesta} la aritmética del último paso reproduce la respuesta enunciada a la precisión con que está escrita. Es una reproducción aritmética; no valida el modelo físico ni sus hipótesis.

\noindent\textbf{Limitaciones}
\begin{itemize}
\item Modelo educativo; la aplicabilidad a una máquina real requiere revisión experta.
\end{itemize}

\noindent\textbf{Errores frecuentes}
\begin{itemize}
\item Usar una cotangente con el mismo ángulo desde meridiana.
\end{itemize}
% studium:end-record eyJpZCI6ICJQUkItMDExNSIsICJzZWN0aW9uIjogImMxNSJ9
\end{studiumworked}

\begin{studiumworked}{Problema resuelto C15-02-P}


\noindent\textbf{Enunciado}
\(C15-02-P \cdot\) Curvatura hacia atrás. Ejercicio de transferencia. Un rodete tiene ángulo relativo desde la dirección meridiana negativo. Halle el trabajo con deslizamiento definido como cociente de remolinos. Antes de consultar la solución, prediga el efecto de cambiar el primer dato respecto del caso resuelto de esta familia. Justifique la tendencia y explique qué restricción podría convertirse en limitante. Datos: \(U=477.12\, \mathrm{m}\,\mathrm{s}^{-1}\); \(C_{m,s}=139\, \mathrm{m}\,\mathrm{s}^{-1}\); \(\lambda_s=-0.5061455\, \mathrm{rad}\); \(\sigma_s=0.89\).

\noindent\textbf{Datos}
\begin{itemize}
\item $U = 477.12\,\mathrm{m}\,\mathrm{s}^{-1}$ — velocidad periférica salida
\item $C_{m,s} = 139\,\mathrm{m}\,\mathrm{s}^{-1}$ — velocidad meridiana salida
\item $\lambda_s = -0.50614548307836\,\mathrm{rad}$ — ángulo relativo ideal desde meridiana
\item $\sigma_s = 0.89$ — factor de deslizamiento
\end{itemize}

\noindent\textbf{Incógnita}
$w_c$, en $\mathrm{J\,kg^{-1}}$; también las magnitudes intermedias desarrolladas.

\noindent\textbf{Modelo físico}

\noindent\textit{Curvatura hacia atrás}

\[
\fitmath{w_c=\sigma_s U \left(C_{m,s} \tan{\left(\lambda_s \right)} + U\right)}
\]

\noindent\textbf{Hipótesis}
\begin{itemize}
\item Factor prescrito positivo; salida ideal con remolino positivo.
\end{itemize}

\noindent\textbf{Desarrollo paso a paso}

\begin{enumerate}

\item El metal ideal determina la componente tangencial absoluta.

\[\fitmath{\mathrm{ctid}=\mathrm{uu}+\mathrm{cm}\cdot \tan\left(\mathrm{lam}\right)}\]

\[\fitmath{\mathrm{ctid}=(477.12)+(139)\cdot \tan\left((-0.5061454831)\right)\approx 400.071\,\mathrm{m/s}}\]

\item Se aplica el cociente de remolinos definido.

\[\fitmath{\mathrm{ct}=\mathrm{slip}\cdot \left(\mathrm{uu}+\mathrm{cm}\cdot \tan\left(\mathrm{lam}\right)\right)}\]

\[\fitmath{\mathrm{ct}=(0.89)\cdot \left((477.12)+(139)\cdot \tan\left((-0.5061454831)\right)\right)\approx 356.0632\,\mathrm{m/s}}\]

\item Sin prerrotación, Euler queda U2 Ctheta2.

\[\fitmath{\mathrm{wc}=\mathrm{uu}\cdot \mathrm{slip}\cdot \left(\mathrm{uu}+\mathrm{cm}\cdot \tan\left(\mathrm{lam}\right)\right)}\]

\[\fitmath{\mathrm{wc}=(477.12)\cdot (0.89)\cdot \left((477.12)+(139)\cdot \tan\left((-0.5061454831)\right)\right)\approx 169884.9\,\mathrm{J/kg}}\]

\end{enumerate}


\noindent\textbf{Sustitución numérica}

\[
\fitmath{\left(0.89\right) \left(477.12\,\mathrm{m}\,\mathrm{s}^{-1}\right) \left(\left(139\,\mathrm{m}\,\mathrm{s}^{-1}\right) \tan{\left(\left(-0.50614548307836\,\mathrm{rad}\right) \right)} + \left(477.12\,\mathrm{m}\,\mathrm{s}^{-1}\right)\right)}
\]

\noindent\textbf{Resultado}

\[
\fitmath{w_c=1.698848870\times10^{5}\,\mathrm{J}\,\mathrm{kg}^{-1}}
\]

\noindent\textbf{Interpretación}
La curvatura hacia atrás reduce el remolino ideal para estos datos. No intercambie esta definición con $1-dC_\theta/U$.

\noindent\textbf{Comprobación de los pasos}
\begin{itemize}
\item Paso 2: $400.071$ — reproducido
\item Paso 4: $356.063$ — reproducido
\item Paso 6: $169885$ — reproducido
\end{itemize}

\noindent\textbf{Comprobación de la respuesta} la aritmética del último paso reproduce la respuesta enunciada a la precisión con que está escrita. Es una reproducción aritmética; no valida el modelo físico ni sus hipótesis.

\noindent\textbf{Limitaciones}
\begin{itemize}
\item Modelo educativo; la aplicabilidad a una máquina real requiere revisión experta.
\end{itemize}

\noindent\textbf{Errores frecuentes}
\begin{itemize}
\item Usar una cotangente con el mismo ángulo desde meridiana.
\end{itemize}
% studium:end-record eyJpZCI6ICJQUkItMDExNiIsICJzZWN0aW9uIjogImMxNSJ9
\end{studiumworked}

\begin{studiumworked}{Problema resuelto C15-03-R}


\noindent\textbf{Enunciado}
\(C15-03-R \cdot\) Inductor y Mach de punta. Un inductor sin prerrotación recibe gas con temperatura estática dada. Halle la velocidad relativa de punta y su Mach. Datos: \(\Omega=1920\, \mathrm{s}^{-1}\); \(r_t=0.102\, \mathrm{m}\); \(C_x=148\, \mathrm{m}\,\mathrm{s}^{-1}\); \(T_e=283\, \mathrm{K}\); \(\gamma=1.4\); \(R_g=287\, \mathrm{J}\,\mathrm{kg}^{-1}\,\mathrm{K}^{-1}\).

\noindent\textbf{Datos}
\begin{itemize}
\item $\Omega = 1920\,\mathrm{s}^{-1}$ — velocidad angular
\item $r_t = 0.102\,\mathrm{m}$ — radio punta inductor
\item $C_x = 148\,\mathrm{m}\,\mathrm{s}^{-1}$ — velocidad axial
\item $T_e = 283\,\mathrm{K}$ — temperatura estática
\item $\gamma = 1.4$ — razón de calores
\item $R_g = 287\,\mathrm{J}\,\mathrm{kg}^{-1}\,\mathrm{K}^{-1}$ — constante gas
\end{itemize}

\noindent\textbf{Incógnita}
$\mathrm{Ma}_r$ (adimensional); también las magnitudes intermedias desarrolladas.

\noindent\textbf{Modelo físico}

\noindent\textit{Inductor y Mach de punta}

\[
\fitmath{\mathrm{Ma}_r=\frac{\sqrt{\Omega^{2} r_t^{2} + C_x^{2}}}{\sqrt{R_g T_e \gamma}}}
\]

\noindent\textbf{Hipótesis}
\begin{itemize}
\item Sin prerrotación y velocidad axial uniforme.
\end{itemize}

\noindent\textbf{Desarrollo paso a paso}

\begin{enumerate}

\item La punta tiene mayor velocidad periférica que la raíz.

\[\fitmath{\mathrm{Ut}=\Omega\cdot \mathrm{rt}}\]

\[\fitmath{\mathrm{Ut}=(1920)\cdot (0.102)\approx 195.84\,\mathrm{m/s}}\]

\item Se forma el módulo relativo de entrada.

\[\fitmath{W=\sqrt{{\mathrm{cx}}^{2}+{\left(\Omega\cdot \mathrm{rt}\right)}^{2}}}\]

\[\fitmath{W=\sqrt{{(148)}^{2}+{\left((1920)\cdot (0.102)\right)}^{2}}\approx 245.4736\,\mathrm{m/s}}\]

\item El sonido corresponde al estado estático.

\[\fitmath{\mathrm{Mr}=\frac{\sqrt{{\mathrm{cx}}^{2}+{\left(\Omega\cdot \mathrm{rt}\right)}^{2}}}{\sqrt{\gamma\cdot R_g\cdot \mathrm{Te}}}}\]

\[\fitmath{\mathrm{Mr}=\frac{\sqrt{{(148)}^{2}+{\left((1920)\cdot (0.102)\right)}^{2}}}{\sqrt{(1.4)\cdot (287)\cdot (283)}}\approx 0.7279586\,}\]

\end{enumerate}


\noindent\textbf{Sustitución numérica}

\[
\fitmath{\frac{\sqrt{\left(\left(1920\,\mathrm{s}^{-1}\right)\right)^{2} \left(0.102\,\mathrm{m}\right)^{2} + \left(\left(148\,\mathrm{m}\,\mathrm{s}^{-1}\right)\right)^{2}}}{\sqrt{\left(287\,\mathrm{J}\,\mathrm{kg}^{-1}\,\mathrm{K}^{-1}\right) \left(283\,\mathrm{K}\right) \left(1.4\right)}}}
\]

\noindent\textbf{Resultado}

\[
\fitmath{\mathrm{Ma}_r=7.279585817\times10^{-1}}
\]

\noindent\textbf{Interpretación}
La temperatura estática se comparte entre referencias; la velocidad relativa aumenta por la rotación.

\noindent\textbf{Comprobación de los pasos}
\begin{itemize}
\item Paso 2: $195.84$ — reproducido
\item Paso 4: $245.474$ — reproducido
\item Paso 6: $0.727959$ — reproducido
\end{itemize}

\noindent\textbf{Comprobación de la respuesta} la aritmética del último paso reproduce la respuesta enunciada a la precisión con que está escrita. Es una reproducción aritmética; no valida el modelo físico ni sus hipótesis.

\noindent\textbf{Limitaciones}
\begin{itemize}
\item Modelo educativo; la aplicabilidad a una máquina real requiere revisión experta.
\end{itemize}

\noindent\textbf{Errores frecuentes}
\begin{itemize}
\item Usar radio exterior del rodete en el ojo del inductor.
\end{itemize}
% studium:end-record eyJpZCI6ICJQUkItMDExNyIsICJzZWN0aW9uIjogImMxNSJ9
\end{studiumworked}

\begin{studiumworked}{Problema resuelto C15-03-P}


\noindent\textbf{Enunciado}
\(C15-03-P \cdot\) Inductor y Mach de punta. Ejercicio de transferencia. Un inductor sin prerrotación recibe gas con temperatura estática dada. Halle la velocidad relativa de punta y su Mach. Antes de consultar la solución, prediga el efecto de cambiar el primer dato respecto del caso resuelto de esta familia. Justifique la tendencia y explique qué restricción podría convertirse en limitante. Datos: \(\Omega=2044.8\, \mathrm{s}^{-1}\); \(r_t=0.102\, \mathrm{m}\); \(C_x=148\, \mathrm{m}\,\mathrm{s}^{-1}\); \(T_e=283\, \mathrm{K}\); \(\gamma=1.4\); \(R_g=287\, \mathrm{J}\,\mathrm{kg}^{-1}\,\mathrm{K}^{-1}\).

\noindent\textbf{Datos}
\begin{itemize}
\item $\Omega = 2044.8\,\mathrm{s}^{-1}$ — velocidad angular
\item $r_t = 0.102\,\mathrm{m}$ — radio punta inductor
\item $C_x = 148\,\mathrm{m}\,\mathrm{s}^{-1}$ — velocidad axial
\item $T_e = 283\,\mathrm{K}$ — temperatura estática
\item $\gamma = 1.4$ — razón de calores
\item $R_g = 287\,\mathrm{J}\,\mathrm{kg}^{-1}\,\mathrm{K}^{-1}$ — constante gas
\end{itemize}

\noindent\textbf{Incógnita}
$\mathrm{Ma}_r$ (adimensional); también las magnitudes intermedias desarrolladas.

\noindent\textbf{Modelo físico}

\noindent\textit{Inductor y Mach de punta}

\[
\fitmath{\mathrm{Ma}_r=\frac{\sqrt{\Omega^{2} r_t^{2} + C_x^{2}}}{\sqrt{R_g T_e \gamma}}}
\]

\noindent\textbf{Hipótesis}
\begin{itemize}
\item Sin prerrotación y velocidad axial uniforme.
\end{itemize}

\noindent\textbf{Desarrollo paso a paso}

\begin{enumerate}

\item La punta tiene mayor velocidad periférica que la raíz.

\[\fitmath{\mathrm{Ut}=\Omega\cdot \mathrm{rt}}\]

\[\fitmath{\mathrm{Ut}=(2044.8)\cdot (0.102)\approx 208.5696\,\mathrm{m/s}}\]

\item Se forma el módulo relativo de entrada.

\[\fitmath{W=\sqrt{{\mathrm{cx}}^{2}+{\left(\Omega\cdot \mathrm{rt}\right)}^{2}}}\]

\[\fitmath{W=\sqrt{{(148)}^{2}+{\left((2044.8)\cdot (0.102)\right)}^{2}}\approx 255.7446\,\mathrm{m/s}}\]

\item El sonido corresponde al estado estático.

\[\fitmath{\mathrm{Mr}=\frac{\sqrt{{\mathrm{cx}}^{2}+{\left(\Omega\cdot \mathrm{rt}\right)}^{2}}}{\sqrt{\gamma\cdot R_g\cdot \mathrm{Te}}}}\]

\[\fitmath{\mathrm{Mr}=\frac{\sqrt{{(148)}^{2}+{\left((2044.8)\cdot (0.102)\right)}^{2}}}{\sqrt{(1.4)\cdot (287)\cdot (283)}}\approx 0.7584173\,}\]

\end{enumerate}


\noindent\textbf{Sustitución numérica}

\[
\fitmath{\frac{\sqrt{\left(\left(2044.8\,\mathrm{s}^{-1}\right)\right)^{2} \left(0.102\,\mathrm{m}\right)^{2} + \left(\left(148\,\mathrm{m}\,\mathrm{s}^{-1}\right)\right)^{2}}}{\sqrt{\left(287\,\mathrm{J}\,\mathrm{kg}^{-1}\,\mathrm{K}^{-1}\right) \left(283\,\mathrm{K}\right) \left(1.4\right)}}}
\]

\noindent\textbf{Resultado}

\[
\fitmath{\mathrm{Ma}_r=7.584172706\times10^{-1}}
\]

\noindent\textbf{Interpretación}
La temperatura estática se comparte entre referencias; la velocidad relativa aumenta por la rotación.

\noindent\textbf{Comprobación de los pasos}
\begin{itemize}
\item Paso 2: $208.57$ — reproducido
\item Paso 4: $255.745$ — reproducido
\item Paso 6: $0.758417$ — reproducido
\end{itemize}

\noindent\textbf{Comprobación de la respuesta} la aritmética del último paso reproduce la respuesta enunciada a la precisión con que está escrita. Es una reproducción aritmética; no valida el modelo físico ni sus hipótesis.

\noindent\textbf{Limitaciones}
\begin{itemize}
\item Modelo educativo; la aplicabilidad a una máquina real requiere revisión experta.
\end{itemize}

\noindent\textbf{Errores frecuentes}
\begin{itemize}
\item Usar radio exterior del rodete en el ojo del inductor.
\end{itemize}
% studium:end-record eyJpZCI6ICJQUkItMDExOCIsICJzZWN0aW9uIjogImMxNSJ9
\end{studiumworked}

\begin{studiumworked}{Problema resuelto C15-04-R}


\noindent\textbf{Enunciado}
\(C15-04-R \cdot\) Prerrotación y trabajo. Manteniendo la descarga real, determine cuánto trabajo aporta el rodete con prerrotación positiva. Datos: \(\Omega=2080\, \mathrm{s}^{-1}\); \(r_e=0.079\, \mathrm{m}\); \(r_s=0.204\, \mathrm{m}\); \(C_{\theta e}=47\, \mathrm{m}\,\mathrm{s}^{-1}\); \(C_{\theta s}=351\, \mathrm{m}\,\mathrm{s}^{-1}\).

\noindent\textbf{Datos}
\begin{itemize}
\item $\Omega = 2080\,\mathrm{s}^{-1}$ — velocidad angular
\item $r_e = 0.079\,\mathrm{m}$ — radio entrada
\item $r_s = 0.204\,\mathrm{m}$ — radio salida
\item $C_{\theta e} = 47\,\mathrm{m}\,\mathrm{s}^{-1}$ — remolino entrada
\item $C_{\theta s} = 351\,\mathrm{m}\,\mathrm{s}^{-1}$ — remolino real salida
\end{itemize}

\noindent\textbf{Incógnita}
$w_c$, en $\mathrm{J\,kg^{-1}}$; también las magnitudes intermedias desarrolladas.

\noindent\textbf{Modelo físico}

\noindent\textit{Prerrotación y trabajo}

\[
\fitmath{w_c=\Omega \left(- C_{\theta e} r_e + C_{\theta s} r_s\right)}
\]

\noindent\textbf{Hipótesis}
\begin{itemize}
\item Euler general con valores medios compatibles.
\end{itemize}

\noindent\textbf{Desarrollo paso a paso}

\begin{enumerate}

\item Se calcula la contribución de descarga.

\[\fitmath{\mathrm{exit}=\Omega\cdot \mathrm{rs}\cdot \mathrm{cs}}\]

\[\fitmath{\mathrm{exit}=(2080)\cdot (0.204)\cdot (351)\approx 148936.3\,\mathrm{J/kg}}\]

\item Se resta la contribución de entrada.

\[\fitmath{\mathrm{wc}=\Omega\cdot \left(\mathrm{rs}\cdot \mathrm{cs}-\mathrm{re}\cdot \mathrm{ce}\right)}\]

\[\fitmath{\mathrm{wc}=(2080)\cdot \left((0.204)\cdot (351)-(0.079)\cdot (47)\right)\approx 141213.3\,\mathrm{J/kg}}\]

\end{enumerate}


\noindent\textbf{Sustitución numérica}

\[
\fitmath{\left(2080\,\mathrm{s}^{-1}\right) \left(- \left(47\,\mathrm{m}\,\mathrm{s}^{-1}\right) \left(0.079\,\mathrm{m}\right) + \left(351\,\mathrm{m}\,\mathrm{s}^{-1}\right) \left(0.204\,\mathrm{m}\right)\right)}
\]

\noindent\textbf{Resultado}

\[
\fitmath{w_c=1.412132800\times10^{5}\,\mathrm{J}\,\mathrm{kg}^{-1}}
\]

\noindent\textbf{Interpretación}
La prerrotación positiva reduce trabajo para descarga fijada. Su efecto en incidencia debe evaluarse aparte.

\noindent\textbf{Comprobación de los pasos}
\begin{itemize}
\item Paso 2: $148936$ — reproducido
\item Paso 4: $141213$ — reproducido
\end{itemize}

\noindent\textbf{Comprobación de la respuesta} la aritmética del último paso reproduce la respuesta enunciada a la precisión con que está escrita. Es una reproducción aritmética; no valida el modelo físico ni sus hipótesis.

\noindent\textbf{Limitaciones}
\begin{itemize}
\item Modelo educativo; la aplicabilidad a una máquina real requiere revisión experta.
\end{itemize}

\noindent\textbf{Errores frecuentes}
\begin{itemize}
\item Suponer prerrotación nula aunque el enunciado la proporcione.
\end{itemize}
% studium:end-record eyJpZCI6ICJQUkItMDExOSIsICJzZWN0aW9uIjogImMxNSJ9
\end{studiumworked}

\begin{studiumworked}{Problema resuelto C15-04-P}


\noindent\textbf{Enunciado}
\(C15-04-P \cdot\) Prerrotación y trabajo. Ejercicio de transferencia. Manteniendo la descarga real, determine cuánto trabajo aporta el rodete con prerrotación positiva. Antes de consultar la solución, prediga el efecto de cambiar el primer dato respecto del caso resuelto de esta familia. Justifique la tendencia y explique qué restricción podría convertirse en limitante. Datos: \(\Omega=2215.2\, \mathrm{s}^{-1}\); \(r_e=0.079\, \mathrm{m}\); \(r_s=0.204\, \mathrm{m}\); \(C_{\theta e}=47\, \mathrm{m}\,\mathrm{s}^{-1}\); \(C_{\theta s}=351\, \mathrm{m}\,\mathrm{s}^{-1}\).

\noindent\textbf{Datos}
\begin{itemize}
\item $\Omega = 2215.2\,\mathrm{s}^{-1}$ — velocidad angular
\item $r_e = 0.079\,\mathrm{m}$ — radio entrada
\item $r_s = 0.204\,\mathrm{m}$ — radio salida
\item $C_{\theta e} = 47\,\mathrm{m}\,\mathrm{s}^{-1}$ — remolino entrada
\item $C_{\theta s} = 351\,\mathrm{m}\,\mathrm{s}^{-1}$ — remolino real salida
\end{itemize}

\noindent\textbf{Incógnita}
$w_c$, en $\mathrm{J\,kg^{-1}}$; también las magnitudes intermedias desarrolladas.

\noindent\textbf{Modelo físico}

\noindent\textit{Prerrotación y trabajo}

\[
\fitmath{w_c=\Omega \left(- C_{\theta e} r_e + C_{\theta s} r_s\right)}
\]

\noindent\textbf{Hipótesis}
\begin{itemize}
\item Euler general con valores medios compatibles.
\end{itemize}

\noindent\textbf{Desarrollo paso a paso}

\begin{enumerate}

\item Se calcula la contribución de descarga.

\[\fitmath{\mathrm{exit}=\Omega\cdot \mathrm{rs}\cdot \mathrm{cs}}\]

\[\fitmath{\mathrm{exit}=(2215.2)\cdot (0.204)\cdot (351)\approx 158617.2\,\mathrm{J/kg}}\]

\item Se resta la contribución de entrada.

\[\fitmath{\mathrm{wc}=\Omega\cdot \left(\mathrm{rs}\cdot \mathrm{cs}-\mathrm{re}\cdot \mathrm{ce}\right)}\]

\[\fitmath{\mathrm{wc}=(2215.2)\cdot \left((0.204)\cdot (351)-(0.079)\cdot (47)\right)\approx 150392.1\,\mathrm{J/kg}}\]

\end{enumerate}


\noindent\textbf{Sustitución numérica}

\[
\fitmath{\left(2215.2\,\mathrm{s}^{-1}\right) \left(- \left(47\,\mathrm{m}\,\mathrm{s}^{-1}\right) \left(0.079\,\mathrm{m}\right) + \left(351\,\mathrm{m}\,\mathrm{s}^{-1}\right) \left(0.204\,\mathrm{m}\right)\right)}
\]

\noindent\textbf{Resultado}

\[
\fitmath{w_c=1.503921432\times10^{5}\,\mathrm{J}\,\mathrm{kg}^{-1}}
\]

\noindent\textbf{Interpretación}
La prerrotación positiva reduce trabajo para descarga fijada. Su efecto en incidencia debe evaluarse aparte.

\noindent\textbf{Comprobación de los pasos}
\begin{itemize}
\item Paso 2: $158617$ — reproducido
\item Paso 4: $150392$ — reproducido
\end{itemize}

\noindent\textbf{Comprobación de la respuesta} la aritmética del último paso reproduce la respuesta enunciada a la precisión con que está escrita. Es una reproducción aritmética; no valida el modelo físico ni sus hipótesis.

\noindent\textbf{Limitaciones}
\begin{itemize}
\item Modelo educativo; la aplicabilidad a una máquina real requiere revisión experta.
\end{itemize}

\noindent\textbf{Errores frecuentes}
\begin{itemize}
\item Suponer prerrotación nula aunque el enunciado la proporcione.
\end{itemize}
% studium:end-record eyJpZCI6ICJQUkItMDEyMCIsICJzZWN0aW9uIjogImMxNSJ9
\end{studiumworked}

\chapter{Difusores radiales, recuperación y geometría de paso}
\begin{studiumworked}{Problema resuelto C16-01-R}


\noindent\textbf{Enunciado}
\(C16-01-R \cdot\) Difusor sin álabes ideal. En un difusor de ancho y densidad constantes, calcule las componentes de salida y el incremento ideal de presión. Datos: \(r_e=0.21\, \mathrm{m}\); \(r_s=0.32\, \mathrm{m}\); \(C_r=128\, \mathrm{m}\,\mathrm{s}^{-1}\); \(C_\theta=337\, \mathrm{m}\,\mathrm{s}^{-1}\); \(\rho=2.9\, \mathrm{kg}\,\mathrm{m}^{-3}\).

\noindent\textbf{Datos}
\begin{itemize}
\item $r_e = 0.21\,\mathrm{m}$ — radio entrada
\item $r_s = 0.32\,\mathrm{m}$ — radio salida
\item $C_r = 128\,\mathrm{m}\,\mathrm{s}^{-1}$ — componente radial entrada
\item $C_\theta = 337\,\mathrm{m}\,\mathrm{s}^{-1}$ — componente tangencial entrada
\item $\rho = 2.9\,\mathrm{kg}\,\mathrm{m}^{-3}$ — densidad uniforme
\end{itemize}

\noindent\textbf{Incógnita}
$\Delta p$, en $\mathrm{Pa}$; también las magnitudes intermedias desarrolladas.

\noindent\textbf{Modelo físico}

\noindent\textit{Difusor sin álabes ideal}

\[
\fitmath{\Delta p=\frac{\rho \left(C_r^{2} + C_\theta^{2}\right) \left(- \frac{r_e^{2}}{r_s^{2}} + 1\right)}{2}}
\]

\noindent\textbf{Hipótesis}
\begin{itemize}
\item Adiabático ideal, ancho y densidad constantes, sin par.
\end{itemize}

\noindent\textbf{Desarrollo paso a paso}

\begin{enumerate}

\item Componente radial de salida

\[\fitmath{\mathrm{crs}=\frac{\mathrm{cr}\cdot \mathrm{re}}{\mathrm{rs}}}\]

\[\fitmath{\mathrm{crs}=\frac{(128)\cdot (0.21)}{(0.32)}\approx 84\,\mathrm{m/s}}\]

\item Continuidad y ausencia de par hacen escalar ambas componentes con r inverso.

\[\fitmath{\mathrm{cts}=\frac{\mathrm{ct}\cdot \mathrm{re}}{\mathrm{rs}}}\]

\[\fitmath{\mathrm{cts}=\frac{(337)\cdot (0.21)}{(0.32)}\approx 221.1562\,\mathrm{m/s}}\]

\item Se evalúa el incremento de presión por deceleración ideal.

\[\fitmath{\mathrm{dp}=\frac{\rho\cdot \left({\mathrm{cr}}^{2}+{\mathrm{ct}}^{2}\right)\cdot \left(1-{\left(\frac{\mathrm{re}}{\mathrm{rs}}\right)}^{2}\right)}{2}}\]

\[\fitmath{\mathrm{dp}=\frac{(2.9)\cdot \left({(128)}^{2}+{(337)}^{2}\right)\cdot \left(1-{\left(\frac{(0.21)}{(0.32)}\right)}^{2}\right)}{2}\approx 107281\,\mathrm{Pa}}\]

\end{enumerate}


\noindent\textbf{Sustitución numérica}

\[
\fitmath{\frac{\left(2.9\,\mathrm{kg}\,\mathrm{m}^{-3}\right) \left(\left(\left(128\,\mathrm{m}\,\mathrm{s}^{-1}\right)\right)^{2} + \left(\left(337\,\mathrm{m}\,\mathrm{s}^{-1}\right)\right)^{2}\right) \left(- \frac{\left(0.21\,\mathrm{m}\right)^{2}}{\left(0.32\,\mathrm{m}\right)^{2}} + 1\right)}{2}}
\]

\noindent\textbf{Resultado}

\[
\fitmath{\Delta p=1.072810240\times10^{5}\,\mathrm{Pa}}
\]

\noindent\textbf{Componente radial de salida}
\[
\fitmath{C_{rs}=C_{re}r_e/r_s=84\,\mathrm{m/s}}
\]

\noindent\textbf{Comprobación de continuidad.} $r_sC_{rs}=r_eC_{re}$ con ancho y densidad constantes; junto con $r_sC_{\theta s}=r_eC_{\theta e}$ reproduce la recuperación ideal de presión.

\noindent\textbf{Interpretación}
La presión aumenta y el momento angular se conserva. Es una integración incompresible ideal, no una curva ensayada.

\noindent\textbf{Comprobación de los pasos}
\begin{itemize}
\item Paso 2: $221.156$ — reproducido
\item Paso 4: $107281$ — reproducido
\end{itemize}

\noindent\textbf{Comprobación de la respuesta} la aritmética del último paso reproduce la respuesta enunciada a la precisión con que está escrita. Es una reproducción aritmética; no valida el modelo físico ni sus hipótesis.

\noindent\textbf{Limitaciones}
\begin{itemize}
\item Modelo educativo; la aplicabilidad a una máquina real requiere revisión experta.
\end{itemize}

\noindent\textbf{Errores frecuentes}
\begin{itemize}
\item Conservar $C_\theta$ en lugar de $rC_\theta$.
\end{itemize}
% studium:end-record eyJpZCI6ICJQUkItMDEyMSIsICJzZWN0aW9uIjogImMxNiJ9
\end{studiumworked}

\begin{studiumworked}{Problema resuelto C16-01-P}


\noindent\textbf{Enunciado}
\(C16-01-P \cdot\) Difusor sin álabes ideal. Ejercicio de transferencia. En un difusor de ancho y densidad constantes, calcule las componentes de salida y el incremento ideal de presión. Antes de consultar la solución, prediga el efecto de cambiar el primer dato respecto del caso resuelto de esta familia. Justifique la tendencia y explique qué restricción podría convertirse en limitante. Datos: \(r_e=0.22365\, \mathrm{m}\); \(r_s=0.32\, \mathrm{m}\); \(C_r=128\, \mathrm{m}\,\mathrm{s}^{-1}\); \(C_\theta=337\, \mathrm{m}\,\mathrm{s}^{-1}\); \(\rho=2.9\, \mathrm{kg}\,\mathrm{m}^{-3}\).

\noindent\textbf{Datos}
\begin{itemize}
\item $r_e = 0.22365\,\mathrm{m}$ — radio entrada
\item $r_s = 0.32\,\mathrm{m}$ — radio salida
\item $C_r = 128\,\mathrm{m}\,\mathrm{s}^{-1}$ — componente radial entrada
\item $C_\theta = 337\,\mathrm{m}\,\mathrm{s}^{-1}$ — componente tangencial entrada
\item $\rho = 2.9\,\mathrm{kg}\,\mathrm{m}^{-3}$ — densidad uniforme
\end{itemize}

\noindent\textbf{Incógnita}
$\Delta p$, en $\mathrm{Pa}$; también las magnitudes intermedias desarrolladas.

\noindent\textbf{Modelo físico}

\noindent\textit{Difusor sin álabes ideal}

\[
\fitmath{\Delta p=\frac{\rho \left(C_r^{2} + C_\theta^{2}\right) \left(- \frac{r_e^{2}}{r_s^{2}} + 1\right)}{2}}
\]

\noindent\textbf{Hipótesis}
\begin{itemize}
\item Adiabático ideal, ancho y densidad constantes, sin par.
\end{itemize}

\noindent\textbf{Desarrollo paso a paso}

\begin{enumerate}

\item Componente radial de salida

\[\fitmath{\mathrm{crs}=\frac{\mathrm{cr}\cdot \mathrm{re}}{\mathrm{rs}}}\]

\[\fitmath{\mathrm{crs}=\frac{(128)\cdot (0.22365)}{(0.32)}\approx 89.46\,\mathrm{m/s}}\]

\item Continuidad y ausencia de par hacen escalar ambas componentes con r inverso.

\[\fitmath{\mathrm{cts}=\frac{\mathrm{ct}\cdot \mathrm{re}}{\mathrm{rs}}}\]

\[\fitmath{\mathrm{cts}=\frac{(337)\cdot (0.22365)}{(0.32)}\approx 235.5314\,\mathrm{m/s}}\]

\item Se evalúa el incremento de presión por deceleración ideal.

\[\fitmath{\mathrm{dp}=\frac{\rho\cdot \left({\mathrm{cr}}^{2}+{\mathrm{ct}}^{2}\right)\cdot \left(1-{\left(\frac{\mathrm{re}}{\mathrm{rs}}\right)}^{2}\right)}{2}}\]

\[\fitmath{\mathrm{dp}=\frac{(2.9)\cdot \left({(128)}^{2}+{(337)}^{2}\right)\cdot \left(1-{\left(\frac{(0.22365)}{(0.32)}\right)}^{2}\right)}{2}\approx 96388.55\,\mathrm{Pa}}\]

\end{enumerate}


\noindent\textbf{Sustitución numérica}

\[
\fitmath{\frac{\left(2.9\,\mathrm{kg}\,\mathrm{m}^{-3}\right) \left(\left(\left(128\,\mathrm{m}\,\mathrm{s}^{-1}\right)\right)^{2} + \left(\left(337\,\mathrm{m}\,\mathrm{s}^{-1}\right)\right)^{2}\right) \left(- \frac{\left(0.22365\,\mathrm{m}\right)^{2}}{\left(0.32\,\mathrm{m}\right)^{2}} + 1\right)}{2}}
\]

\noindent\textbf{Resultado}

\[
\fitmath{\Delta p=9.638855435\times10^{4}\,\mathrm{Pa}}
\]

\noindent\textbf{Componente radial de salida}
\[
\fitmath{C_{rs}=C_{re}r_e/r_s=89.46\,\mathrm{m/s}}
\]

\noindent\textbf{Comprobación de continuidad.} $r_sC_{rs}=r_eC_{re}$ con ancho y densidad constantes; junto con $r_sC_{\theta s}=r_eC_{\theta e}$ reproduce la recuperación ideal de presión.

\noindent\textbf{Interpretación}
La presión aumenta y el momento angular se conserva. Es una integración incompresible ideal, no una curva ensayada.

\noindent\textbf{Comprobación de los pasos}
\begin{itemize}
\item Paso 2: $235.531$ — reproducido
\item Paso 4: $96388.6$ — reproducido
\end{itemize}

\noindent\textbf{Comprobación de la respuesta} la aritmética del último paso reproduce la respuesta enunciada a la precisión con que está escrita. Es una reproducción aritmética; no valida el modelo físico ni sus hipótesis.

\noindent\textbf{Limitaciones}
\begin{itemize}
\item Modelo educativo; la aplicabilidad a una máquina real requiere revisión experta.
\end{itemize}

\noindent\textbf{Errores frecuentes}
\begin{itemize}
\item Conservar $C_\theta$ en lugar de $rC_\theta$.
\end{itemize}
% studium:end-record eyJpZCI6ICJQUkItMDEyMiIsICJzZWN0aW9uIjogImMxNiJ9
\end{studiumworked}

\begin{studiumworked}{Problema resuelto C16-02-R}


\noindent\textbf{Enunciado}
\(C16-02-R \cdot\) Recuperación con pérdida prescrita. Calcule el coeficiente de recuperación estática y el incremento de presión con una pérdida normalizada dada. Datos: \(C_e=364\, \mathrm{m}\,\mathrm{s}^{-1}\); \(C_s=216\, \mathrm{m}\,\mathrm{s}^{-1}\); \(\zeta_D=0.13\); \(\rho=3.1\, \mathrm{kg}\,\mathrm{m}^{-3}\).

\noindent\textbf{Datos}
\begin{itemize}
\item $C_e = 364\,\mathrm{m}\,\mathrm{s}^{-1}$ — velocidad entrada
\item $C_s = 216\,\mathrm{m}\,\mathrm{s}^{-1}$ — velocidad salida
\item $\zeta_D = 0.13$ — pérdida normalizada con dinámica de entrada
\item $\rho = 3.1\,\mathrm{kg}\,\mathrm{m}^{-3}$ — densidad
\end{itemize}

\noindent\textbf{Incógnita}
$\Delta p$, en $\mathrm{Pa}$; también las magnitudes intermedias desarrolladas.

\noindent\textbf{Modelo físico}

\noindent\textit{Recuperación con pérdida prescrita}

\[
\fitmath{\Delta p=\frac{\rho C_e^{2} \left(- \zeta_D + 1 - \frac{C_s^{2}}{C_e^{2}}\right)}{2}}
\]

\noindent\textbf{Hipótesis}
\begin{itemize}
\item Modelo incompresible, pérdida definida con velocidad entrada.
\end{itemize}

\noindent\textbf{Desarrollo paso a paso}

\begin{enumerate}

\item La recuperación resta energía residual y pérdida.

\[\fitmath{\mathrm{Cp}=1-{\left(\frac{\mathrm{vs}}{\mathrm{ve}}\right)}^{2}-\mathrm{loss}}\]

\[\fitmath{\mathrm{Cp}=1-{\left(\frac{(216)}{(364)}\right)}^{2}-(0.13)\approx 0.5178686\,}\]

\item Se dimensionaliza con la presión dinámica de entrada.

\[\fitmath{\mathrm{dp}=\frac{\rho\cdot {\mathrm{ve}}^{2}\cdot \left(1-{\left(\frac{\mathrm{vs}}{\mathrm{ve}}\right)}^{2}-\mathrm{loss}\right)}{2}}\]

\[\fitmath{\mathrm{dp}=\frac{(3.1)\cdot {(364)}^{2}\cdot \left(1-{\left(\frac{(216)}{(364)}\right)}^{2}-(0.13)\right)}{2}\approx 106354.1\,\mathrm{Pa}}\]

\end{enumerate}


\noindent\textbf{Sustitución numérica}

\[
\fitmath{\frac{\left(3.1\,\mathrm{kg}\,\mathrm{m}^{-3}\right) \left(\left(364\,\mathrm{m}\,\mathrm{s}^{-1}\right)\right)^{2} \left(- \left(0.13\right) + 1 - \frac{\left(\left(216\,\mathrm{m}\,\mathrm{s}^{-1}\right)\right)^{2}}{\left(\left(364\,\mathrm{m}\,\mathrm{s}^{-1}\right)\right)^{2}}\right)}{2}}
\]

\noindent\textbf{Resultado}

\[
\fitmath{\Delta p=1.063540560\times10^{5}\,\mathrm{Pa}}
\]

\noindent\textbf{Interpretación}
La pérdida total y la velocidad residual reducen recuperación por mecanismos distintos.

\noindent\textbf{Comprobación de los pasos}
\begin{itemize}
\item Paso 2: $0.517869$ — reproducido
\item Paso 4: $106354$ — reproducido
\end{itemize}

\noindent\textbf{Comprobación de la respuesta} la aritmética del último paso reproduce la respuesta enunciada a la precisión con que está escrita. Es una reproducción aritmética; no valida el modelo físico ni sus hipótesis.

\noindent\textbf{Limitaciones}
\begin{itemize}
\item Modelo educativo; la aplicabilidad a una máquina real requiere revisión experta.
\end{itemize}

\noindent\textbf{Errores frecuentes}
\begin{itemize}
\item Restar una pérdida normalizada con otro denominador.
\end{itemize}
% studium:end-record eyJpZCI6ICJQUkItMDEyMyIsICJzZWN0aW9uIjogImMxNiJ9
\end{studiumworked}

\begin{studiumworked}{Problema resuelto C16-02-P}


\noindent\textbf{Enunciado}
\(C16-02-P \cdot\) Recuperación con pérdida prescrita. Ejercicio de transferencia. Calcule el coeficiente de recuperación estática y el incremento de presión con una pérdida normalizada dada. Antes de consultar la solución, prediga el efecto de cambiar el primer dato respecto del caso resuelto de esta familia. Justifique la tendencia y explique qué restricción podría convertirse en limitante. Datos: \(C_e=387.66\, \mathrm{m}\,\mathrm{s}^{-1}\); \(C_s=216\, \mathrm{m}\,\mathrm{s}^{-1}\); \(\zeta_D=0.13\); \(\rho=3.1\, \mathrm{kg}\,\mathrm{m}^{-3}\).

\noindent\textbf{Datos}
\begin{itemize}
\item $C_e = 387.66\,\mathrm{m}\,\mathrm{s}^{-1}$ — velocidad entrada
\item $C_s = 216\,\mathrm{m}\,\mathrm{s}^{-1}$ — velocidad salida
\item $\zeta_D = 0.13$ — pérdida normalizada con dinámica de entrada
\item $\rho = 3.1\,\mathrm{kg}\,\mathrm{m}^{-3}$ — densidad
\end{itemize}

\noindent\textbf{Incógnita}
$\Delta p$, en $\mathrm{Pa}$; también las magnitudes intermedias desarrolladas.

\noindent\textbf{Modelo físico}

\noindent\textit{Recuperación con pérdida prescrita}

\[
\fitmath{\Delta p=\frac{\rho C_e^{2} \left(- \zeta_D + 1 - \frac{C_s^{2}}{C_e^{2}}\right)}{2}}
\]

\noindent\textbf{Hipótesis}
\begin{itemize}
\item Modelo incompresible, pérdida definida con velocidad entrada.
\end{itemize}

\noindent\textbf{Desarrollo paso a paso}

\begin{enumerate}

\item La recuperación resta energía residual y pérdida.

\[\fitmath{\mathrm{Cp}=1-{\left(\frac{\mathrm{vs}}{\mathrm{ve}}\right)}^{2}-\mathrm{loss}}\]

\[\fitmath{\mathrm{Cp}=1-{\left(\frac{(216)}{(387.66)}\right)}^{2}-(0.13)\approx 0.5595401\,}\]

\item Se dimensionaliza con la presión dinámica de entrada.

\[\fitmath{\mathrm{dp}=\frac{\rho\cdot {\mathrm{ve}}^{2}\cdot \left(1-{\left(\frac{\mathrm{vs}}{\mathrm{ve}}\right)}^{2}-\mathrm{loss}\right)}{2}}\]

\[\fitmath{\mathrm{dp}=\frac{(3.1)\cdot {(387.66)}^{2}\cdot \left(1-{\left(\frac{(216)}{(387.66)}\right)}^{2}-(0.13)\right)}{2}\approx 130336.2\,\mathrm{Pa}}\]

\end{enumerate}


\noindent\textbf{Sustitución numérica}

\[
\fitmath{\frac{\left(3.1\,\mathrm{kg}\,\mathrm{m}^{-3}\right) \left(\left(387.66\,\mathrm{m}\,\mathrm{s}^{-1}\right)\right)^{2} \left(- \left(0.13\right) + 1 - \frac{\left(\left(216\,\mathrm{m}\,\mathrm{s}^{-1}\right)\right)^{2}}{\left(\left(387.66\,\mathrm{m}\,\mathrm{s}^{-1}\right)\right)^{2}}\right)}{2}}
\]

\noindent\textbf{Resultado}

\[
\fitmath{\Delta p=1.303361516\times10^{5}\,\mathrm{Pa}}
\]

\noindent\textbf{Interpretación}
La pérdida total y la velocidad residual reducen recuperación por mecanismos distintos.

\noindent\textbf{Comprobación de los pasos}
\begin{itemize}
\item Paso 2: $0.55954$ — reproducido
\item Paso 4: $130336$ — reproducido
\end{itemize}

\noindent\textbf{Comprobación de la respuesta} la aritmética del último paso reproduce la respuesta enunciada a la precisión con que está escrita. Es una reproducción aritmética; no valida el modelo físico ni sus hipótesis.

\noindent\textbf{Limitaciones}
\begin{itemize}
\item Modelo educativo; la aplicabilidad a una máquina real requiere revisión experta.
\end{itemize}

\noindent\textbf{Errores frecuentes}
\begin{itemize}
\item Restar una pérdida normalizada con otro denominador.
\end{itemize}
% studium:end-record eyJpZCI6ICJQUkItMDEyNCIsICJzZWN0aW9uIjogImMxNiJ9
\end{studiumworked}

\begin{studiumworked}{Problema resuelto C16-03-R}


\noindent\textbf{Enunciado}
\(C16-03-R \cdot\) Ancho del paso cilíndrico. Dimensione el ancho geométrico de salida radial con el caudal y área disponible dados. Datos: \(\dot m=8.7\, \mathrm{kg}\,\mathrm{s}^{-1}\); \(r=0.198\, \mathrm{m}\); \(\rho=3.4\, \mathrm{kg}\,\mathrm{m}^{-3}\); \(C_r=126\, \mathrm{m}\,\mathrm{s}^{-1}\); \(B=0.94\).

\noindent\textbf{Datos}
\begin{itemize}
\item $\dot m = 8.7\,\mathrm{kg}\,\mathrm{s}^{-1}$ — caudal
\item $r = 0.198\,\mathrm{m}$ — radio salida
\item $\rho = 3.4\,\mathrm{kg}\,\mathrm{m}^{-3}$ — densidad estática
\item $C_r = 126\,\mathrm{m}\,\mathrm{s}^{-1}$ — velocidad radial
\item $B = 0.94$ — fracción área disponible
\end{itemize}

\noindent\textbf{Incógnita}
$b$, en $\mathrm{m}$; también las magnitudes intermedias desarrolladas.

\noindent\textbf{Modelo físico}

\noindent\textit{Ancho del paso cilíndrico}

\[
\fitmath{b=\frac{\dot m}{2 \pi B C_r r \rho}}
\]

\noindent\textbf{Hipótesis}
\begin{itemize}
\item Paso cilíndrico uniforme y bloqueo explícito.
\end{itemize}

\noindent\textbf{Desarrollo paso a paso}

\begin{enumerate}

\item La continuidad fija el área cilíndrica efectiva.

\[\fitmath{\mathrm{ae}=\frac{\mathrm{md}}{\rho\cdot \mathrm{cr}}}\]

\[\fitmath{\mathrm{ae}=\frac{(8.7)}{(3.4)\cdot (126)}\approx 0.02030812\,\mathrm{m^{2}}}\]

\item La circunferencia y el bloqueo fijan el ancho.

\[\fitmath{\mathrm{width}=\frac{\mathrm{md}}{2\cdot \pi\cdot \mathrm{rad}\cdot \rho\cdot \mathrm{cr}\cdot \mathrm{bb}}}\]

\[\fitmath{\mathrm{width}=\frac{(8.7)}{2\cdot \pi\cdot (0.198)\cdot (3.4)\cdot (126)\cdot (0.94)}\approx 0.01736588\,\mathrm{m}}\]

\end{enumerate}


\noindent\textbf{Sustitución numérica}

\[
\fitmath{\frac{\left(8.7\,\mathrm{kg}\,\mathrm{s}^{-1}\right)}{2 \pi \left(0.94\right) \left(126\,\mathrm{m}\,\mathrm{s}^{-1}\right) \left(0.198\,\mathrm{m}\right) \left(3.4\,\mathrm{kg}\,\mathrm{m}^{-3}\right)}}
\]

\noindent\textbf{Resultado}

\[
\fitmath{b=1.736588330\times10^{-2}\,\mathrm{m}}
\]

\noindent\textbf{Interpretación}
Se usa componente radial, aunque el módulo absoluto sea mucho mayor por remolino.

\noindent\textbf{Comprobación de los pasos}
\begin{itemize}
\item Paso 2: $0.0203081$ — reproducido
\item Paso 4: $0.0173659$ — reproducido
\end{itemize}

\noindent\textbf{Comprobación de la respuesta} la aritmética del último paso reproduce la respuesta enunciada a la precisión con que está escrita. Es una reproducción aritmética; no valida el modelo físico ni sus hipótesis.

\noindent\textbf{Limitaciones}
\begin{itemize}
\item Modelo educativo; la aplicabilidad a una máquina real requiere revisión experta.
\end{itemize}

\noindent\textbf{Errores frecuentes}
\begin{itemize}
\item Usar el módulo absoluto en la continuidad radial.
\end{itemize}
% studium:end-record eyJpZCI6ICJQUkItMDEyNSIsICJzZWN0aW9uIjogImMxNiJ9
\end{studiumworked}

\begin{studiumworked}{Problema resuelto C16-03-P}


\noindent\textbf{Enunciado}
\(C16-03-P \cdot\) Ancho del paso cilíndrico. Ejercicio de transferencia. Dimensione el ancho geométrico de salida radial con el caudal y área disponible dados. Antes de consultar la solución, prediga el efecto de cambiar el primer dato respecto del caso resuelto de esta familia. Justifique la tendencia y explique qué restricción podría convertirse en limitante. Datos: \(\dot m=9.2655\, \mathrm{kg}\,\mathrm{s}^{-1}\); \(r=0.198\, \mathrm{m}\); \(\rho=3.4\, \mathrm{kg}\,\mathrm{m}^{-3}\); \(C_r=126\, \mathrm{m}\,\mathrm{s}^{-1}\); \(B=0.94\).

\noindent\textbf{Datos}
\begin{itemize}
\item $\dot m = 9.2655\,\mathrm{kg}\,\mathrm{s}^{-1}$ — caudal
\item $r = 0.198\,\mathrm{m}$ — radio salida
\item $\rho = 3.4\,\mathrm{kg}\,\mathrm{m}^{-3}$ — densidad estática
\item $C_r = 126\,\mathrm{m}\,\mathrm{s}^{-1}$ — velocidad radial
\item $B = 0.94$ — fracción área disponible
\end{itemize}

\noindent\textbf{Incógnita}
$b$, en $\mathrm{m}$; también las magnitudes intermedias desarrolladas.

\noindent\textbf{Modelo físico}

\noindent\textit{Ancho del paso cilíndrico}

\[
\fitmath{b=\frac{\dot m}{2 \pi B C_r r \rho}}
\]

\noindent\textbf{Hipótesis}
\begin{itemize}
\item Paso cilíndrico uniforme y bloqueo explícito.
\end{itemize}

\noindent\textbf{Desarrollo paso a paso}

\begin{enumerate}

\item La continuidad fija el área cilíndrica efectiva.

\[\fitmath{\mathrm{ae}=\frac{\mathrm{md}}{\rho\cdot \mathrm{cr}}}\]

\[\fitmath{\mathrm{ae}=\frac{(9.2655)}{(3.4)\cdot (126)}\approx 0.02162815\,\mathrm{m^{2}}}\]

\item La circunferencia y el bloqueo fijan el ancho.

\[\fitmath{\mathrm{width}=\frac{\mathrm{md}}{2\cdot \pi\cdot \mathrm{rad}\cdot \rho\cdot \mathrm{cr}\cdot \mathrm{bb}}}\]

\[\fitmath{\mathrm{width}=\frac{(9.2655)}{2\cdot \pi\cdot (0.198)\cdot (3.4)\cdot (126)\cdot (0.94)}\approx 0.01849467\,\mathrm{m}}\]

\end{enumerate}


\noindent\textbf{Sustitución numérica}

\[
\fitmath{\frac{\left(9.2655\,\mathrm{kg}\,\mathrm{s}^{-1}\right)}{2 \pi \left(0.94\right) \left(126\,\mathrm{m}\,\mathrm{s}^{-1}\right) \left(0.198\,\mathrm{m}\right) \left(3.4\,\mathrm{kg}\,\mathrm{m}^{-3}\right)}}
\]

\noindent\textbf{Resultado}

\[
\fitmath{b=1.849466572\times10^{-2}\,\mathrm{m}}
\]

\noindent\textbf{Interpretación}
Se usa componente radial, aunque el módulo absoluto sea mucho mayor por remolino.

\noindent\textbf{Comprobación de los pasos}
\begin{itemize}
\item Paso 2: $0.0216282$ — reproducido
\item Paso 4: $0.0184947$ — reproducido
\end{itemize}

\noindent\textbf{Comprobación de la respuesta} la aritmética del último paso reproduce la respuesta enunciada a la precisión con que está escrita. Es una reproducción aritmética; no valida el modelo físico ni sus hipótesis.

\noindent\textbf{Limitaciones}
\begin{itemize}
\item Modelo educativo; la aplicabilidad a una máquina real requiere revisión experta.
\end{itemize}

\noindent\textbf{Errores frecuentes}
\begin{itemize}
\item Usar el módulo absoluto en la continuidad radial.
\end{itemize}
% studium:end-record eyJpZCI6ICJQUkItMDEyNiIsICJzZWN0aW9uIjogImMxNiJ9
\end{studiumworked}

\begin{studiumworked}{Problema resuelto C16-04-R}


\noindent\textbf{Enunciado}
\(C16-04-R \cdot\) Estado estático de descarga. La descarga del rodete tiene estado total y componentes radial y tangencial dadas. Halle temperatura estática y densidad para dimensionar el difusor. Datos: \(T_0=472\, \mathrm{K}\); \(p_0=342000\, \mathrm{Pa}\); \(C_r=141\, \mathrm{m}\,\mathrm{s}^{-1}\); \(C_\theta=328\, \mathrm{m}\,\mathrm{s}^{-1}\); \(c_p=1004.5\, \mathrm{J}\,\mathrm{kg}^{-1}\,\mathrm{K}^{-1}\); \(\gamma=1.4\); \(R_g=287\, \mathrm{J}\,\mathrm{kg}^{-1}\,\mathrm{K}^{-1}\).

\noindent\textbf{Datos}
\begin{itemize}
\item $T_0 = 472\,\mathrm{K}$ — temperatura total
\item $p_0 = 342000\,\mathrm{Pa}$ — presión total
\item $C_r = 141\,\mathrm{m}\,\mathrm{s}^{-1}$ — componente radial
\item $C_\theta = 328\,\mathrm{m}\,\mathrm{s}^{-1}$ — componente tangencial
\item $c_p = 1004.5\,\mathrm{J}\,\mathrm{kg}^{-1}\,\mathrm{K}^{-1}$ — calor específico
\item $\gamma = 1.4$ — razón de calores
\item $R_g = 287\,\mathrm{J}\,\mathrm{kg}^{-1}\,\mathrm{K}^{-1}$ — constante gas
\end{itemize}

\noindent\textbf{Incógnita}
$\rho$, en $\mathrm{kg\,m^{-3}}$; también las magnitudes intermedias desarrolladas.

\noindent\textbf{Modelo físico}

\noindent\textit{Estado estático de descarga}

\[
\fitmath{\rho=\frac{p_0 \left(\frac{T_0 - \frac{C_r^{2} + C_\theta^{2}}{2 c_p}}{T_0}\right)^{\frac{\gamma}{\gamma - 1}}}{R_g \left(T_0 - \frac{C_r^{2} + C_\theta^{2}}{2 c_p}\right)}}
\]

\noindent\textbf{Hipótesis}
\begin{itemize}
\item Gas perfecto constante; componentes ortogonales.
\end{itemize}

\noindent\textbf{Desarrollo paso a paso}

\begin{enumerate}

\item Ambas componentes contribuyen a energía cinética.

\[\fitmath{T=T_0-\frac{{\mathrm{cr}}^{2}+{\mathrm{ct}}^{2}}{2\cdot \mathrm{cp}}}\]

\[\fitmath{T=(472)-\frac{{(141)}^{2}+{(328)}^{2}}{2\cdot (1004.5)}\approx 408.553\,\mathrm{K}}\]

\item La presión estática deriva del remanso local.

\[\fitmath{p=p_0\cdot {\left(\frac{T_0-\frac{{\mathrm{cr}}^{2}+{\mathrm{ct}}^{2}}{2\cdot \mathrm{cp}}}{T_0}\right)}^{\frac{\gamma}{\gamma-1}}}\]

\[\fitmath{p=(342000)\cdot {\left(\frac{(472)-\frac{{(141)}^{2}+{(328)}^{2}}{2\cdot (1004.5)}}{(472)}\right)}^{\frac{(1.4)}{(1.4)-1}}\approx 206347.2\,\mathrm{Pa}}\]

\item Se calcula densidad a partir del mismo estado estático.

\[\fitmath{\rho=\frac{p_0\cdot {\left(\frac{T_0-\frac{{\mathrm{cr}}^{2}+{\mathrm{ct}}^{2}}{2\cdot \mathrm{cp}}}{T_0}\right)}^{\frac{\gamma}{\gamma-1}}}{R_g\cdot \left(T_0-\frac{{\mathrm{cr}}^{2}+{\mathrm{ct}}^{2}}{2\cdot \mathrm{cp}}\right)}}\]

\[\fitmath{\rho=\frac{(342000)\cdot {\left(\frac{(472)-\frac{{(141)}^{2}+{(328)}^{2}}{2\cdot (1004.5)}}{(472)}\right)}^{\frac{(1.4)}{(1.4)-1}}}{(287)\cdot \left((472)-\frac{{(141)}^{2}+{(328)}^{2}}{2\cdot (1004.5)}\right)}\approx 1.75982\,\mathrm{kg/m^{3}}}\]

\end{enumerate}


\noindent\textbf{Sustitución numérica}

\[
\fitmath{\frac{\left(342000\,\mathrm{Pa}\right) \left(\frac{\left(472\,\mathrm{K}\right) - \frac{\left(\left(141\,\mathrm{m}\,\mathrm{s}^{-1}\right)\right)^{2} + \left(\left(328\,\mathrm{m}\,\mathrm{s}^{-1}\right)\right)^{2}}{2 \left(1004.5\,\mathrm{J}\,\mathrm{kg}^{-1}\,\mathrm{K}^{-1}\right)}}{\left(472\,\mathrm{K}\right)}\right)^{\frac{\left(1.4\right)}{\left(1.4\right) - 1}}}{\left(287\,\mathrm{J}\,\mathrm{kg}^{-1}\,\mathrm{K}^{-1}\right) \left(\left(472\,\mathrm{K}\right) - \frac{\left(\left(141\,\mathrm{m}\,\mathrm{s}^{-1}\right)\right)^{2} + \left(\left(328\,\mathrm{m}\,\mathrm{s}^{-1}\right)\right)^{2}}{2 \left(1004.5\,\mathrm{J}\,\mathrm{kg}^{-1}\,\mathrm{K}^{-1}\right)}\right)}}
\]

\noindent\textbf{Resultado}

\[
\fitmath{\rho=1.759820021\times10^{0}\,\mathrm{kg}\,\mathrm{m}^{-3}}
\]

\noindent\textbf{Interpretación}
Usar solo Cr para temperatura sobreestima la densidad disponible. La continuidad sí utiliza Cr.

\noindent\textbf{Comprobación de los pasos}
\begin{itemize}
\item Paso 2: $408.553$ — reproducido
\item Paso 4: $206347$ — reproducido
\item Paso 6: $1.75982$ — reproducido
\end{itemize}

\noindent\textbf{Comprobación de la respuesta} la aritmética del último paso reproduce la respuesta enunciada a la precisión con que está escrita. Es una reproducción aritmética; no valida el modelo físico ni sus hipótesis.

\noindent\textbf{Limitaciones}
\begin{itemize}
\item Modelo educativo; la aplicabilidad a una máquina real requiere revisión experta.
\end{itemize}

\noindent\textbf{Errores frecuentes}
\begin{itemize}
\item Usar distintas definiciones de velocidad en energía sin justificarlo.
\end{itemize}
% studium:end-record eyJpZCI6ICJQUkItMDEyNyIsICJzZWN0aW9uIjogImMxNiJ9
\end{studiumworked}

\begin{studiumworked}{Problema resuelto C16-04-P}


\noindent\textbf{Enunciado}
\(C16-04-P \cdot\) Estado estático de descarga. Ejercicio de transferencia. La descarga del rodete tiene estado total y componentes radial y tangencial dadas. Halle temperatura estática y densidad para dimensionar el difusor. Antes de consultar la solución, prediga el efecto de cambiar el primer dato respecto del caso resuelto de esta familia. Justifique la tendencia y explique qué restricción podría convertirse en limitante. Datos: \(T_0=502.68\, \mathrm{K}\); \(p_0=342000\, \mathrm{Pa}\); \(C_r=141\, \mathrm{m}\,\mathrm{s}^{-1}\); \(C_\theta=328\, \mathrm{m}\,\mathrm{s}^{-1}\); \(c_p=1004.5\, \mathrm{J}\,\mathrm{kg}^{-1}\,\mathrm{K}^{-1}\); \(\gamma=1.4\); \(R_g=287\, \mathrm{J}\,\mathrm{kg}^{-1}\,\mathrm{K}^{-1}\).

\noindent\textbf{Datos}
\begin{itemize}
\item $T_0 = 502.68\,\mathrm{K}$ — temperatura total
\item $p_0 = 342000\,\mathrm{Pa}$ — presión total
\item $C_r = 141\,\mathrm{m}\,\mathrm{s}^{-1}$ — componente radial
\item $C_\theta = 328\,\mathrm{m}\,\mathrm{s}^{-1}$ — componente tangencial
\item $c_p = 1004.5\,\mathrm{J}\,\mathrm{kg}^{-1}\,\mathrm{K}^{-1}$ — calor específico
\item $\gamma = 1.4$ — razón de calores
\item $R_g = 287\,\mathrm{J}\,\mathrm{kg}^{-1}\,\mathrm{K}^{-1}$ — constante gas
\end{itemize}

\noindent\textbf{Incógnita}
$\rho$, en $\mathrm{kg\,m^{-3}}$; también las magnitudes intermedias desarrolladas.

\noindent\textbf{Modelo físico}

\noindent\textit{Estado estático de descarga}

\[
\fitmath{\rho=\frac{p_0 \left(\frac{T_0 - \frac{C_r^{2} + C_\theta^{2}}{2 c_p}}{T_0}\right)^{\frac{\gamma}{\gamma - 1}}}{R_g \left(T_0 - \frac{C_r^{2} + C_\theta^{2}}{2 c_p}\right)}}
\]

\noindent\textbf{Hipótesis}
\begin{itemize}
\item Gas perfecto constante; componentes ortogonales.
\end{itemize}

\noindent\textbf{Desarrollo paso a paso}

\begin{enumerate}

\item Ambas componentes contribuyen a energía cinética.

\[\fitmath{T=T_0-\frac{{\mathrm{cr}}^{2}+{\mathrm{ct}}^{2}}{2\cdot \mathrm{cp}}}\]

\[\fitmath{T=(502.68)-\frac{{(141)}^{2}+{(328)}^{2}}{2\cdot (1004.5)}\approx 439.233\,\mathrm{K}}\]

\item La presión estática deriva del remanso local.

\[\fitmath{p=p_0\cdot {\left(\frac{T_0-\frac{{\mathrm{cr}}^{2}+{\mathrm{ct}}^{2}}{2\cdot \mathrm{cp}}}{T_0}\right)}^{\frac{\gamma}{\gamma-1}}}\]

\[\fitmath{p=(342000)\cdot {\left(\frac{(502.68)-\frac{{(141)}^{2}+{(328)}^{2}}{2\cdot (1004.5)}}{(502.68)}\right)}^{\frac{(1.4)}{(1.4)-1}}\approx 213274\,\mathrm{Pa}}\]

\item Se calcula densidad a partir del mismo estado estático.

\[\fitmath{\rho=\frac{p_0\cdot {\left(\frac{T_0-\frac{{\mathrm{cr}}^{2}+{\mathrm{ct}}^{2}}{2\cdot \mathrm{cp}}}{T_0}\right)}^{\frac{\gamma}{\gamma-1}}}{R_g\cdot \left(T_0-\frac{{\mathrm{cr}}^{2}+{\mathrm{ct}}^{2}}{2\cdot \mathrm{cp}}\right)}}\]

\[\fitmath{\rho=\frac{(342000)\cdot {\left(\frac{(502.68)-\frac{{(141)}^{2}+{(328)}^{2}}{2\cdot (1004.5)}}{(502.68)}\right)}^{\frac{(1.4)}{(1.4)-1}}}{(287)\cdot \left((502.68)-\frac{{(141)}^{2}+{(328)}^{2}}{2\cdot (1004.5)}\right)}\approx 1.691847\,\mathrm{kg/m^{3}}}\]

\end{enumerate}


\noindent\textbf{Sustitución numérica}

\[
\fitmath{\frac{\left(342000\,\mathrm{Pa}\right) \left(\frac{\left(502.68\,\mathrm{K}\right) - \frac{\left(\left(141\,\mathrm{m}\,\mathrm{s}^{-1}\right)\right)^{2} + \left(\left(328\,\mathrm{m}\,\mathrm{s}^{-1}\right)\right)^{2}}{2 \left(1004.5\,\mathrm{J}\,\mathrm{kg}^{-1}\,\mathrm{K}^{-1}\right)}}{\left(502.68\,\mathrm{K}\right)}\right)^{\frac{\left(1.4\right)}{\left(1.4\right) - 1}}}{\left(287\,\mathrm{J}\,\mathrm{kg}^{-1}\,\mathrm{K}^{-1}\right) \left(\left(502.68\,\mathrm{K}\right) - \frac{\left(\left(141\,\mathrm{m}\,\mathrm{s}^{-1}\right)\right)^{2} + \left(\left(328\,\mathrm{m}\,\mathrm{s}^{-1}\right)\right)^{2}}{2 \left(1004.5\,\mathrm{J}\,\mathrm{kg}^{-1}\,\mathrm{K}^{-1}\right)}\right)}}
\]

\noindent\textbf{Resultado}

\[
\fitmath{\rho=1.691846742\times10^{0}\,\mathrm{kg}\,\mathrm{m}^{-3}}
\]

\noindent\textbf{Interpretación}
Usar solo Cr para temperatura sobreestima la densidad disponible. La continuidad sí utiliza Cr.

\noindent\textbf{Comprobación de los pasos}
\begin{itemize}
\item Paso 2: $439.233$ — reproducido
\item Paso 4: $213274$ — reproducido
\item Paso 6: $1.69185$ — reproducido
\end{itemize}

\noindent\textbf{Comprobación de la respuesta} la aritmética del último paso reproduce la respuesta enunciada a la precisión con que está escrita. Es una reproducción aritmética; no valida el modelo físico ni sus hipótesis.

\noindent\textbf{Limitaciones}
\begin{itemize}
\item Modelo educativo; la aplicabilidad a una máquina real requiere revisión experta.
\end{itemize}

\noindent\textbf{Errores frecuentes}
\begin{itemize}
\item Usar distintas definiciones de velocidad en energía sin justificarlo.
\end{itemize}
% studium:end-record eyJpZCI6ICJQUkItMDEyOCIsICJzZWN0aW9uIjogImMxNiJ9
\end{studiumworked}

\chapter{Prestaciones y límites de un compresor radial}
\begin{studiumworked}{Problema resuelto C17-01-R}


\noindent\textbf{Enunciado}
\(C17-01-R \cdot\) Velocidad exigida por trabajo. Un rodete de salida radial debe suministrar trabajo dado con factor de deslizamiento prescrito. Calcule la velocidad periférica y el radio para el giro indicado. Datos: \(w_c=169000\, \mathrm{J}\,\mathrm{kg}^{-1}\); \(\sigma_s=0.86\); \(\Omega=2140\, \mathrm{s}^{-1}\).

\noindent\textbf{Datos}
\begin{itemize}
\item $w_c = 169000\,\mathrm{J}\,\mathrm{kg}^{-1}$ — trabajo requerido
\item $\sigma_s = 0.86$ — factor de deslizamiento
\item $\Omega = 2140\,\mathrm{s}^{-1}$ — velocidad angular
\end{itemize}

\noindent\textbf{Incógnita}
$r$, en $\mathrm{m}$; también las magnitudes intermedias desarrolladas.

\noindent\textbf{Modelo físico}

\noindent\textit{Velocidad exigida por trabajo}

\[
\fitmath{r=\frac{\sqrt{\frac{w_c}{\sigma_s}}}{\Omega}}
\]

\noindent\textbf{Hipótesis}
\begin{itemize}
\item Salida radial, sin prerrotación, slip constante.
\end{itemize}

\noindent\textbf{Desarrollo paso a paso}

\begin{enumerate}

\item La salida radial sin prerrotación da wc=slip U2 al cuadrado.

\[\fitmath{\mathrm{uu}=\sqrt{\frac{\mathrm{wc}}{\mathrm{slip}}}}\]

\[\fitmath{\mathrm{uu}=\sqrt{\frac{(169000)}{(0.86)}}\approx 443.2963\,\mathrm{m/s}}\]

\item La geometría requerida sigue de U=Omega r.

\[\fitmath{\mathrm{rad}=\frac{\sqrt{\frac{\mathrm{wc}}{\mathrm{slip}}}}{\Omega}}\]

\[\fitmath{\mathrm{rad}=\frac{\sqrt{\frac{(169000)}{(0.86)}}}{(2140)}\approx 0.2071478\,\mathrm{m}}\]

\end{enumerate}


\noindent\textbf{Sustitución numérica}

\[
\fitmath{\frac{\sqrt{\frac{\left(169000\,\mathrm{J}\,\mathrm{kg}^{-1}\right)}{\left(0.86\right)}}}{\left(2140\,\mathrm{s}^{-1}\right)}}
\]

\noindent\textbf{Resultado}

\[
\fitmath{r=2.071478134\times10^{-1}\,\mathrm{m}}
\]

\noindent\textbf{Interpretación}
El tamaño y el giro están ligados al trabajo. Las tensiones y el Mach deben comprobarse después.

\noindent\textbf{Comprobación de los pasos}
\begin{itemize}
\item Paso 2: $443.296$ — reproducido
\item Paso 4: $0.207148$ — reproducido
\end{itemize}

\noindent\textbf{Comprobación de la respuesta} la aritmética del último paso reproduce la respuesta enunciada a la precisión con que está escrita. Es una reproducción aritmética; no valida el modelo físico ni sus hipótesis.

\noindent\textbf{Limitaciones}
\begin{itemize}
\item Modelo educativo; la aplicabilidad a una máquina real requiere revisión experta.
\end{itemize}

\noindent\textbf{Errores frecuentes}
\begin{itemize}
\item Sustituir eficiencia isentrópica por factor de deslizamiento.
\end{itemize}
% studium:end-record eyJpZCI6ICJQUkItMDEyOSIsICJzZWN0aW9uIjogImMxNyJ9
\end{studiumworked}

\begin{studiumworked}{Problema resuelto C17-01-P}


\noindent\textbf{Enunciado}
\(C17-01-P \cdot\) Velocidad exigida por trabajo. Ejercicio de transferencia. Un rodete de salida radial debe suministrar trabajo dado con factor de deslizamiento prescrito. Calcule la velocidad periférica y el radio para el giro indicado. Antes de consultar la solución, prediga el efecto de cambiar el primer dato respecto del caso resuelto de esta familia. Justifique la tendencia y explique qué restricción podría convertirse en limitante. Datos: \(w_c=179985\, \mathrm{J}\,\mathrm{kg}^{-1}\); \(\sigma_s=0.86\); \(\Omega=2140\, \mathrm{s}^{-1}\).

\noindent\textbf{Datos}
\begin{itemize}
\item $w_c = 179985\,\mathrm{J}\,\mathrm{kg}^{-1}$ — trabajo requerido
\item $\sigma_s = 0.86$ — factor de deslizamiento
\item $\Omega = 2140\,\mathrm{s}^{-1}$ — velocidad angular
\end{itemize}

\noindent\textbf{Incógnita}
$r$, en $\mathrm{m}$; también las magnitudes intermedias desarrolladas.

\noindent\textbf{Modelo físico}

\noindent\textit{Velocidad exigida por trabajo}

\[
\fitmath{r=\frac{\sqrt{\frac{w_c}{\sigma_s}}}{\Omega}}
\]

\noindent\textbf{Hipótesis}
\begin{itemize}
\item Salida radial, sin prerrotación, slip constante.
\end{itemize}

\noindent\textbf{Desarrollo paso a paso}

\begin{enumerate}

\item La salida radial sin prerrotación da wc=slip U2 al cuadrado.

\[\fitmath{\mathrm{uu}=\sqrt{\frac{\mathrm{wc}}{\mathrm{slip}}}}\]

\[\fitmath{\mathrm{uu}=\sqrt{\frac{(179985)}{(0.86)}}\approx 457.4766\,\mathrm{m/s}}\]

\item La geometría requerida sigue de U=Omega r.

\[\fitmath{\mathrm{rad}=\frac{\sqrt{\frac{\mathrm{wc}}{\mathrm{slip}}}}{\Omega}}\]

\[\fitmath{\mathrm{rad}=\frac{\sqrt{\frac{(179985)}{(0.86)}}}{(2140)}\approx 0.2137741\,\mathrm{m}}\]

\end{enumerate}


\noindent\textbf{Sustitución numérica}

\[
\fitmath{\frac{\sqrt{\frac{\left(179985\,\mathrm{J}\,\mathrm{kg}^{-1}\right)}{\left(0.86\right)}}}{\left(2140\,\mathrm{s}^{-1}\right)}}
\]

\noindent\textbf{Resultado}

\[
\fitmath{r=2.137741347\times10^{-1}\,\mathrm{m}}
\]

\noindent\textbf{Interpretación}
El tamaño y el giro están ligados al trabajo. Las tensiones y el Mach deben comprobarse después.

\noindent\textbf{Comprobación de los pasos}
\begin{itemize}
\item Paso 2: $457.477$ — reproducido
\item Paso 4: $0.213774$ — reproducido
\end{itemize}

\noindent\textbf{Comprobación de la respuesta} la aritmética del último paso reproduce la respuesta enunciada a la precisión con que está escrita. Es una reproducción aritmética; no valida el modelo físico ni sus hipótesis.

\noindent\textbf{Limitaciones}
\begin{itemize}
\item Modelo educativo; la aplicabilidad a una máquina real requiere revisión experta.
\end{itemize}

\noindent\textbf{Errores frecuentes}
\begin{itemize}
\item Sustituir eficiencia isentrópica por factor de deslizamiento.
\end{itemize}
% studium:end-record eyJpZCI6ICJQUkItMDEzMCIsICJzZWN0aW9uIjogImMxNyJ9
\end{studiumworked}

\begin{studiumworked}{Problema resuelto C17-02-R}


\noindent\textbf{Enunciado}
\(C17-02-R \cdot\) Presión y sensibilidad de eficiencia. Para trabajo de Euler fijado, calcule la relación de presión que produce el rendimiento suministrado. Datos: \(w_c=178000\, \mathrm{J}\,\mathrm{kg}^{-1}\); \(T_{0e}=297\, \mathrm{K}\); \(\eta=0.81\); \(c_p=1004.5\, \mathrm{J}\,\mathrm{kg}^{-1}\,\mathrm{K}^{-1}\); \(\gamma=1.4\).

\noindent\textbf{Datos}
\begin{itemize}
\item $w_c = 178000\,\mathrm{J}\,\mathrm{kg}^{-1}$ — trabajo de Euler comunicado
\item $T_{0e} = 297\,\mathrm{K}$ — temperatura total entrada
\item $\eta = 0.81$ — eficiencia total-total
\item $c_p = 1004.5\,\mathrm{J}\,\mathrm{kg}^{-1}\,\mathrm{K}^{-1}$ — calor específico
\item $\gamma = 1.4$ — razón de calores
\end{itemize}

\noindent\textbf{Incógnita}
$\pi_c$ (adimensional); también las magnitudes intermedias desarrolladas.

\noindent\textbf{Modelo físico}

\noindent\textit{Presión y sensibilidad de eficiencia}

\[
\fitmath{\pi_c=\left(1 + \frac{\eta w_c}{T_{0e} c_p}\right)^{\frac{\gamma}{\gamma - 1}}}
\]

\noindent\textbf{Hipótesis}
\begin{itemize}
\item Gas perfecto y trabajo fijado, sin refrigeración.
\end{itemize}

\noindent\textbf{Desarrollo paso a paso}

\begin{enumerate}

\item El trabajo real fija incremento total independientemente de eta.

\[\fitmath{\mathrm{dt}=\frac{\mathrm{wc}}{\mathrm{cp}}}\]

\[\fitmath{\mathrm{dt}=\frac{(178000)}{(1004.5)}\approx 177.2026\,\mathrm{K}}\]

\item La eficiencia fija el incremento isentrópico equivalente.

\[\fitmath{\mathrm{pr}={\left(1+\frac{\eta\cdot \mathrm{wc}}{\mathrm{cp}\cdot \mathrm{Tin}}\right)}^{\frac{\gamma}{\gamma-1}}}\]

\[\fitmath{\mathrm{pr}={\left(1+\frac{(0.81)\cdot (178000)}{(1004.5)\cdot (297)}\right)}^{\frac{(1.4)}{(1.4)-1}}\approx 3.974484\,}\]

\end{enumerate}


\noindent\textbf{Sustitución numérica}

\[
\fitmath{\left(1 + \frac{\left(0.81\right) \left(178000\,\mathrm{J}\,\mathrm{kg}^{-1}\right)}{\left(297\,\mathrm{K}\right) \left(1004.5\,\mathrm{J}\,\mathrm{kg}^{-1}\,\mathrm{K}^{-1}\right)}\right)^{\frac{\left(1.4\right)}{\left(1.4\right) - 1}}}
\]

\noindent\textbf{Resultado}

\[
\fitmath{\pi_c=3.974484195\times10^{0}}
\]

\noindent\textbf{Interpretación}
Variar eta a trabajo fijo cambia la presión, mientras el aumento real de temperatura permanece fijo.

\noindent\textbf{Comprobación de los pasos}
\begin{itemize}
\item Paso 2: $177.203$ — reproducido
\item Paso 4: $3.97448$ — reproducido
\end{itemize}

\noindent\textbf{Comprobación de la respuesta} la aritmética del último paso reproduce la respuesta enunciada a la precisión con que está escrita. Es una reproducción aritmética; no valida el modelo físico ni sus hipótesis.

\noindent\textbf{Limitaciones}
\begin{itemize}
\item Modelo educativo; la aplicabilidad a una máquina real requiere revisión experta.
\end{itemize}

\noindent\textbf{Errores frecuentes}
\begin{itemize}
\item Reducir el calentamiento real multiplicando por eficiencia.
\end{itemize}
% studium:end-record eyJpZCI6ICJQUkItMDEzMSIsICJzZWN0aW9uIjogImMxNyJ9
\end{studiumworked}

\begin{studiumworked}{Problema resuelto C17-02-P}


\noindent\textbf{Enunciado}
\(C17-02-P \cdot\) Presión y sensibilidad de eficiencia. Ejercicio de transferencia. Para trabajo de Euler fijado, calcule la relación de presión que produce el rendimiento suministrado. Antes de consultar la solución, prediga el efecto de cambiar el primer dato respecto del caso resuelto de esta familia. Justifique la tendencia y explique qué restricción podría convertirse en limitante. Datos: \(w_c=189570\, \mathrm{J}\,\mathrm{kg}^{-1}\); \(T_{0e}=297\, \mathrm{K}\); \(\eta=0.81\); \(c_p=1004.5\, \mathrm{J}\,\mathrm{kg}^{-1}\,\mathrm{K}^{-1}\); \(\gamma=1.4\).

\noindent\textbf{Datos}
\begin{itemize}
\item $w_c = 189570\,\mathrm{J}\,\mathrm{kg}^{-1}$ — trabajo de Euler comunicado
\item $T_{0e} = 297\,\mathrm{K}$ — temperatura total entrada
\item $\eta = 0.81$ — eficiencia total-total
\item $c_p = 1004.5\,\mathrm{J}\,\mathrm{kg}^{-1}\,\mathrm{K}^{-1}$ — calor específico
\item $\gamma = 1.4$ — razón de calores
\end{itemize}

\noindent\textbf{Incógnita}
$\pi_c$ (adimensional); también las magnitudes intermedias desarrolladas.

\noindent\textbf{Modelo físico}

\noindent\textit{Presión y sensibilidad de eficiencia}

\[
\fitmath{\pi_c=\left(1 + \frac{\eta w_c}{T_{0e} c_p}\right)^{\frac{\gamma}{\gamma - 1}}}
\]

\noindent\textbf{Hipótesis}
\begin{itemize}
\item Gas perfecto y trabajo fijado, sin refrigeración.
\end{itemize}

\noindent\textbf{Desarrollo paso a paso}

\begin{enumerate}

\item El trabajo real fija incremento total independientemente de eta.

\[\fitmath{\mathrm{dt}=\frac{\mathrm{wc}}{\mathrm{cp}}}\]

\[\fitmath{\mathrm{dt}=\frac{(189570)}{(1004.5)}\approx 188.7208\,\mathrm{K}}\]

\item La eficiencia fija el incremento isentrópico equivalente.

\[\fitmath{\mathrm{pr}={\left(1+\frac{\eta\cdot \mathrm{wc}}{\mathrm{cp}\cdot \mathrm{Tin}}\right)}^{\frac{\gamma}{\gamma-1}}}\]

\[\fitmath{\mathrm{pr}={\left(1+\frac{(0.81)\cdot (189570)}{(1004.5)\cdot (297)}\right)}^{\frac{(1.4)}{(1.4)-1}}\approx 4.276969\,}\]

\end{enumerate}


\noindent\textbf{Sustitución numérica}

\[
\fitmath{\left(1 + \frac{\left(0.81\right) \left(189570\,\mathrm{J}\,\mathrm{kg}^{-1}\right)}{\left(297\,\mathrm{K}\right) \left(1004.5\,\mathrm{J}\,\mathrm{kg}^{-1}\,\mathrm{K}^{-1}\right)}\right)^{\frac{\left(1.4\right)}{\left(1.4\right) - 1}}}
\]

\noindent\textbf{Resultado}

\[
\fitmath{\pi_c=4.276969338\times10^{0}}
\]

\noindent\textbf{Interpretación}
Variar eta a trabajo fijo cambia la presión, mientras el aumento real de temperatura permanece fijo.

\noindent\textbf{Comprobación de los pasos}
\begin{itemize}
\item Paso 2: $188.721$ — reproducido
\item Paso 4: $4.27697$ — reproducido
\end{itemize}

\noindent\textbf{Comprobación de la respuesta} la aritmética del último paso reproduce la respuesta enunciada a la precisión con que está escrita. Es una reproducción aritmética; no valida el modelo físico ni sus hipótesis.

\noindent\textbf{Limitaciones}
\begin{itemize}
\item Modelo educativo; la aplicabilidad a una máquina real requiere revisión experta.
\end{itemize}

\noindent\textbf{Errores frecuentes}
\begin{itemize}
\item Reducir el calentamiento real multiplicando por eficiencia.
\end{itemize}
% studium:end-record eyJpZCI6ICJQUkItMDEzMiIsICJzZWN0aW9uIjogImMxNyJ9
\end{studiumworked}

\begin{studiumworked}{Problema resuelto C17-03-R}


\noindent\textbf{Enunciado}
\(C17-03-R \cdot\) Pendiente cinemática con caudal. Un modelo de rodete mantiene densidad, ancho, slip y giro constantes. Calcule la derivada del trabajo respecto al caudal. Datos: \(U=434\, \mathrm{m}\,\mathrm{s}^{-1}\); \(\sigma_s=0.88\); \(\lambda_s=-0.5410521\, \mathrm{rad}\); \(r=0.203\, \mathrm{m}\); \(b=0.017\, \mathrm{m}\); \(\rho=3.2\, \mathrm{kg}\,\mathrm{m}^{-3}\).

\noindent\textbf{Datos}
\begin{itemize}
\item $U = 434\,\mathrm{m}\,\mathrm{s}^{-1}$ — velocidad periférica
\item $\sigma_s = 0.88$ — factor de deslizamiento
\item $\lambda_s = -0.54105206811824\,\mathrm{rad}$ — ángulo ideal desde meridiana
\item $r = 0.203\,\mathrm{m}$ — radio salida
\item $b = 0.017\,\mathrm{m}$ — ancho salida
\item $\rho = 3.2\,\mathrm{kg}\,\mathrm{m}^{-3}$ — densidad salida
\end{itemize}

\noindent\textbf{Incógnita}
$\mathrm{d}w_c/\mathrm{d}\dot{m}$, en $\mathrm{J\,s\,kg^{-2}}$; también las magnitudes intermedias desarrolladas.

\noindent\textbf{Modelo físico}

\noindent\textit{Pendiente cinemática con caudal}

\[
\fitmath{\frac{\mathrm{d}w_c}{\mathrm{d}\dot{m}}=\frac{\sigma_s U \tan{\left(\lambda_s \right)}}{2 \pi r \rho b}}
\]

\noindent\textbf{Hipótesis}
\begin{itemize}
\item Densidad, ancho, slip, radio y giro fijos; sin prerrotación.
\end{itemize}

\noindent\textbf{Desarrollo paso a paso}

\begin{enumerate}

\item La continuidad da la derivada de velocidad meridiana por caudal.

\[\fitmath{\mathrm{dcm}=\frac{1}{2\cdot \pi\cdot \mathrm{rad}\cdot \mathrm{width}\cdot \rho}}\]

\[\fitmath{\mathrm{dcm}=\frac{1}{2\cdot \pi\cdot (0.203)\cdot (0.017)\cdot (3.2)}\approx 14.41203\,\mathrm{m/kg}}\]

\item La derivada del modelo de Euler incorpora el ángulo.

\[\fitmath{\mathrm{slope}=\frac{\mathrm{uu}\cdot \mathrm{slip}\cdot \tan\left(\mathrm{lam}\right)}{2\cdot \pi\cdot \mathrm{rad}\cdot \mathrm{width}\cdot \rho}}\]

\[\fitmath{\mathrm{slope}=\frac{(434)\cdot (0.88)\cdot \tan\left((-0.5410520681)\right)}{2\cdot \pi\cdot (0.203)\cdot (0.017)\cdot (3.2)}\approx -3307.283\,\mathrm{J\,s/kg^{2}}}\]

\end{enumerate}


\noindent\textbf{Sustitución numérica}

\[
\fitmath{\frac{\left(0.88\right) \left(434\,\mathrm{m}\,\mathrm{s}^{-1}\right) \tan{\left(\left(-0.54105206811824\,\mathrm{rad}\right) \right)}}{2 \pi \left(0.203\,\mathrm{m}\right) \left(3.2\,\mathrm{kg}\,\mathrm{m}^{-3}\right) \left(0.017\,\mathrm{m}\right)}}
\]

\noindent\textbf{Resultado}

\[
\fitmath{\frac{\mathrm{d}w_c}{\mathrm{d}\dot{m}}=-3.307282833\times10^{3}\,\mathrm{J}\,\mathrm{s}\,\mathrm{kg}^{-2}}
\]

\noindent\textbf{Interpretación}
La pendiente negativa proviene de curvatura hacia atrás. No es una predicción de la línea de bombeo.

\noindent\textbf{Comprobación de los pasos}
\begin{itemize}
\item Paso 2: $14.412$ — reproducido
\item Paso 4: $-3307.28$ — reproducido
\end{itemize}

\noindent\textbf{Comprobación de la respuesta} la aritmética del último paso reproduce la respuesta enunciada a la precisión con que está escrita. Es una reproducción aritmética; no valida el modelo físico ni sus hipótesis.

\noindent\textbf{Limitaciones}
\begin{itemize}
\item Modelo educativo; la aplicabilidad a una máquina real requiere revisión experta.
\end{itemize}

\noindent\textbf{Errores frecuentes}
\begin{itemize}
\item Presentar la pendiente de Euler como estabilidad global.
\end{itemize}
% studium:end-record eyJpZCI6ICJQUkItMDEzMyIsICJzZWN0aW9uIjogImMxNyJ9
\end{studiumworked}

\begin{studiumworked}{Problema resuelto C17-03-P}


\noindent\textbf{Enunciado}
\(C17-03-P \cdot\) Pendiente cinemática con caudal. Ejercicio de transferencia. Un modelo de rodete mantiene densidad, ancho, slip y giro constantes. Calcule la derivada del trabajo respecto al caudal. Antes de consultar la solución, prediga el efecto de cambiar el primer dato respecto del caso resuelto de esta familia. Justifique la tendencia y explique qué restricción podría convertirse en limitante. Datos: \(U=462.21\, \mathrm{m}\,\mathrm{s}^{-1}\); \(\sigma_s=0.88\); \(\lambda_s=-0.5410521\, \mathrm{rad}\); \(r=0.203\, \mathrm{m}\); \(b=0.017\, \mathrm{m}\); \(\rho=3.2\, \mathrm{kg}\,\mathrm{m}^{-3}\).

\noindent\textbf{Datos}
\begin{itemize}
\item $U = 462.21\,\mathrm{m}\,\mathrm{s}^{-1}$ — velocidad periférica
\item $\sigma_s = 0.88$ — factor de deslizamiento
\item $\lambda_s = -0.54105206811824\,\mathrm{rad}$ — ángulo ideal desde meridiana
\item $r = 0.203\,\mathrm{m}$ — radio salida
\item $b = 0.017\,\mathrm{m}$ — ancho salida
\item $\rho = 3.2\,\mathrm{kg}\,\mathrm{m}^{-3}$ — densidad salida
\end{itemize}

\noindent\textbf{Incógnita}
$\mathrm{d}w_c/\mathrm{d}\dot{m}$, en $\mathrm{J\,s\,kg^{-2}}$; también las magnitudes intermedias desarrolladas.

\noindent\textbf{Modelo físico}

\noindent\textit{Pendiente cinemática con caudal}

\[
\fitmath{\frac{\mathrm{d}w_c}{\mathrm{d}\dot{m}}=\frac{\sigma_s U \tan{\left(\lambda_s \right)}}{2 \pi r \rho b}}
\]

\noindent\textbf{Hipótesis}
\begin{itemize}
\item Densidad, ancho, slip, radio y giro fijos; sin prerrotación.
\end{itemize}

\noindent\textbf{Desarrollo paso a paso}

\begin{enumerate}

\item La continuidad da la derivada de velocidad meridiana por caudal.

\[\fitmath{\mathrm{dcm}=\frac{1}{2\cdot \pi\cdot \mathrm{rad}\cdot \mathrm{width}\cdot \rho}}\]

\[\fitmath{\mathrm{dcm}=\frac{1}{2\cdot \pi\cdot (0.203)\cdot (0.017)\cdot (3.2)}\approx 14.41203\,\mathrm{m/kg}}\]

\item La derivada del modelo de Euler incorpora el ángulo.

\[\fitmath{\mathrm{slope}=\frac{\mathrm{uu}\cdot \mathrm{slip}\cdot \tan\left(\mathrm{lam}\right)}{2\cdot \pi\cdot \mathrm{rad}\cdot \mathrm{width}\cdot \rho}}\]

\[\fitmath{\mathrm{slope}=\frac{(462.21)\cdot (0.88)\cdot \tan\left((-0.5410520681)\right)}{2\cdot \pi\cdot (0.203)\cdot (0.017)\cdot (3.2)}\approx -3522.256\,\mathrm{J\,s/kg^{2}}}\]

\end{enumerate}


\noindent\textbf{Sustitución numérica}

\[
\fitmath{\frac{\left(0.88\right) \left(462.21\,\mathrm{m}\,\mathrm{s}^{-1}\right) \tan{\left(\left(-0.54105206811824\,\mathrm{rad}\right) \right)}}{2 \pi \left(0.203\,\mathrm{m}\right) \left(3.2\,\mathrm{kg}\,\mathrm{m}^{-3}\right) \left(0.017\,\mathrm{m}\right)}}
\]

\noindent\textbf{Resultado}

\[
\fitmath{\frac{\mathrm{d}w_c}{\mathrm{d}\dot{m}}=-3.522256217\times10^{3}\,\mathrm{J}\,\mathrm{s}\,\mathrm{kg}^{-2}}
\]

\noindent\textbf{Interpretación}
La pendiente negativa proviene de curvatura hacia atrás. No es una predicción de la línea de bombeo.

\noindent\textbf{Comprobación de los pasos}
\begin{itemize}
\item Paso 2: $14.412$ — reproducido
\item Paso 4: $-3522.26$ — reproducido
\end{itemize}

\noindent\textbf{Comprobación de la respuesta} la aritmética del último paso reproduce la respuesta enunciada a la precisión con que está escrita. Es una reproducción aritmética; no valida el modelo físico ni sus hipótesis.

\noindent\textbf{Limitaciones}
\begin{itemize}
\item Modelo educativo; la aplicabilidad a una máquina real requiere revisión experta.
\end{itemize}

\noindent\textbf{Errores frecuentes}
\begin{itemize}
\item Presentar la pendiente de Euler como estabilidad global.
\end{itemize}
% studium:end-record eyJpZCI6ICJQUkItMDEzNCIsICJzZWN0aW9uIjogImMxNyJ9
\end{studiumworked}

\begin{studiumworked}{Problema resuelto C17-04-R}


\noindent\textbf{Enunciado}
\(C17-04-R \cdot\) Balance inverso de rendimiento. De un ensayo educativo se conocen aumento total de temperatura y relación de presión. Calcule la eficiencia total-total y juzgue admisibilidad. Datos: \(T_{0e}=303\, \mathrm{K}\); \(T_{0s}=469\, \mathrm{K}\); \(\pi_c=3.85\); \(\gamma=1.4\).

\noindent\textbf{Datos}
\begin{itemize}
\item $T_{0e} = 303\,\mathrm{K}$ — temperatura total entrada
\item $T_{0s} = 469\,\mathrm{K}$ — temperatura total salida
\item $\pi_c = 3.85$ — relación de presión total
\item $\gamma = 1.4$ — razón de calores
\end{itemize}

\noindent\textbf{Incógnita}
$\eta$ (adimensional); también las magnitudes intermedias desarrolladas.

\noindent\textbf{Modelo físico}

\noindent\textit{Balance inverso de rendimiento}

\[
\fitmath{\eta=\frac{T_{0e} \left(\pi_c^{\frac{\gamma - 1}{\gamma}} - 1\right)}{- T_{0e} + T_{0s}}}
\]

\noindent\textbf{Hipótesis}
\begin{itemize}
\item Gas perfecto constante, sin calor externo; mediciones medias representativas.
\end{itemize}

\noindent\textbf{Desarrollo paso a paso}

\begin{enumerate}

\item Se calcula la temperatura isentrópica equivalente.

\[\fitmath{\mathrm{Tis}=\mathrm{Tin}\cdot {\mathrm{pr}}^{\frac{\gamma-1}{\gamma}}}\]

\[\fitmath{\mathrm{Tis}=(303)\cdot {(3.85)}^{\frac{(1.4)-1}{(1.4)}}\approx 445.3661\,\mathrm{K}}\]

\item Se divide aumento ideal por aumento real.

\[\fitmath{\mathrm{eta}=\frac{\mathrm{Tin}\cdot \left({\mathrm{pr}}^{\frac{\gamma-1}{\gamma}}-1\right)}{\mathrm{Tout}-\mathrm{Tin}}}\]

\[\fitmath{\mathrm{eta}=\frac{(303)\cdot \left({(3.85)}^{\frac{(1.4)-1}{(1.4)}}-1\right)}{(469)-(303)}\approx 0.8576269\,}\]

\end{enumerate}


\noindent\textbf{Sustitución numérica}

\[
\fitmath{\frac{\left(303\,\mathrm{K}\right) \left(\left(3.85\right)^{\frac{\left(1.4\right) - 1}{\left(1.4\right)}} - 1\right)}{- \left(303\,\mathrm{K}\right) + \left(469\,\mathrm{K}\right)}}
\]

\noindent\textbf{Resultado}

\[
\fitmath{\eta=8.576268981\times10^{-1}}
\]

\noindent\textbf{Interpretación}
El valor debe estar entre cero y uno bajo las hipótesis adiabáticas. Un valor admisible no prueba precisión experimental.

\noindent\textbf{Comprobación de los pasos}
\begin{itemize}
\item Paso 2: $445.366$ — reproducido
\item Paso 4: $0.857627$ — reproducido
\end{itemize}

\noindent\textbf{Comprobación de la respuesta} la aritmética del último paso reproduce la respuesta enunciada a la precisión con que está escrita. Es una reproducción aritmética; no valida el modelo físico ni sus hipótesis.

\noindent\textbf{Limitaciones}
\begin{itemize}
\item Modelo educativo; la aplicabilidad a una máquina real requiere revisión experta.
\end{itemize}

\noindent\textbf{Errores frecuentes}
\begin{itemize}
\item Confundir presión estática medida con total sin corrección.
\end{itemize}
% studium:end-record eyJpZCI6ICJQUkItMDEzNSIsICJzZWN0aW9uIjogImMxNyJ9
\end{studiumworked}

\begin{studiumworked}{Problema resuelto C17-04-P}


\noindent\textbf{Enunciado}
\(C17-04-P \cdot\) Balance inverso de rendimiento. Ejercicio de transferencia. De un ensayo educativo se conocen aumento total de temperatura y relación de presión. Calcule la eficiencia total-total y juzgue admisibilidad. Antes de consultar la solución, prediga el efecto de cambiar el primer dato respecto del caso resuelto de esta familia. Justifique la tendencia y explique qué restricción podría convertirse en limitante. Datos: \(T_{0e}=322.695\, \mathrm{K}\); \(T_{0s}=469\, \mathrm{K}\); \(\pi_c=3.85\); \(\gamma=1.4\).

\noindent\textbf{Datos}
\begin{itemize}
\item $T_{0e} = 322.695\,\mathrm{K}$ — temperatura total entrada
\item $T_{0s} = 469\,\mathrm{K}$ — temperatura total salida
\item $\pi_c = 3.85$ — relación de presión total
\item $\gamma = 1.4$ — razón de calores
\end{itemize}

\noindent\textbf{Incógnita}
$\eta$ (adimensional); también las magnitudes intermedias desarrolladas.

\noindent\textbf{Modelo físico}

\noindent\textit{Balance inverso de rendimiento}

\[
\fitmath{\eta=\frac{T_{0e} \left(\pi_c^{\frac{\gamma - 1}{\gamma}} - 1\right)}{- T_{0e} + T_{0s}}}
\]

\noindent\textbf{Hipótesis}
\begin{itemize}
\item Gas perfecto constante, sin calor externo; mediciones medias representativas.
\end{itemize}

\noindent\textbf{Desarrollo paso a paso}

\begin{enumerate}

\item Se calcula la temperatura isentrópica equivalente.

\[\fitmath{\mathrm{Tis}=\mathrm{Tin}\cdot {\mathrm{pr}}^{\frac{\gamma-1}{\gamma}}}\]

\[\fitmath{\mathrm{Tis}=(322.695)\cdot {(3.85)}^{\frac{(1.4)-1}{(1.4)}}\approx 474.3149\,\mathrm{K}}\]

\item Se divide aumento ideal por aumento real.

\[\fitmath{\eta=\frac{\mathrm{Tin}\cdot \left({\mathrm{pr}}^{\frac{\gamma-1}{\gamma}}-1\right)}{\mathrm{Tout}-\mathrm{Tin}}}\]

\[\fitmath{\eta=\frac{(322.695)\cdot \left({(3.85)}^{\frac{(1.4)-1}{(1.4)}}-1\right)}{(469)-(322.695)}\approx 1.036327\,}\]

\end{enumerate}


\noindent\textbf{Sustitución numérica}

\[
\fitmath{\frac{\left(322.695\,\mathrm{K}\right) \left(\left(3.85\right)^{\frac{\left(1.4\right) - 1}{\left(1.4\right)}} - 1\right)}{- \left(322.695\,\mathrm{K}\right) + \left(469\,\mathrm{K}\right)}}
\]

\noindent\textbf{Resultado}

\[
\fitmath{\eta=1.036327257\times10^{0}}
\]

\noindent\textbf{Interpretación}
La eficiencia calculada supera uno. Estos datos no son compatibles con el modelo adiabático: revise identificación estática-total, calibración y calor externo. El resultado es un diagnóstico de inconsistencia, no una prestación admisible.

\noindent\textbf{Comprobación de los pasos}
\begin{itemize}
\item Paso 2: $474.315$ — reproducido
\item Paso 4: $1.03633$ — reproducido
\end{itemize}

\noindent\textbf{Comprobación de la respuesta} la aritmética del último paso reproduce la respuesta enunciada a la precisión con que está escrita. Es una reproducción aritmética; no valida el modelo físico ni sus hipótesis.

\noindent\textbf{Limitaciones}
\begin{itemize}
\item Modelo educativo; la aplicabilidad a una máquina real requiere revisión experta.
\end{itemize}

\noindent\textbf{Errores frecuentes}
\begin{itemize}
\item Confundir presión estática medida con total sin corrección.
\end{itemize}
% studium:end-record eyJpZCI6ICJQUkItMDE1MSIsICJzZWN0aW9uIjogImMxNyJ9
\end{studiumworked}

\chapter{Práctica reproducible: secciones a distintos radios}
\begin{studiumworked}{Problema resuelto C18-01-R}


\noindent\textbf{Enunciado}
\(C18-01-R \cdot\) Rotor compresor a radio elegido. Aplique la plantilla radial educativa a la entrada de un rotor compresor. Calcule U, beta y ángulo metálico con incidencia prescrita. Datos: \(r=0.271\, \mathrm{m}\); \(\Omega=1010\, \mathrm{s}^{-1}\); \(K_e=14\, \mathrm{m}^{2}\,\mathrm{s}^{-1}\); \(C_x=159\, \mathrm{m}\,\mathrm{s}^{-1}\); \(i=-0.01745329\, \mathrm{rad}\).

\noindent\textbf{Datos}
\begin{itemize}
\item $r = 0.271\,\mathrm{m}$ — radio elegido
\item $\Omega = 1010\,\mathrm{s}^{-1}$ — velocidad angular
\item $K_e = 14\,\mathrm{m}^{2}\,\mathrm{s}^{-1}$ — constante remolino de entrada
\item $C_x = 159\,\mathrm{m}\,\mathrm{s}^{-1}$ — velocidad axial
\item $i = -0.017453292519943\,\mathrm{rad}$ — incidencia algebraica prescrita
\end{itemize}

\noindent\textbf{Incógnita}
$\kappa_e$, en $\mathrm{rad}$; también las magnitudes intermedias desarrolladas.

\noindent\textbf{Modelo físico}

\noindent\textit{Rotor compresor a radio elegido}

\[
\fitmath{\kappa_e=- i + \operatorname{atan}{\left(\frac{- \Omega r + \frac{K_e}{r}}{C_x} \right)}}
\]

\noindent\textbf{Hipótesis}
\begin{itemize}
\item Vórtice libre por estación, incidencia prescrita.
\end{itemize}

\noindent\textbf{Desarrollo paso a paso}

\begin{enumerate}

\item Velocidad periférica

\[\fitmath{U=\Omega\cdot \mathrm{rad}}\]

\[\fitmath{U=(1010)\cdot (0.271)\approx 273.71\,\mathrm{m/s}}\]

\item La velocidad tangencial absoluta sigue vórtice libre.

\[\fitmath{\mathrm{ct}=\frac{\mathrm{kv}}{\mathrm{rad}}}\]

\[\fitmath{\mathrm{ct}=\frac{(14)}{(0.271)}\approx 51.66052\,\mathrm{m/s}}\]

\item Se forma el ángulo del gas en referencia rotor.

\[\fitmath{\beta=\arctan\left(\frac{\frac{\mathrm{kv}}{\mathrm{rad}}-\Omega\cdot \mathrm{rad}}{\mathrm{cx}}\right)}\]

\[\fitmath{\beta=\arctan\left(\frac{\frac{(14)}{(0.271)}-(1010)\cdot (0.271)}{(159)}\right)\approx -0.9493752\,\mathrm{rad}}\]

\item Se despeja metal a partir de incidencia definida.

\[\fitmath{\mathrm{kap}=\arctan\left(\frac{\frac{\mathrm{kv}}{\mathrm{rad}}-\Omega\cdot \mathrm{rad}}{\mathrm{cx}}\right)-\mathrm{inc}}\]

\[\fitmath{\mathrm{kap}=\arctan\left(\frac{\frac{(14)}{(0.271)}-(1010)\cdot (0.271)}{(159)}\right)-(-0.01745329252)\approx -0.9319219\,\mathrm{rad}}\]

\end{enumerate}


\noindent\textbf{Sustitución numérica}

\[
\fitmath{- \left(-0.017453292519943\,\mathrm{rad}\right) + \operatorname{atan}{\left(\frac{- \left(1010\,\mathrm{s}^{-1}\right) \left(0.271\,\mathrm{m}\right) + \frac{\left(14\,\mathrm{m}^{2}\,\mathrm{s}^{-1}\right)}{\left(0.271\,\mathrm{m}\right)}}{\left(159\,\mathrm{m}\,\mathrm{s}^{-1}\right)} \right)}}
\]

\noindent\textbf{Resultado}

\[
\fitmath{\kappa_e=-9.319219110\times10^{-1}\,\mathrm{rad}}
\]

\noindent\textbf{Velocidad periférica}
\[
\fitmath{U=\Omega r=273.71\,\mathrm{m/s}}
\]

\noindent\textbf{Comprobación cinemática.} Insertar este $U$ en $\beta=\arctan[(C_\theta-U)/C_x]$ reproduce el ángulo desarrollado.

\noindent\textbf{Interpretación}
La misma tabla debe actualizarse al variar r u Omega. El metal es una propuesta geométrica educativa.

\noindent\textbf{Comprobación de los pasos}
\begin{itemize}
\item Paso 2: $51.6605$ — reproducido
\item Paso 4: $-0.949375$ — reproducido
\item Paso 6: $-0.931922$ — reproducido
\end{itemize}

\noindent\textbf{Comprobación de la respuesta} la aritmética del último paso reproduce la respuesta enunciada a la precisión con que está escrita. Es una reproducción aritmética; no valida el modelo físico ni sus hipótesis.

\noindent\textbf{Limitaciones}
\begin{itemize}
\item Modelo educativo; la aplicabilidad a una máquina real requiere revisión experta.
\end{itemize}

\noindent\textbf{Errores frecuentes}
\begin{itemize}
\item Usar alfa para construir un rotor.
\end{itemize}
% studium:end-record eyJpZCI6ICJQUkItMDEzNyIsICJzZWN0aW9uIjogImMxOCJ9
\end{studiumworked}

\begin{studiumworked}{Problema resuelto C18-01-P}


\noindent\textbf{Enunciado}
\(C18-01-P \cdot\) Rotor compresor a radio elegido. Ejercicio de transferencia. Aplique la plantilla radial educativa a la entrada de un rotor compresor. Calcule U, beta y ángulo metálico con incidencia prescrita. Antes de consultar la solución, prediga el efecto de cambiar el primer dato respecto del caso resuelto de esta familia. Justifique la tendencia y explique qué restricción podría convertirse en limitante. Datos: \(r=0.288615\, \mathrm{m}\); \(\Omega=1010\, \mathrm{s}^{-1}\); \(K_e=14\, \mathrm{m}^{2}\,\mathrm{s}^{-1}\); \(C_x=159\, \mathrm{m}\,\mathrm{s}^{-1}\); \(i=-0.01745329\, \mathrm{rad}\).

\noindent\textbf{Datos}
\begin{itemize}
\item $r = 0.288615\,\mathrm{m}$ — radio elegido
\item $\Omega = 1010\,\mathrm{s}^{-1}$ — velocidad angular
\item $K_e = 14\,\mathrm{m}^{2}\,\mathrm{s}^{-1}$ — constante remolino de entrada
\item $C_x = 159\,\mathrm{m}\,\mathrm{s}^{-1}$ — velocidad axial
\item $i = -0.017453292519943\,\mathrm{rad}$ — incidencia algebraica prescrita
\end{itemize}

\noindent\textbf{Incógnita}
$\kappa_e$, en $\mathrm{rad}$; también las magnitudes intermedias desarrolladas.

\noindent\textbf{Modelo físico}

\noindent\textit{Rotor compresor a radio elegido}

\[
\fitmath{\kappa_e=- i + \operatorname{atan}{\left(\frac{- \Omega r + \frac{K_e}{r}}{C_x} \right)}}
\]

\noindent\textbf{Hipótesis}
\begin{itemize}
\item Vórtice libre por estación, incidencia prescrita.
\end{itemize}

\noindent\textbf{Desarrollo paso a paso}

\begin{enumerate}

\item Velocidad periférica

\[\fitmath{U=\Omega\cdot \mathrm{rad}}\]

\[\fitmath{U=(1010)\cdot (0.288615)\approx 291.5011\,\mathrm{m/s}}\]

\item La velocidad tangencial absoluta sigue vórtice libre.

\[\fitmath{\mathrm{ct}=\frac{\mathrm{kv}}{\mathrm{rad}}}\]

\[\fitmath{\mathrm{ct}=\frac{(14)}{(0.288615)}\approx 48.50753\,\mathrm{m/s}}\]

\item Se forma el ángulo del gas en referencia rotor.

\[\fitmath{\beta=\arctan\left(\frac{\frac{\mathrm{kv}}{\mathrm{rad}}-\Omega\cdot \mathrm{rad}}{\mathrm{cx}}\right)}\]

\[\fitmath{\beta=\arctan\left(\frac{\frac{(14)}{(0.288615)}-(1010)\cdot (0.288615)}{(159)}\right)\approx -0.9913775\,\mathrm{rad}}\]

\item Se despeja metal a partir de incidencia definida.

\[\fitmath{\mathrm{kap}=\arctan\left(\frac{\frac{\mathrm{kv}}{\mathrm{rad}}-\Omega\cdot \mathrm{rad}}{\mathrm{cx}}\right)-\mathrm{inc}}\]

\[\fitmath{\mathrm{kap}=\arctan\left(\frac{\frac{(14)}{(0.288615)}-(1010)\cdot (0.288615)}{(159)}\right)-(-0.01745329252)\approx -0.9739242\,\mathrm{rad}}\]

\end{enumerate}


\noindent\textbf{Sustitución numérica}

\[
\fitmath{- \left(-0.017453292519943\,\mathrm{rad}\right) + \operatorname{atan}{\left(\frac{- \left(1010\,\mathrm{s}^{-1}\right) \left(0.288615\,\mathrm{m}\right) + \frac{\left(14\,\mathrm{m}^{2}\,\mathrm{s}^{-1}\right)}{\left(0.288615\,\mathrm{m}\right)}}{\left(159\,\mathrm{m}\,\mathrm{s}^{-1}\right)} \right)}}
\]

\noindent\textbf{Resultado}

\[
\fitmath{\kappa_e=-9.739241832\times10^{-1}\,\mathrm{rad}}
\]

\noindent\textbf{Velocidad periférica}
\[
\fitmath{U=\Omega r=291.50115\,\mathrm{m/s}}
\]

\noindent\textbf{Comprobación cinemática.} Insertar este $U$ en $\beta=\arctan[(C_\theta-U)/C_x]$ reproduce el ángulo desarrollado.

\noindent\textbf{Interpretación}
La misma tabla debe actualizarse al variar r u Omega. El metal es una propuesta geométrica educativa.

\noindent\textbf{Comprobación de los pasos}
\begin{itemize}
\item Paso 2: $48.5075$ — reproducido
\item Paso 4: $-0.991377$ — reproducido
\item Paso 6: $-0.973924$ — reproducido
\end{itemize}

\noindent\textbf{Comprobación de la respuesta} la aritmética del último paso reproduce la respuesta enunciada a la precisión con que está escrita. Es una reproducción aritmética; no valida el modelo físico ni sus hipótesis.

\noindent\textbf{Limitaciones}
\begin{itemize}
\item Modelo educativo; la aplicabilidad a una máquina real requiere revisión experta.
\end{itemize}

\noindent\textbf{Errores frecuentes}
\begin{itemize}
\item Usar alfa para construir un rotor.
\end{itemize}
% studium:end-record eyJpZCI6ICJQUkItMDEzOCIsICJzZWN0aW9uIjogImMxOCJ9
\end{studiumworked}

\begin{studiumworked}{Problema resuelto C18-02-R}


\noindent\textbf{Enunciado}
\(C18-02-R \cdot\) Estator compresor y orientación de salida. Con remolino de salida prescrito, halle el ángulo absoluto y el ángulo metálico para un estator compresor. Datos: \(r=0.282\, \mathrm{m}\); \(K_s=9\, \mathrm{m}^{2}\,\mathrm{s}^{-1}\); \(C_x=162\, \mathrm{m}\,\mathrm{s}^{-1}\); \(\delta=0.03839724\, \mathrm{rad}\).

\noindent\textbf{Datos}
\begin{itemize}
\item $r = 0.282\,\mathrm{m}$ — radio
\item $K_s = 9\,\mathrm{m}^{2}\,\mathrm{s}^{-1}$ — constante remolino salida
\item $C_x = 162\,\mathrm{m}\,\mathrm{s}^{-1}$ — velocidad axial
\item $\delta = 0.038397243543875\,\mathrm{rad}$ — desviación algebraica prescrita
\end{itemize}

\noindent\textbf{Incógnita}
$\kappa_e$, en $\mathrm{rad}$; también las magnitudes intermedias desarrolladas.

\noindent\textbf{Modelo físico}

\noindent\textit{Estator compresor y orientación de salida}

\[
\fitmath{\kappa_e=- \delta + \operatorname{atan}{\left(\frac{K_s}{C_x r} \right)}}
\]

\noindent\textbf{Hipótesis}
\begin{itemize}
\item Velocidades absolutas prescritas, desviación como dato.
\end{itemize}

\noindent\textbf{Desarrollo paso a paso}

\begin{enumerate}

\item La fila fija utiliza velocidad absoluta.

\[\fitmath{\alpha=\arctan\left(\frac{\mathrm{kv}}{\mathrm{rad}\cdot \mathrm{cx}}\right)}\]

\[\fitmath{\alpha=\arctan\left(\frac{(9)}{(0.282)\cdot (162)}\right)\approx 0.1945146\,\mathrm{rad}}\]

\item La definición delta=gas-metal permite despejar metal.

\[\fitmath{\mathrm{kap}=\arctan\left(\frac{\mathrm{kv}}{\mathrm{rad}\cdot \mathrm{cx}}\right)-\mathrm{dev}}\]

\[\fitmath{\mathrm{kap}=\arctan\left(\frac{(9)}{(0.282)\cdot (162)}\right)-(0.03839724354)\approx 0.1561174\,\mathrm{rad}}\]

\end{enumerate}


\noindent\textbf{Sustitución numérica}

\[
\fitmath{- \left(0.038397243543875\,\mathrm{rad}\right) + \operatorname{atan}{\left(\frac{\left(9\,\mathrm{m}^{2}\,\mathrm{s}^{-1}\right)}{\left(162\,\mathrm{m}\,\mathrm{s}^{-1}\right) \left(0.282\,\mathrm{m}\right)} \right)}}
\]

\noindent\textbf{Resultado}

\[
\fitmath{\kappa_e=1.561173538\times10^{-1}\,\mathrm{rad}}
\]

\noindent\textbf{Interpretación}
No aparece Omega porque la fila está fija; solo cambiaría el resultado si el cambio de operación modifica los datos absolutos.

\noindent\textbf{Comprobación de los pasos}
\begin{itemize}
\item Paso 2: $0.194515$ — reproducido
\item Paso 4: $0.156117$ — reproducido
\end{itemize}

\noindent\textbf{Comprobación de la respuesta} la aritmética del último paso reproduce la respuesta enunciada a la precisión con que está escrita. Es una reproducción aritmética; no valida el modelo físico ni sus hipótesis.

\noindent\textbf{Limitaciones}
\begin{itemize}
\item Modelo educativo; la aplicabilidad a una máquina real requiere revisión experta.
\end{itemize}

\noindent\textbf{Errores frecuentes}
\begin{itemize}
\item Restar U a una fila fija.
\end{itemize}
% studium:end-record eyJpZCI6ICJQUkItMDEzOSIsICJzZWN0aW9uIjogImMxOCJ9
\end{studiumworked}

\begin{studiumworked}{Problema resuelto C18-02-P}


\noindent\textbf{Enunciado}
\(C18-02-P \cdot\) Estator compresor y orientación de salida. Ejercicio de transferencia. Con remolino de salida prescrito, halle el ángulo absoluto y el ángulo metálico para un estator compresor. Antes de consultar la solución, prediga el efecto de cambiar el primer dato respecto del caso resuelto de esta familia. Justifique la tendencia y explique qué restricción podría convertirse en limitante. Datos: \(r=0.30033\, \mathrm{m}\); \(K_s=9\, \mathrm{m}^{2}\,\mathrm{s}^{-1}\); \(C_x=162\, \mathrm{m}\,\mathrm{s}^{-1}\); \(\delta=0.03839724\, \mathrm{rad}\).

\noindent\textbf{Datos}
\begin{itemize}
\item $r = 0.30033\,\mathrm{m}$ — radio
\item $K_s = 9\,\mathrm{m}^{2}\,\mathrm{s}^{-1}$ — constante remolino salida
\item $C_x = 162\,\mathrm{m}\,\mathrm{s}^{-1}$ — velocidad axial
\item $\delta = 0.038397243543875\,\mathrm{rad}$ — desviación algebraica prescrita
\end{itemize}

\noindent\textbf{Incógnita}
$\kappa_e$, en $\mathrm{rad}$; también las magnitudes intermedias desarrolladas.

\noindent\textbf{Modelo físico}

\noindent\textit{Estator compresor y orientación de salida}

\[
\fitmath{\kappa_e=- \delta + \operatorname{atan}{\left(\frac{K_s}{C_x r} \right)}}
\]

\noindent\textbf{Hipótesis}
\begin{itemize}
\item Velocidades absolutas prescritas, desviación como dato.
\end{itemize}

\noindent\textbf{Desarrollo paso a paso}

\begin{enumerate}

\item La fila fija utiliza velocidad absoluta.

\[\fitmath{\alpha=\arctan\left(\frac{\mathrm{kv}}{\mathrm{rad}\cdot \mathrm{cx}}\right)}\]

\[\fitmath{\alpha=\arctan\left(\frac{(9)}{(0.30033)\cdot (162)}\right)\approx 0.1829141\,\mathrm{rad}}\]

\item La definición delta=gas-metal permite despejar metal.

\[\fitmath{\mathrm{kap}=\arctan\left(\frac{\mathrm{kv}}{\mathrm{rad}\cdot \mathrm{cx}}\right)-\mathrm{dev}}\]

\[\fitmath{\mathrm{kap}=\arctan\left(\frac{(9)}{(0.30033)\cdot (162)}\right)-(0.03839724354)\approx 0.1445168\,\mathrm{rad}}\]

\end{enumerate}


\noindent\textbf{Sustitución numérica}

\[
\fitmath{- \left(0.038397243543875\,\mathrm{rad}\right) + \operatorname{atan}{\left(\frac{\left(9\,\mathrm{m}^{2}\,\mathrm{s}^{-1}\right)}{\left(162\,\mathrm{m}\,\mathrm{s}^{-1}\right) \left(0.30033\,\mathrm{m}\right)} \right)}}
\]

\noindent\textbf{Resultado}

\[
\fitmath{\kappa_e=1.445168334\times10^{-1}\,\mathrm{rad}}
\]

\noindent\textbf{Interpretación}
No aparece Omega porque la fila está fija; solo cambiaría el resultado si el cambio de operación modifica los datos absolutos.

\noindent\textbf{Comprobación de los pasos}
\begin{itemize}
\item Paso 2: $0.182914$ — reproducido
\item Paso 4: $0.144517$ — reproducido
\end{itemize}

\noindent\textbf{Comprobación de la respuesta} la aritmética del último paso reproduce la respuesta enunciada a la precisión con que está escrita. Es una reproducción aritmética; no valida el modelo físico ni sus hipótesis.

\noindent\textbf{Limitaciones}
\begin{itemize}
\item Modelo educativo; la aplicabilidad a una máquina real requiere revisión experta.
\end{itemize}

\noindent\textbf{Errores frecuentes}
\begin{itemize}
\item Restar U a una fila fija.
\end{itemize}
% studium:end-record eyJpZCI6ICJQUkItMDE0MCIsICJzZWN0aW9uIjogImMxOCJ9
\end{studiumworked}

\begin{studiumworked}{Problema resuelto C18-03-R}


\noindent\textbf{Enunciado}
\(C18-03-R \cdot\) Estator turbina y calado educativo. Una sección educativa del estator de turbina tiene pendientes metálicas dadas. Con d igual al promedio de pendientes, calcule calado y cuerda real. Datos: \(m_e=0.12\); \(m_s=1.38\); \(c_x=0.046\, \mathrm{m}\).

\noindent\textbf{Datos}
\begin{itemize}
\item $m_e = 0.12$ — pendiente metal entrada
\item $m_s = 1.38$ — pendiente metal salida
\item $c_x = 0.046\,\mathrm{m}$ — cuerda axial
\end{itemize}

\noindent\textbf{Incógnita}
$c$, en $\mathrm{m}$; también las magnitudes intermedias desarrolladas.

\noindent\textbf{Modelo físico}

\noindent\textit{Estator turbina y calado educativo}

\[
\fitmath{c=c_x \sqrt{\left(\frac{m_e}{2} + \frac{m_s}{2}\right)^{2} + 1}}
\]

\noindent\textbf{Hipótesis}
\begin{itemize}
\item Línea media Hermite educativa; x monótona.
\end{itemize}

\noindent\textbf{Desarrollo paso a paso}

\begin{enumerate}

\item El desplazamiento normalizado se prescribe como promedio de pendientes.

\[\fitmath{\mathrm{dd}=\frac{\mathrm{me}+\mathrm{ms}}{2}}\]

\[\fitmath{\mathrm{dd}=\frac{(0.12)+(1.38)}{2}\approx 0.75\,}\]

\item El calado de cuerda es atan d.

\[\fitmath{\chi=\arctan\left(\frac{\mathrm{me}+\mathrm{ms}}{2}\right)}\]

\[\fitmath{\chi=\arctan\left(\frac{(0.12)+(1.38)}{2}\right)\approx 0.6435011\,\mathrm{rad}}\]

\item La cuerda real se deduce del desplazamiento final.

\[\fitmath{\mathrm{chord}=\mathrm{cxg}\cdot \sqrt{1+{\left(\frac{\mathrm{me}+\mathrm{ms}}{2}\right)}^{2}}}\]

\[\fitmath{\mathrm{chord}=(0.046)\cdot \sqrt{1+{\left(\frac{(0.12)+(1.38)}{2}\right)}^{2}}\approx 0.0575\,\mathrm{m}}\]

\end{enumerate}


\noindent\textbf{Sustitución numérica}

\[
\fitmath{\left(0.046\,\mathrm{m}\right) \sqrt{\left(\frac{\left(0.12\right)}{2} + \frac{\left(1.38\right)}{2}\right)^{2} + 1}}
\]

\noindent\textbf{Resultado}

\[
\fitmath{c=5.750000000\times10^{-2}\,\mathrm{m}}
\]

\noindent\textbf{Interpretación}
Promediar pendientes no equivale a promediar ángulos. La ley es geométrica y no una correlación de mínima pérdida.

\noindent\textbf{Comprobación de los pasos}
\begin{itemize}
\item Paso 2: $0.75$ — reproducido
\item Paso 4: $0.643501$ — reproducido
\item Paso 6: $0.0575$ — reproducido
\end{itemize}

\noindent\textbf{Comprobación de la respuesta} la aritmética del último paso reproduce la respuesta enunciada a la precisión con que está escrita. Es una reproducción aritmética; no valida el modelo físico ni sus hipótesis.

\noindent\textbf{Limitaciones}
\begin{itemize}
\item Modelo educativo; la aplicabilidad a una máquina real requiere revisión experta.
\end{itemize}

\noindent\textbf{Errores frecuentes}
\begin{itemize}
\item Presentar esta sección como un perfil de fabricante.
\end{itemize}
% studium:end-record eyJpZCI6ICJQUkItMDE0MSIsICJzZWN0aW9uIjogImMxOCJ9
\end{studiumworked}

\begin{studiumworked}{Problema resuelto C18-03-P}


\noindent\textbf{Enunciado}
\(C18-03-P \cdot\) Estator turbina y calado educativo. Ejercicio de transferencia. Una sección educativa del estator de turbina tiene pendientes metálicas dadas. Con d igual al promedio de pendientes, calcule calado y cuerda real. Antes de consultar la solución, prediga el efecto de cambiar el primer dato respecto del caso resuelto de esta familia. Justifique la tendencia y explique qué restricción podría convertirse en limitante. Datos: \(m_e=0.1278\); \(m_s=1.38\); \(c_x=0.046\, \mathrm{m}\).

\noindent\textbf{Datos}
\begin{itemize}
\item $m_e = 0.1278$ — pendiente metal entrada
\item $m_s = 1.38$ — pendiente metal salida
\item $c_x = 0.046\,\mathrm{m}$ — cuerda axial
\end{itemize}

\noindent\textbf{Incógnita}
$c$, en $\mathrm{m}$; también las magnitudes intermedias desarrolladas.

\noindent\textbf{Modelo físico}

\noindent\textit{Estator turbina y calado educativo}

\[
\fitmath{c=c_x \sqrt{\left(\frac{m_e}{2} + \frac{m_s}{2}\right)^{2} + 1}}
\]

\noindent\textbf{Hipótesis}
\begin{itemize}
\item Línea media Hermite educativa; x monótona.
\end{itemize}

\noindent\textbf{Desarrollo paso a paso}

\begin{enumerate}

\item El desplazamiento normalizado se prescribe como promedio de pendientes.

\[\fitmath{\mathrm{dd}=\frac{\mathrm{me}+\mathrm{ms}}{2}}\]

\[\fitmath{\mathrm{dd}=\frac{(0.1278)+(1.38)}{2}\approx 0.7539\,}\]

\item El calado de cuerda es atan d.

\[\fitmath{\mathrm{chi}=\arctan\left(\frac{\mathrm{me}+\mathrm{ms}}{2}\right)}\]

\[\fitmath{\mathrm{chi}=\arctan\left(\frac{(0.1278)+(1.38)}{2}\right)\approx 0.6459924\,\mathrm{rad}}\]

\item La cuerda real se deduce del desplazamiento final.

\[\fitmath{\mathrm{chord}=\mathrm{cxg}\cdot \sqrt{1+{\left(\frac{\mathrm{me}+\mathrm{ms}}{2}\right)}^{2}}}\]

\[\fitmath{\mathrm{chord}=(0.046)\cdot \sqrt{1+{\left(\frac{(0.1278)+(1.38)}{2}\right)}^{2}}\approx 0.05760782\,\mathrm{m}}\]

\end{enumerate}


\noindent\textbf{Sustitución numérica}

\[
\fitmath{\left(0.046\,\mathrm{m}\right) \sqrt{\left(\frac{\left(0.1278\right)}{2} + \frac{\left(1.38\right)}{2}\right)^{2} + 1}}
\]

\noindent\textbf{Resultado}

\[
\fitmath{c=5.760781878\times10^{-2}\,\mathrm{m}}
\]

\noindent\textbf{Interpretación}
Promediar pendientes no equivale a promediar ángulos. La ley es geométrica y no una correlación de mínima pérdida.

\noindent\textbf{Comprobación de los pasos}
\begin{itemize}
\item Paso 2: $0.7539$ — reproducido
\item Paso 4: $0.645992$ — reproducido
\item Paso 6: $0.0576078$ — reproducido
\end{itemize}

\noindent\textbf{Comprobación de la respuesta} la aritmética del último paso reproduce la respuesta enunciada a la precisión con que está escrita. Es una reproducción aritmética; no valida el modelo físico ni sus hipótesis.

\noindent\textbf{Limitaciones}
\begin{itemize}
\item Modelo educativo; la aplicabilidad a una máquina real requiere revisión experta.
\end{itemize}

\noindent\textbf{Errores frecuentes}
\begin{itemize}
\item Presentar esta sección como un perfil de fabricante.
\end{itemize}
% studium:end-record eyJpZCI6ICJQUkItMDE0MiIsICJzZWN0aW9uIjogImMxOCJ9
\end{studiumworked}

\begin{studiumworked}{Problema resuelto C18-04-R}


\noindent\textbf{Enunciado}
\(C18-04-R \cdot\) Rotor turbina y espesor por radio. En la salida del rotor de turbina, calcule beta y la coordenada transversal de la superficie superior en un punto de pendiente conocida. Datos: \(r=0.294\, \mathrm{m}\); \(\Omega=1120\, \mathrm{s}^{-1}\); \(K_s=-7\, \mathrm{m}^{2}\,\mathrm{s}^{-1}\); \(C_x=201\, \mathrm{m}\,\mathrm{s}^{-1}\); \(y=0.014\, \mathrm{m}\); \(t_n=0.002\, \mathrm{m}\); \(m_y=-0.62\).

\noindent\textbf{Datos}
\begin{itemize}
\item $r = 0.294\,\mathrm{m}$ — radio
\item $\Omega = 1120\,\mathrm{s}^{-1}$ — velocidad angular
\item $K_s = -7\,\mathrm{m}^{2}\,\mathrm{s}^{-1}$ — constante remolino salida
\item $C_x = 201\,\mathrm{m}\,\mathrm{s}^{-1}$ — velocidad axial
\item $y = 0.014\,\mathrm{m}$ — ordenada línea media en el punto
\item $t_n = 0.002\,\mathrm{m}$ — semiespesor normal
\item $m_y = -0.62$ — pendiente de línea media en el punto
\end{itemize}

\noindent\textbf{Incógnita}
$y_s$, en $\mathrm{m}$; también las magnitudes intermedias desarrolladas.

\noindent\textbf{Modelo físico}

\noindent\textit{Rotor turbina y espesor por radio}

\[
\fitmath{y_s=\frac{t_n}{\sqrt{m_y^{2} + 1}} + y}
\]

\noindent\textbf{Hipótesis}
\begin{itemize}
\item Vórtice libre y geometría normal local educativa.
\end{itemize}

\noindent\textbf{Desarrollo paso a paso}

\begin{enumerate}

\item El ángulo del gas se forma en referencia móvil.

\[\fitmath{\beta=\arctan\left(\frac{\frac{\mathrm{kv}}{\mathrm{rad}}-\Omega\cdot \mathrm{rad}}{\mathrm{cx}}\right)}\]

\[\fitmath{\beta=\arctan\left(\frac{\frac{(-7)}{(0.294)}-(1120)\cdot (0.294)}{(201)}\right)\approx -1.053286\,\mathrm{rad}}\]

\item La superficie superior desplaza la ordenada según la normal.

\[\fitmath{\mathrm{ys}=\mathrm{yy}+\frac{\mathrm{tt}}{\sqrt{1+{\mathrm{slope}}^{2}}}}\]

\[\fitmath{\mathrm{ys}=(0.014)+\frac{(0.002)}{\sqrt{1+{(-0.62)}^{2}}}\approx 0.01569981\,\mathrm{m}}\]

\end{enumerate}


\noindent\textbf{Sustitución numérica}

\[
\fitmath{\frac{\left(0.002\,\mathrm{m}\right)}{\sqrt{\left(-0.62\right)^{2} + 1}} + \left(0.014\,\mathrm{m}\right)}
\]

\noindent\textbf{Resultado}

\[
\fitmath{y_s=1.569980538\times10^{-2}\,\mathrm{m}}
\]

\noindent\textbf{Interpretación}
El dibujo completo debe desplazar también x. Esta salida tiene remolino negativo definido, sin cambiar el sentido de giro.

\noindent\textbf{Comprobación de los pasos}
\begin{itemize}
\item Paso 2: $-1.05329$ — reproducido
\item Paso 4: $0.0156998$ — reproducido
\end{itemize}

\noindent\textbf{Comprobación de la respuesta} la aritmética del último paso reproduce la respuesta enunciada a la precisión con que está escrita. Es una reproducción aritmética; no valida el modelo físico ni sus hipótesis.

\noindent\textbf{Limitaciones}
\begin{itemize}
\item Modelo educativo; la aplicabilidad a una máquina real requiere revisión experta.
\end{itemize}

\noindent\textbf{Errores frecuentes}
\begin{itemize}
\item Cambiar de signo U para representar una turbina.
\end{itemize}
% studium:end-record eyJpZCI6ICJQUkItMDE0MyIsICJzZWN0aW9uIjogImMxOCJ9
\end{studiumworked}

\begin{studiumworked}{Problema resuelto C18-04-P}


\noindent\textbf{Enunciado}
\(C18-04-P \cdot\) Rotor turbina y espesor por radio. Ejercicio de transferencia. En la salida del rotor de turbina, calcule beta y la coordenada transversal de la superficie superior en un punto de pendiente conocida. Antes de consultar la solución, prediga el efecto de cambiar el primer dato respecto del caso resuelto de esta familia. Justifique la tendencia y explique qué restricción podría convertirse en limitante. Datos: \(r=0.31311\, \mathrm{m}\); \(\Omega=1120\, \mathrm{s}^{-1}\); \(K_s=-7\, \mathrm{m}^{2}\,\mathrm{s}^{-1}\); \(C_x=201\, \mathrm{m}\,\mathrm{s}^{-1}\); \(y=0.014\, \mathrm{m}\); \(t_n=0.002\, \mathrm{m}\); \(m_y=-0.62\).

\noindent\textbf{Datos}
\begin{itemize}
\item $r = 0.31311\,\mathrm{m}$ — radio
\item $\Omega = 1120\,\mathrm{s}^{-1}$ — velocidad angular
\item $K_s = -7\,\mathrm{m}^{2}\,\mathrm{s}^{-1}$ — constante remolino salida
\item $C_x = 201\,\mathrm{m}\,\mathrm{s}^{-1}$ — velocidad axial
\item $y = 0.014\,\mathrm{m}$ — ordenada línea media en el punto
\item $t_n = 0.002\,\mathrm{m}$ — semiespesor normal
\item $m_y = -0.62$ — pendiente de línea media en el punto
\end{itemize}

\noindent\textbf{Incógnita}
$y_s$, en $\mathrm{m}$; también las magnitudes intermedias desarrolladas.

\noindent\textbf{Modelo físico}

\noindent\textit{Rotor turbina y espesor por radio}

\[
\fitmath{y_s=\frac{t_n}{\sqrt{m_y^{2} + 1}} + y}
\]

\noindent\textbf{Hipótesis}
\begin{itemize}
\item Vórtice libre y geometría normal local educativa.
\end{itemize}

\noindent\textbf{Desarrollo paso a paso}

\begin{enumerate}

\item El ángulo del gas se forma en referencia móvil.

\[\fitmath{\beta=\arctan\left(\frac{\frac{\mathrm{kv}}{\mathrm{rad}}-\Omega\cdot \mathrm{rad}}{\mathrm{cx}}\right)}\]

\[\fitmath{\beta=\arctan\left(\frac{\frac{(-7)}{(0.31311)}-(1120)\cdot (0.31311)}{(201)}\right)\approx -1.07658\,\mathrm{rad}}\]

\item La superficie superior desplaza la ordenada según la normal.

\[\fitmath{\mathrm{ys}=\mathrm{yy}+\frac{\mathrm{tt}}{\sqrt{1+{\mathrm{slope}}^{2}}}}\]

\[\fitmath{\mathrm{ys}=(0.014)+\frac{(0.002)}{\sqrt{1+{(-0.62)}^{2}}}\approx 0.01569981\,\mathrm{m}}\]

\end{enumerate}


\noindent\textbf{Sustitución numérica}

\[
\fitmath{\frac{\left(0.002\,\mathrm{m}\right)}{\sqrt{\left(-0.62\right)^{2} + 1}} + \left(0.014\,\mathrm{m}\right)}
\]

\noindent\textbf{Resultado}

\[
\fitmath{y_s=1.569980538\times10^{-2}\,\mathrm{m}}
\]

\noindent\textbf{Interpretación}
El dibujo completo debe desplazar también x. Esta salida tiene remolino negativo definido, sin cambiar el sentido de giro.

\noindent\textbf{Comprobación de los pasos}
\begin{itemize}
\item Paso 2: $-1.07658$ — reproducido
\item Paso 4: $0.0156998$ — reproducido
\end{itemize}

\noindent\textbf{Comprobación de la respuesta} la aritmética del último paso reproduce la respuesta enunciada a la precisión con que está escrita. Es una reproducción aritmética; no valida el modelo físico ni sus hipótesis.

\noindent\textbf{Limitaciones}
\begin{itemize}
\item Modelo educativo; la aplicabilidad a una máquina real requiere revisión experta.
\end{itemize}

\noindent\textbf{Errores frecuentes}
\begin{itemize}
\item Cambiar de signo U para representar una turbina.
\end{itemize}
% studium:end-record eyJpZCI6ICJQUkItMDE0NCIsICJzZWN0aW9uIjogImMxOCJ9
\end{studiumworked}

\chapter{Simulacros originales y estrategia de resolución}
\begin{studiumworked}{Problema resuelto C19-01-P}


\noindent\textbf{Enunciado}
\(C19-01-P \cdot A1 \cdot\) Cierre de compresor axial. Un rotor compresor al radio medio recibe Cx y remolinos dados. Obtenga trabajo, temperatura total final, presión total final y área geométrica de entrada. Resultado final: área de entrada. Datos: \(T_{0e}=302\, \mathrm{K}\); \(p_{0e}=118000\, \mathrm{Pa}\); \(C_x=156\, \mathrm{m}\,\mathrm{s}^{-1}\); \(C_{\theta e}=38\, \mathrm{m}\,\mathrm{s}^{-1}\); \(C_{\theta s}=147\, \mathrm{m}\,\mathrm{s}^{-1}\); \(U=291\, \mathrm{m}\,\mathrm{s}^{-1}\); \(\dot m=21\, \mathrm{kg}\,\mathrm{s}^{-1}\); \(\eta=0.87\); \(B=0.96\); \(c_p=1004.5\, \mathrm{J}\,\mathrm{kg}^{-1}\,\mathrm{K}^{-1}\); \(\gamma=1.4\); \(R_g=287\, \mathrm{J}\,\mathrm{kg}^{-1}\,\mathrm{K}^{-1}\).

\noindent\textbf{Datos}
\begin{itemize}
\item $T_{0e} = 302\,\mathrm{K}$ — temperatura total entrada
\item $p_{0e} = 118000\,\mathrm{Pa}$ — presión total entrada
\item $C_x = 156\,\mathrm{m}\,\mathrm{s}^{-1}$ — componente axial
\item $C_{\theta e} = 38\,\mathrm{m}\,\mathrm{s}^{-1}$ — remolino entrada rotor
\item $C_{\theta s} = 147\,\mathrm{m}\,\mathrm{s}^{-1}$ — remolino salida rotor
\item $U = 291\,\mathrm{m}\,\mathrm{s}^{-1}$ — velocidad periférica
\item $\dot m = 21\,\mathrm{kg}\,\mathrm{s}^{-1}$ — caudal
\item $\eta = 0.87$ — eficiencia total-total
\item $B = 0.96$ — fracción disponible
\item $c_p = 1004.5\,\mathrm{J}\,\mathrm{kg}^{-1}\,\mathrm{K}^{-1}$ — calor específico
\item $\gamma = 1.4$ — razón de calores
\item $R_g = 287\,\mathrm{J}\,\mathrm{kg}^{-1}\,\mathrm{K}^{-1}$ — constante gas
\end{itemize}

\noindent\textbf{Incógnita}
$A$, en $\mathrm{m^{2}}$; también las magnitudes intermedias desarrolladas.

\noindent\textbf{Modelo físico}

\noindent\textit{A1 \(\cdot\) Cierre de compresor axial}

\[
\fitmath{A=\frac{R_g \dot m \left(\frac{T_{0e} - \frac{C_{\theta e}^{2} + C_x^{2}}{2 c_p}}{T_{0e}}\right)^{- \frac{\gamma}{\gamma - 1}} \left(T_{0e} - \frac{C_{\theta e}^{2} + C_x^{2}}{2 c_p}\right)}{B C_x p_{0e}}}
\]

\noindent\textbf{Hipótesis}
\begin{itemize}
\item Gas perfecto, radio medio y Cx constantes, escalón adiabático.
\end{itemize}

\noindent\textbf{Desarrollo paso a paso}

\begin{enumerate}

\item Euler fija el trabajo.

\[\fitmath{\mathrm{wc}=\mathrm{uu}\cdot \left(\mathrm{cs}-\mathrm{ce}\right)}\]

\[\fitmath{\mathrm{wc}=(291)\cdot \left((147)-(38)\right)\approx 31719\,\mathrm{J/kg}}\]

\item El balance de energía fija la temperatura total de salida.

\[\fitmath{\mathrm{Tout}=\mathrm{Tin}+\frac{\mathrm{uu}\cdot \left(\mathrm{cs}-\mathrm{ce}\right)}{\mathrm{cp}}}\]

\[\fitmath{\mathrm{Tout}=(302)+\frac{(291)\cdot \left((147)-(38)\right)}{(1004.5)}\approx 333.5769\,\mathrm{K}}\]

\item La eficiencia determina presión total final.

\[\fitmath{\mathrm{pout}=\mathrm{pin}\cdot {\left(1+\frac{\eta\cdot \mathrm{uu}\cdot \left(\mathrm{cs}-\mathrm{ce}\right)}{\mathrm{cp}\cdot \mathrm{Tin}}\right)}^{\frac{\gamma}{\gamma-1}}}\]

\[\fitmath{\mathrm{pout}=(118000)\cdot {\left(1+\frac{(0.87)\cdot (291)\cdot \left((147)-(38)\right)}{(1004.5)\cdot (302)}\right)}^{\frac{(1.4)}{(1.4)-1}}\approx 160037.6\,\mathrm{Pa}}\]

\item La temperatura estática de entrada usa el módulo absoluto.

\[\fitmath{\mathrm{Te}=\mathrm{Tin}-\frac{{\mathrm{cx}}^{2}+{\mathrm{ce}}^{2}}{2\cdot \mathrm{cp}}}\]

\[\fitmath{\mathrm{Te}=(302)-\frac{{(156)}^{2}+{(38)}^{2}}{2\cdot (1004.5)}\approx 289.1677\,\mathrm{K}}\]

\item La continuidad utiliza densidad estática y componente axial.

\[\fitmath{\mathrm{area}=\frac{\mathrm{md}\cdot R_g\cdot \left(\mathrm{Tin}-\frac{{\mathrm{cx}}^{2}+{\mathrm{ce}}^{2}}{2\cdot \mathrm{cp}}\right)}{\mathrm{cx}\cdot \mathrm{bb}\cdot \mathrm{pin}\cdot {\left(\frac{\mathrm{Tin}-\frac{{\mathrm{cx}}^{2}+{\mathrm{ce}}^{2}}{2\cdot \mathrm{cp}}}{\mathrm{Tin}}\right)}^{\frac{\gamma}{\gamma-1}}}}\]

\[\fitmath{\mathrm{area}=\frac{(21)\cdot (287)\cdot \left((302)-\frac{{(156)}^{2}+{(38)}^{2}}{2\cdot (1004.5)}\right)}{(156)\cdot (0.96)\cdot (118000)\cdot {\left(\frac{(302)-\frac{{(156)}^{2}+{(38)}^{2}}{2\cdot (1004.5)}}{(302)}\right)}^{\frac{(1.4)}{(1.4)-1}}}\approx 0.1148082\,\mathrm{m^{2}}}\]

\end{enumerate}


\noindent\textbf{Sustitución numérica}

\[
\fitmath{\frac{\left(287\,\mathrm{J}\,\mathrm{kg}^{-1}\,\mathrm{K}^{-1}\right) \left(21\,\mathrm{kg}\,\mathrm{s}^{-1}\right) \left(\frac{\left(302\,\mathrm{K}\right) - \frac{\left(\left(38\,\mathrm{m}\,\mathrm{s}^{-1}\right)\right)^{2} + \left(\left(156\,\mathrm{m}\,\mathrm{s}^{-1}\right)\right)^{2}}{2 \left(1004.5\,\mathrm{J}\,\mathrm{kg}^{-1}\,\mathrm{K}^{-1}\right)}}{\left(302\,\mathrm{K}\right)}\right)^{- \frac{\left(1.4\right)}{\left(1.4\right) - 1}} \left(\left(302\,\mathrm{K}\right) - \frac{\left(\left(38\,\mathrm{m}\,\mathrm{s}^{-1}\right)\right)^{2} + \left(\left(156\,\mathrm{m}\,\mathrm{s}^{-1}\right)\right)^{2}}{2 \left(1004.5\,\mathrm{J}\,\mathrm{kg}^{-1}\,\mathrm{K}^{-1}\right)}\right)}{\left(0.96\right) \left(156\,\mathrm{m}\,\mathrm{s}^{-1}\right) \left(118000\,\mathrm{Pa}\right)}}
\]

\noindent\textbf{Resultado}

\[
\fitmath{A=1.148082328\times10^{-1}\,\mathrm{m}^{2}}
\]

\noindent\textbf{Interpretación}
Reconstruya masa con rho estática y B A. La presión final no determina la pérdida de cada fila por separado.

\noindent\textbf{Comprobación de los pasos}
\begin{itemize}
\item Paso 2: $31719$ — reproducido
\item Paso 4: $333.577$ — reproducido
\item Paso 6: $160038$ — reproducido
\item Paso 8: $289.168$ — reproducido
\item Paso 10: $0.114808$ — reproducido
\end{itemize}

\noindent\textbf{Comprobación de la respuesta} la aritmética del último paso reproduce la respuesta enunciada a la precisión con que está escrita. Es una reproducción aritmética; no valida el modelo físico ni sus hipótesis.

\noindent\textbf{Limitaciones}
\begin{itemize}
\item Modelo educativo; la aplicabilidad a una máquina real requiere revisión experta.
\end{itemize}

\noindent\textbf{Errores frecuentes}
\begin{itemize}
\item Usar densidad total para el anillo.
\end{itemize}
% studium:end-record eyJpZCI6ICJQUkItMDE0NSIsICJzZWN0aW9uIjogImMxOSJ9
\begin{center}\begin{tikzpicture}\draw[thick] (0,0)--(7,0);\draw[thick] (0,2.5)--(7,2.5);\node[draw,fill=orange!20,minimum width=2cm,minimum height=1.5cm] at (3.5,1.25) {Rotor};\draw[->,blue,thick] (.2,1.2)--(2,1.2) node[midway,above] {$C_x,C_{\theta e}$};\draw[->,blue,thick] (5,1.2)--(6.8,1.2) node[midway,above] {$C_x,C_{\theta s}$};\draw[->,orange,thick] (3.5,2.8)--(3.5,3.6) node[above] {$U>0$};\node[below,align=center] at (3.5,-0.7) {$w_c=3.1719\times10^{4}\,\mathrm{J}\,\mathrm{kg}^{-1},\quad T_{0s}=3.3358\times10^{2}\,\mathrm{K}$};\node[below,align=center] at (3.5,-1.2999999999999998) {$p_{0s}=1.6004\times10^{5}\,\mathrm{Pa},\quad T_e=2.8917\times10^{2}\,\mathrm{K}$};\node[below,align=center] at (3.5,-1.9) {$A=1.1481\times10^{-1}\,\mathrm{m}^{2}$};\end{tikzpicture}\end{center}
\end{studiumworked}

\begin{studiumworked}{Problema resuelto C19-02-P}


\noindent\textbf{Enunciado}
\(C19-02-P \cdot A2 \cdot\) Turbina con descarga axial. Diseñe un escalón repetitivo de turbina que descarga axialmente. Calcule carga, trabajo, temperatura total de salida y presión total de salida. Datos: \(T_{0e}=1490\, \mathrm{K}\); \(p_{0e}=940000\, \mathrm{Pa}\); \(\mathcal R=0.48\); \(U=314\, \mathrm{m}\,\mathrm{s}^{-1}\); \(\eta=0.895\); \(c_p=1156.697\, \mathrm{J}\,\mathrm{kg}^{-1}\,\mathrm{K}^{-1}\); \(\gamma=1.33\).

\noindent\textbf{Datos}
\begin{itemize}
\item $T_{0e} = 1490\,\mathrm{K}$ — temperatura total entrada
\item $p_{0e} = 940000\,\mathrm{Pa}$ — presión total entrada
\item $\mathcal R = 0.48$ — reacción
\item $U = 314\,\mathrm{m}\,\mathrm{s}^{-1}$ — velocidad periférica
\item $\eta = 0.895$ — eficiencia total-total
\item $c_p = 1156.696969697\,\mathrm{J}\,\mathrm{kg}^{-1}\,\mathrm{K}^{-1}$ — calor específico
\item $\gamma = 1.33$ — razón de calores
\end{itemize}

\noindent\textbf{Incógnita}
$p_{0s}$, en $\mathrm{Pa}$; también las magnitudes intermedias desarrolladas.

\noindent\textbf{Modelo físico}

\noindent\textit{A2 \(\cdot\) Turbina con descarga axial}

\[
\fitmath{p_{0s}=p_{0e} \left(1 - \frac{U^{2} \left(2 - 2 \mathcal R\right)}{T_{0e} c_p \eta}\right)^{\frac{\gamma}{\gamma - 1}}}
\]

\noindent\textbf{Hipótesis}
\begin{itemize}
\item Repetitivo, mismo radio y Cx constantes, sin refrigeración.
\end{itemize}

\noindent\textbf{Desarrollo paso a paso}

\begin{enumerate}

\item La descarga axial liga carga y reacción.

\[\fitmath{\mathrm{psi}=2\cdot \left(1-\mathrm{rea}\right)}\]

\[\fitmath{\mathrm{psi}=2\cdot \left(1-(0.48)\right)\approx 1.04\,}\]

\item Euler da el trabajo extraído.

\[\fitmath{\mathrm{wt}=2\cdot \left(1-\mathrm{rea}\right)\cdot {\mathrm{uu}}^{2}}\]

\[\fitmath{\mathrm{wt}=2\cdot \left(1-(0.48)\right)\cdot {(314)}^{2}\approx 102539.8\,\mathrm{J/kg}}\]

\item El remanso disminuye por el trabajo.

\[\fitmath{\mathrm{Tout}=\mathrm{Tin}-\frac{2\cdot \left(1-\mathrm{rea}\right)\cdot {\mathrm{uu}}^{2}}{\mathrm{cp}}}\]

\[\fitmath{\mathrm{Tout}=(1490)-\frac{2\cdot \left(1-(0.48)\right)\cdot {(314)}^{2}}{(1156.69697)}\approx 1401.351\,\mathrm{K}}\]

\item El estado ideal equivalente permite hallar presión total.

\[\fitmath{\mathrm{pout}=\mathrm{pin}\cdot {\left(1-\frac{2\cdot \left(1-\mathrm{rea}\right)\cdot {\mathrm{uu}}^{2}}{\mathrm{eta}\cdot \mathrm{cp}\cdot \mathrm{Tin}}\right)}^{\frac{\gamma}{\gamma-1}}}\]

\[\fitmath{\mathrm{pout}=(940000)\cdot {\left(1-\frac{2\cdot \left(1-(0.48)\right)\cdot {(314)}^{2}}{(0.895)\cdot (1156.69697)\cdot (1490)}\right)}^{\frac{(1.33)}{(1.33)-1}}\approx 712401.5\,\mathrm{Pa}}\]

\end{enumerate}


\noindent\textbf{Sustitución numérica}

\[
\fitmath{\left(940000\,\mathrm{Pa}\right) \left(1 - \frac{\left(\left(314\,\mathrm{m}\,\mathrm{s}^{-1}\right)\right)^{2} \left(2 - 2 \left(0.48\right)\right)}{\left(1490\,\mathrm{K}\right) \left(1156.696969697\,\mathrm{J}\,\mathrm{kg}^{-1}\,\mathrm{K}^{-1}\right) \left(0.895\right)}\right)^{\frac{\left(1.33\right)}{\left(1.33\right) - 1}}}
\]

\noindent\textbf{Resultado}

\[
\fitmath{p_{0s}=7.124015323\times10^{5}\,\mathrm{Pa}}
\]

\noindent\textbf{Interpretación}
La razón de presiones es menor que uno y la base isentrópica debe permanecer positiva. No se ha calculado la presión estática sin una velocidad de salida.

\noindent\textbf{Comprobación de los pasos}
\begin{itemize}
\item Paso 2: $1.04$ — reproducido
\item Paso 4: $102540$ — reproducido
\item Paso 6: $1401.35$ — reproducido
\item Paso 8: $712402$ — reproducido
\end{itemize}

\noindent\textbf{Comprobación de la respuesta} la aritmética del último paso reproduce la respuesta enunciada a la precisión con que está escrita. Es una reproducción aritmética; no valida el modelo físico ni sus hipótesis.

\noindent\textbf{Limitaciones}
\begin{itemize}
\item Modelo educativo; la aplicabilidad a una máquina real requiere revisión experta.
\end{itemize}

\noindent\textbf{Errores frecuentes}
\begin{itemize}
\item Imponer una carga incompatible con reacción y descarga axial.
\end{itemize}
% studium:end-record eyJpZCI6ICJQUkItMDE0NiIsICJzZWN0aW9uIjogImMxOSJ9
\begin{center}\begin{tikzpicture}\node[draw,fill=blue!15,minimum height=1.5cm] (s) at (2,1) {Estator};\node[draw,fill=orange!20,minimum height=1.5cm] (r) at (5,1) {Rotor};\draw[->] (0,1)--(s);\draw[->] (s)--(r) node[midway,above] {2};\draw[->] (r)--(7,1) node[right] {$C_{\theta3}=0$};\node at (0,1.6) {1};\node at (7,.4) {3};\draw[->,orange,thick] (5,2)--(5,2.8) node[above] {$U$};\node[below,align=center] at (3.5,-0.7) {$\psi=1.0400\times10^{0}\,,\quad w_t=1.0254\times10^{5}\,\mathrm{J}\,\mathrm{kg}^{-1}$};\node[below,align=center] at (3.5,-1.2999999999999998) {$T_{0s}=1.4014\times10^{3}\,\mathrm{K},\quad p_{0s}=7.1240\times10^{5}\,\mathrm{Pa}$};\end{tikzpicture}\end{center}
\end{studiumworked}

\begin{studiumworked}{Problema resuelto C19-03-P}


\noindent\textbf{Enunciado}
\(C19-03-P \cdot A3 \cdot\) Distribución radial y tracción. Un escalón compresor tiene leyes libres en ambas estaciones. Calcule trabajo radialmente uniforme, reacción en raíz y tensión centrífuga nominal de pala uniforme. Resultado final: tensión. Datos: \(\Omega=1090\, \mathrm{s}^{-1}\); \(K_e=14\, \mathrm{m}^{2}\,\mathrm{s}^{-1}\); \(K_s=44\, \mathrm{m}^{2}\,\mathrm{s}^{-1}\); \(r_h=0.218\, \mathrm{m}\); \(r_t=0.302\, \mathrm{m}\); \(\rho=8050\, \mathrm{kg}\,\mathrm{m}^{-3}\).

\noindent\textbf{Datos}
\begin{itemize}
\item $\Omega = 1090\,\mathrm{s}^{-1}$ — velocidad angular
\item $K_e = 14\,\mathrm{m}^{2}\,\mathrm{s}^{-1}$ — constante entrada
\item $K_s = 44\,\mathrm{m}^{2}\,\mathrm{s}^{-1}$ — constante salida
\item $r_h = 0.218\,\mathrm{m}$ — radio raíz
\item $r_t = 0.302\,\mathrm{m}$ — radio punta
\item $\rho = 8050\,\mathrm{kg}\,\mathrm{m}^{-3}$ — densidad material
\end{itemize}

\noindent\textbf{Incógnita}
$\sigma$, en $\mathrm{Pa}$; también las magnitudes intermedias desarrolladas.

\noindent\textbf{Modelo físico}

\noindent\textit{A3 \(\cdot\) Distribución radial y tracción}

\[
\fitmath{\sigma=\frac{\Omega^{2} \rho \left(- r_h^{2} + r_t^{2}\right)}{2}}
\]

\noindent\textbf{Hipótesis}
\begin{itemize}
\item Reacción reducida de escalón repetitivo; pala mecánica uniforme.
\end{itemize}

\noindent\textbf{Desarrollo paso a paso}

\begin{enumerate}

\item Euler hace uniforme el trabajo para estas leyes libres.

\[\fitmath{\mathrm{wc}=\Omega\cdot \left(\mathrm{ks}-\mathrm{ke}\right)}\]

\[\fitmath{\mathrm{wc}=(1090)\cdot \left((44)-(14)\right)\approx 32700\,\mathrm{J/kg}}\]

\item La reacción se evalúa con U local en raíz.

\[\fitmath{\mathrm{rea}=1-\frac{\mathrm{ke}+\mathrm{ks}}{2\cdot \Omega\cdot {\mathrm{rh}}^{2}}}\]

\[\fitmath{\mathrm{rea}=1-\frac{(14)+(44)}{2\cdot (1090)\cdot {(0.218)}^{2}}\approx 0.440167\,}\]

\item La tracción nominal integra toda la masa exterior.

\[\fitmath{\mathrm{stress}=\frac{\rho\cdot {\Omega}^{2}\cdot \left({\mathrm{rt}}^{2}-{\mathrm{rh}}^{2}\right)}{2}}\]

\[\fitmath{\mathrm{stress}=\frac{(8050)\cdot {(1090)}^{2}\cdot \left({(0.302)}^{2}-{(0.218)}^{2}\right)}{2}\approx 2.088822\times10^{8}\,\mathrm{Pa}}\]

\end{enumerate}


\noindent\textbf{Sustitución numérica}

\[
\fitmath{\frac{\left(\left(1090\,\mathrm{s}^{-1}\right)\right)^{2} \left(8050\,\mathrm{kg}\,\mathrm{m}^{-3}\right) \left(- \left(0.218\,\mathrm{m}\right)^{2} + \left(0.302\,\mathrm{m}\right)^{2}\right)}{2}}
\]

\noindent\textbf{Resultado}

\[
\fitmath{\sigma=2.088822372\times10^{8}\,\mathrm{Pa}}
\]

\noindent\textbf{Interpretación}
La admisibilidad aerodinámica y mecánica son comprobaciones distintas. No se ha suministrado una curva de vida ni concentración de raíz.

\noindent\textbf{Comprobación de los pasos}
\begin{itemize}
\item Paso 2: $32700$ — reproducido
\item Paso 4: $0.440167$ — reproducido
\item Paso 6: $2.08882\times 10^{8}$ — reproducido
\end{itemize}

\noindent\textbf{Comprobación de la respuesta} la aritmética del último paso reproduce la respuesta enunciada a la precisión con que está escrita. Es una reproducción aritmética; no valida el modelo físico ni sus hipótesis.

\noindent\textbf{Limitaciones}
\begin{itemize}
\item Modelo educativo; la aplicabilidad a una máquina real requiere revisión experta.
\end{itemize}

\noindent\textbf{Errores frecuentes}
\begin{itemize}
\item Usar la reacción media en toda la pala.
\end{itemize}
% studium:end-record eyJpZCI6ICJQUkItMDE0NyIsICJzZWN0aW9uIjogImMxOSJ9
\begin{center}\begin{tikzpicture}\draw[fill=gray!15,thick] (1,0) rectangle (2,3.7);\draw[thick] (.5,0)--(2.5,0);\foreach \x in {.6,.9,...,2.4}{\draw (\x,0)--(\x-.25,-.25);}\draw[->,orange] (1.5,2)--(1.5,4.5) node[above] {Inercia centrífuga};\draw[<->] (3.2,0)--(3.2,3.7) node[midway,right] {$r_t-r_h$};\node[left] at (1,0) {$r_h$};\node[left] at (1,3.7) {$r_t$};\node at (5.5,2.8) {$rC_\theta=K_v$};\node at (5.5,1.7) {$U=\Omega r$};\node[below,align=center] at (3.5,-0.7) {$w_c=3.2700\times10^{4}\,\mathrm{J}\,\mathrm{kg}^{-1},\quad \mathcal R=4.4017\times10^{-1}\,$};\node[below,align=center] at (3.5,-1.2999999999999998) {$\sigma=2.0888\times10^{8}\,\mathrm{Pa}$};\end{tikzpicture}\end{center}
\end{studiumworked}

\begin{studiumworked}{Problema resuelto C19-04-P}


\noindent\textbf{Enunciado}
\(C19-04-P \cdot B1 \cdot\) Rodete radial y prestación. Un rodete de salida radial sin prerrotación tiene radio y giro dados. Con slip y eficiencia prescritos, halle trabajo, temperatura total y relación de presión. Datos: \(\Omega=2070\, \mathrm{s}^{-1}\); \(r=0.206\, \mathrm{m}\); \(\sigma_s=0.875\); \(\eta=0.825\); \(T_{0e}=296\, \mathrm{K}\); \(c_p=1004.5\, \mathrm{J}\,\mathrm{kg}^{-1}\,\mathrm{K}^{-1}\); \(\gamma=1.4\).

\noindent\textbf{Datos}
\begin{itemize}
\item $\Omega = 2070\,\mathrm{s}^{-1}$ — velocidad angular
\item $r = 0.206\,\mathrm{m}$ — radio salida
\item $\sigma_s = 0.875$ — factor $C_\theta/C_{\theta,\mathrm{ideal}}$
\item $\eta = 0.825$ — eficiencia total-total
\item $T_{0e} = 296\,\mathrm{K}$ — temperatura total entrada
\item $c_p = 1004.5\,\mathrm{J}\,\mathrm{kg}^{-1}\,\mathrm{K}^{-1}$ — calor específico
\item $\gamma = 1.4$ — razón de calores
\end{itemize}

\noindent\textbf{Incógnita}
$\pi_c$ (adimensional); también las magnitudes intermedias desarrolladas.

\noindent\textbf{Modelo físico}

\noindent\textit{B1 \(\cdot\) Rodete radial y prestación}

\[
\fitmath{\pi_c=\left(\frac{\Omega^{2} \eta r^{2} \sigma_s}{T_{0e} c_p} + 1\right)^{\frac{\gamma}{\gamma - 1}}}
\]

\noindent\textbf{Hipótesis}
\begin{itemize}
\item Gas perfecto, salida radial ideal antes de slip; sin prerrotación.
\end{itemize}

\noindent\textbf{Desarrollo paso a paso}

\begin{enumerate}

\item La velocidad periférica depende del radio de salida.

\[\fitmath{\mathrm{uu}=\Omega\cdot \mathrm{rad}}\]

\[\fitmath{\mathrm{uu}=(2070)\cdot (0.206)\approx 426.42\,\mathrm{m/s}}\]

\item El deslizamiento modifica el remolino y el trabajo.

\[\fitmath{\mathrm{wc}=\mathrm{slip}\cdot {\left(\Omega\cdot \mathrm{rad}\right)}^{2}}\]

\[\fitmath{\mathrm{wc}=(0.875)\cdot {\left((2070)\cdot (0.206)\right)}^{2}\approx 159104.8\,\mathrm{J/kg}}\]

\item La temperatura total aumenta según el trabajo real.

\[\fitmath{\mathrm{Tout}=\mathrm{Tin}+\frac{\mathrm{slip}\cdot {\left(\Omega\cdot \mathrm{rad}\right)}^{2}}{\mathrm{cp}}}\]

\[\fitmath{\mathrm{Tout}=(296)+\frac{(0.875)\cdot {\left((2070)\cdot (0.206)\right)}^{2}}{(1004.5)}\approx 454.392\,\mathrm{K}}\]

\item La eficiencia enlaza trabajo con razón de presión.

\[\fitmath{\mathrm{pr}={\left(1+\frac{\mathrm{eta}\cdot \mathrm{slip}\cdot {\left(\Omega\cdot \mathrm{rad}\right)}^{2}}{\mathrm{cp}\cdot \mathrm{Tin}}\right)}^{\frac{\gamma}{\gamma-1}}}\]

\[\fitmath{\mathrm{pr}={\left(1+\frac{(0.825)\cdot (0.875)\cdot {\left((2070)\cdot (0.206)\right)}^{2}}{(1004.5)\cdot (296)}\right)}^{\frac{(1.4)}{(1.4)-1}}\approx 3.595949\,}\]

\end{enumerate}


\noindent\textbf{Sustitución numérica}

\[
\fitmath{\left(\frac{\left(\left(2070\,\mathrm{s}^{-1}\right)\right)^{2} \left(0.825\right) \left(0.206\,\mathrm{m}\right)^{2} \left(0.875\right)}{\left(296\,\mathrm{K}\right) \left(1004.5\,\mathrm{J}\,\mathrm{kg}^{-1}\,\mathrm{K}^{-1}\right)} + 1\right)^{\frac{\left(1.4\right)}{\left(1.4\right) - 1}}}
\]

\noindent\textbf{Resultado}

\[
\fitmath{\pi_c=3.595948824\times10^{0}}
\]

\noindent\textbf{Interpretación}
El resultado no es un mapa. Deben revisarse Mach de inductor, difusión, giro y tensión para juzgar realizabilidad.

\noindent\textbf{Comprobación de los pasos}
\begin{itemize}
\item Paso 2: $426.42$ — reproducido
\item Paso 4: $159105$ — reproducido
\item Paso 6: $454.392$ — reproducido
\item Paso 8: $3.59595$ — reproducido
\end{itemize}

\noindent\textbf{Comprobación de la respuesta} la aritmética del último paso reproduce la respuesta enunciada a la precisión con que está escrita. Es una reproducción aritmética; no valida el modelo físico ni sus hipótesis.

\noindent\textbf{Limitaciones}
\begin{itemize}
\item Modelo educativo; la aplicabilidad a una máquina real requiere revisión experta.
\end{itemize}

\noindent\textbf{Errores frecuentes}
\begin{itemize}
\item Confundir slip y rendimiento isentrópico.
\end{itemize}
% studium:end-record eyJpZCI6ICJQUkItMDE0OCIsICJzZWN0aW9uIjogImMxOSJ9
\begin{center}\begin{tikzpicture}\draw[thick] (2.5,1.7) circle (1.5);\draw[thick] (2.5,1.7) circle (.45);\foreach \a in {0,45,...,315}{\draw[blue,rotate around={\a:(2.5,1.7)}] (2.95,1.7)--(4,1.7);}\draw[->,orange,thick] (3,3.6) arc (70:140:1.8);\node at (2.5,4) {$\Omega>0$};\draw[->] (4,1.7)--(6,1.7) node[right] {$C_{r2}$};\draw[->] (4,1.7)--(4,2.8) node[right] {$C_{\theta2}$};\node at (5.7,.5) {Prerrotación nula};\node[below,align=center] at (3.5,-0.7) {$U=4.2642\times10^{2}\,\mathrm{m}\,\mathrm{s}^{-1},\quad w_c=1.5910\times10^{5}\,\mathrm{J}\,\mathrm{kg}^{-1}$};\node[below,align=center] at (3.5,-1.2999999999999998) {$T_{0s}=4.5439\times10^{2}\,\mathrm{K},\quad \pi_c=3.5959\times10^{0}\,$};\end{tikzpicture}\end{center}
\end{studiumworked}

\begin{studiumworked}{Problema resuelto C19-05-P}


\noindent\textbf{Enunciado}
\(C19-05-P \cdot B2 \cdot\) Cascada y elección entera. Una fila compresora tiene velocidades relativas y cambio tangencial dados. Se ha elegido un número entero de álabes; calcule paso, solidez real y difusión. Datos: \(r=0.279\, \mathrm{m}\); \(c=0.054\, \mathrm{m}\); \(Z=43\); \(W_e=296\, \mathrm{m}\,\mathrm{s}^{-1}\); \(W_s=224\, \mathrm{m}\,\mathrm{s}^{-1}\); \(|\Delta W_\theta|=103\, \mathrm{m}\,\mathrm{s}^{-1}\).

\noindent\textbf{Datos}
\begin{itemize}
\item $r = 0.279\,\mathrm{m}$ — radio
\item $c = 0.054\,\mathrm{m}$ — cuerda real
\item $Z = 43$ — número de álabes
\item $W_e = 296\,\mathrm{m}\,\mathrm{s}^{-1}$ — módulo relativo entrada
\item $W_s = 224\,\mathrm{m}\,\mathrm{s}^{-1}$ — módulo relativo salida
\item $|\Delta W_\theta| = 103\,\mathrm{m}\,\mathrm{s}^{-1}$ — magnitud de cambio tangencial
\end{itemize}

\noindent\textbf{Incógnita}
$D$ (adimensional); también las magnitudes intermedias desarrolladas.

\noindent\textbf{Modelo físico}

\noindent\textit{B2 \(\cdot\) Cascada y elección entera}

\[
\fitmath{D=1 - \frac{W_s}{W_e} + \frac{\pi |\Delta W_\theta| r}{c W_e Z}}
\]

\noindent\textbf{Hipótesis}
\begin{itemize}
\item Indicador de difusión definido en referencia relativa.
\end{itemize}

\noindent\textbf{Desarrollo paso a paso}

\begin{enumerate}

\item El entero fija el paso real.

\[\fitmath{\mathrm{ss}=\frac{2\cdot \pi\cdot \mathrm{rad}}{\mathrm{zz}}}\]

\[\fitmath{\mathrm{ss}=\frac{2\cdot \pi\cdot (0.279)}{(43)}\approx 0.04076764\,\mathrm{m}}\]

\item Se calcula la solidez alcanzada.

\[\fitmath{\mathrm{sol}=\frac{\mathrm{chord}\cdot \mathrm{zz}}{2\cdot \pi\cdot \mathrm{rad}}}\]

\[\fitmath{\mathrm{sol}=\frac{(0.054)\cdot (43)}{2\cdot \pi\cdot (0.279)}\approx 1.32458\,}\]

\item La difusión utiliza esa solidez, no una anterior al redondeo.

\[\fitmath{D=1-\frac{\mathrm{vs}}{\mathrm{ve}}+\frac{\mathrm{dv}}{2\cdot \frac{\mathrm{chord}\cdot \mathrm{zz}}{2\cdot \pi\cdot \mathrm{rad}}\cdot \mathrm{ve}}}\]

\[\fitmath{D=1-\frac{(224)}{(296)}+\frac{(103)}{2\cdot \frac{(0.054)\cdot (43)}{2\cdot \pi\cdot (0.279)}\cdot (296)}\approx 0.3745955\,}\]

\end{enumerate}


\noindent\textbf{Sustitución numérica}

\[
\fitmath{1 - \frac{\left(224\,\mathrm{m}\,\mathrm{s}^{-1}\right)}{\left(296\,\mathrm{m}\,\mathrm{s}^{-1}\right)} + \frac{\pi \left(103\,\mathrm{m}\,\mathrm{s}^{-1}\right) \left(0.279\,\mathrm{m}\right)}{\left(0.054\,\mathrm{m}\right) \left(296\,\mathrm{m}\,\mathrm{s}^{-1}\right) \left(43\right)}}
\]

\noindent\textbf{Resultado}

\[
\fitmath{D=3.745954502\times10^{-1}}
\]

\noindent\textbf{Interpretación}
Compare con un límite solo si se ha suministrado su rango de validez. D no determina una eficiencia del perfil.

\noindent\textbf{Comprobación de los pasos}
\begin{itemize}
\item Paso 2: $0.0407676$ — reproducido
\item Paso 4: $1.32458$ — reproducido
\item Paso 6: $0.374595$ — reproducido
\end{itemize}

\noindent\textbf{Comprobación de la respuesta} la aritmética del último paso reproduce la respuesta enunciada a la precisión con que está escrita. Es una reproducción aritmética; no valida el modelo físico ni sus hipótesis.

\noindent\textbf{Limitaciones}
\begin{itemize}
\item Modelo educativo; la aplicabilidad a una máquina real requiere revisión experta.
\end{itemize}

\noindent\textbf{Errores frecuentes}
\begin{itemize}
\item Mantener solidez objetivo tras redondear Z.
\end{itemize}
% studium:end-record eyJpZCI6ICJQUkItMDE0OSIsICJzZWN0aW9uIjogImMxOSJ9
\begin{center}\begin{tikzpicture}\foreach \y in {0,1.8,3.6}{\draw[blue,thick] (0,\y) .. controls (1,\y+.8) and (2.8,\y+.8) .. (4,\y+.3);\draw[blue,thick] (0,\y) .. controls (1,\y+.2) and (2.8,\y+.3) .. (4,\y+.3);}\draw[<->] (-.5,0)--(-.5,1.8) node[midway,left] {$s$};\draw[<->] (0,4.6)--(4,4.9) node[midway,above] {$c$};\node[align=center] at (6,2) {$Z_b$ entero\\$V_e,V_s,|\Delta V_\theta|$};\node[below,align=center] at (3.5,-0.7) {$s=4.0768\times10^{-2}\,\mathrm{m},\quad \sigma=1.3246\times10^{0}\,$};\node[below,align=center] at (3.5,-1.2999999999999998) {$D=3.7460\times10^{-1}\,$};\end{tikzpicture}\end{center}
\end{studiumworked}

\begin{studiumworked}{Problema resuelto C19-06-P}


\noindent\textbf{Enunciado}
\(C19-06-P \cdot B3 \cdot\) Balance térmico integrado. Una pared se refrigera por convección interior. Calcule flujo por área, temperatura caliente y caudal mínimo energético para el incremento permitido. Resultado final: caudal. Datos: \(T_{aw}=1580\, \mathrm{K}\); \(T_c=720\, \mathrm{K}\); \(h_g=1540\, \mathrm{W}\,\mathrm{m}^{-2}\,\mathrm{K}^{-1}\); \(h_c=2320\, \mathrm{W}\,\mathrm{m}^{-2}\,\mathrm{K}^{-1}\); \(t=0.0012\, \mathrm{m}\); \(k_m=23\, \mathrm{W}\,\mathrm{m}^{-1}\,\mathrm{K}^{-1}\); \(A=0.028\, \mathrm{m}^{2}\); \(c_p=1080\, \mathrm{J}\,\mathrm{kg}^{-1}\,\mathrm{K}^{-1}\); \(\Delta T_c=165\, \mathrm{K}\).

\noindent\textbf{Datos}
\begin{itemize}
\item $T_{aw} = 1580\,\mathrm{K}$ — temperatura recuperación exterior
\item $T_c = 720\,\mathrm{K}$ — temperatura refrigerante local media del modelo
\item $h_g = 1540\,\mathrm{W}\,\mathrm{m}^{-2}\,\mathrm{K}^{-1}$ — coeficiente exterior
\item $h_c = 2320\,\mathrm{W}\,\mathrm{m}^{-2}\,\mathrm{K}^{-1}$ — coeficiente interior
\item $t = 0.0012\,\mathrm{m}$ — espesor metal
\item $k_m = 23\,\mathrm{W}\,\mathrm{m}^{-1}\,\mathrm{K}^{-1}$ — conductividad
\item $A = 0.028\,\mathrm{m}^{2}$ — superficie intercambio
\item $c_p = 1080\,\mathrm{J}\,\mathrm{kg}^{-1}\,\mathrm{K}^{-1}$ — calor específico refrigerante
\item $\Delta T_c = 165\,\mathrm{K}$ — incremento permitido
\end{itemize}

\noindent\textbf{Incógnita}
$\dot{m}_c$, en $\mathrm{kg\,s^{-1}}$; también las magnitudes intermedias desarrolladas.

\noindent\textbf{Modelo físico}

\noindent\textit{B3 \(\cdot\) Balance térmico integrado}

\[
\fitmath{\dot{m}_c=\frac{A \left(- T_c + T_{aw}\right)}{c_p \Delta T_c \left(\frac{1}{h_g} + \frac{1}{h_c} + \frac{t}{k_m}\right)}}
\]

\noindent\textbf{Hipótesis}
\begin{itemize}
\item Red estacionaria unidimensional, propiedades constantes, sin radiación.
\end{itemize}

\noindent\textbf{Desarrollo paso a paso}

\begin{enumerate}

\item Se suma la resistencia térmica por área.

\[\fitmath{\mathrm{Res}=\frac{1}{\mathrm{hg}}+\frac{\mathrm{th}}{\mathrm{cond}}+\frac{1}{\mathrm{hc}}}\]

\[\fitmath{\mathrm{Res}=\frac{1}{(1540)}+\frac{(0.0012)}{(23)}+\frac{1}{(2320)}\approx 0.001132559\,\mathrm{m^{2}\,K/W}}\]

\item Se calcula el flujo estacionario.

\[\fitmath{\mathrm{flux}=\frac{\mathrm{Tg}-\mathrm{Tc}}{\frac{1}{\mathrm{hg}}+\frac{\mathrm{th}}{\mathrm{cond}}+\frac{1}{\mathrm{hc}}}}\]

\[\fitmath{\mathrm{flux}=\frac{(1580)-(720)}{\frac{1}{(1540)}+\frac{(0.0012)}{(23)}+\frac{1}{(2320)}}\approx 759342.3\,\mathrm{W/m^{2}}}\]

\item La cara caliente sigue del salto convectivo exterior.

\[\fitmath{\mathrm{Tw}=\mathrm{Tg}-\frac{\mathrm{Tg}-\mathrm{Tc}}{\mathrm{hg}\cdot \left(\frac{1}{\mathrm{hg}}+\frac{\mathrm{th}}{\mathrm{cond}}+\frac{1}{\mathrm{hc}}\right)}}\]

\[\fitmath{\mathrm{Tw}=(1580)-\frac{(1580)-(720)}{(1540)\cdot \left(\frac{1}{(1540)}+\frac{(0.0012)}{(23)}+\frac{1}{(2320)}\right)}\approx 1086.921\,\mathrm{K}}\]

\item La energía absorbida determina el mínimo de caudal.

\[\fitmath{\mathrm{mc}=\frac{\mathrm{area}\cdot \left(\mathrm{Tg}-\mathrm{Tc}\right)}{\left(\frac{1}{\mathrm{hg}}+\frac{\mathrm{th}}{\mathrm{cond}}+\frac{1}{\mathrm{hc}}\right)\cdot \mathrm{cp}\cdot \mathrm{dt}}}\]

\[\fitmath{\mathrm{mc}=\frac{(0.028)\cdot \left((1580)-(720)\right)}{\left(\frac{1}{(1540)}+\frac{(0.0012)}{(23)}+\frac{1}{(2320)}\right)\cdot (1080)\cdot (165)}\approx 0.119313\,\mathrm{kg/s}}\]

\end{enumerate}


\noindent\textbf{Sustitución numérica}

\[
\fitmath{\frac{\left(0.028\,\mathrm{m}^{2}\right) \left(- \left(720\,\mathrm{K}\right) + \left(1580\,\mathrm{K}\right)\right)}{\left(1080\,\mathrm{J}\,\mathrm{kg}^{-1}\,\mathrm{K}^{-1}\right) \left(165\,\mathrm{K}\right) \left(\frac{1}{\left(1540\,\mathrm{W}\,\mathrm{m}^{-2}\,\mathrm{K}^{-1}\right)} + \frac{1}{\left(2320\,\mathrm{W}\,\mathrm{m}^{-2}\,\mathrm{K}^{-1}\right)} + \frac{\left(0.0012\,\mathrm{m}\right)}{\left(23\,\mathrm{W}\,\mathrm{m}^{-1}\,\mathrm{K}^{-1}\right)}\right)}}
\]

\noindent\textbf{Resultado}

\[
\fitmath{\dot{m}_c=1.193130451\times10^{-1}\,\mathrm{kg}\,\mathrm{s}^{-1}}
\]

\noindent\textbf{Interpretación}
La temperatura Tc es un valor medio local prescrito para esta red. El mínimo energético no garantiza coeficiente interior ni presión de suministro.

\noindent\textbf{Comprobación de los pasos}
\begin{itemize}
\item Paso 2: $0.00113256$ — reproducido
\item Paso 4: $759342$ — reproducido
\item Paso 6: $1086.92$ — reproducido
\item Paso 8: $0.119313$ — reproducido
\end{itemize}

\noindent\textbf{Comprobación de la respuesta} la aritmética del último paso reproduce la respuesta enunciada a la precisión con que está escrita. Es una reproducción aritmética; no valida el modelo físico ni sus hipótesis.

\noindent\textbf{Limitaciones}
\begin{itemize}
\item Modelo educativo; la aplicabilidad a una máquina real requiere revisión experta.
\end{itemize}

\noindent\textbf{Errores frecuentes}
\begin{itemize}
\item Tratar el balance energético como diseño de refrigeración seguro.
\end{itemize}
% studium:end-record eyJpZCI6ICJQUkItMDE1MCIsICJzZWN0aW9uIjogImMxOSJ9
\begin{center}\begin{tikzpicture}\node (a) at (0,1) {$T_{aw}$};\node[draw,minimum height=1cm] (b) at (1.8,1) {$1/h_g$};\node[draw,minimum height=1cm] (c) at (3.8,1) {$t/k_m$};\node[draw,minimum height=1cm] (d) at (5.8,1) {$1/h_c$};\node (e) at (7.3,1) {$T_c$};\draw[->] (a)--(b);\draw[->] (b)--(c);\draw[->] (c)--(d);\draw[->] (d)--(e);\node at (2.8,1.8) {$T_{w,g}$};\draw[->,red,thick] (1,2.5)--(6.5,2.5) node[midway,above] {$q^{\prime\prime}$};\node[below,align=center] at (3.5,-0.7) {$R_{\mathrm{th}}=1.1326\times10^{-3}\,\mathrm{m}^{2}\,\mathrm{K}\,\mathrm{W}^{-1},\quad q''=7.5934\times10^{5}\,\mathrm{W}\,\mathrm{m}^{-2}$};\node[below,align=center] at (3.5,-1.2999999999999998) {$T_w=1.0869\times10^{3}\,\mathrm{K},\quad \dot{m}_c=1.1931\times10^{-1}\,\mathrm{kg}\,\mathrm{s}^{-1}$};\end{tikzpicture}\end{center}
\end{studiumworked}

\backmatter
\begin{thebibliography}{99}
\bibitem{guia} Universidad de León. \emph{Guía docente 2026/2027, 0710340: Turbomáquinas}. Documento de cinco páginas aportado por el usuario. Temario p. 2, metodología y evaluación pp. 3--4. Lectura del documento; no se afirma una validación criptográfica de su firma.
\bibitem{ule} Universidad de León. \emph{Material docente de Turbomáquinas}: presentación 2026/2027; enunciado del trabajo de hoja de cálculo; Generalidades; Diseño de escalón de compresor axial; Turbinas axiales; Esfuerzos y refrigeración; Compresores radiales. Siete PDF aportados, 144 páginas. Las referencias a diapositivas usan la página del PDF.
\bibitem{mit24} Martinez-Sanchez, M. \emph{16.50 Introduction to Propulsion Systems}, Spring 2012, lecciones 24, 26, 27, 28, 29, 30 y 32. MIT OpenCourseWare. Notas del curso basadas en el trabajo docente de Jack L. Kerrebrock. Consultadas para contraste de principios y modelos; imágenes de terceros excluidas. \url{https://ocw.mit.edu/courses/16-50-introduction-to-propulsion-systems-spring-2012/pages/lecture-notes/}. Licencia y condiciones: \url{https://ocw.mit.edu/terms/}.
\bibitem{hill} Hill, P. G.; Peterson, C. R. \emph{Mechanics and Thermodynamics of Propulsion}, 2.ª ed. La guía 2026/2027 identifica Pearson India Education, 2010, ISBN 978-81-317-2951-9. El catálogo editorial confirma título, autores, edición e ISBN; su lista de 2020 indica 2009 como año de esa referencia. Se conserva la fecha de la guía con esta discrepancia pendiente de comprobar en el ejemplar. La edición internacional de 1992, Addison-Wesley, ISBN 978-0-201-14659-2, es la citada en las diapositivas. Metadatos contrastados; libros no consultados íntegramente. \url{https://www.pearson.com/en-us/subject-catalog/p/mechanics-and-thermodynamics-of-propulsion/P200000003466/9780201146592}.
\bibitem{mattingly} Mattingly, J. D.; Heiser, W. H.; Pratt, D. T. \emph{Aircraft Engine Design}, 2.ª ed. AIAA, Reston, 2002. ISBN 978-1-56347-538-2. DOI: \url{https://doi.org/10.2514/4.861444}. Bibliografía recomendada por la guía; identidad contrastada con el registro editorial en Crossref. Texto comercial no consultado íntegramente.
\bibitem{oates} Oates, G. C. (ed.). \emph{Aerothermodynamics of Aircraft Engine Components}. AIAA, New York, 1985. ISBN 978-0-915928-97-2. DOI: \url{https://doi.org/10.2514/4.861338}. Bibliografía de la guía; metadatos contrastados con Crossref. Oates figura como editor. Texto comercial no consultado íntegramente.
\bibitem{lewis} Lewis, R. I. \emph{Turbomachinery Performance Analysis}. Arnold, London, 1996. ISBN 978-0-340-63191-1, según la guía. Autor, título, editorial y año contrastados en la página del autor: \url{https://www.staff.ncl.ac.uk/r.i.lewis/arnold.htm}. Texto no consultado íntegramente.
\bibitem{powers} \emph{Journal of Propulsion and Power}. AIAA. ISSN electrónico 1533-3876. La entrada sin título asociada a Joseph M. Powers en la guía se identifica por este ISSN; corresponde a una revista, no a un libro de Powers. \url{https://doi.org/10.2514/jpp}. No se atribuye una fórmula a un artículo no seleccionado y leído.
\bibitem{cumpsty} Cumpsty, N. A. \emph{Compressor Aerodynamics}. Reimpresión de Krieger, 2004, de la edición Longman de 1989. ISBN 978-1-57524-247-7. Referencia auxiliar de la presentación. Metadatos y disponibilidad abierta contrastados en Imperial College London: \url{https://hdl.handle.net/10044/1/128140}. Se consultó el registro bibliográfico, no el texto completo; no se usa como testigo de una fórmula concreta.
\bibitem{rr} Rolls-Royce. \emph{The Jet Engine}, 5.ª ed. Wiley, 2015. ISBN 978-1-119-06599-9. Referencia auxiliar de la presentación. Datos contrastados en el catálogo editorial: \url{https://www.wiley-vch.de/en/areas-interest/engineering/the-jet-engine-978-1-119-06599-9}. Texto no consultado íntegramente.
\bibitem{wuolah} \emph{Exámenes} y \emph{Problemas}, documentos estudiantiles anónimos distribuidos por Wuolah, 13 y 11 páginas. Referencias privadas no oficiales; únicamente análisis de competencias, sin reproducción de contenidos protegidos ni aceptación de resultados manuscritos.
\end{thebibliography}
\end{document}
